% !TeX program = lualatex
% ===================================================================
% sumario-longo.tex — mede o sumário sem hifenização (issue #35, "Fora do
% escopo"): com a hifenização desligada, a entrada longa de capítulo pode
% bater no número da página. Gera 3 x 16 capítulos, cada um com um prefixo
% (de 3 a 18 palavras) de uma frase longa: frase 1 de palavras médias,
% frase 2 com números, frase 3 só de palavras muito longas. Conta no log
% os "Overfull \hbox" do sumário: o sumário é a única parte do documento
% sem parágrafo de corpo, então todo Overfull "in paragraph" antes do
% primeiro \chapter é entrada que passou da margem.
%
% Uso (de uma pasta com livros_crp_acessivel_book.sty e crpsp-leg.sty no
% caminho; as fontes NEWJUNE por OSFONTDIR):
%   lualatex '\def\var{R}% !TeX program = lualatex
% ===================================================================
% sumario-longo.tex — mede o sumário sem hifenização (issue #35, "Fora do
% escopo"): com a hifenização desligada, a entrada longa de capítulo pode
% bater no número da página. Gera 3 x 16 capítulos, cada um com um prefixo
% (de 3 a 18 palavras) de uma frase longa: frase 1 de palavras médias,
% frase 2 com números, frase 3 só de palavras muito longas. Conta no log
% os "Overfull \hbox" do sumário: o sumário é a única parte do documento
% sem parágrafo de corpo, então todo Overfull "in paragraph" antes do
% primeiro \chapter é entrada que passou da margem.
%
% Uso (de uma pasta com livros_crp_acessivel_book.sty e crpsp-leg.sty no
% caminho; as fontes NEWJUNE por OSFONTDIR):
%   lualatex '\def\var{R}% !TeX program = lualatex
% ===================================================================
% sumario-longo.tex — mede o sumário sem hifenização (issue #35, "Fora do
% escopo"): com a hifenização desligada, a entrada longa de capítulo pode
% bater no número da página. Gera 3 x 16 capítulos, cada um com um prefixo
% (de 3 a 18 palavras) de uma frase longa: frase 1 de palavras médias,
% frase 2 com números, frase 3 só de palavras muito longas. Conta no log
% os "Overfull \hbox" do sumário: o sumário é a única parte do documento
% sem parágrafo de corpo, então todo Overfull "in paragraph" antes do
% primeiro \chapter é entrada que passou da margem.
%
% Uso (de uma pasta com livros_crp_acessivel_book.sty e crpsp-leg.sty no
% caminho; as fontes NEWJUNE por OSFONTDIR):
%   lualatex '\def\var{R}% !TeX program = lualatex
% ===================================================================
% sumario-longo.tex — mede o sumário sem hifenização (issue #35, "Fora do
% escopo"): com a hifenização desligada, a entrada longa de capítulo pode
% bater no número da página. Gera 3 x 16 capítulos, cada um com um prefixo
% (de 3 a 18 palavras) de uma frase longa: frase 1 de palavras médias,
% frase 2 com números, frase 3 só de palavras muito longas. Conta no log
% os "Overfull \hbox" do sumário: o sumário é a única parte do documento
% sem parágrafo de corpo, então todo Overfull "in paragraph" antes do
% primeiro \chapter é entrada que passou da margem.
%
% Uso (de uma pasta com livros_crp_acessivel_book.sty e crpsp-leg.sty no
% caminho; as fontes NEWJUNE por OSFONTDIR):
%   lualatex '\def\var{R}\input{sumario-longo}'   % R: rígido (até 0.5.x)
%   lualatex '\def\var{B}\input{sumario-longo}'   % B: plus 4em (o do guia)
%   lualatex '\def\var{C}\input{sumario-longo}'   % C: plus 4em + \tolerance 10000
%   lualatex '\def\var{S}\input{sumario-longo}'   % S: o que a .sty define
% Cada variante só redefine o \l@chapter depois da carga da .sty.
% ===================================================================
\DocumentMetadata{lang=pt-BR, pdfstandard=ua-2, pdfversion=2.0, tagging=on}
\documentclass[11pt,a5paper,twoside,openright]{book}
\usepackage{livros_crp_acessivel_book}
\hypersetup{pdftitle={Sumario longo}, pdfauthor={CRP SP}}
\makeatletter
\providecommand\var{S}
\def\crp@vR{R}\def\crp@vB{B}\def\crp@vC{C}
\def\crp@lchap#1{%  #1 = acréscimo à margem direita (rightskip)
  \renewcommand*\l@chapter[2]{%
    \ifnum \c@tocdepth >\m@ne
        \addpenalty{-\@highpenalty}\addvspace{0.5em \@plus\p@}%
        \setlength\@tempdima{1.5em}%
        \begingroup
            \parindent \z@ \rightskip \@pnumwidth #1
            \parfillskip -\@pnumwidth
            \leavevmode \bfseries
            \advance\leftskip\@tempdima \hskip -\leftskip
            ##1\nobreak\hfil
            \nobreak\hb@xt@\@pnumwidth{\hss ##2\kern-\p@\kern\p@}\par
            \penalty\@highpenalty
        \endgroup
    \fi}}
\ifx\var\crp@vR \crp@lchap{}\fi
\ifx\var\crp@vB \crp@lchap{\@plus 4em}\fi
\ifx\var\crp@vC \crp@lchap{\@plus 4em \tolerance=10000 \hbadness=10000 }\fi
\makeatother
\begin{document}
\frontmatter
\tableofcontents
\mainmatter
\chapter{Deficiência psicossocial e}
\chapter{Deficiência psicossocial e transtorno}
\chapter{Deficiência psicossocial e transtorno do}
\chapter{Deficiência psicossocial e transtorno do espectro}
\chapter{Deficiência psicossocial e transtorno do espectro autista:}
\chapter{Deficiência psicossocial e transtorno do espectro autista: orientações}
\chapter{Deficiência psicossocial e transtorno do espectro autista: orientações para}
\chapter{Deficiência psicossocial e transtorno do espectro autista: orientações para a}
\chapter{Deficiência psicossocial e transtorno do espectro autista: orientações para a realização}
\chapter{Deficiência psicossocial e transtorno do espectro autista: orientações para a realização de}
\chapter{Deficiência psicossocial e transtorno do espectro autista: orientações para a realização de eventos}
\chapter{Deficiência psicossocial e transtorno do espectro autista: orientações para a realização de eventos e}
\chapter{Deficiência psicossocial e transtorno do espectro autista: orientações para a realização de eventos e apresentações}
\chapter{Deficiência psicossocial e transtorno do espectro autista: orientações para a realização de eventos e apresentações acessíveis}
\chapter{Deficiência psicossocial e transtorno do espectro autista: orientações para a realização de eventos e apresentações acessíveis em}
\chapter{Deficiência psicossocial e transtorno do espectro autista: orientações para a realização de eventos e apresentações acessíveis em todo}
\chapter{Comunicação alternativa e}
\chapter{Comunicação alternativa e aumentativa}
\chapter{Comunicação alternativa e aumentativa na}
\chapter{Comunicação alternativa e aumentativa na participação}
\chapter{Comunicação alternativa e aumentativa na participação social}
\chapter{Comunicação alternativa e aumentativa na participação social de}
\chapter{Comunicação alternativa e aumentativa na participação social de pessoas}
\chapter{Comunicação alternativa e aumentativa na participação social de pessoas com}
\chapter{Comunicação alternativa e aumentativa na participação social de pessoas com deficiência}
\chapter{Comunicação alternativa e aumentativa na participação social de pessoas com deficiência intelectual}
\chapter{Comunicação alternativa e aumentativa na participação social de pessoas com deficiência intelectual em}
\chapter{Comunicação alternativa e aumentativa na participação social de pessoas com deficiência intelectual em eventos}
\chapter{Comunicação alternativa e aumentativa na participação social de pessoas com deficiência intelectual em eventos públicos}
\chapter{Comunicação alternativa e aumentativa na participação social de pessoas com deficiência intelectual em eventos públicos de}
\chapter{Comunicação alternativa e aumentativa na participação social de pessoas com deficiência intelectual em eventos públicos de grande}
\chapter{Comunicação alternativa e aumentativa na participação social de pessoas com deficiência intelectual em eventos públicos de grande porte}
\chapter{Psicopedagogicamente inconstitucionalíssimo otorrinolaringologista}
\chapter{Psicopedagogicamente inconstitucionalíssimo otorrinolaringologista socioeducativo}
\chapter{Psicopedagogicamente inconstitucionalíssimo otorrinolaringologista socioeducativo ético-profissional}
\chapter{Psicopedagogicamente inconstitucionalíssimo otorrinolaringologista socioeducativo ético-profissional sociodemográfico}
\chapter{Psicopedagogicamente inconstitucionalíssimo otorrinolaringologista socioeducativo ético-profissional sociodemográfico interdisciplinaridade}
\chapter{Psicopedagogicamente inconstitucionalíssimo otorrinolaringologista socioeducativo ético-profissional sociodemográfico interdisciplinaridade neuropsicológica}
\chapter{Psicopedagogicamente inconstitucionalíssimo otorrinolaringologista socioeducativo ético-profissional sociodemográfico interdisciplinaridade neuropsicológica avaliação}
\chapter{Psicopedagogicamente inconstitucionalíssimo otorrinolaringologista socioeducativo ético-profissional sociodemográfico interdisciplinaridade neuropsicológica avaliação}
\chapter{Psicopedagogicamente inconstitucionalíssimo otorrinolaringologista socioeducativo ético-profissional sociodemográfico interdisciplinaridade neuropsicológica avaliação}
\chapter{Psicopedagogicamente inconstitucionalíssimo otorrinolaringologista socioeducativo ético-profissional sociodemográfico interdisciplinaridade neuropsicológica avaliação}
\chapter{Psicopedagogicamente inconstitucionalíssimo otorrinolaringologista socioeducativo ético-profissional sociodemográfico interdisciplinaridade neuropsicológica avaliação}
\chapter{Psicopedagogicamente inconstitucionalíssimo otorrinolaringologista socioeducativo ético-profissional sociodemográfico interdisciplinaridade neuropsicológica avaliação}
\chapter{Psicopedagogicamente inconstitucionalíssimo otorrinolaringologista socioeducativo ético-profissional sociodemográfico interdisciplinaridade neuropsicológica avaliação}
\chapter{Psicopedagogicamente inconstitucionalíssimo otorrinolaringologista socioeducativo ético-profissional sociodemográfico interdisciplinaridade neuropsicológica avaliação}
\chapter{Psicopedagogicamente inconstitucionalíssimo otorrinolaringologista socioeducativo ético-profissional sociodemográfico interdisciplinaridade neuropsicológica avaliação}
\chapter{Psicopedagogicamente inconstitucionalíssimo otorrinolaringologista socioeducativo ético-profissional sociodemográfico interdisciplinaridade neuropsicológica avaliação}
\end{document}
'   % R: rígido (até 0.5.x)
%   lualatex '\def\var{B}% !TeX program = lualatex
% ===================================================================
% sumario-longo.tex — mede o sumário sem hifenização (issue #35, "Fora do
% escopo"): com a hifenização desligada, a entrada longa de capítulo pode
% bater no número da página. Gera 3 x 16 capítulos, cada um com um prefixo
% (de 3 a 18 palavras) de uma frase longa: frase 1 de palavras médias,
% frase 2 com números, frase 3 só de palavras muito longas. Conta no log
% os "Overfull \hbox" do sumário: o sumário é a única parte do documento
% sem parágrafo de corpo, então todo Overfull "in paragraph" antes do
% primeiro \chapter é entrada que passou da margem.
%
% Uso (de uma pasta com livros_crp_acessivel_book.sty e crpsp-leg.sty no
% caminho; as fontes NEWJUNE por OSFONTDIR):
%   lualatex '\def\var{R}\input{sumario-longo}'   % R: rígido (até 0.5.x)
%   lualatex '\def\var{B}\input{sumario-longo}'   % B: plus 4em (o do guia)
%   lualatex '\def\var{C}\input{sumario-longo}'   % C: plus 4em + \tolerance 10000
%   lualatex '\def\var{S}\input{sumario-longo}'   % S: o que a .sty define
% Cada variante só redefine o \l@chapter depois da carga da .sty.
% ===================================================================
\DocumentMetadata{lang=pt-BR, pdfstandard=ua-2, pdfversion=2.0, tagging=on}
\documentclass[11pt,a5paper,twoside,openright]{book}
\usepackage{livros_crp_acessivel_book}
\hypersetup{pdftitle={Sumario longo}, pdfauthor={CRP SP}}
\makeatletter
\providecommand\var{S}
\def\crp@vR{R}\def\crp@vB{B}\def\crp@vC{C}
\def\crp@lchap#1{%  #1 = acréscimo à margem direita (rightskip)
  \renewcommand*\l@chapter[2]{%
    \ifnum \c@tocdepth >\m@ne
        \addpenalty{-\@highpenalty}\addvspace{0.5em \@plus\p@}%
        \setlength\@tempdima{1.5em}%
        \begingroup
            \parindent \z@ \rightskip \@pnumwidth #1
            \parfillskip -\@pnumwidth
            \leavevmode \bfseries
            \advance\leftskip\@tempdima \hskip -\leftskip
            ##1\nobreak\hfil
            \nobreak\hb@xt@\@pnumwidth{\hss ##2\kern-\p@\kern\p@}\par
            \penalty\@highpenalty
        \endgroup
    \fi}}
\ifx\var\crp@vR \crp@lchap{}\fi
\ifx\var\crp@vB \crp@lchap{\@plus 4em}\fi
\ifx\var\crp@vC \crp@lchap{\@plus 4em \tolerance=10000 \hbadness=10000 }\fi
\makeatother
\begin{document}
\frontmatter
\tableofcontents
\mainmatter
\chapter{Deficiência psicossocial e}
\chapter{Deficiência psicossocial e transtorno}
\chapter{Deficiência psicossocial e transtorno do}
\chapter{Deficiência psicossocial e transtorno do espectro}
\chapter{Deficiência psicossocial e transtorno do espectro autista:}
\chapter{Deficiência psicossocial e transtorno do espectro autista: orientações}
\chapter{Deficiência psicossocial e transtorno do espectro autista: orientações para}
\chapter{Deficiência psicossocial e transtorno do espectro autista: orientações para a}
\chapter{Deficiência psicossocial e transtorno do espectro autista: orientações para a realização}
\chapter{Deficiência psicossocial e transtorno do espectro autista: orientações para a realização de}
\chapter{Deficiência psicossocial e transtorno do espectro autista: orientações para a realização de eventos}
\chapter{Deficiência psicossocial e transtorno do espectro autista: orientações para a realização de eventos e}
\chapter{Deficiência psicossocial e transtorno do espectro autista: orientações para a realização de eventos e apresentações}
\chapter{Deficiência psicossocial e transtorno do espectro autista: orientações para a realização de eventos e apresentações acessíveis}
\chapter{Deficiência psicossocial e transtorno do espectro autista: orientações para a realização de eventos e apresentações acessíveis em}
\chapter{Deficiência psicossocial e transtorno do espectro autista: orientações para a realização de eventos e apresentações acessíveis em todo}
\chapter{Comunicação alternativa e}
\chapter{Comunicação alternativa e aumentativa}
\chapter{Comunicação alternativa e aumentativa na}
\chapter{Comunicação alternativa e aumentativa na participação}
\chapter{Comunicação alternativa e aumentativa na participação social}
\chapter{Comunicação alternativa e aumentativa na participação social de}
\chapter{Comunicação alternativa e aumentativa na participação social de pessoas}
\chapter{Comunicação alternativa e aumentativa na participação social de pessoas com}
\chapter{Comunicação alternativa e aumentativa na participação social de pessoas com deficiência}
\chapter{Comunicação alternativa e aumentativa na participação social de pessoas com deficiência intelectual}
\chapter{Comunicação alternativa e aumentativa na participação social de pessoas com deficiência intelectual em}
\chapter{Comunicação alternativa e aumentativa na participação social de pessoas com deficiência intelectual em eventos}
\chapter{Comunicação alternativa e aumentativa na participação social de pessoas com deficiência intelectual em eventos públicos}
\chapter{Comunicação alternativa e aumentativa na participação social de pessoas com deficiência intelectual em eventos públicos de}
\chapter{Comunicação alternativa e aumentativa na participação social de pessoas com deficiência intelectual em eventos públicos de grande}
\chapter{Comunicação alternativa e aumentativa na participação social de pessoas com deficiência intelectual em eventos públicos de grande porte}
\chapter{Psicopedagogicamente inconstitucionalíssimo otorrinolaringologista}
\chapter{Psicopedagogicamente inconstitucionalíssimo otorrinolaringologista socioeducativo}
\chapter{Psicopedagogicamente inconstitucionalíssimo otorrinolaringologista socioeducativo ético-profissional}
\chapter{Psicopedagogicamente inconstitucionalíssimo otorrinolaringologista socioeducativo ético-profissional sociodemográfico}
\chapter{Psicopedagogicamente inconstitucionalíssimo otorrinolaringologista socioeducativo ético-profissional sociodemográfico interdisciplinaridade}
\chapter{Psicopedagogicamente inconstitucionalíssimo otorrinolaringologista socioeducativo ético-profissional sociodemográfico interdisciplinaridade neuropsicológica}
\chapter{Psicopedagogicamente inconstitucionalíssimo otorrinolaringologista socioeducativo ético-profissional sociodemográfico interdisciplinaridade neuropsicológica avaliação}
\chapter{Psicopedagogicamente inconstitucionalíssimo otorrinolaringologista socioeducativo ético-profissional sociodemográfico interdisciplinaridade neuropsicológica avaliação}
\chapter{Psicopedagogicamente inconstitucionalíssimo otorrinolaringologista socioeducativo ético-profissional sociodemográfico interdisciplinaridade neuropsicológica avaliação}
\chapter{Psicopedagogicamente inconstitucionalíssimo otorrinolaringologista socioeducativo ético-profissional sociodemográfico interdisciplinaridade neuropsicológica avaliação}
\chapter{Psicopedagogicamente inconstitucionalíssimo otorrinolaringologista socioeducativo ético-profissional sociodemográfico interdisciplinaridade neuropsicológica avaliação}
\chapter{Psicopedagogicamente inconstitucionalíssimo otorrinolaringologista socioeducativo ético-profissional sociodemográfico interdisciplinaridade neuropsicológica avaliação}
\chapter{Psicopedagogicamente inconstitucionalíssimo otorrinolaringologista socioeducativo ético-profissional sociodemográfico interdisciplinaridade neuropsicológica avaliação}
\chapter{Psicopedagogicamente inconstitucionalíssimo otorrinolaringologista socioeducativo ético-profissional sociodemográfico interdisciplinaridade neuropsicológica avaliação}
\chapter{Psicopedagogicamente inconstitucionalíssimo otorrinolaringologista socioeducativo ético-profissional sociodemográfico interdisciplinaridade neuropsicológica avaliação}
\chapter{Psicopedagogicamente inconstitucionalíssimo otorrinolaringologista socioeducativo ético-profissional sociodemográfico interdisciplinaridade neuropsicológica avaliação}
\end{document}
'   % B: plus 4em (o do guia)
%   lualatex '\def\var{C}% !TeX program = lualatex
% ===================================================================
% sumario-longo.tex — mede o sumário sem hifenização (issue #35, "Fora do
% escopo"): com a hifenização desligada, a entrada longa de capítulo pode
% bater no número da página. Gera 3 x 16 capítulos, cada um com um prefixo
% (de 3 a 18 palavras) de uma frase longa: frase 1 de palavras médias,
% frase 2 com números, frase 3 só de palavras muito longas. Conta no log
% os "Overfull \hbox" do sumário: o sumário é a única parte do documento
% sem parágrafo de corpo, então todo Overfull "in paragraph" antes do
% primeiro \chapter é entrada que passou da margem.
%
% Uso (de uma pasta com livros_crp_acessivel_book.sty e crpsp-leg.sty no
% caminho; as fontes NEWJUNE por OSFONTDIR):
%   lualatex '\def\var{R}\input{sumario-longo}'   % R: rígido (até 0.5.x)
%   lualatex '\def\var{B}\input{sumario-longo}'   % B: plus 4em (o do guia)
%   lualatex '\def\var{C}\input{sumario-longo}'   % C: plus 4em + \tolerance 10000
%   lualatex '\def\var{S}\input{sumario-longo}'   % S: o que a .sty define
% Cada variante só redefine o \l@chapter depois da carga da .sty.
% ===================================================================
\DocumentMetadata{lang=pt-BR, pdfstandard=ua-2, pdfversion=2.0, tagging=on}
\documentclass[11pt,a5paper,twoside,openright]{book}
\usepackage{livros_crp_acessivel_book}
\hypersetup{pdftitle={Sumario longo}, pdfauthor={CRP SP}}
\makeatletter
\providecommand\var{S}
\def\crp@vR{R}\def\crp@vB{B}\def\crp@vC{C}
\def\crp@lchap#1{%  #1 = acréscimo à margem direita (rightskip)
  \renewcommand*\l@chapter[2]{%
    \ifnum \c@tocdepth >\m@ne
        \addpenalty{-\@highpenalty}\addvspace{0.5em \@plus\p@}%
        \setlength\@tempdima{1.5em}%
        \begingroup
            \parindent \z@ \rightskip \@pnumwidth #1
            \parfillskip -\@pnumwidth
            \leavevmode \bfseries
            \advance\leftskip\@tempdima \hskip -\leftskip
            ##1\nobreak\hfil
            \nobreak\hb@xt@\@pnumwidth{\hss ##2\kern-\p@\kern\p@}\par
            \penalty\@highpenalty
        \endgroup
    \fi}}
\ifx\var\crp@vR \crp@lchap{}\fi
\ifx\var\crp@vB \crp@lchap{\@plus 4em}\fi
\ifx\var\crp@vC \crp@lchap{\@plus 4em \tolerance=10000 \hbadness=10000 }\fi
\makeatother
\begin{document}
\frontmatter
\tableofcontents
\mainmatter
\chapter{Deficiência psicossocial e}
\chapter{Deficiência psicossocial e transtorno}
\chapter{Deficiência psicossocial e transtorno do}
\chapter{Deficiência psicossocial e transtorno do espectro}
\chapter{Deficiência psicossocial e transtorno do espectro autista:}
\chapter{Deficiência psicossocial e transtorno do espectro autista: orientações}
\chapter{Deficiência psicossocial e transtorno do espectro autista: orientações para}
\chapter{Deficiência psicossocial e transtorno do espectro autista: orientações para a}
\chapter{Deficiência psicossocial e transtorno do espectro autista: orientações para a realização}
\chapter{Deficiência psicossocial e transtorno do espectro autista: orientações para a realização de}
\chapter{Deficiência psicossocial e transtorno do espectro autista: orientações para a realização de eventos}
\chapter{Deficiência psicossocial e transtorno do espectro autista: orientações para a realização de eventos e}
\chapter{Deficiência psicossocial e transtorno do espectro autista: orientações para a realização de eventos e apresentações}
\chapter{Deficiência psicossocial e transtorno do espectro autista: orientações para a realização de eventos e apresentações acessíveis}
\chapter{Deficiência psicossocial e transtorno do espectro autista: orientações para a realização de eventos e apresentações acessíveis em}
\chapter{Deficiência psicossocial e transtorno do espectro autista: orientações para a realização de eventos e apresentações acessíveis em todo}
\chapter{Comunicação alternativa e}
\chapter{Comunicação alternativa e aumentativa}
\chapter{Comunicação alternativa e aumentativa na}
\chapter{Comunicação alternativa e aumentativa na participação}
\chapter{Comunicação alternativa e aumentativa na participação social}
\chapter{Comunicação alternativa e aumentativa na participação social de}
\chapter{Comunicação alternativa e aumentativa na participação social de pessoas}
\chapter{Comunicação alternativa e aumentativa na participação social de pessoas com}
\chapter{Comunicação alternativa e aumentativa na participação social de pessoas com deficiência}
\chapter{Comunicação alternativa e aumentativa na participação social de pessoas com deficiência intelectual}
\chapter{Comunicação alternativa e aumentativa na participação social de pessoas com deficiência intelectual em}
\chapter{Comunicação alternativa e aumentativa na participação social de pessoas com deficiência intelectual em eventos}
\chapter{Comunicação alternativa e aumentativa na participação social de pessoas com deficiência intelectual em eventos públicos}
\chapter{Comunicação alternativa e aumentativa na participação social de pessoas com deficiência intelectual em eventos públicos de}
\chapter{Comunicação alternativa e aumentativa na participação social de pessoas com deficiência intelectual em eventos públicos de grande}
\chapter{Comunicação alternativa e aumentativa na participação social de pessoas com deficiência intelectual em eventos públicos de grande porte}
\chapter{Psicopedagogicamente inconstitucionalíssimo otorrinolaringologista}
\chapter{Psicopedagogicamente inconstitucionalíssimo otorrinolaringologista socioeducativo}
\chapter{Psicopedagogicamente inconstitucionalíssimo otorrinolaringologista socioeducativo ético-profissional}
\chapter{Psicopedagogicamente inconstitucionalíssimo otorrinolaringologista socioeducativo ético-profissional sociodemográfico}
\chapter{Psicopedagogicamente inconstitucionalíssimo otorrinolaringologista socioeducativo ético-profissional sociodemográfico interdisciplinaridade}
\chapter{Psicopedagogicamente inconstitucionalíssimo otorrinolaringologista socioeducativo ético-profissional sociodemográfico interdisciplinaridade neuropsicológica}
\chapter{Psicopedagogicamente inconstitucionalíssimo otorrinolaringologista socioeducativo ético-profissional sociodemográfico interdisciplinaridade neuropsicológica avaliação}
\chapter{Psicopedagogicamente inconstitucionalíssimo otorrinolaringologista socioeducativo ético-profissional sociodemográfico interdisciplinaridade neuropsicológica avaliação}
\chapter{Psicopedagogicamente inconstitucionalíssimo otorrinolaringologista socioeducativo ético-profissional sociodemográfico interdisciplinaridade neuropsicológica avaliação}
\chapter{Psicopedagogicamente inconstitucionalíssimo otorrinolaringologista socioeducativo ético-profissional sociodemográfico interdisciplinaridade neuropsicológica avaliação}
\chapter{Psicopedagogicamente inconstitucionalíssimo otorrinolaringologista socioeducativo ético-profissional sociodemográfico interdisciplinaridade neuropsicológica avaliação}
\chapter{Psicopedagogicamente inconstitucionalíssimo otorrinolaringologista socioeducativo ético-profissional sociodemográfico interdisciplinaridade neuropsicológica avaliação}
\chapter{Psicopedagogicamente inconstitucionalíssimo otorrinolaringologista socioeducativo ético-profissional sociodemográfico interdisciplinaridade neuropsicológica avaliação}
\chapter{Psicopedagogicamente inconstitucionalíssimo otorrinolaringologista socioeducativo ético-profissional sociodemográfico interdisciplinaridade neuropsicológica avaliação}
\chapter{Psicopedagogicamente inconstitucionalíssimo otorrinolaringologista socioeducativo ético-profissional sociodemográfico interdisciplinaridade neuropsicológica avaliação}
\chapter{Psicopedagogicamente inconstitucionalíssimo otorrinolaringologista socioeducativo ético-profissional sociodemográfico interdisciplinaridade neuropsicológica avaliação}
\end{document}
'   % C: plus 4em + \tolerance 10000
%   lualatex '\def\var{S}% !TeX program = lualatex
% ===================================================================
% sumario-longo.tex — mede o sumário sem hifenização (issue #35, "Fora do
% escopo"): com a hifenização desligada, a entrada longa de capítulo pode
% bater no número da página. Gera 3 x 16 capítulos, cada um com um prefixo
% (de 3 a 18 palavras) de uma frase longa: frase 1 de palavras médias,
% frase 2 com números, frase 3 só de palavras muito longas. Conta no log
% os "Overfull \hbox" do sumário: o sumário é a única parte do documento
% sem parágrafo de corpo, então todo Overfull "in paragraph" antes do
% primeiro \chapter é entrada que passou da margem.
%
% Uso (de uma pasta com livros_crp_acessivel_book.sty e crpsp-leg.sty no
% caminho; as fontes NEWJUNE por OSFONTDIR):
%   lualatex '\def\var{R}\input{sumario-longo}'   % R: rígido (até 0.5.x)
%   lualatex '\def\var{B}\input{sumario-longo}'   % B: plus 4em (o do guia)
%   lualatex '\def\var{C}\input{sumario-longo}'   % C: plus 4em + \tolerance 10000
%   lualatex '\def\var{S}\input{sumario-longo}'   % S: o que a .sty define
% Cada variante só redefine o \l@chapter depois da carga da .sty.
% ===================================================================
\DocumentMetadata{lang=pt-BR, pdfstandard=ua-2, pdfversion=2.0, tagging=on}
\documentclass[11pt,a5paper,twoside,openright]{book}
\usepackage{livros_crp_acessivel_book}
\hypersetup{pdftitle={Sumario longo}, pdfauthor={CRP SP}}
\makeatletter
\providecommand\var{S}
\def\crp@vR{R}\def\crp@vB{B}\def\crp@vC{C}
\def\crp@lchap#1{%  #1 = acréscimo à margem direita (rightskip)
  \renewcommand*\l@chapter[2]{%
    \ifnum \c@tocdepth >\m@ne
        \addpenalty{-\@highpenalty}\addvspace{0.5em \@plus\p@}%
        \setlength\@tempdima{1.5em}%
        \begingroup
            \parindent \z@ \rightskip \@pnumwidth #1
            \parfillskip -\@pnumwidth
            \leavevmode \bfseries
            \advance\leftskip\@tempdima \hskip -\leftskip
            ##1\nobreak\hfil
            \nobreak\hb@xt@\@pnumwidth{\hss ##2\kern-\p@\kern\p@}\par
            \penalty\@highpenalty
        \endgroup
    \fi}}
\ifx\var\crp@vR \crp@lchap{}\fi
\ifx\var\crp@vB \crp@lchap{\@plus 4em}\fi
\ifx\var\crp@vC \crp@lchap{\@plus 4em \tolerance=10000 \hbadness=10000 }\fi
\makeatother
\begin{document}
\frontmatter
\tableofcontents
\mainmatter
\chapter{Deficiência psicossocial e}
\chapter{Deficiência psicossocial e transtorno}
\chapter{Deficiência psicossocial e transtorno do}
\chapter{Deficiência psicossocial e transtorno do espectro}
\chapter{Deficiência psicossocial e transtorno do espectro autista:}
\chapter{Deficiência psicossocial e transtorno do espectro autista: orientações}
\chapter{Deficiência psicossocial e transtorno do espectro autista: orientações para}
\chapter{Deficiência psicossocial e transtorno do espectro autista: orientações para a}
\chapter{Deficiência psicossocial e transtorno do espectro autista: orientações para a realização}
\chapter{Deficiência psicossocial e transtorno do espectro autista: orientações para a realização de}
\chapter{Deficiência psicossocial e transtorno do espectro autista: orientações para a realização de eventos}
\chapter{Deficiência psicossocial e transtorno do espectro autista: orientações para a realização de eventos e}
\chapter{Deficiência psicossocial e transtorno do espectro autista: orientações para a realização de eventos e apresentações}
\chapter{Deficiência psicossocial e transtorno do espectro autista: orientações para a realização de eventos e apresentações acessíveis}
\chapter{Deficiência psicossocial e transtorno do espectro autista: orientações para a realização de eventos e apresentações acessíveis em}
\chapter{Deficiência psicossocial e transtorno do espectro autista: orientações para a realização de eventos e apresentações acessíveis em todo}
\chapter{Comunicação alternativa e}
\chapter{Comunicação alternativa e aumentativa}
\chapter{Comunicação alternativa e aumentativa na}
\chapter{Comunicação alternativa e aumentativa na participação}
\chapter{Comunicação alternativa e aumentativa na participação social}
\chapter{Comunicação alternativa e aumentativa na participação social de}
\chapter{Comunicação alternativa e aumentativa na participação social de pessoas}
\chapter{Comunicação alternativa e aumentativa na participação social de pessoas com}
\chapter{Comunicação alternativa e aumentativa na participação social de pessoas com deficiência}
\chapter{Comunicação alternativa e aumentativa na participação social de pessoas com deficiência intelectual}
\chapter{Comunicação alternativa e aumentativa na participação social de pessoas com deficiência intelectual em}
\chapter{Comunicação alternativa e aumentativa na participação social de pessoas com deficiência intelectual em eventos}
\chapter{Comunicação alternativa e aumentativa na participação social de pessoas com deficiência intelectual em eventos públicos}
\chapter{Comunicação alternativa e aumentativa na participação social de pessoas com deficiência intelectual em eventos públicos de}
\chapter{Comunicação alternativa e aumentativa na participação social de pessoas com deficiência intelectual em eventos públicos de grande}
\chapter{Comunicação alternativa e aumentativa na participação social de pessoas com deficiência intelectual em eventos públicos de grande porte}
\chapter{Psicopedagogicamente inconstitucionalíssimo otorrinolaringologista}
\chapter{Psicopedagogicamente inconstitucionalíssimo otorrinolaringologista socioeducativo}
\chapter{Psicopedagogicamente inconstitucionalíssimo otorrinolaringologista socioeducativo ético-profissional}
\chapter{Psicopedagogicamente inconstitucionalíssimo otorrinolaringologista socioeducativo ético-profissional sociodemográfico}
\chapter{Psicopedagogicamente inconstitucionalíssimo otorrinolaringologista socioeducativo ético-profissional sociodemográfico interdisciplinaridade}
\chapter{Psicopedagogicamente inconstitucionalíssimo otorrinolaringologista socioeducativo ético-profissional sociodemográfico interdisciplinaridade neuropsicológica}
\chapter{Psicopedagogicamente inconstitucionalíssimo otorrinolaringologista socioeducativo ético-profissional sociodemográfico interdisciplinaridade neuropsicológica avaliação}
\chapter{Psicopedagogicamente inconstitucionalíssimo otorrinolaringologista socioeducativo ético-profissional sociodemográfico interdisciplinaridade neuropsicológica avaliação}
\chapter{Psicopedagogicamente inconstitucionalíssimo otorrinolaringologista socioeducativo ético-profissional sociodemográfico interdisciplinaridade neuropsicológica avaliação}
\chapter{Psicopedagogicamente inconstitucionalíssimo otorrinolaringologista socioeducativo ético-profissional sociodemográfico interdisciplinaridade neuropsicológica avaliação}
\chapter{Psicopedagogicamente inconstitucionalíssimo otorrinolaringologista socioeducativo ético-profissional sociodemográfico interdisciplinaridade neuropsicológica avaliação}
\chapter{Psicopedagogicamente inconstitucionalíssimo otorrinolaringologista socioeducativo ético-profissional sociodemográfico interdisciplinaridade neuropsicológica avaliação}
\chapter{Psicopedagogicamente inconstitucionalíssimo otorrinolaringologista socioeducativo ético-profissional sociodemográfico interdisciplinaridade neuropsicológica avaliação}
\chapter{Psicopedagogicamente inconstitucionalíssimo otorrinolaringologista socioeducativo ético-profissional sociodemográfico interdisciplinaridade neuropsicológica avaliação}
\chapter{Psicopedagogicamente inconstitucionalíssimo otorrinolaringologista socioeducativo ético-profissional sociodemográfico interdisciplinaridade neuropsicológica avaliação}
\chapter{Psicopedagogicamente inconstitucionalíssimo otorrinolaringologista socioeducativo ético-profissional sociodemográfico interdisciplinaridade neuropsicológica avaliação}
\end{document}
'   % S: o que a .sty define
% Cada variante só redefine o \l@chapter depois da carga da .sty.
% ===================================================================
\DocumentMetadata{lang=pt-BR, pdfstandard=ua-2, pdfversion=2.0, tagging=on}
\documentclass[11pt,a5paper,twoside,openright]{book}
\usepackage{livros_crp_acessivel_book}
\hypersetup{pdftitle={Sumario longo}, pdfauthor={CRP SP}}
\makeatletter
\providecommand\var{S}
\def\crp@vR{R}\def\crp@vB{B}\def\crp@vC{C}
\def\crp@lchap#1{%  #1 = acréscimo à margem direita (rightskip)
  \renewcommand*\l@chapter[2]{%
    \ifnum \c@tocdepth >\m@ne
        \addpenalty{-\@highpenalty}\addvspace{0.5em \@plus\p@}%
        \setlength\@tempdima{1.5em}%
        \begingroup
            \parindent \z@ \rightskip \@pnumwidth #1
            \parfillskip -\@pnumwidth
            \leavevmode \bfseries
            \advance\leftskip\@tempdima \hskip -\leftskip
            ##1\nobreak\hfil
            \nobreak\hb@xt@\@pnumwidth{\hss ##2\kern-\p@\kern\p@}\par
            \penalty\@highpenalty
        \endgroup
    \fi}}
\ifx\var\crp@vR \crp@lchap{}\fi
\ifx\var\crp@vB \crp@lchap{\@plus 4em}\fi
\ifx\var\crp@vC \crp@lchap{\@plus 4em \tolerance=10000 \hbadness=10000 }\fi
\makeatother
\begin{document}
\frontmatter
\tableofcontents
\mainmatter
\chapter{Deficiência psicossocial e}
\chapter{Deficiência psicossocial e transtorno}
\chapter{Deficiência psicossocial e transtorno do}
\chapter{Deficiência psicossocial e transtorno do espectro}
\chapter{Deficiência psicossocial e transtorno do espectro autista:}
\chapter{Deficiência psicossocial e transtorno do espectro autista: orientações}
\chapter{Deficiência psicossocial e transtorno do espectro autista: orientações para}
\chapter{Deficiência psicossocial e transtorno do espectro autista: orientações para a}
\chapter{Deficiência psicossocial e transtorno do espectro autista: orientações para a realização}
\chapter{Deficiência psicossocial e transtorno do espectro autista: orientações para a realização de}
\chapter{Deficiência psicossocial e transtorno do espectro autista: orientações para a realização de eventos}
\chapter{Deficiência psicossocial e transtorno do espectro autista: orientações para a realização de eventos e}
\chapter{Deficiência psicossocial e transtorno do espectro autista: orientações para a realização de eventos e apresentações}
\chapter{Deficiência psicossocial e transtorno do espectro autista: orientações para a realização de eventos e apresentações acessíveis}
\chapter{Deficiência psicossocial e transtorno do espectro autista: orientações para a realização de eventos e apresentações acessíveis em}
\chapter{Deficiência psicossocial e transtorno do espectro autista: orientações para a realização de eventos e apresentações acessíveis em todo}
\chapter{Comunicação alternativa e}
\chapter{Comunicação alternativa e aumentativa}
\chapter{Comunicação alternativa e aumentativa na}
\chapter{Comunicação alternativa e aumentativa na participação}
\chapter{Comunicação alternativa e aumentativa na participação social}
\chapter{Comunicação alternativa e aumentativa na participação social de}
\chapter{Comunicação alternativa e aumentativa na participação social de pessoas}
\chapter{Comunicação alternativa e aumentativa na participação social de pessoas com}
\chapter{Comunicação alternativa e aumentativa na participação social de pessoas com deficiência}
\chapter{Comunicação alternativa e aumentativa na participação social de pessoas com deficiência intelectual}
\chapter{Comunicação alternativa e aumentativa na participação social de pessoas com deficiência intelectual em}
\chapter{Comunicação alternativa e aumentativa na participação social de pessoas com deficiência intelectual em eventos}
\chapter{Comunicação alternativa e aumentativa na participação social de pessoas com deficiência intelectual em eventos públicos}
\chapter{Comunicação alternativa e aumentativa na participação social de pessoas com deficiência intelectual em eventos públicos de}
\chapter{Comunicação alternativa e aumentativa na participação social de pessoas com deficiência intelectual em eventos públicos de grande}
\chapter{Comunicação alternativa e aumentativa na participação social de pessoas com deficiência intelectual em eventos públicos de grande porte}
\chapter{Psicopedagogicamente inconstitucionalíssimo otorrinolaringologista}
\chapter{Psicopedagogicamente inconstitucionalíssimo otorrinolaringologista socioeducativo}
\chapter{Psicopedagogicamente inconstitucionalíssimo otorrinolaringologista socioeducativo ético-profissional}
\chapter{Psicopedagogicamente inconstitucionalíssimo otorrinolaringologista socioeducativo ético-profissional sociodemográfico}
\chapter{Psicopedagogicamente inconstitucionalíssimo otorrinolaringologista socioeducativo ético-profissional sociodemográfico interdisciplinaridade}
\chapter{Psicopedagogicamente inconstitucionalíssimo otorrinolaringologista socioeducativo ético-profissional sociodemográfico interdisciplinaridade neuropsicológica}
\chapter{Psicopedagogicamente inconstitucionalíssimo otorrinolaringologista socioeducativo ético-profissional sociodemográfico interdisciplinaridade neuropsicológica avaliação}
\chapter{Psicopedagogicamente inconstitucionalíssimo otorrinolaringologista socioeducativo ético-profissional sociodemográfico interdisciplinaridade neuropsicológica avaliação}
\chapter{Psicopedagogicamente inconstitucionalíssimo otorrinolaringologista socioeducativo ético-profissional sociodemográfico interdisciplinaridade neuropsicológica avaliação}
\chapter{Psicopedagogicamente inconstitucionalíssimo otorrinolaringologista socioeducativo ético-profissional sociodemográfico interdisciplinaridade neuropsicológica avaliação}
\chapter{Psicopedagogicamente inconstitucionalíssimo otorrinolaringologista socioeducativo ético-profissional sociodemográfico interdisciplinaridade neuropsicológica avaliação}
\chapter{Psicopedagogicamente inconstitucionalíssimo otorrinolaringologista socioeducativo ético-profissional sociodemográfico interdisciplinaridade neuropsicológica avaliação}
\chapter{Psicopedagogicamente inconstitucionalíssimo otorrinolaringologista socioeducativo ético-profissional sociodemográfico interdisciplinaridade neuropsicológica avaliação}
\chapter{Psicopedagogicamente inconstitucionalíssimo otorrinolaringologista socioeducativo ético-profissional sociodemográfico interdisciplinaridade neuropsicológica avaliação}
\chapter{Psicopedagogicamente inconstitucionalíssimo otorrinolaringologista socioeducativo ético-profissional sociodemográfico interdisciplinaridade neuropsicológica avaliação}
\chapter{Psicopedagogicamente inconstitucionalíssimo otorrinolaringologista socioeducativo ético-profissional sociodemográfico interdisciplinaridade neuropsicológica avaliação}
\end{document}
'   % R: rígido (até 0.5.x)
%   lualatex '\def\var{B}% !TeX program = lualatex
% ===================================================================
% sumario-longo.tex — mede o sumário sem hifenização (issue #35, "Fora do
% escopo"): com a hifenização desligada, a entrada longa de capítulo pode
% bater no número da página. Gera 3 x 16 capítulos, cada um com um prefixo
% (de 3 a 18 palavras) de uma frase longa: frase 1 de palavras médias,
% frase 2 com números, frase 3 só de palavras muito longas. Conta no log
% os "Overfull \hbox" do sumário: o sumário é a única parte do documento
% sem parágrafo de corpo, então todo Overfull "in paragraph" antes do
% primeiro \chapter é entrada que passou da margem.
%
% Uso (de uma pasta com livros_crp_acessivel_book.sty e crpsp-leg.sty no
% caminho; as fontes NEWJUNE por OSFONTDIR):
%   lualatex '\def\var{R}% !TeX program = lualatex
% ===================================================================
% sumario-longo.tex — mede o sumário sem hifenização (issue #35, "Fora do
% escopo"): com a hifenização desligada, a entrada longa de capítulo pode
% bater no número da página. Gera 3 x 16 capítulos, cada um com um prefixo
% (de 3 a 18 palavras) de uma frase longa: frase 1 de palavras médias,
% frase 2 com números, frase 3 só de palavras muito longas. Conta no log
% os "Overfull \hbox" do sumário: o sumário é a única parte do documento
% sem parágrafo de corpo, então todo Overfull "in paragraph" antes do
% primeiro \chapter é entrada que passou da margem.
%
% Uso (de uma pasta com livros_crp_acessivel_book.sty e crpsp-leg.sty no
% caminho; as fontes NEWJUNE por OSFONTDIR):
%   lualatex '\def\var{R}\input{sumario-longo}'   % R: rígido (até 0.5.x)
%   lualatex '\def\var{B}\input{sumario-longo}'   % B: plus 4em (o do guia)
%   lualatex '\def\var{C}\input{sumario-longo}'   % C: plus 4em + \tolerance 10000
%   lualatex '\def\var{S}\input{sumario-longo}'   % S: o que a .sty define
% Cada variante só redefine o \l@chapter depois da carga da .sty.
% ===================================================================
\DocumentMetadata{lang=pt-BR, pdfstandard=ua-2, pdfversion=2.0, tagging=on}
\documentclass[11pt,a5paper,twoside,openright]{book}
\usepackage{livros_crp_acessivel_book}
\hypersetup{pdftitle={Sumario longo}, pdfauthor={CRP SP}}
\makeatletter
\providecommand\var{S}
\def\crp@vR{R}\def\crp@vB{B}\def\crp@vC{C}
\def\crp@lchap#1{%  #1 = acréscimo à margem direita (rightskip)
  \renewcommand*\l@chapter[2]{%
    \ifnum \c@tocdepth >\m@ne
        \addpenalty{-\@highpenalty}\addvspace{0.5em \@plus\p@}%
        \setlength\@tempdima{1.5em}%
        \begingroup
            \parindent \z@ \rightskip \@pnumwidth #1
            \parfillskip -\@pnumwidth
            \leavevmode \bfseries
            \advance\leftskip\@tempdima \hskip -\leftskip
            ##1\nobreak\hfil
            \nobreak\hb@xt@\@pnumwidth{\hss ##2\kern-\p@\kern\p@}\par
            \penalty\@highpenalty
        \endgroup
    \fi}}
\ifx\var\crp@vR \crp@lchap{}\fi
\ifx\var\crp@vB \crp@lchap{\@plus 4em}\fi
\ifx\var\crp@vC \crp@lchap{\@plus 4em \tolerance=10000 \hbadness=10000 }\fi
\makeatother
\begin{document}
\frontmatter
\tableofcontents
\mainmatter
\chapter{Deficiência psicossocial e}
\chapter{Deficiência psicossocial e transtorno}
\chapter{Deficiência psicossocial e transtorno do}
\chapter{Deficiência psicossocial e transtorno do espectro}
\chapter{Deficiência psicossocial e transtorno do espectro autista:}
\chapter{Deficiência psicossocial e transtorno do espectro autista: orientações}
\chapter{Deficiência psicossocial e transtorno do espectro autista: orientações para}
\chapter{Deficiência psicossocial e transtorno do espectro autista: orientações para a}
\chapter{Deficiência psicossocial e transtorno do espectro autista: orientações para a realização}
\chapter{Deficiência psicossocial e transtorno do espectro autista: orientações para a realização de}
\chapter{Deficiência psicossocial e transtorno do espectro autista: orientações para a realização de eventos}
\chapter{Deficiência psicossocial e transtorno do espectro autista: orientações para a realização de eventos e}
\chapter{Deficiência psicossocial e transtorno do espectro autista: orientações para a realização de eventos e apresentações}
\chapter{Deficiência psicossocial e transtorno do espectro autista: orientações para a realização de eventos e apresentações acessíveis}
\chapter{Deficiência psicossocial e transtorno do espectro autista: orientações para a realização de eventos e apresentações acessíveis em}
\chapter{Deficiência psicossocial e transtorno do espectro autista: orientações para a realização de eventos e apresentações acessíveis em todo}
\chapter{Comunicação alternativa e}
\chapter{Comunicação alternativa e aumentativa}
\chapter{Comunicação alternativa e aumentativa na}
\chapter{Comunicação alternativa e aumentativa na participação}
\chapter{Comunicação alternativa e aumentativa na participação social}
\chapter{Comunicação alternativa e aumentativa na participação social de}
\chapter{Comunicação alternativa e aumentativa na participação social de pessoas}
\chapter{Comunicação alternativa e aumentativa na participação social de pessoas com}
\chapter{Comunicação alternativa e aumentativa na participação social de pessoas com deficiência}
\chapter{Comunicação alternativa e aumentativa na participação social de pessoas com deficiência intelectual}
\chapter{Comunicação alternativa e aumentativa na participação social de pessoas com deficiência intelectual em}
\chapter{Comunicação alternativa e aumentativa na participação social de pessoas com deficiência intelectual em eventos}
\chapter{Comunicação alternativa e aumentativa na participação social de pessoas com deficiência intelectual em eventos públicos}
\chapter{Comunicação alternativa e aumentativa na participação social de pessoas com deficiência intelectual em eventos públicos de}
\chapter{Comunicação alternativa e aumentativa na participação social de pessoas com deficiência intelectual em eventos públicos de grande}
\chapter{Comunicação alternativa e aumentativa na participação social de pessoas com deficiência intelectual em eventos públicos de grande porte}
\chapter{Psicopedagogicamente inconstitucionalíssimo otorrinolaringologista}
\chapter{Psicopedagogicamente inconstitucionalíssimo otorrinolaringologista socioeducativo}
\chapter{Psicopedagogicamente inconstitucionalíssimo otorrinolaringologista socioeducativo ético-profissional}
\chapter{Psicopedagogicamente inconstitucionalíssimo otorrinolaringologista socioeducativo ético-profissional sociodemográfico}
\chapter{Psicopedagogicamente inconstitucionalíssimo otorrinolaringologista socioeducativo ético-profissional sociodemográfico interdisciplinaridade}
\chapter{Psicopedagogicamente inconstitucionalíssimo otorrinolaringologista socioeducativo ético-profissional sociodemográfico interdisciplinaridade neuropsicológica}
\chapter{Psicopedagogicamente inconstitucionalíssimo otorrinolaringologista socioeducativo ético-profissional sociodemográfico interdisciplinaridade neuropsicológica avaliação}
\chapter{Psicopedagogicamente inconstitucionalíssimo otorrinolaringologista socioeducativo ético-profissional sociodemográfico interdisciplinaridade neuropsicológica avaliação}
\chapter{Psicopedagogicamente inconstitucionalíssimo otorrinolaringologista socioeducativo ético-profissional sociodemográfico interdisciplinaridade neuropsicológica avaliação}
\chapter{Psicopedagogicamente inconstitucionalíssimo otorrinolaringologista socioeducativo ético-profissional sociodemográfico interdisciplinaridade neuropsicológica avaliação}
\chapter{Psicopedagogicamente inconstitucionalíssimo otorrinolaringologista socioeducativo ético-profissional sociodemográfico interdisciplinaridade neuropsicológica avaliação}
\chapter{Psicopedagogicamente inconstitucionalíssimo otorrinolaringologista socioeducativo ético-profissional sociodemográfico interdisciplinaridade neuropsicológica avaliação}
\chapter{Psicopedagogicamente inconstitucionalíssimo otorrinolaringologista socioeducativo ético-profissional sociodemográfico interdisciplinaridade neuropsicológica avaliação}
\chapter{Psicopedagogicamente inconstitucionalíssimo otorrinolaringologista socioeducativo ético-profissional sociodemográfico interdisciplinaridade neuropsicológica avaliação}
\chapter{Psicopedagogicamente inconstitucionalíssimo otorrinolaringologista socioeducativo ético-profissional sociodemográfico interdisciplinaridade neuropsicológica avaliação}
\chapter{Psicopedagogicamente inconstitucionalíssimo otorrinolaringologista socioeducativo ético-profissional sociodemográfico interdisciplinaridade neuropsicológica avaliação}
\end{document}
'   % R: rígido (até 0.5.x)
%   lualatex '\def\var{B}% !TeX program = lualatex
% ===================================================================
% sumario-longo.tex — mede o sumário sem hifenização (issue #35, "Fora do
% escopo"): com a hifenização desligada, a entrada longa de capítulo pode
% bater no número da página. Gera 3 x 16 capítulos, cada um com um prefixo
% (de 3 a 18 palavras) de uma frase longa: frase 1 de palavras médias,
% frase 2 com números, frase 3 só de palavras muito longas. Conta no log
% os "Overfull \hbox" do sumário: o sumário é a única parte do documento
% sem parágrafo de corpo, então todo Overfull "in paragraph" antes do
% primeiro \chapter é entrada que passou da margem.
%
% Uso (de uma pasta com livros_crp_acessivel_book.sty e crpsp-leg.sty no
% caminho; as fontes NEWJUNE por OSFONTDIR):
%   lualatex '\def\var{R}\input{sumario-longo}'   % R: rígido (até 0.5.x)
%   lualatex '\def\var{B}\input{sumario-longo}'   % B: plus 4em (o do guia)
%   lualatex '\def\var{C}\input{sumario-longo}'   % C: plus 4em + \tolerance 10000
%   lualatex '\def\var{S}\input{sumario-longo}'   % S: o que a .sty define
% Cada variante só redefine o \l@chapter depois da carga da .sty.
% ===================================================================
\DocumentMetadata{lang=pt-BR, pdfstandard=ua-2, pdfversion=2.0, tagging=on}
\documentclass[11pt,a5paper,twoside,openright]{book}
\usepackage{livros_crp_acessivel_book}
\hypersetup{pdftitle={Sumario longo}, pdfauthor={CRP SP}}
\makeatletter
\providecommand\var{S}
\def\crp@vR{R}\def\crp@vB{B}\def\crp@vC{C}
\def\crp@lchap#1{%  #1 = acréscimo à margem direita (rightskip)
  \renewcommand*\l@chapter[2]{%
    \ifnum \c@tocdepth >\m@ne
        \addpenalty{-\@highpenalty}\addvspace{0.5em \@plus\p@}%
        \setlength\@tempdima{1.5em}%
        \begingroup
            \parindent \z@ \rightskip \@pnumwidth #1
            \parfillskip -\@pnumwidth
            \leavevmode \bfseries
            \advance\leftskip\@tempdima \hskip -\leftskip
            ##1\nobreak\hfil
            \nobreak\hb@xt@\@pnumwidth{\hss ##2\kern-\p@\kern\p@}\par
            \penalty\@highpenalty
        \endgroup
    \fi}}
\ifx\var\crp@vR \crp@lchap{}\fi
\ifx\var\crp@vB \crp@lchap{\@plus 4em}\fi
\ifx\var\crp@vC \crp@lchap{\@plus 4em \tolerance=10000 \hbadness=10000 }\fi
\makeatother
\begin{document}
\frontmatter
\tableofcontents
\mainmatter
\chapter{Deficiência psicossocial e}
\chapter{Deficiência psicossocial e transtorno}
\chapter{Deficiência psicossocial e transtorno do}
\chapter{Deficiência psicossocial e transtorno do espectro}
\chapter{Deficiência psicossocial e transtorno do espectro autista:}
\chapter{Deficiência psicossocial e transtorno do espectro autista: orientações}
\chapter{Deficiência psicossocial e transtorno do espectro autista: orientações para}
\chapter{Deficiência psicossocial e transtorno do espectro autista: orientações para a}
\chapter{Deficiência psicossocial e transtorno do espectro autista: orientações para a realização}
\chapter{Deficiência psicossocial e transtorno do espectro autista: orientações para a realização de}
\chapter{Deficiência psicossocial e transtorno do espectro autista: orientações para a realização de eventos}
\chapter{Deficiência psicossocial e transtorno do espectro autista: orientações para a realização de eventos e}
\chapter{Deficiência psicossocial e transtorno do espectro autista: orientações para a realização de eventos e apresentações}
\chapter{Deficiência psicossocial e transtorno do espectro autista: orientações para a realização de eventos e apresentações acessíveis}
\chapter{Deficiência psicossocial e transtorno do espectro autista: orientações para a realização de eventos e apresentações acessíveis em}
\chapter{Deficiência psicossocial e transtorno do espectro autista: orientações para a realização de eventos e apresentações acessíveis em todo}
\chapter{Comunicação alternativa e}
\chapter{Comunicação alternativa e aumentativa}
\chapter{Comunicação alternativa e aumentativa na}
\chapter{Comunicação alternativa e aumentativa na participação}
\chapter{Comunicação alternativa e aumentativa na participação social}
\chapter{Comunicação alternativa e aumentativa na participação social de}
\chapter{Comunicação alternativa e aumentativa na participação social de pessoas}
\chapter{Comunicação alternativa e aumentativa na participação social de pessoas com}
\chapter{Comunicação alternativa e aumentativa na participação social de pessoas com deficiência}
\chapter{Comunicação alternativa e aumentativa na participação social de pessoas com deficiência intelectual}
\chapter{Comunicação alternativa e aumentativa na participação social de pessoas com deficiência intelectual em}
\chapter{Comunicação alternativa e aumentativa na participação social de pessoas com deficiência intelectual em eventos}
\chapter{Comunicação alternativa e aumentativa na participação social de pessoas com deficiência intelectual em eventos públicos}
\chapter{Comunicação alternativa e aumentativa na participação social de pessoas com deficiência intelectual em eventos públicos de}
\chapter{Comunicação alternativa e aumentativa na participação social de pessoas com deficiência intelectual em eventos públicos de grande}
\chapter{Comunicação alternativa e aumentativa na participação social de pessoas com deficiência intelectual em eventos públicos de grande porte}
\chapter{Psicopedagogicamente inconstitucionalíssimo otorrinolaringologista}
\chapter{Psicopedagogicamente inconstitucionalíssimo otorrinolaringologista socioeducativo}
\chapter{Psicopedagogicamente inconstitucionalíssimo otorrinolaringologista socioeducativo ético-profissional}
\chapter{Psicopedagogicamente inconstitucionalíssimo otorrinolaringologista socioeducativo ético-profissional sociodemográfico}
\chapter{Psicopedagogicamente inconstitucionalíssimo otorrinolaringologista socioeducativo ético-profissional sociodemográfico interdisciplinaridade}
\chapter{Psicopedagogicamente inconstitucionalíssimo otorrinolaringologista socioeducativo ético-profissional sociodemográfico interdisciplinaridade neuropsicológica}
\chapter{Psicopedagogicamente inconstitucionalíssimo otorrinolaringologista socioeducativo ético-profissional sociodemográfico interdisciplinaridade neuropsicológica avaliação}
\chapter{Psicopedagogicamente inconstitucionalíssimo otorrinolaringologista socioeducativo ético-profissional sociodemográfico interdisciplinaridade neuropsicológica avaliação}
\chapter{Psicopedagogicamente inconstitucionalíssimo otorrinolaringologista socioeducativo ético-profissional sociodemográfico interdisciplinaridade neuropsicológica avaliação}
\chapter{Psicopedagogicamente inconstitucionalíssimo otorrinolaringologista socioeducativo ético-profissional sociodemográfico interdisciplinaridade neuropsicológica avaliação}
\chapter{Psicopedagogicamente inconstitucionalíssimo otorrinolaringologista socioeducativo ético-profissional sociodemográfico interdisciplinaridade neuropsicológica avaliação}
\chapter{Psicopedagogicamente inconstitucionalíssimo otorrinolaringologista socioeducativo ético-profissional sociodemográfico interdisciplinaridade neuropsicológica avaliação}
\chapter{Psicopedagogicamente inconstitucionalíssimo otorrinolaringologista socioeducativo ético-profissional sociodemográfico interdisciplinaridade neuropsicológica avaliação}
\chapter{Psicopedagogicamente inconstitucionalíssimo otorrinolaringologista socioeducativo ético-profissional sociodemográfico interdisciplinaridade neuropsicológica avaliação}
\chapter{Psicopedagogicamente inconstitucionalíssimo otorrinolaringologista socioeducativo ético-profissional sociodemográfico interdisciplinaridade neuropsicológica avaliação}
\chapter{Psicopedagogicamente inconstitucionalíssimo otorrinolaringologista socioeducativo ético-profissional sociodemográfico interdisciplinaridade neuropsicológica avaliação}
\end{document}
'   % B: plus 4em (o do guia)
%   lualatex '\def\var{C}% !TeX program = lualatex
% ===================================================================
% sumario-longo.tex — mede o sumário sem hifenização (issue #35, "Fora do
% escopo"): com a hifenização desligada, a entrada longa de capítulo pode
% bater no número da página. Gera 3 x 16 capítulos, cada um com um prefixo
% (de 3 a 18 palavras) de uma frase longa: frase 1 de palavras médias,
% frase 2 com números, frase 3 só de palavras muito longas. Conta no log
% os "Overfull \hbox" do sumário: o sumário é a única parte do documento
% sem parágrafo de corpo, então todo Overfull "in paragraph" antes do
% primeiro \chapter é entrada que passou da margem.
%
% Uso (de uma pasta com livros_crp_acessivel_book.sty e crpsp-leg.sty no
% caminho; as fontes NEWJUNE por OSFONTDIR):
%   lualatex '\def\var{R}\input{sumario-longo}'   % R: rígido (até 0.5.x)
%   lualatex '\def\var{B}\input{sumario-longo}'   % B: plus 4em (o do guia)
%   lualatex '\def\var{C}\input{sumario-longo}'   % C: plus 4em + \tolerance 10000
%   lualatex '\def\var{S}\input{sumario-longo}'   % S: o que a .sty define
% Cada variante só redefine o \l@chapter depois da carga da .sty.
% ===================================================================
\DocumentMetadata{lang=pt-BR, pdfstandard=ua-2, pdfversion=2.0, tagging=on}
\documentclass[11pt,a5paper,twoside,openright]{book}
\usepackage{livros_crp_acessivel_book}
\hypersetup{pdftitle={Sumario longo}, pdfauthor={CRP SP}}
\makeatletter
\providecommand\var{S}
\def\crp@vR{R}\def\crp@vB{B}\def\crp@vC{C}
\def\crp@lchap#1{%  #1 = acréscimo à margem direita (rightskip)
  \renewcommand*\l@chapter[2]{%
    \ifnum \c@tocdepth >\m@ne
        \addpenalty{-\@highpenalty}\addvspace{0.5em \@plus\p@}%
        \setlength\@tempdima{1.5em}%
        \begingroup
            \parindent \z@ \rightskip \@pnumwidth #1
            \parfillskip -\@pnumwidth
            \leavevmode \bfseries
            \advance\leftskip\@tempdima \hskip -\leftskip
            ##1\nobreak\hfil
            \nobreak\hb@xt@\@pnumwidth{\hss ##2\kern-\p@\kern\p@}\par
            \penalty\@highpenalty
        \endgroup
    \fi}}
\ifx\var\crp@vR \crp@lchap{}\fi
\ifx\var\crp@vB \crp@lchap{\@plus 4em}\fi
\ifx\var\crp@vC \crp@lchap{\@plus 4em \tolerance=10000 \hbadness=10000 }\fi
\makeatother
\begin{document}
\frontmatter
\tableofcontents
\mainmatter
\chapter{Deficiência psicossocial e}
\chapter{Deficiência psicossocial e transtorno}
\chapter{Deficiência psicossocial e transtorno do}
\chapter{Deficiência psicossocial e transtorno do espectro}
\chapter{Deficiência psicossocial e transtorno do espectro autista:}
\chapter{Deficiência psicossocial e transtorno do espectro autista: orientações}
\chapter{Deficiência psicossocial e transtorno do espectro autista: orientações para}
\chapter{Deficiência psicossocial e transtorno do espectro autista: orientações para a}
\chapter{Deficiência psicossocial e transtorno do espectro autista: orientações para a realização}
\chapter{Deficiência psicossocial e transtorno do espectro autista: orientações para a realização de}
\chapter{Deficiência psicossocial e transtorno do espectro autista: orientações para a realização de eventos}
\chapter{Deficiência psicossocial e transtorno do espectro autista: orientações para a realização de eventos e}
\chapter{Deficiência psicossocial e transtorno do espectro autista: orientações para a realização de eventos e apresentações}
\chapter{Deficiência psicossocial e transtorno do espectro autista: orientações para a realização de eventos e apresentações acessíveis}
\chapter{Deficiência psicossocial e transtorno do espectro autista: orientações para a realização de eventos e apresentações acessíveis em}
\chapter{Deficiência psicossocial e transtorno do espectro autista: orientações para a realização de eventos e apresentações acessíveis em todo}
\chapter{Comunicação alternativa e}
\chapter{Comunicação alternativa e aumentativa}
\chapter{Comunicação alternativa e aumentativa na}
\chapter{Comunicação alternativa e aumentativa na participação}
\chapter{Comunicação alternativa e aumentativa na participação social}
\chapter{Comunicação alternativa e aumentativa na participação social de}
\chapter{Comunicação alternativa e aumentativa na participação social de pessoas}
\chapter{Comunicação alternativa e aumentativa na participação social de pessoas com}
\chapter{Comunicação alternativa e aumentativa na participação social de pessoas com deficiência}
\chapter{Comunicação alternativa e aumentativa na participação social de pessoas com deficiência intelectual}
\chapter{Comunicação alternativa e aumentativa na participação social de pessoas com deficiência intelectual em}
\chapter{Comunicação alternativa e aumentativa na participação social de pessoas com deficiência intelectual em eventos}
\chapter{Comunicação alternativa e aumentativa na participação social de pessoas com deficiência intelectual em eventos públicos}
\chapter{Comunicação alternativa e aumentativa na participação social de pessoas com deficiência intelectual em eventos públicos de}
\chapter{Comunicação alternativa e aumentativa na participação social de pessoas com deficiência intelectual em eventos públicos de grande}
\chapter{Comunicação alternativa e aumentativa na participação social de pessoas com deficiência intelectual em eventos públicos de grande porte}
\chapter{Psicopedagogicamente inconstitucionalíssimo otorrinolaringologista}
\chapter{Psicopedagogicamente inconstitucionalíssimo otorrinolaringologista socioeducativo}
\chapter{Psicopedagogicamente inconstitucionalíssimo otorrinolaringologista socioeducativo ético-profissional}
\chapter{Psicopedagogicamente inconstitucionalíssimo otorrinolaringologista socioeducativo ético-profissional sociodemográfico}
\chapter{Psicopedagogicamente inconstitucionalíssimo otorrinolaringologista socioeducativo ético-profissional sociodemográfico interdisciplinaridade}
\chapter{Psicopedagogicamente inconstitucionalíssimo otorrinolaringologista socioeducativo ético-profissional sociodemográfico interdisciplinaridade neuropsicológica}
\chapter{Psicopedagogicamente inconstitucionalíssimo otorrinolaringologista socioeducativo ético-profissional sociodemográfico interdisciplinaridade neuropsicológica avaliação}
\chapter{Psicopedagogicamente inconstitucionalíssimo otorrinolaringologista socioeducativo ético-profissional sociodemográfico interdisciplinaridade neuropsicológica avaliação}
\chapter{Psicopedagogicamente inconstitucionalíssimo otorrinolaringologista socioeducativo ético-profissional sociodemográfico interdisciplinaridade neuropsicológica avaliação}
\chapter{Psicopedagogicamente inconstitucionalíssimo otorrinolaringologista socioeducativo ético-profissional sociodemográfico interdisciplinaridade neuropsicológica avaliação}
\chapter{Psicopedagogicamente inconstitucionalíssimo otorrinolaringologista socioeducativo ético-profissional sociodemográfico interdisciplinaridade neuropsicológica avaliação}
\chapter{Psicopedagogicamente inconstitucionalíssimo otorrinolaringologista socioeducativo ético-profissional sociodemográfico interdisciplinaridade neuropsicológica avaliação}
\chapter{Psicopedagogicamente inconstitucionalíssimo otorrinolaringologista socioeducativo ético-profissional sociodemográfico interdisciplinaridade neuropsicológica avaliação}
\chapter{Psicopedagogicamente inconstitucionalíssimo otorrinolaringologista socioeducativo ético-profissional sociodemográfico interdisciplinaridade neuropsicológica avaliação}
\chapter{Psicopedagogicamente inconstitucionalíssimo otorrinolaringologista socioeducativo ético-profissional sociodemográfico interdisciplinaridade neuropsicológica avaliação}
\chapter{Psicopedagogicamente inconstitucionalíssimo otorrinolaringologista socioeducativo ético-profissional sociodemográfico interdisciplinaridade neuropsicológica avaliação}
\end{document}
'   % C: plus 4em + \tolerance 10000
%   lualatex '\def\var{S}% !TeX program = lualatex
% ===================================================================
% sumario-longo.tex — mede o sumário sem hifenização (issue #35, "Fora do
% escopo"): com a hifenização desligada, a entrada longa de capítulo pode
% bater no número da página. Gera 3 x 16 capítulos, cada um com um prefixo
% (de 3 a 18 palavras) de uma frase longa: frase 1 de palavras médias,
% frase 2 com números, frase 3 só de palavras muito longas. Conta no log
% os "Overfull \hbox" do sumário: o sumário é a única parte do documento
% sem parágrafo de corpo, então todo Overfull "in paragraph" antes do
% primeiro \chapter é entrada que passou da margem.
%
% Uso (de uma pasta com livros_crp_acessivel_book.sty e crpsp-leg.sty no
% caminho; as fontes NEWJUNE por OSFONTDIR):
%   lualatex '\def\var{R}\input{sumario-longo}'   % R: rígido (até 0.5.x)
%   lualatex '\def\var{B}\input{sumario-longo}'   % B: plus 4em (o do guia)
%   lualatex '\def\var{C}\input{sumario-longo}'   % C: plus 4em + \tolerance 10000
%   lualatex '\def\var{S}\input{sumario-longo}'   % S: o que a .sty define
% Cada variante só redefine o \l@chapter depois da carga da .sty.
% ===================================================================
\DocumentMetadata{lang=pt-BR, pdfstandard=ua-2, pdfversion=2.0, tagging=on}
\documentclass[11pt,a5paper,twoside,openright]{book}
\usepackage{livros_crp_acessivel_book}
\hypersetup{pdftitle={Sumario longo}, pdfauthor={CRP SP}}
\makeatletter
\providecommand\var{S}
\def\crp@vR{R}\def\crp@vB{B}\def\crp@vC{C}
\def\crp@lchap#1{%  #1 = acréscimo à margem direita (rightskip)
  \renewcommand*\l@chapter[2]{%
    \ifnum \c@tocdepth >\m@ne
        \addpenalty{-\@highpenalty}\addvspace{0.5em \@plus\p@}%
        \setlength\@tempdima{1.5em}%
        \begingroup
            \parindent \z@ \rightskip \@pnumwidth #1
            \parfillskip -\@pnumwidth
            \leavevmode \bfseries
            \advance\leftskip\@tempdima \hskip -\leftskip
            ##1\nobreak\hfil
            \nobreak\hb@xt@\@pnumwidth{\hss ##2\kern-\p@\kern\p@}\par
            \penalty\@highpenalty
        \endgroup
    \fi}}
\ifx\var\crp@vR \crp@lchap{}\fi
\ifx\var\crp@vB \crp@lchap{\@plus 4em}\fi
\ifx\var\crp@vC \crp@lchap{\@plus 4em \tolerance=10000 \hbadness=10000 }\fi
\makeatother
\begin{document}
\frontmatter
\tableofcontents
\mainmatter
\chapter{Deficiência psicossocial e}
\chapter{Deficiência psicossocial e transtorno}
\chapter{Deficiência psicossocial e transtorno do}
\chapter{Deficiência psicossocial e transtorno do espectro}
\chapter{Deficiência psicossocial e transtorno do espectro autista:}
\chapter{Deficiência psicossocial e transtorno do espectro autista: orientações}
\chapter{Deficiência psicossocial e transtorno do espectro autista: orientações para}
\chapter{Deficiência psicossocial e transtorno do espectro autista: orientações para a}
\chapter{Deficiência psicossocial e transtorno do espectro autista: orientações para a realização}
\chapter{Deficiência psicossocial e transtorno do espectro autista: orientações para a realização de}
\chapter{Deficiência psicossocial e transtorno do espectro autista: orientações para a realização de eventos}
\chapter{Deficiência psicossocial e transtorno do espectro autista: orientações para a realização de eventos e}
\chapter{Deficiência psicossocial e transtorno do espectro autista: orientações para a realização de eventos e apresentações}
\chapter{Deficiência psicossocial e transtorno do espectro autista: orientações para a realização de eventos e apresentações acessíveis}
\chapter{Deficiência psicossocial e transtorno do espectro autista: orientações para a realização de eventos e apresentações acessíveis em}
\chapter{Deficiência psicossocial e transtorno do espectro autista: orientações para a realização de eventos e apresentações acessíveis em todo}
\chapter{Comunicação alternativa e}
\chapter{Comunicação alternativa e aumentativa}
\chapter{Comunicação alternativa e aumentativa na}
\chapter{Comunicação alternativa e aumentativa na participação}
\chapter{Comunicação alternativa e aumentativa na participação social}
\chapter{Comunicação alternativa e aumentativa na participação social de}
\chapter{Comunicação alternativa e aumentativa na participação social de pessoas}
\chapter{Comunicação alternativa e aumentativa na participação social de pessoas com}
\chapter{Comunicação alternativa e aumentativa na participação social de pessoas com deficiência}
\chapter{Comunicação alternativa e aumentativa na participação social de pessoas com deficiência intelectual}
\chapter{Comunicação alternativa e aumentativa na participação social de pessoas com deficiência intelectual em}
\chapter{Comunicação alternativa e aumentativa na participação social de pessoas com deficiência intelectual em eventos}
\chapter{Comunicação alternativa e aumentativa na participação social de pessoas com deficiência intelectual em eventos públicos}
\chapter{Comunicação alternativa e aumentativa na participação social de pessoas com deficiência intelectual em eventos públicos de}
\chapter{Comunicação alternativa e aumentativa na participação social de pessoas com deficiência intelectual em eventos públicos de grande}
\chapter{Comunicação alternativa e aumentativa na participação social de pessoas com deficiência intelectual em eventos públicos de grande porte}
\chapter{Psicopedagogicamente inconstitucionalíssimo otorrinolaringologista}
\chapter{Psicopedagogicamente inconstitucionalíssimo otorrinolaringologista socioeducativo}
\chapter{Psicopedagogicamente inconstitucionalíssimo otorrinolaringologista socioeducativo ético-profissional}
\chapter{Psicopedagogicamente inconstitucionalíssimo otorrinolaringologista socioeducativo ético-profissional sociodemográfico}
\chapter{Psicopedagogicamente inconstitucionalíssimo otorrinolaringologista socioeducativo ético-profissional sociodemográfico interdisciplinaridade}
\chapter{Psicopedagogicamente inconstitucionalíssimo otorrinolaringologista socioeducativo ético-profissional sociodemográfico interdisciplinaridade neuropsicológica}
\chapter{Psicopedagogicamente inconstitucionalíssimo otorrinolaringologista socioeducativo ético-profissional sociodemográfico interdisciplinaridade neuropsicológica avaliação}
\chapter{Psicopedagogicamente inconstitucionalíssimo otorrinolaringologista socioeducativo ético-profissional sociodemográfico interdisciplinaridade neuropsicológica avaliação}
\chapter{Psicopedagogicamente inconstitucionalíssimo otorrinolaringologista socioeducativo ético-profissional sociodemográfico interdisciplinaridade neuropsicológica avaliação}
\chapter{Psicopedagogicamente inconstitucionalíssimo otorrinolaringologista socioeducativo ético-profissional sociodemográfico interdisciplinaridade neuropsicológica avaliação}
\chapter{Psicopedagogicamente inconstitucionalíssimo otorrinolaringologista socioeducativo ético-profissional sociodemográfico interdisciplinaridade neuropsicológica avaliação}
\chapter{Psicopedagogicamente inconstitucionalíssimo otorrinolaringologista socioeducativo ético-profissional sociodemográfico interdisciplinaridade neuropsicológica avaliação}
\chapter{Psicopedagogicamente inconstitucionalíssimo otorrinolaringologista socioeducativo ético-profissional sociodemográfico interdisciplinaridade neuropsicológica avaliação}
\chapter{Psicopedagogicamente inconstitucionalíssimo otorrinolaringologista socioeducativo ético-profissional sociodemográfico interdisciplinaridade neuropsicológica avaliação}
\chapter{Psicopedagogicamente inconstitucionalíssimo otorrinolaringologista socioeducativo ético-profissional sociodemográfico interdisciplinaridade neuropsicológica avaliação}
\chapter{Psicopedagogicamente inconstitucionalíssimo otorrinolaringologista socioeducativo ético-profissional sociodemográfico interdisciplinaridade neuropsicológica avaliação}
\end{document}
'   % S: o que a .sty define
% Cada variante só redefine o \l@chapter depois da carga da .sty.
% ===================================================================
\DocumentMetadata{lang=pt-BR, pdfstandard=ua-2, pdfversion=2.0, tagging=on}
\documentclass[11pt,a5paper,twoside,openright]{book}
\usepackage{livros_crp_acessivel_book}
\hypersetup{pdftitle={Sumario longo}, pdfauthor={CRP SP}}
\makeatletter
\providecommand\var{S}
\def\crp@vR{R}\def\crp@vB{B}\def\crp@vC{C}
\def\crp@lchap#1{%  #1 = acréscimo à margem direita (rightskip)
  \renewcommand*\l@chapter[2]{%
    \ifnum \c@tocdepth >\m@ne
        \addpenalty{-\@highpenalty}\addvspace{0.5em \@plus\p@}%
        \setlength\@tempdima{1.5em}%
        \begingroup
            \parindent \z@ \rightskip \@pnumwidth #1
            \parfillskip -\@pnumwidth
            \leavevmode \bfseries
            \advance\leftskip\@tempdima \hskip -\leftskip
            ##1\nobreak\hfil
            \nobreak\hb@xt@\@pnumwidth{\hss ##2\kern-\p@\kern\p@}\par
            \penalty\@highpenalty
        \endgroup
    \fi}}
\ifx\var\crp@vR \crp@lchap{}\fi
\ifx\var\crp@vB \crp@lchap{\@plus 4em}\fi
\ifx\var\crp@vC \crp@lchap{\@plus 4em \tolerance=10000 \hbadness=10000 }\fi
\makeatother
\begin{document}
\frontmatter
\tableofcontents
\mainmatter
\chapter{Deficiência psicossocial e}
\chapter{Deficiência psicossocial e transtorno}
\chapter{Deficiência psicossocial e transtorno do}
\chapter{Deficiência psicossocial e transtorno do espectro}
\chapter{Deficiência psicossocial e transtorno do espectro autista:}
\chapter{Deficiência psicossocial e transtorno do espectro autista: orientações}
\chapter{Deficiência psicossocial e transtorno do espectro autista: orientações para}
\chapter{Deficiência psicossocial e transtorno do espectro autista: orientações para a}
\chapter{Deficiência psicossocial e transtorno do espectro autista: orientações para a realização}
\chapter{Deficiência psicossocial e transtorno do espectro autista: orientações para a realização de}
\chapter{Deficiência psicossocial e transtorno do espectro autista: orientações para a realização de eventos}
\chapter{Deficiência psicossocial e transtorno do espectro autista: orientações para a realização de eventos e}
\chapter{Deficiência psicossocial e transtorno do espectro autista: orientações para a realização de eventos e apresentações}
\chapter{Deficiência psicossocial e transtorno do espectro autista: orientações para a realização de eventos e apresentações acessíveis}
\chapter{Deficiência psicossocial e transtorno do espectro autista: orientações para a realização de eventos e apresentações acessíveis em}
\chapter{Deficiência psicossocial e transtorno do espectro autista: orientações para a realização de eventos e apresentações acessíveis em todo}
\chapter{Comunicação alternativa e}
\chapter{Comunicação alternativa e aumentativa}
\chapter{Comunicação alternativa e aumentativa na}
\chapter{Comunicação alternativa e aumentativa na participação}
\chapter{Comunicação alternativa e aumentativa na participação social}
\chapter{Comunicação alternativa e aumentativa na participação social de}
\chapter{Comunicação alternativa e aumentativa na participação social de pessoas}
\chapter{Comunicação alternativa e aumentativa na participação social de pessoas com}
\chapter{Comunicação alternativa e aumentativa na participação social de pessoas com deficiência}
\chapter{Comunicação alternativa e aumentativa na participação social de pessoas com deficiência intelectual}
\chapter{Comunicação alternativa e aumentativa na participação social de pessoas com deficiência intelectual em}
\chapter{Comunicação alternativa e aumentativa na participação social de pessoas com deficiência intelectual em eventos}
\chapter{Comunicação alternativa e aumentativa na participação social de pessoas com deficiência intelectual em eventos públicos}
\chapter{Comunicação alternativa e aumentativa na participação social de pessoas com deficiência intelectual em eventos públicos de}
\chapter{Comunicação alternativa e aumentativa na participação social de pessoas com deficiência intelectual em eventos públicos de grande}
\chapter{Comunicação alternativa e aumentativa na participação social de pessoas com deficiência intelectual em eventos públicos de grande porte}
\chapter{Psicopedagogicamente inconstitucionalíssimo otorrinolaringologista}
\chapter{Psicopedagogicamente inconstitucionalíssimo otorrinolaringologista socioeducativo}
\chapter{Psicopedagogicamente inconstitucionalíssimo otorrinolaringologista socioeducativo ético-profissional}
\chapter{Psicopedagogicamente inconstitucionalíssimo otorrinolaringologista socioeducativo ético-profissional sociodemográfico}
\chapter{Psicopedagogicamente inconstitucionalíssimo otorrinolaringologista socioeducativo ético-profissional sociodemográfico interdisciplinaridade}
\chapter{Psicopedagogicamente inconstitucionalíssimo otorrinolaringologista socioeducativo ético-profissional sociodemográfico interdisciplinaridade neuropsicológica}
\chapter{Psicopedagogicamente inconstitucionalíssimo otorrinolaringologista socioeducativo ético-profissional sociodemográfico interdisciplinaridade neuropsicológica avaliação}
\chapter{Psicopedagogicamente inconstitucionalíssimo otorrinolaringologista socioeducativo ético-profissional sociodemográfico interdisciplinaridade neuropsicológica avaliação}
\chapter{Psicopedagogicamente inconstitucionalíssimo otorrinolaringologista socioeducativo ético-profissional sociodemográfico interdisciplinaridade neuropsicológica avaliação}
\chapter{Psicopedagogicamente inconstitucionalíssimo otorrinolaringologista socioeducativo ético-profissional sociodemográfico interdisciplinaridade neuropsicológica avaliação}
\chapter{Psicopedagogicamente inconstitucionalíssimo otorrinolaringologista socioeducativo ético-profissional sociodemográfico interdisciplinaridade neuropsicológica avaliação}
\chapter{Psicopedagogicamente inconstitucionalíssimo otorrinolaringologista socioeducativo ético-profissional sociodemográfico interdisciplinaridade neuropsicológica avaliação}
\chapter{Psicopedagogicamente inconstitucionalíssimo otorrinolaringologista socioeducativo ético-profissional sociodemográfico interdisciplinaridade neuropsicológica avaliação}
\chapter{Psicopedagogicamente inconstitucionalíssimo otorrinolaringologista socioeducativo ético-profissional sociodemográfico interdisciplinaridade neuropsicológica avaliação}
\chapter{Psicopedagogicamente inconstitucionalíssimo otorrinolaringologista socioeducativo ético-profissional sociodemográfico interdisciplinaridade neuropsicológica avaliação}
\chapter{Psicopedagogicamente inconstitucionalíssimo otorrinolaringologista socioeducativo ético-profissional sociodemográfico interdisciplinaridade neuropsicológica avaliação}
\end{document}
'   % B: plus 4em (o do guia)
%   lualatex '\def\var{C}% !TeX program = lualatex
% ===================================================================
% sumario-longo.tex — mede o sumário sem hifenização (issue #35, "Fora do
% escopo"): com a hifenização desligada, a entrada longa de capítulo pode
% bater no número da página. Gera 3 x 16 capítulos, cada um com um prefixo
% (de 3 a 18 palavras) de uma frase longa: frase 1 de palavras médias,
% frase 2 com números, frase 3 só de palavras muito longas. Conta no log
% os "Overfull \hbox" do sumário: o sumário é a única parte do documento
% sem parágrafo de corpo, então todo Overfull "in paragraph" antes do
% primeiro \chapter é entrada que passou da margem.
%
% Uso (de uma pasta com livros_crp_acessivel_book.sty e crpsp-leg.sty no
% caminho; as fontes NEWJUNE por OSFONTDIR):
%   lualatex '\def\var{R}% !TeX program = lualatex
% ===================================================================
% sumario-longo.tex — mede o sumário sem hifenização (issue #35, "Fora do
% escopo"): com a hifenização desligada, a entrada longa de capítulo pode
% bater no número da página. Gera 3 x 16 capítulos, cada um com um prefixo
% (de 3 a 18 palavras) de uma frase longa: frase 1 de palavras médias,
% frase 2 com números, frase 3 só de palavras muito longas. Conta no log
% os "Overfull \hbox" do sumário: o sumário é a única parte do documento
% sem parágrafo de corpo, então todo Overfull "in paragraph" antes do
% primeiro \chapter é entrada que passou da margem.
%
% Uso (de uma pasta com livros_crp_acessivel_book.sty e crpsp-leg.sty no
% caminho; as fontes NEWJUNE por OSFONTDIR):
%   lualatex '\def\var{R}\input{sumario-longo}'   % R: rígido (até 0.5.x)
%   lualatex '\def\var{B}\input{sumario-longo}'   % B: plus 4em (o do guia)
%   lualatex '\def\var{C}\input{sumario-longo}'   % C: plus 4em + \tolerance 10000
%   lualatex '\def\var{S}\input{sumario-longo}'   % S: o que a .sty define
% Cada variante só redefine o \l@chapter depois da carga da .sty.
% ===================================================================
\DocumentMetadata{lang=pt-BR, pdfstandard=ua-2, pdfversion=2.0, tagging=on}
\documentclass[11pt,a5paper,twoside,openright]{book}
\usepackage{livros_crp_acessivel_book}
\hypersetup{pdftitle={Sumario longo}, pdfauthor={CRP SP}}
\makeatletter
\providecommand\var{S}
\def\crp@vR{R}\def\crp@vB{B}\def\crp@vC{C}
\def\crp@lchap#1{%  #1 = acréscimo à margem direita (rightskip)
  \renewcommand*\l@chapter[2]{%
    \ifnum \c@tocdepth >\m@ne
        \addpenalty{-\@highpenalty}\addvspace{0.5em \@plus\p@}%
        \setlength\@tempdima{1.5em}%
        \begingroup
            \parindent \z@ \rightskip \@pnumwidth #1
            \parfillskip -\@pnumwidth
            \leavevmode \bfseries
            \advance\leftskip\@tempdima \hskip -\leftskip
            ##1\nobreak\hfil
            \nobreak\hb@xt@\@pnumwidth{\hss ##2\kern-\p@\kern\p@}\par
            \penalty\@highpenalty
        \endgroup
    \fi}}
\ifx\var\crp@vR \crp@lchap{}\fi
\ifx\var\crp@vB \crp@lchap{\@plus 4em}\fi
\ifx\var\crp@vC \crp@lchap{\@plus 4em \tolerance=10000 \hbadness=10000 }\fi
\makeatother
\begin{document}
\frontmatter
\tableofcontents
\mainmatter
\chapter{Deficiência psicossocial e}
\chapter{Deficiência psicossocial e transtorno}
\chapter{Deficiência psicossocial e transtorno do}
\chapter{Deficiência psicossocial e transtorno do espectro}
\chapter{Deficiência psicossocial e transtorno do espectro autista:}
\chapter{Deficiência psicossocial e transtorno do espectro autista: orientações}
\chapter{Deficiência psicossocial e transtorno do espectro autista: orientações para}
\chapter{Deficiência psicossocial e transtorno do espectro autista: orientações para a}
\chapter{Deficiência psicossocial e transtorno do espectro autista: orientações para a realização}
\chapter{Deficiência psicossocial e transtorno do espectro autista: orientações para a realização de}
\chapter{Deficiência psicossocial e transtorno do espectro autista: orientações para a realização de eventos}
\chapter{Deficiência psicossocial e transtorno do espectro autista: orientações para a realização de eventos e}
\chapter{Deficiência psicossocial e transtorno do espectro autista: orientações para a realização de eventos e apresentações}
\chapter{Deficiência psicossocial e transtorno do espectro autista: orientações para a realização de eventos e apresentações acessíveis}
\chapter{Deficiência psicossocial e transtorno do espectro autista: orientações para a realização de eventos e apresentações acessíveis em}
\chapter{Deficiência psicossocial e transtorno do espectro autista: orientações para a realização de eventos e apresentações acessíveis em todo}
\chapter{Comunicação alternativa e}
\chapter{Comunicação alternativa e aumentativa}
\chapter{Comunicação alternativa e aumentativa na}
\chapter{Comunicação alternativa e aumentativa na participação}
\chapter{Comunicação alternativa e aumentativa na participação social}
\chapter{Comunicação alternativa e aumentativa na participação social de}
\chapter{Comunicação alternativa e aumentativa na participação social de pessoas}
\chapter{Comunicação alternativa e aumentativa na participação social de pessoas com}
\chapter{Comunicação alternativa e aumentativa na participação social de pessoas com deficiência}
\chapter{Comunicação alternativa e aumentativa na participação social de pessoas com deficiência intelectual}
\chapter{Comunicação alternativa e aumentativa na participação social de pessoas com deficiência intelectual em}
\chapter{Comunicação alternativa e aumentativa na participação social de pessoas com deficiência intelectual em eventos}
\chapter{Comunicação alternativa e aumentativa na participação social de pessoas com deficiência intelectual em eventos públicos}
\chapter{Comunicação alternativa e aumentativa na participação social de pessoas com deficiência intelectual em eventos públicos de}
\chapter{Comunicação alternativa e aumentativa na participação social de pessoas com deficiência intelectual em eventos públicos de grande}
\chapter{Comunicação alternativa e aumentativa na participação social de pessoas com deficiência intelectual em eventos públicos de grande porte}
\chapter{Psicopedagogicamente inconstitucionalíssimo otorrinolaringologista}
\chapter{Psicopedagogicamente inconstitucionalíssimo otorrinolaringologista socioeducativo}
\chapter{Psicopedagogicamente inconstitucionalíssimo otorrinolaringologista socioeducativo ético-profissional}
\chapter{Psicopedagogicamente inconstitucionalíssimo otorrinolaringologista socioeducativo ético-profissional sociodemográfico}
\chapter{Psicopedagogicamente inconstitucionalíssimo otorrinolaringologista socioeducativo ético-profissional sociodemográfico interdisciplinaridade}
\chapter{Psicopedagogicamente inconstitucionalíssimo otorrinolaringologista socioeducativo ético-profissional sociodemográfico interdisciplinaridade neuropsicológica}
\chapter{Psicopedagogicamente inconstitucionalíssimo otorrinolaringologista socioeducativo ético-profissional sociodemográfico interdisciplinaridade neuropsicológica avaliação}
\chapter{Psicopedagogicamente inconstitucionalíssimo otorrinolaringologista socioeducativo ético-profissional sociodemográfico interdisciplinaridade neuropsicológica avaliação}
\chapter{Psicopedagogicamente inconstitucionalíssimo otorrinolaringologista socioeducativo ético-profissional sociodemográfico interdisciplinaridade neuropsicológica avaliação}
\chapter{Psicopedagogicamente inconstitucionalíssimo otorrinolaringologista socioeducativo ético-profissional sociodemográfico interdisciplinaridade neuropsicológica avaliação}
\chapter{Psicopedagogicamente inconstitucionalíssimo otorrinolaringologista socioeducativo ético-profissional sociodemográfico interdisciplinaridade neuropsicológica avaliação}
\chapter{Psicopedagogicamente inconstitucionalíssimo otorrinolaringologista socioeducativo ético-profissional sociodemográfico interdisciplinaridade neuropsicológica avaliação}
\chapter{Psicopedagogicamente inconstitucionalíssimo otorrinolaringologista socioeducativo ético-profissional sociodemográfico interdisciplinaridade neuropsicológica avaliação}
\chapter{Psicopedagogicamente inconstitucionalíssimo otorrinolaringologista socioeducativo ético-profissional sociodemográfico interdisciplinaridade neuropsicológica avaliação}
\chapter{Psicopedagogicamente inconstitucionalíssimo otorrinolaringologista socioeducativo ético-profissional sociodemográfico interdisciplinaridade neuropsicológica avaliação}
\chapter{Psicopedagogicamente inconstitucionalíssimo otorrinolaringologista socioeducativo ético-profissional sociodemográfico interdisciplinaridade neuropsicológica avaliação}
\end{document}
'   % R: rígido (até 0.5.x)
%   lualatex '\def\var{B}% !TeX program = lualatex
% ===================================================================
% sumario-longo.tex — mede o sumário sem hifenização (issue #35, "Fora do
% escopo"): com a hifenização desligada, a entrada longa de capítulo pode
% bater no número da página. Gera 3 x 16 capítulos, cada um com um prefixo
% (de 3 a 18 palavras) de uma frase longa: frase 1 de palavras médias,
% frase 2 com números, frase 3 só de palavras muito longas. Conta no log
% os "Overfull \hbox" do sumário: o sumário é a única parte do documento
% sem parágrafo de corpo, então todo Overfull "in paragraph" antes do
% primeiro \chapter é entrada que passou da margem.
%
% Uso (de uma pasta com livros_crp_acessivel_book.sty e crpsp-leg.sty no
% caminho; as fontes NEWJUNE por OSFONTDIR):
%   lualatex '\def\var{R}\input{sumario-longo}'   % R: rígido (até 0.5.x)
%   lualatex '\def\var{B}\input{sumario-longo}'   % B: plus 4em (o do guia)
%   lualatex '\def\var{C}\input{sumario-longo}'   % C: plus 4em + \tolerance 10000
%   lualatex '\def\var{S}\input{sumario-longo}'   % S: o que a .sty define
% Cada variante só redefine o \l@chapter depois da carga da .sty.
% ===================================================================
\DocumentMetadata{lang=pt-BR, pdfstandard=ua-2, pdfversion=2.0, tagging=on}
\documentclass[11pt,a5paper,twoside,openright]{book}
\usepackage{livros_crp_acessivel_book}
\hypersetup{pdftitle={Sumario longo}, pdfauthor={CRP SP}}
\makeatletter
\providecommand\var{S}
\def\crp@vR{R}\def\crp@vB{B}\def\crp@vC{C}
\def\crp@lchap#1{%  #1 = acréscimo à margem direita (rightskip)
  \renewcommand*\l@chapter[2]{%
    \ifnum \c@tocdepth >\m@ne
        \addpenalty{-\@highpenalty}\addvspace{0.5em \@plus\p@}%
        \setlength\@tempdima{1.5em}%
        \begingroup
            \parindent \z@ \rightskip \@pnumwidth #1
            \parfillskip -\@pnumwidth
            \leavevmode \bfseries
            \advance\leftskip\@tempdima \hskip -\leftskip
            ##1\nobreak\hfil
            \nobreak\hb@xt@\@pnumwidth{\hss ##2\kern-\p@\kern\p@}\par
            \penalty\@highpenalty
        \endgroup
    \fi}}
\ifx\var\crp@vR \crp@lchap{}\fi
\ifx\var\crp@vB \crp@lchap{\@plus 4em}\fi
\ifx\var\crp@vC \crp@lchap{\@plus 4em \tolerance=10000 \hbadness=10000 }\fi
\makeatother
\begin{document}
\frontmatter
\tableofcontents
\mainmatter
\chapter{Deficiência psicossocial e}
\chapter{Deficiência psicossocial e transtorno}
\chapter{Deficiência psicossocial e transtorno do}
\chapter{Deficiência psicossocial e transtorno do espectro}
\chapter{Deficiência psicossocial e transtorno do espectro autista:}
\chapter{Deficiência psicossocial e transtorno do espectro autista: orientações}
\chapter{Deficiência psicossocial e transtorno do espectro autista: orientações para}
\chapter{Deficiência psicossocial e transtorno do espectro autista: orientações para a}
\chapter{Deficiência psicossocial e transtorno do espectro autista: orientações para a realização}
\chapter{Deficiência psicossocial e transtorno do espectro autista: orientações para a realização de}
\chapter{Deficiência psicossocial e transtorno do espectro autista: orientações para a realização de eventos}
\chapter{Deficiência psicossocial e transtorno do espectro autista: orientações para a realização de eventos e}
\chapter{Deficiência psicossocial e transtorno do espectro autista: orientações para a realização de eventos e apresentações}
\chapter{Deficiência psicossocial e transtorno do espectro autista: orientações para a realização de eventos e apresentações acessíveis}
\chapter{Deficiência psicossocial e transtorno do espectro autista: orientações para a realização de eventos e apresentações acessíveis em}
\chapter{Deficiência psicossocial e transtorno do espectro autista: orientações para a realização de eventos e apresentações acessíveis em todo}
\chapter{Comunicação alternativa e}
\chapter{Comunicação alternativa e aumentativa}
\chapter{Comunicação alternativa e aumentativa na}
\chapter{Comunicação alternativa e aumentativa na participação}
\chapter{Comunicação alternativa e aumentativa na participação social}
\chapter{Comunicação alternativa e aumentativa na participação social de}
\chapter{Comunicação alternativa e aumentativa na participação social de pessoas}
\chapter{Comunicação alternativa e aumentativa na participação social de pessoas com}
\chapter{Comunicação alternativa e aumentativa na participação social de pessoas com deficiência}
\chapter{Comunicação alternativa e aumentativa na participação social de pessoas com deficiência intelectual}
\chapter{Comunicação alternativa e aumentativa na participação social de pessoas com deficiência intelectual em}
\chapter{Comunicação alternativa e aumentativa na participação social de pessoas com deficiência intelectual em eventos}
\chapter{Comunicação alternativa e aumentativa na participação social de pessoas com deficiência intelectual em eventos públicos}
\chapter{Comunicação alternativa e aumentativa na participação social de pessoas com deficiência intelectual em eventos públicos de}
\chapter{Comunicação alternativa e aumentativa na participação social de pessoas com deficiência intelectual em eventos públicos de grande}
\chapter{Comunicação alternativa e aumentativa na participação social de pessoas com deficiência intelectual em eventos públicos de grande porte}
\chapter{Psicopedagogicamente inconstitucionalíssimo otorrinolaringologista}
\chapter{Psicopedagogicamente inconstitucionalíssimo otorrinolaringologista socioeducativo}
\chapter{Psicopedagogicamente inconstitucionalíssimo otorrinolaringologista socioeducativo ético-profissional}
\chapter{Psicopedagogicamente inconstitucionalíssimo otorrinolaringologista socioeducativo ético-profissional sociodemográfico}
\chapter{Psicopedagogicamente inconstitucionalíssimo otorrinolaringologista socioeducativo ético-profissional sociodemográfico interdisciplinaridade}
\chapter{Psicopedagogicamente inconstitucionalíssimo otorrinolaringologista socioeducativo ético-profissional sociodemográfico interdisciplinaridade neuropsicológica}
\chapter{Psicopedagogicamente inconstitucionalíssimo otorrinolaringologista socioeducativo ético-profissional sociodemográfico interdisciplinaridade neuropsicológica avaliação}
\chapter{Psicopedagogicamente inconstitucionalíssimo otorrinolaringologista socioeducativo ético-profissional sociodemográfico interdisciplinaridade neuropsicológica avaliação}
\chapter{Psicopedagogicamente inconstitucionalíssimo otorrinolaringologista socioeducativo ético-profissional sociodemográfico interdisciplinaridade neuropsicológica avaliação}
\chapter{Psicopedagogicamente inconstitucionalíssimo otorrinolaringologista socioeducativo ético-profissional sociodemográfico interdisciplinaridade neuropsicológica avaliação}
\chapter{Psicopedagogicamente inconstitucionalíssimo otorrinolaringologista socioeducativo ético-profissional sociodemográfico interdisciplinaridade neuropsicológica avaliação}
\chapter{Psicopedagogicamente inconstitucionalíssimo otorrinolaringologista socioeducativo ético-profissional sociodemográfico interdisciplinaridade neuropsicológica avaliação}
\chapter{Psicopedagogicamente inconstitucionalíssimo otorrinolaringologista socioeducativo ético-profissional sociodemográfico interdisciplinaridade neuropsicológica avaliação}
\chapter{Psicopedagogicamente inconstitucionalíssimo otorrinolaringologista socioeducativo ético-profissional sociodemográfico interdisciplinaridade neuropsicológica avaliação}
\chapter{Psicopedagogicamente inconstitucionalíssimo otorrinolaringologista socioeducativo ético-profissional sociodemográfico interdisciplinaridade neuropsicológica avaliação}
\chapter{Psicopedagogicamente inconstitucionalíssimo otorrinolaringologista socioeducativo ético-profissional sociodemográfico interdisciplinaridade neuropsicológica avaliação}
\end{document}
'   % B: plus 4em (o do guia)
%   lualatex '\def\var{C}% !TeX program = lualatex
% ===================================================================
% sumario-longo.tex — mede o sumário sem hifenização (issue #35, "Fora do
% escopo"): com a hifenização desligada, a entrada longa de capítulo pode
% bater no número da página. Gera 3 x 16 capítulos, cada um com um prefixo
% (de 3 a 18 palavras) de uma frase longa: frase 1 de palavras médias,
% frase 2 com números, frase 3 só de palavras muito longas. Conta no log
% os "Overfull \hbox" do sumário: o sumário é a única parte do documento
% sem parágrafo de corpo, então todo Overfull "in paragraph" antes do
% primeiro \chapter é entrada que passou da margem.
%
% Uso (de uma pasta com livros_crp_acessivel_book.sty e crpsp-leg.sty no
% caminho; as fontes NEWJUNE por OSFONTDIR):
%   lualatex '\def\var{R}\input{sumario-longo}'   % R: rígido (até 0.5.x)
%   lualatex '\def\var{B}\input{sumario-longo}'   % B: plus 4em (o do guia)
%   lualatex '\def\var{C}\input{sumario-longo}'   % C: plus 4em + \tolerance 10000
%   lualatex '\def\var{S}\input{sumario-longo}'   % S: o que a .sty define
% Cada variante só redefine o \l@chapter depois da carga da .sty.
% ===================================================================
\DocumentMetadata{lang=pt-BR, pdfstandard=ua-2, pdfversion=2.0, tagging=on}
\documentclass[11pt,a5paper,twoside,openright]{book}
\usepackage{livros_crp_acessivel_book}
\hypersetup{pdftitle={Sumario longo}, pdfauthor={CRP SP}}
\makeatletter
\providecommand\var{S}
\def\crp@vR{R}\def\crp@vB{B}\def\crp@vC{C}
\def\crp@lchap#1{%  #1 = acréscimo à margem direita (rightskip)
  \renewcommand*\l@chapter[2]{%
    \ifnum \c@tocdepth >\m@ne
        \addpenalty{-\@highpenalty}\addvspace{0.5em \@plus\p@}%
        \setlength\@tempdima{1.5em}%
        \begingroup
            \parindent \z@ \rightskip \@pnumwidth #1
            \parfillskip -\@pnumwidth
            \leavevmode \bfseries
            \advance\leftskip\@tempdima \hskip -\leftskip
            ##1\nobreak\hfil
            \nobreak\hb@xt@\@pnumwidth{\hss ##2\kern-\p@\kern\p@}\par
            \penalty\@highpenalty
        \endgroup
    \fi}}
\ifx\var\crp@vR \crp@lchap{}\fi
\ifx\var\crp@vB \crp@lchap{\@plus 4em}\fi
\ifx\var\crp@vC \crp@lchap{\@plus 4em \tolerance=10000 \hbadness=10000 }\fi
\makeatother
\begin{document}
\frontmatter
\tableofcontents
\mainmatter
\chapter{Deficiência psicossocial e}
\chapter{Deficiência psicossocial e transtorno}
\chapter{Deficiência psicossocial e transtorno do}
\chapter{Deficiência psicossocial e transtorno do espectro}
\chapter{Deficiência psicossocial e transtorno do espectro autista:}
\chapter{Deficiência psicossocial e transtorno do espectro autista: orientações}
\chapter{Deficiência psicossocial e transtorno do espectro autista: orientações para}
\chapter{Deficiência psicossocial e transtorno do espectro autista: orientações para a}
\chapter{Deficiência psicossocial e transtorno do espectro autista: orientações para a realização}
\chapter{Deficiência psicossocial e transtorno do espectro autista: orientações para a realização de}
\chapter{Deficiência psicossocial e transtorno do espectro autista: orientações para a realização de eventos}
\chapter{Deficiência psicossocial e transtorno do espectro autista: orientações para a realização de eventos e}
\chapter{Deficiência psicossocial e transtorno do espectro autista: orientações para a realização de eventos e apresentações}
\chapter{Deficiência psicossocial e transtorno do espectro autista: orientações para a realização de eventos e apresentações acessíveis}
\chapter{Deficiência psicossocial e transtorno do espectro autista: orientações para a realização de eventos e apresentações acessíveis em}
\chapter{Deficiência psicossocial e transtorno do espectro autista: orientações para a realização de eventos e apresentações acessíveis em todo}
\chapter{Comunicação alternativa e}
\chapter{Comunicação alternativa e aumentativa}
\chapter{Comunicação alternativa e aumentativa na}
\chapter{Comunicação alternativa e aumentativa na participação}
\chapter{Comunicação alternativa e aumentativa na participação social}
\chapter{Comunicação alternativa e aumentativa na participação social de}
\chapter{Comunicação alternativa e aumentativa na participação social de pessoas}
\chapter{Comunicação alternativa e aumentativa na participação social de pessoas com}
\chapter{Comunicação alternativa e aumentativa na participação social de pessoas com deficiência}
\chapter{Comunicação alternativa e aumentativa na participação social de pessoas com deficiência intelectual}
\chapter{Comunicação alternativa e aumentativa na participação social de pessoas com deficiência intelectual em}
\chapter{Comunicação alternativa e aumentativa na participação social de pessoas com deficiência intelectual em eventos}
\chapter{Comunicação alternativa e aumentativa na participação social de pessoas com deficiência intelectual em eventos públicos}
\chapter{Comunicação alternativa e aumentativa na participação social de pessoas com deficiência intelectual em eventos públicos de}
\chapter{Comunicação alternativa e aumentativa na participação social de pessoas com deficiência intelectual em eventos públicos de grande}
\chapter{Comunicação alternativa e aumentativa na participação social de pessoas com deficiência intelectual em eventos públicos de grande porte}
\chapter{Psicopedagogicamente inconstitucionalíssimo otorrinolaringologista}
\chapter{Psicopedagogicamente inconstitucionalíssimo otorrinolaringologista socioeducativo}
\chapter{Psicopedagogicamente inconstitucionalíssimo otorrinolaringologista socioeducativo ético-profissional}
\chapter{Psicopedagogicamente inconstitucionalíssimo otorrinolaringologista socioeducativo ético-profissional sociodemográfico}
\chapter{Psicopedagogicamente inconstitucionalíssimo otorrinolaringologista socioeducativo ético-profissional sociodemográfico interdisciplinaridade}
\chapter{Psicopedagogicamente inconstitucionalíssimo otorrinolaringologista socioeducativo ético-profissional sociodemográfico interdisciplinaridade neuropsicológica}
\chapter{Psicopedagogicamente inconstitucionalíssimo otorrinolaringologista socioeducativo ético-profissional sociodemográfico interdisciplinaridade neuropsicológica avaliação}
\chapter{Psicopedagogicamente inconstitucionalíssimo otorrinolaringologista socioeducativo ético-profissional sociodemográfico interdisciplinaridade neuropsicológica avaliação}
\chapter{Psicopedagogicamente inconstitucionalíssimo otorrinolaringologista socioeducativo ético-profissional sociodemográfico interdisciplinaridade neuropsicológica avaliação}
\chapter{Psicopedagogicamente inconstitucionalíssimo otorrinolaringologista socioeducativo ético-profissional sociodemográfico interdisciplinaridade neuropsicológica avaliação}
\chapter{Psicopedagogicamente inconstitucionalíssimo otorrinolaringologista socioeducativo ético-profissional sociodemográfico interdisciplinaridade neuropsicológica avaliação}
\chapter{Psicopedagogicamente inconstitucionalíssimo otorrinolaringologista socioeducativo ético-profissional sociodemográfico interdisciplinaridade neuropsicológica avaliação}
\chapter{Psicopedagogicamente inconstitucionalíssimo otorrinolaringologista socioeducativo ético-profissional sociodemográfico interdisciplinaridade neuropsicológica avaliação}
\chapter{Psicopedagogicamente inconstitucionalíssimo otorrinolaringologista socioeducativo ético-profissional sociodemográfico interdisciplinaridade neuropsicológica avaliação}
\chapter{Psicopedagogicamente inconstitucionalíssimo otorrinolaringologista socioeducativo ético-profissional sociodemográfico interdisciplinaridade neuropsicológica avaliação}
\chapter{Psicopedagogicamente inconstitucionalíssimo otorrinolaringologista socioeducativo ético-profissional sociodemográfico interdisciplinaridade neuropsicológica avaliação}
\end{document}
'   % C: plus 4em + \tolerance 10000
%   lualatex '\def\var{S}% !TeX program = lualatex
% ===================================================================
% sumario-longo.tex — mede o sumário sem hifenização (issue #35, "Fora do
% escopo"): com a hifenização desligada, a entrada longa de capítulo pode
% bater no número da página. Gera 3 x 16 capítulos, cada um com um prefixo
% (de 3 a 18 palavras) de uma frase longa: frase 1 de palavras médias,
% frase 2 com números, frase 3 só de palavras muito longas. Conta no log
% os "Overfull \hbox" do sumário: o sumário é a única parte do documento
% sem parágrafo de corpo, então todo Overfull "in paragraph" antes do
% primeiro \chapter é entrada que passou da margem.
%
% Uso (de uma pasta com livros_crp_acessivel_book.sty e crpsp-leg.sty no
% caminho; as fontes NEWJUNE por OSFONTDIR):
%   lualatex '\def\var{R}\input{sumario-longo}'   % R: rígido (até 0.5.x)
%   lualatex '\def\var{B}\input{sumario-longo}'   % B: plus 4em (o do guia)
%   lualatex '\def\var{C}\input{sumario-longo}'   % C: plus 4em + \tolerance 10000
%   lualatex '\def\var{S}\input{sumario-longo}'   % S: o que a .sty define
% Cada variante só redefine o \l@chapter depois da carga da .sty.
% ===================================================================
\DocumentMetadata{lang=pt-BR, pdfstandard=ua-2, pdfversion=2.0, tagging=on}
\documentclass[11pt,a5paper,twoside,openright]{book}
\usepackage{livros_crp_acessivel_book}
\hypersetup{pdftitle={Sumario longo}, pdfauthor={CRP SP}}
\makeatletter
\providecommand\var{S}
\def\crp@vR{R}\def\crp@vB{B}\def\crp@vC{C}
\def\crp@lchap#1{%  #1 = acréscimo à margem direita (rightskip)
  \renewcommand*\l@chapter[2]{%
    \ifnum \c@tocdepth >\m@ne
        \addpenalty{-\@highpenalty}\addvspace{0.5em \@plus\p@}%
        \setlength\@tempdima{1.5em}%
        \begingroup
            \parindent \z@ \rightskip \@pnumwidth #1
            \parfillskip -\@pnumwidth
            \leavevmode \bfseries
            \advance\leftskip\@tempdima \hskip -\leftskip
            ##1\nobreak\hfil
            \nobreak\hb@xt@\@pnumwidth{\hss ##2\kern-\p@\kern\p@}\par
            \penalty\@highpenalty
        \endgroup
    \fi}}
\ifx\var\crp@vR \crp@lchap{}\fi
\ifx\var\crp@vB \crp@lchap{\@plus 4em}\fi
\ifx\var\crp@vC \crp@lchap{\@plus 4em \tolerance=10000 \hbadness=10000 }\fi
\makeatother
\begin{document}
\frontmatter
\tableofcontents
\mainmatter
\chapter{Deficiência psicossocial e}
\chapter{Deficiência psicossocial e transtorno}
\chapter{Deficiência psicossocial e transtorno do}
\chapter{Deficiência psicossocial e transtorno do espectro}
\chapter{Deficiência psicossocial e transtorno do espectro autista:}
\chapter{Deficiência psicossocial e transtorno do espectro autista: orientações}
\chapter{Deficiência psicossocial e transtorno do espectro autista: orientações para}
\chapter{Deficiência psicossocial e transtorno do espectro autista: orientações para a}
\chapter{Deficiência psicossocial e transtorno do espectro autista: orientações para a realização}
\chapter{Deficiência psicossocial e transtorno do espectro autista: orientações para a realização de}
\chapter{Deficiência psicossocial e transtorno do espectro autista: orientações para a realização de eventos}
\chapter{Deficiência psicossocial e transtorno do espectro autista: orientações para a realização de eventos e}
\chapter{Deficiência psicossocial e transtorno do espectro autista: orientações para a realização de eventos e apresentações}
\chapter{Deficiência psicossocial e transtorno do espectro autista: orientações para a realização de eventos e apresentações acessíveis}
\chapter{Deficiência psicossocial e transtorno do espectro autista: orientações para a realização de eventos e apresentações acessíveis em}
\chapter{Deficiência psicossocial e transtorno do espectro autista: orientações para a realização de eventos e apresentações acessíveis em todo}
\chapter{Comunicação alternativa e}
\chapter{Comunicação alternativa e aumentativa}
\chapter{Comunicação alternativa e aumentativa na}
\chapter{Comunicação alternativa e aumentativa na participação}
\chapter{Comunicação alternativa e aumentativa na participação social}
\chapter{Comunicação alternativa e aumentativa na participação social de}
\chapter{Comunicação alternativa e aumentativa na participação social de pessoas}
\chapter{Comunicação alternativa e aumentativa na participação social de pessoas com}
\chapter{Comunicação alternativa e aumentativa na participação social de pessoas com deficiência}
\chapter{Comunicação alternativa e aumentativa na participação social de pessoas com deficiência intelectual}
\chapter{Comunicação alternativa e aumentativa na participação social de pessoas com deficiência intelectual em}
\chapter{Comunicação alternativa e aumentativa na participação social de pessoas com deficiência intelectual em eventos}
\chapter{Comunicação alternativa e aumentativa na participação social de pessoas com deficiência intelectual em eventos públicos}
\chapter{Comunicação alternativa e aumentativa na participação social de pessoas com deficiência intelectual em eventos públicos de}
\chapter{Comunicação alternativa e aumentativa na participação social de pessoas com deficiência intelectual em eventos públicos de grande}
\chapter{Comunicação alternativa e aumentativa na participação social de pessoas com deficiência intelectual em eventos públicos de grande porte}
\chapter{Psicopedagogicamente inconstitucionalíssimo otorrinolaringologista}
\chapter{Psicopedagogicamente inconstitucionalíssimo otorrinolaringologista socioeducativo}
\chapter{Psicopedagogicamente inconstitucionalíssimo otorrinolaringologista socioeducativo ético-profissional}
\chapter{Psicopedagogicamente inconstitucionalíssimo otorrinolaringologista socioeducativo ético-profissional sociodemográfico}
\chapter{Psicopedagogicamente inconstitucionalíssimo otorrinolaringologista socioeducativo ético-profissional sociodemográfico interdisciplinaridade}
\chapter{Psicopedagogicamente inconstitucionalíssimo otorrinolaringologista socioeducativo ético-profissional sociodemográfico interdisciplinaridade neuropsicológica}
\chapter{Psicopedagogicamente inconstitucionalíssimo otorrinolaringologista socioeducativo ético-profissional sociodemográfico interdisciplinaridade neuropsicológica avaliação}
\chapter{Psicopedagogicamente inconstitucionalíssimo otorrinolaringologista socioeducativo ético-profissional sociodemográfico interdisciplinaridade neuropsicológica avaliação}
\chapter{Psicopedagogicamente inconstitucionalíssimo otorrinolaringologista socioeducativo ético-profissional sociodemográfico interdisciplinaridade neuropsicológica avaliação}
\chapter{Psicopedagogicamente inconstitucionalíssimo otorrinolaringologista socioeducativo ético-profissional sociodemográfico interdisciplinaridade neuropsicológica avaliação}
\chapter{Psicopedagogicamente inconstitucionalíssimo otorrinolaringologista socioeducativo ético-profissional sociodemográfico interdisciplinaridade neuropsicológica avaliação}
\chapter{Psicopedagogicamente inconstitucionalíssimo otorrinolaringologista socioeducativo ético-profissional sociodemográfico interdisciplinaridade neuropsicológica avaliação}
\chapter{Psicopedagogicamente inconstitucionalíssimo otorrinolaringologista socioeducativo ético-profissional sociodemográfico interdisciplinaridade neuropsicológica avaliação}
\chapter{Psicopedagogicamente inconstitucionalíssimo otorrinolaringologista socioeducativo ético-profissional sociodemográfico interdisciplinaridade neuropsicológica avaliação}
\chapter{Psicopedagogicamente inconstitucionalíssimo otorrinolaringologista socioeducativo ético-profissional sociodemográfico interdisciplinaridade neuropsicológica avaliação}
\chapter{Psicopedagogicamente inconstitucionalíssimo otorrinolaringologista socioeducativo ético-profissional sociodemográfico interdisciplinaridade neuropsicológica avaliação}
\end{document}
'   % S: o que a .sty define
% Cada variante só redefine o \l@chapter depois da carga da .sty.
% ===================================================================
\DocumentMetadata{lang=pt-BR, pdfstandard=ua-2, pdfversion=2.0, tagging=on}
\documentclass[11pt,a5paper,twoside,openright]{book}
\usepackage{livros_crp_acessivel_book}
\hypersetup{pdftitle={Sumario longo}, pdfauthor={CRP SP}}
\makeatletter
\providecommand\var{S}
\def\crp@vR{R}\def\crp@vB{B}\def\crp@vC{C}
\def\crp@lchap#1{%  #1 = acréscimo à margem direita (rightskip)
  \renewcommand*\l@chapter[2]{%
    \ifnum \c@tocdepth >\m@ne
        \addpenalty{-\@highpenalty}\addvspace{0.5em \@plus\p@}%
        \setlength\@tempdima{1.5em}%
        \begingroup
            \parindent \z@ \rightskip \@pnumwidth #1
            \parfillskip -\@pnumwidth
            \leavevmode \bfseries
            \advance\leftskip\@tempdima \hskip -\leftskip
            ##1\nobreak\hfil
            \nobreak\hb@xt@\@pnumwidth{\hss ##2\kern-\p@\kern\p@}\par
            \penalty\@highpenalty
        \endgroup
    \fi}}
\ifx\var\crp@vR \crp@lchap{}\fi
\ifx\var\crp@vB \crp@lchap{\@plus 4em}\fi
\ifx\var\crp@vC \crp@lchap{\@plus 4em \tolerance=10000 \hbadness=10000 }\fi
\makeatother
\begin{document}
\frontmatter
\tableofcontents
\mainmatter
\chapter{Deficiência psicossocial e}
\chapter{Deficiência psicossocial e transtorno}
\chapter{Deficiência psicossocial e transtorno do}
\chapter{Deficiência psicossocial e transtorno do espectro}
\chapter{Deficiência psicossocial e transtorno do espectro autista:}
\chapter{Deficiência psicossocial e transtorno do espectro autista: orientações}
\chapter{Deficiência psicossocial e transtorno do espectro autista: orientações para}
\chapter{Deficiência psicossocial e transtorno do espectro autista: orientações para a}
\chapter{Deficiência psicossocial e transtorno do espectro autista: orientações para a realização}
\chapter{Deficiência psicossocial e transtorno do espectro autista: orientações para a realização de}
\chapter{Deficiência psicossocial e transtorno do espectro autista: orientações para a realização de eventos}
\chapter{Deficiência psicossocial e transtorno do espectro autista: orientações para a realização de eventos e}
\chapter{Deficiência psicossocial e transtorno do espectro autista: orientações para a realização de eventos e apresentações}
\chapter{Deficiência psicossocial e transtorno do espectro autista: orientações para a realização de eventos e apresentações acessíveis}
\chapter{Deficiência psicossocial e transtorno do espectro autista: orientações para a realização de eventos e apresentações acessíveis em}
\chapter{Deficiência psicossocial e transtorno do espectro autista: orientações para a realização de eventos e apresentações acessíveis em todo}
\chapter{Comunicação alternativa e}
\chapter{Comunicação alternativa e aumentativa}
\chapter{Comunicação alternativa e aumentativa na}
\chapter{Comunicação alternativa e aumentativa na participação}
\chapter{Comunicação alternativa e aumentativa na participação social}
\chapter{Comunicação alternativa e aumentativa na participação social de}
\chapter{Comunicação alternativa e aumentativa na participação social de pessoas}
\chapter{Comunicação alternativa e aumentativa na participação social de pessoas com}
\chapter{Comunicação alternativa e aumentativa na participação social de pessoas com deficiência}
\chapter{Comunicação alternativa e aumentativa na participação social de pessoas com deficiência intelectual}
\chapter{Comunicação alternativa e aumentativa na participação social de pessoas com deficiência intelectual em}
\chapter{Comunicação alternativa e aumentativa na participação social de pessoas com deficiência intelectual em eventos}
\chapter{Comunicação alternativa e aumentativa na participação social de pessoas com deficiência intelectual em eventos públicos}
\chapter{Comunicação alternativa e aumentativa na participação social de pessoas com deficiência intelectual em eventos públicos de}
\chapter{Comunicação alternativa e aumentativa na participação social de pessoas com deficiência intelectual em eventos públicos de grande}
\chapter{Comunicação alternativa e aumentativa na participação social de pessoas com deficiência intelectual em eventos públicos de grande porte}
\chapter{Psicopedagogicamente inconstitucionalíssimo otorrinolaringologista}
\chapter{Psicopedagogicamente inconstitucionalíssimo otorrinolaringologista socioeducativo}
\chapter{Psicopedagogicamente inconstitucionalíssimo otorrinolaringologista socioeducativo ético-profissional}
\chapter{Psicopedagogicamente inconstitucionalíssimo otorrinolaringologista socioeducativo ético-profissional sociodemográfico}
\chapter{Psicopedagogicamente inconstitucionalíssimo otorrinolaringologista socioeducativo ético-profissional sociodemográfico interdisciplinaridade}
\chapter{Psicopedagogicamente inconstitucionalíssimo otorrinolaringologista socioeducativo ético-profissional sociodemográfico interdisciplinaridade neuropsicológica}
\chapter{Psicopedagogicamente inconstitucionalíssimo otorrinolaringologista socioeducativo ético-profissional sociodemográfico interdisciplinaridade neuropsicológica avaliação}
\chapter{Psicopedagogicamente inconstitucionalíssimo otorrinolaringologista socioeducativo ético-profissional sociodemográfico interdisciplinaridade neuropsicológica avaliação}
\chapter{Psicopedagogicamente inconstitucionalíssimo otorrinolaringologista socioeducativo ético-profissional sociodemográfico interdisciplinaridade neuropsicológica avaliação}
\chapter{Psicopedagogicamente inconstitucionalíssimo otorrinolaringologista socioeducativo ético-profissional sociodemográfico interdisciplinaridade neuropsicológica avaliação}
\chapter{Psicopedagogicamente inconstitucionalíssimo otorrinolaringologista socioeducativo ético-profissional sociodemográfico interdisciplinaridade neuropsicológica avaliação}
\chapter{Psicopedagogicamente inconstitucionalíssimo otorrinolaringologista socioeducativo ético-profissional sociodemográfico interdisciplinaridade neuropsicológica avaliação}
\chapter{Psicopedagogicamente inconstitucionalíssimo otorrinolaringologista socioeducativo ético-profissional sociodemográfico interdisciplinaridade neuropsicológica avaliação}
\chapter{Psicopedagogicamente inconstitucionalíssimo otorrinolaringologista socioeducativo ético-profissional sociodemográfico interdisciplinaridade neuropsicológica avaliação}
\chapter{Psicopedagogicamente inconstitucionalíssimo otorrinolaringologista socioeducativo ético-profissional sociodemográfico interdisciplinaridade neuropsicológica avaliação}
\chapter{Psicopedagogicamente inconstitucionalíssimo otorrinolaringologista socioeducativo ético-profissional sociodemográfico interdisciplinaridade neuropsicológica avaliação}
\end{document}
'   % C: plus 4em + \tolerance 10000
%   lualatex '\def\var{S}% !TeX program = lualatex
% ===================================================================
% sumario-longo.tex — mede o sumário sem hifenização (issue #35, "Fora do
% escopo"): com a hifenização desligada, a entrada longa de capítulo pode
% bater no número da página. Gera 3 x 16 capítulos, cada um com um prefixo
% (de 3 a 18 palavras) de uma frase longa: frase 1 de palavras médias,
% frase 2 com números, frase 3 só de palavras muito longas. Conta no log
% os "Overfull \hbox" do sumário: o sumário é a única parte do documento
% sem parágrafo de corpo, então todo Overfull "in paragraph" antes do
% primeiro \chapter é entrada que passou da margem.
%
% Uso (de uma pasta com livros_crp_acessivel_book.sty e crpsp-leg.sty no
% caminho; as fontes NEWJUNE por OSFONTDIR):
%   lualatex '\def\var{R}% !TeX program = lualatex
% ===================================================================
% sumario-longo.tex — mede o sumário sem hifenização (issue #35, "Fora do
% escopo"): com a hifenização desligada, a entrada longa de capítulo pode
% bater no número da página. Gera 3 x 16 capítulos, cada um com um prefixo
% (de 3 a 18 palavras) de uma frase longa: frase 1 de palavras médias,
% frase 2 com números, frase 3 só de palavras muito longas. Conta no log
% os "Overfull \hbox" do sumário: o sumário é a única parte do documento
% sem parágrafo de corpo, então todo Overfull "in paragraph" antes do
% primeiro \chapter é entrada que passou da margem.
%
% Uso (de uma pasta com livros_crp_acessivel_book.sty e crpsp-leg.sty no
% caminho; as fontes NEWJUNE por OSFONTDIR):
%   lualatex '\def\var{R}\input{sumario-longo}'   % R: rígido (até 0.5.x)
%   lualatex '\def\var{B}\input{sumario-longo}'   % B: plus 4em (o do guia)
%   lualatex '\def\var{C}\input{sumario-longo}'   % C: plus 4em + \tolerance 10000
%   lualatex '\def\var{S}\input{sumario-longo}'   % S: o que a .sty define
% Cada variante só redefine o \l@chapter depois da carga da .sty.
% ===================================================================
\DocumentMetadata{lang=pt-BR, pdfstandard=ua-2, pdfversion=2.0, tagging=on}
\documentclass[11pt,a5paper,twoside,openright]{book}
\usepackage{livros_crp_acessivel_book}
\hypersetup{pdftitle={Sumario longo}, pdfauthor={CRP SP}}
\makeatletter
\providecommand\var{S}
\def\crp@vR{R}\def\crp@vB{B}\def\crp@vC{C}
\def\crp@lchap#1{%  #1 = acréscimo à margem direita (rightskip)
  \renewcommand*\l@chapter[2]{%
    \ifnum \c@tocdepth >\m@ne
        \addpenalty{-\@highpenalty}\addvspace{0.5em \@plus\p@}%
        \setlength\@tempdima{1.5em}%
        \begingroup
            \parindent \z@ \rightskip \@pnumwidth #1
            \parfillskip -\@pnumwidth
            \leavevmode \bfseries
            \advance\leftskip\@tempdima \hskip -\leftskip
            ##1\nobreak\hfil
            \nobreak\hb@xt@\@pnumwidth{\hss ##2\kern-\p@\kern\p@}\par
            \penalty\@highpenalty
        \endgroup
    \fi}}
\ifx\var\crp@vR \crp@lchap{}\fi
\ifx\var\crp@vB \crp@lchap{\@plus 4em}\fi
\ifx\var\crp@vC \crp@lchap{\@plus 4em \tolerance=10000 \hbadness=10000 }\fi
\makeatother
\begin{document}
\frontmatter
\tableofcontents
\mainmatter
\chapter{Deficiência psicossocial e}
\chapter{Deficiência psicossocial e transtorno}
\chapter{Deficiência psicossocial e transtorno do}
\chapter{Deficiência psicossocial e transtorno do espectro}
\chapter{Deficiência psicossocial e transtorno do espectro autista:}
\chapter{Deficiência psicossocial e transtorno do espectro autista: orientações}
\chapter{Deficiência psicossocial e transtorno do espectro autista: orientações para}
\chapter{Deficiência psicossocial e transtorno do espectro autista: orientações para a}
\chapter{Deficiência psicossocial e transtorno do espectro autista: orientações para a realização}
\chapter{Deficiência psicossocial e transtorno do espectro autista: orientações para a realização de}
\chapter{Deficiência psicossocial e transtorno do espectro autista: orientações para a realização de eventos}
\chapter{Deficiência psicossocial e transtorno do espectro autista: orientações para a realização de eventos e}
\chapter{Deficiência psicossocial e transtorno do espectro autista: orientações para a realização de eventos e apresentações}
\chapter{Deficiência psicossocial e transtorno do espectro autista: orientações para a realização de eventos e apresentações acessíveis}
\chapter{Deficiência psicossocial e transtorno do espectro autista: orientações para a realização de eventos e apresentações acessíveis em}
\chapter{Deficiência psicossocial e transtorno do espectro autista: orientações para a realização de eventos e apresentações acessíveis em todo}
\chapter{Comunicação alternativa e}
\chapter{Comunicação alternativa e aumentativa}
\chapter{Comunicação alternativa e aumentativa na}
\chapter{Comunicação alternativa e aumentativa na participação}
\chapter{Comunicação alternativa e aumentativa na participação social}
\chapter{Comunicação alternativa e aumentativa na participação social de}
\chapter{Comunicação alternativa e aumentativa na participação social de pessoas}
\chapter{Comunicação alternativa e aumentativa na participação social de pessoas com}
\chapter{Comunicação alternativa e aumentativa na participação social de pessoas com deficiência}
\chapter{Comunicação alternativa e aumentativa na participação social de pessoas com deficiência intelectual}
\chapter{Comunicação alternativa e aumentativa na participação social de pessoas com deficiência intelectual em}
\chapter{Comunicação alternativa e aumentativa na participação social de pessoas com deficiência intelectual em eventos}
\chapter{Comunicação alternativa e aumentativa na participação social de pessoas com deficiência intelectual em eventos públicos}
\chapter{Comunicação alternativa e aumentativa na participação social de pessoas com deficiência intelectual em eventos públicos de}
\chapter{Comunicação alternativa e aumentativa na participação social de pessoas com deficiência intelectual em eventos públicos de grande}
\chapter{Comunicação alternativa e aumentativa na participação social de pessoas com deficiência intelectual em eventos públicos de grande porte}
\chapter{Psicopedagogicamente inconstitucionalíssimo otorrinolaringologista}
\chapter{Psicopedagogicamente inconstitucionalíssimo otorrinolaringologista socioeducativo}
\chapter{Psicopedagogicamente inconstitucionalíssimo otorrinolaringologista socioeducativo ético-profissional}
\chapter{Psicopedagogicamente inconstitucionalíssimo otorrinolaringologista socioeducativo ético-profissional sociodemográfico}
\chapter{Psicopedagogicamente inconstitucionalíssimo otorrinolaringologista socioeducativo ético-profissional sociodemográfico interdisciplinaridade}
\chapter{Psicopedagogicamente inconstitucionalíssimo otorrinolaringologista socioeducativo ético-profissional sociodemográfico interdisciplinaridade neuropsicológica}
\chapter{Psicopedagogicamente inconstitucionalíssimo otorrinolaringologista socioeducativo ético-profissional sociodemográfico interdisciplinaridade neuropsicológica avaliação}
\chapter{Psicopedagogicamente inconstitucionalíssimo otorrinolaringologista socioeducativo ético-profissional sociodemográfico interdisciplinaridade neuropsicológica avaliação}
\chapter{Psicopedagogicamente inconstitucionalíssimo otorrinolaringologista socioeducativo ético-profissional sociodemográfico interdisciplinaridade neuropsicológica avaliação}
\chapter{Psicopedagogicamente inconstitucionalíssimo otorrinolaringologista socioeducativo ético-profissional sociodemográfico interdisciplinaridade neuropsicológica avaliação}
\chapter{Psicopedagogicamente inconstitucionalíssimo otorrinolaringologista socioeducativo ético-profissional sociodemográfico interdisciplinaridade neuropsicológica avaliação}
\chapter{Psicopedagogicamente inconstitucionalíssimo otorrinolaringologista socioeducativo ético-profissional sociodemográfico interdisciplinaridade neuropsicológica avaliação}
\chapter{Psicopedagogicamente inconstitucionalíssimo otorrinolaringologista socioeducativo ético-profissional sociodemográfico interdisciplinaridade neuropsicológica avaliação}
\chapter{Psicopedagogicamente inconstitucionalíssimo otorrinolaringologista socioeducativo ético-profissional sociodemográfico interdisciplinaridade neuropsicológica avaliação}
\chapter{Psicopedagogicamente inconstitucionalíssimo otorrinolaringologista socioeducativo ético-profissional sociodemográfico interdisciplinaridade neuropsicológica avaliação}
\chapter{Psicopedagogicamente inconstitucionalíssimo otorrinolaringologista socioeducativo ético-profissional sociodemográfico interdisciplinaridade neuropsicológica avaliação}
\end{document}
'   % R: rígido (até 0.5.x)
%   lualatex '\def\var{B}% !TeX program = lualatex
% ===================================================================
% sumario-longo.tex — mede o sumário sem hifenização (issue #35, "Fora do
% escopo"): com a hifenização desligada, a entrada longa de capítulo pode
% bater no número da página. Gera 3 x 16 capítulos, cada um com um prefixo
% (de 3 a 18 palavras) de uma frase longa: frase 1 de palavras médias,
% frase 2 com números, frase 3 só de palavras muito longas. Conta no log
% os "Overfull \hbox" do sumário: o sumário é a única parte do documento
% sem parágrafo de corpo, então todo Overfull "in paragraph" antes do
% primeiro \chapter é entrada que passou da margem.
%
% Uso (de uma pasta com livros_crp_acessivel_book.sty e crpsp-leg.sty no
% caminho; as fontes NEWJUNE por OSFONTDIR):
%   lualatex '\def\var{R}\input{sumario-longo}'   % R: rígido (até 0.5.x)
%   lualatex '\def\var{B}\input{sumario-longo}'   % B: plus 4em (o do guia)
%   lualatex '\def\var{C}\input{sumario-longo}'   % C: plus 4em + \tolerance 10000
%   lualatex '\def\var{S}\input{sumario-longo}'   % S: o que a .sty define
% Cada variante só redefine o \l@chapter depois da carga da .sty.
% ===================================================================
\DocumentMetadata{lang=pt-BR, pdfstandard=ua-2, pdfversion=2.0, tagging=on}
\documentclass[11pt,a5paper,twoside,openright]{book}
\usepackage{livros_crp_acessivel_book}
\hypersetup{pdftitle={Sumario longo}, pdfauthor={CRP SP}}
\makeatletter
\providecommand\var{S}
\def\crp@vR{R}\def\crp@vB{B}\def\crp@vC{C}
\def\crp@lchap#1{%  #1 = acréscimo à margem direita (rightskip)
  \renewcommand*\l@chapter[2]{%
    \ifnum \c@tocdepth >\m@ne
        \addpenalty{-\@highpenalty}\addvspace{0.5em \@plus\p@}%
        \setlength\@tempdima{1.5em}%
        \begingroup
            \parindent \z@ \rightskip \@pnumwidth #1
            \parfillskip -\@pnumwidth
            \leavevmode \bfseries
            \advance\leftskip\@tempdima \hskip -\leftskip
            ##1\nobreak\hfil
            \nobreak\hb@xt@\@pnumwidth{\hss ##2\kern-\p@\kern\p@}\par
            \penalty\@highpenalty
        \endgroup
    \fi}}
\ifx\var\crp@vR \crp@lchap{}\fi
\ifx\var\crp@vB \crp@lchap{\@plus 4em}\fi
\ifx\var\crp@vC \crp@lchap{\@plus 4em \tolerance=10000 \hbadness=10000 }\fi
\makeatother
\begin{document}
\frontmatter
\tableofcontents
\mainmatter
\chapter{Deficiência psicossocial e}
\chapter{Deficiência psicossocial e transtorno}
\chapter{Deficiência psicossocial e transtorno do}
\chapter{Deficiência psicossocial e transtorno do espectro}
\chapter{Deficiência psicossocial e transtorno do espectro autista:}
\chapter{Deficiência psicossocial e transtorno do espectro autista: orientações}
\chapter{Deficiência psicossocial e transtorno do espectro autista: orientações para}
\chapter{Deficiência psicossocial e transtorno do espectro autista: orientações para a}
\chapter{Deficiência psicossocial e transtorno do espectro autista: orientações para a realização}
\chapter{Deficiência psicossocial e transtorno do espectro autista: orientações para a realização de}
\chapter{Deficiência psicossocial e transtorno do espectro autista: orientações para a realização de eventos}
\chapter{Deficiência psicossocial e transtorno do espectro autista: orientações para a realização de eventos e}
\chapter{Deficiência psicossocial e transtorno do espectro autista: orientações para a realização de eventos e apresentações}
\chapter{Deficiência psicossocial e transtorno do espectro autista: orientações para a realização de eventos e apresentações acessíveis}
\chapter{Deficiência psicossocial e transtorno do espectro autista: orientações para a realização de eventos e apresentações acessíveis em}
\chapter{Deficiência psicossocial e transtorno do espectro autista: orientações para a realização de eventos e apresentações acessíveis em todo}
\chapter{Comunicação alternativa e}
\chapter{Comunicação alternativa e aumentativa}
\chapter{Comunicação alternativa e aumentativa na}
\chapter{Comunicação alternativa e aumentativa na participação}
\chapter{Comunicação alternativa e aumentativa na participação social}
\chapter{Comunicação alternativa e aumentativa na participação social de}
\chapter{Comunicação alternativa e aumentativa na participação social de pessoas}
\chapter{Comunicação alternativa e aumentativa na participação social de pessoas com}
\chapter{Comunicação alternativa e aumentativa na participação social de pessoas com deficiência}
\chapter{Comunicação alternativa e aumentativa na participação social de pessoas com deficiência intelectual}
\chapter{Comunicação alternativa e aumentativa na participação social de pessoas com deficiência intelectual em}
\chapter{Comunicação alternativa e aumentativa na participação social de pessoas com deficiência intelectual em eventos}
\chapter{Comunicação alternativa e aumentativa na participação social de pessoas com deficiência intelectual em eventos públicos}
\chapter{Comunicação alternativa e aumentativa na participação social de pessoas com deficiência intelectual em eventos públicos de}
\chapter{Comunicação alternativa e aumentativa na participação social de pessoas com deficiência intelectual em eventos públicos de grande}
\chapter{Comunicação alternativa e aumentativa na participação social de pessoas com deficiência intelectual em eventos públicos de grande porte}
\chapter{Psicopedagogicamente inconstitucionalíssimo otorrinolaringologista}
\chapter{Psicopedagogicamente inconstitucionalíssimo otorrinolaringologista socioeducativo}
\chapter{Psicopedagogicamente inconstitucionalíssimo otorrinolaringologista socioeducativo ético-profissional}
\chapter{Psicopedagogicamente inconstitucionalíssimo otorrinolaringologista socioeducativo ético-profissional sociodemográfico}
\chapter{Psicopedagogicamente inconstitucionalíssimo otorrinolaringologista socioeducativo ético-profissional sociodemográfico interdisciplinaridade}
\chapter{Psicopedagogicamente inconstitucionalíssimo otorrinolaringologista socioeducativo ético-profissional sociodemográfico interdisciplinaridade neuropsicológica}
\chapter{Psicopedagogicamente inconstitucionalíssimo otorrinolaringologista socioeducativo ético-profissional sociodemográfico interdisciplinaridade neuropsicológica avaliação}
\chapter{Psicopedagogicamente inconstitucionalíssimo otorrinolaringologista socioeducativo ético-profissional sociodemográfico interdisciplinaridade neuropsicológica avaliação}
\chapter{Psicopedagogicamente inconstitucionalíssimo otorrinolaringologista socioeducativo ético-profissional sociodemográfico interdisciplinaridade neuropsicológica avaliação}
\chapter{Psicopedagogicamente inconstitucionalíssimo otorrinolaringologista socioeducativo ético-profissional sociodemográfico interdisciplinaridade neuropsicológica avaliação}
\chapter{Psicopedagogicamente inconstitucionalíssimo otorrinolaringologista socioeducativo ético-profissional sociodemográfico interdisciplinaridade neuropsicológica avaliação}
\chapter{Psicopedagogicamente inconstitucionalíssimo otorrinolaringologista socioeducativo ético-profissional sociodemográfico interdisciplinaridade neuropsicológica avaliação}
\chapter{Psicopedagogicamente inconstitucionalíssimo otorrinolaringologista socioeducativo ético-profissional sociodemográfico interdisciplinaridade neuropsicológica avaliação}
\chapter{Psicopedagogicamente inconstitucionalíssimo otorrinolaringologista socioeducativo ético-profissional sociodemográfico interdisciplinaridade neuropsicológica avaliação}
\chapter{Psicopedagogicamente inconstitucionalíssimo otorrinolaringologista socioeducativo ético-profissional sociodemográfico interdisciplinaridade neuropsicológica avaliação}
\chapter{Psicopedagogicamente inconstitucionalíssimo otorrinolaringologista socioeducativo ético-profissional sociodemográfico interdisciplinaridade neuropsicológica avaliação}
\end{document}
'   % B: plus 4em (o do guia)
%   lualatex '\def\var{C}% !TeX program = lualatex
% ===================================================================
% sumario-longo.tex — mede o sumário sem hifenização (issue #35, "Fora do
% escopo"): com a hifenização desligada, a entrada longa de capítulo pode
% bater no número da página. Gera 3 x 16 capítulos, cada um com um prefixo
% (de 3 a 18 palavras) de uma frase longa: frase 1 de palavras médias,
% frase 2 com números, frase 3 só de palavras muito longas. Conta no log
% os "Overfull \hbox" do sumário: o sumário é a única parte do documento
% sem parágrafo de corpo, então todo Overfull "in paragraph" antes do
% primeiro \chapter é entrada que passou da margem.
%
% Uso (de uma pasta com livros_crp_acessivel_book.sty e crpsp-leg.sty no
% caminho; as fontes NEWJUNE por OSFONTDIR):
%   lualatex '\def\var{R}\input{sumario-longo}'   % R: rígido (até 0.5.x)
%   lualatex '\def\var{B}\input{sumario-longo}'   % B: plus 4em (o do guia)
%   lualatex '\def\var{C}\input{sumario-longo}'   % C: plus 4em + \tolerance 10000
%   lualatex '\def\var{S}\input{sumario-longo}'   % S: o que a .sty define
% Cada variante só redefine o \l@chapter depois da carga da .sty.
% ===================================================================
\DocumentMetadata{lang=pt-BR, pdfstandard=ua-2, pdfversion=2.0, tagging=on}
\documentclass[11pt,a5paper,twoside,openright]{book}
\usepackage{livros_crp_acessivel_book}
\hypersetup{pdftitle={Sumario longo}, pdfauthor={CRP SP}}
\makeatletter
\providecommand\var{S}
\def\crp@vR{R}\def\crp@vB{B}\def\crp@vC{C}
\def\crp@lchap#1{%  #1 = acréscimo à margem direita (rightskip)
  \renewcommand*\l@chapter[2]{%
    \ifnum \c@tocdepth >\m@ne
        \addpenalty{-\@highpenalty}\addvspace{0.5em \@plus\p@}%
        \setlength\@tempdima{1.5em}%
        \begingroup
            \parindent \z@ \rightskip \@pnumwidth #1
            \parfillskip -\@pnumwidth
            \leavevmode \bfseries
            \advance\leftskip\@tempdima \hskip -\leftskip
            ##1\nobreak\hfil
            \nobreak\hb@xt@\@pnumwidth{\hss ##2\kern-\p@\kern\p@}\par
            \penalty\@highpenalty
        \endgroup
    \fi}}
\ifx\var\crp@vR \crp@lchap{}\fi
\ifx\var\crp@vB \crp@lchap{\@plus 4em}\fi
\ifx\var\crp@vC \crp@lchap{\@plus 4em \tolerance=10000 \hbadness=10000 }\fi
\makeatother
\begin{document}
\frontmatter
\tableofcontents
\mainmatter
\chapter{Deficiência psicossocial e}
\chapter{Deficiência psicossocial e transtorno}
\chapter{Deficiência psicossocial e transtorno do}
\chapter{Deficiência psicossocial e transtorno do espectro}
\chapter{Deficiência psicossocial e transtorno do espectro autista:}
\chapter{Deficiência psicossocial e transtorno do espectro autista: orientações}
\chapter{Deficiência psicossocial e transtorno do espectro autista: orientações para}
\chapter{Deficiência psicossocial e transtorno do espectro autista: orientações para a}
\chapter{Deficiência psicossocial e transtorno do espectro autista: orientações para a realização}
\chapter{Deficiência psicossocial e transtorno do espectro autista: orientações para a realização de}
\chapter{Deficiência psicossocial e transtorno do espectro autista: orientações para a realização de eventos}
\chapter{Deficiência psicossocial e transtorno do espectro autista: orientações para a realização de eventos e}
\chapter{Deficiência psicossocial e transtorno do espectro autista: orientações para a realização de eventos e apresentações}
\chapter{Deficiência psicossocial e transtorno do espectro autista: orientações para a realização de eventos e apresentações acessíveis}
\chapter{Deficiência psicossocial e transtorno do espectro autista: orientações para a realização de eventos e apresentações acessíveis em}
\chapter{Deficiência psicossocial e transtorno do espectro autista: orientações para a realização de eventos e apresentações acessíveis em todo}
\chapter{Comunicação alternativa e}
\chapter{Comunicação alternativa e aumentativa}
\chapter{Comunicação alternativa e aumentativa na}
\chapter{Comunicação alternativa e aumentativa na participação}
\chapter{Comunicação alternativa e aumentativa na participação social}
\chapter{Comunicação alternativa e aumentativa na participação social de}
\chapter{Comunicação alternativa e aumentativa na participação social de pessoas}
\chapter{Comunicação alternativa e aumentativa na participação social de pessoas com}
\chapter{Comunicação alternativa e aumentativa na participação social de pessoas com deficiência}
\chapter{Comunicação alternativa e aumentativa na participação social de pessoas com deficiência intelectual}
\chapter{Comunicação alternativa e aumentativa na participação social de pessoas com deficiência intelectual em}
\chapter{Comunicação alternativa e aumentativa na participação social de pessoas com deficiência intelectual em eventos}
\chapter{Comunicação alternativa e aumentativa na participação social de pessoas com deficiência intelectual em eventos públicos}
\chapter{Comunicação alternativa e aumentativa na participação social de pessoas com deficiência intelectual em eventos públicos de}
\chapter{Comunicação alternativa e aumentativa na participação social de pessoas com deficiência intelectual em eventos públicos de grande}
\chapter{Comunicação alternativa e aumentativa na participação social de pessoas com deficiência intelectual em eventos públicos de grande porte}
\chapter{Psicopedagogicamente inconstitucionalíssimo otorrinolaringologista}
\chapter{Psicopedagogicamente inconstitucionalíssimo otorrinolaringologista socioeducativo}
\chapter{Psicopedagogicamente inconstitucionalíssimo otorrinolaringologista socioeducativo ético-profissional}
\chapter{Psicopedagogicamente inconstitucionalíssimo otorrinolaringologista socioeducativo ético-profissional sociodemográfico}
\chapter{Psicopedagogicamente inconstitucionalíssimo otorrinolaringologista socioeducativo ético-profissional sociodemográfico interdisciplinaridade}
\chapter{Psicopedagogicamente inconstitucionalíssimo otorrinolaringologista socioeducativo ético-profissional sociodemográfico interdisciplinaridade neuropsicológica}
\chapter{Psicopedagogicamente inconstitucionalíssimo otorrinolaringologista socioeducativo ético-profissional sociodemográfico interdisciplinaridade neuropsicológica avaliação}
\chapter{Psicopedagogicamente inconstitucionalíssimo otorrinolaringologista socioeducativo ético-profissional sociodemográfico interdisciplinaridade neuropsicológica avaliação}
\chapter{Psicopedagogicamente inconstitucionalíssimo otorrinolaringologista socioeducativo ético-profissional sociodemográfico interdisciplinaridade neuropsicológica avaliação}
\chapter{Psicopedagogicamente inconstitucionalíssimo otorrinolaringologista socioeducativo ético-profissional sociodemográfico interdisciplinaridade neuropsicológica avaliação}
\chapter{Psicopedagogicamente inconstitucionalíssimo otorrinolaringologista socioeducativo ético-profissional sociodemográfico interdisciplinaridade neuropsicológica avaliação}
\chapter{Psicopedagogicamente inconstitucionalíssimo otorrinolaringologista socioeducativo ético-profissional sociodemográfico interdisciplinaridade neuropsicológica avaliação}
\chapter{Psicopedagogicamente inconstitucionalíssimo otorrinolaringologista socioeducativo ético-profissional sociodemográfico interdisciplinaridade neuropsicológica avaliação}
\chapter{Psicopedagogicamente inconstitucionalíssimo otorrinolaringologista socioeducativo ético-profissional sociodemográfico interdisciplinaridade neuropsicológica avaliação}
\chapter{Psicopedagogicamente inconstitucionalíssimo otorrinolaringologista socioeducativo ético-profissional sociodemográfico interdisciplinaridade neuropsicológica avaliação}
\chapter{Psicopedagogicamente inconstitucionalíssimo otorrinolaringologista socioeducativo ético-profissional sociodemográfico interdisciplinaridade neuropsicológica avaliação}
\end{document}
'   % C: plus 4em + \tolerance 10000
%   lualatex '\def\var{S}% !TeX program = lualatex
% ===================================================================
% sumario-longo.tex — mede o sumário sem hifenização (issue #35, "Fora do
% escopo"): com a hifenização desligada, a entrada longa de capítulo pode
% bater no número da página. Gera 3 x 16 capítulos, cada um com um prefixo
% (de 3 a 18 palavras) de uma frase longa: frase 1 de palavras médias,
% frase 2 com números, frase 3 só de palavras muito longas. Conta no log
% os "Overfull \hbox" do sumário: o sumário é a única parte do documento
% sem parágrafo de corpo, então todo Overfull "in paragraph" antes do
% primeiro \chapter é entrada que passou da margem.
%
% Uso (de uma pasta com livros_crp_acessivel_book.sty e crpsp-leg.sty no
% caminho; as fontes NEWJUNE por OSFONTDIR):
%   lualatex '\def\var{R}\input{sumario-longo}'   % R: rígido (até 0.5.x)
%   lualatex '\def\var{B}\input{sumario-longo}'   % B: plus 4em (o do guia)
%   lualatex '\def\var{C}\input{sumario-longo}'   % C: plus 4em + \tolerance 10000
%   lualatex '\def\var{S}\input{sumario-longo}'   % S: o que a .sty define
% Cada variante só redefine o \l@chapter depois da carga da .sty.
% ===================================================================
\DocumentMetadata{lang=pt-BR, pdfstandard=ua-2, pdfversion=2.0, tagging=on}
\documentclass[11pt,a5paper,twoside,openright]{book}
\usepackage{livros_crp_acessivel_book}
\hypersetup{pdftitle={Sumario longo}, pdfauthor={CRP SP}}
\makeatletter
\providecommand\var{S}
\def\crp@vR{R}\def\crp@vB{B}\def\crp@vC{C}
\def\crp@lchap#1{%  #1 = acréscimo à margem direita (rightskip)
  \renewcommand*\l@chapter[2]{%
    \ifnum \c@tocdepth >\m@ne
        \addpenalty{-\@highpenalty}\addvspace{0.5em \@plus\p@}%
        \setlength\@tempdima{1.5em}%
        \begingroup
            \parindent \z@ \rightskip \@pnumwidth #1
            \parfillskip -\@pnumwidth
            \leavevmode \bfseries
            \advance\leftskip\@tempdima \hskip -\leftskip
            ##1\nobreak\hfil
            \nobreak\hb@xt@\@pnumwidth{\hss ##2\kern-\p@\kern\p@}\par
            \penalty\@highpenalty
        \endgroup
    \fi}}
\ifx\var\crp@vR \crp@lchap{}\fi
\ifx\var\crp@vB \crp@lchap{\@plus 4em}\fi
\ifx\var\crp@vC \crp@lchap{\@plus 4em \tolerance=10000 \hbadness=10000 }\fi
\makeatother
\begin{document}
\frontmatter
\tableofcontents
\mainmatter
\chapter{Deficiência psicossocial e}
\chapter{Deficiência psicossocial e transtorno}
\chapter{Deficiência psicossocial e transtorno do}
\chapter{Deficiência psicossocial e transtorno do espectro}
\chapter{Deficiência psicossocial e transtorno do espectro autista:}
\chapter{Deficiência psicossocial e transtorno do espectro autista: orientações}
\chapter{Deficiência psicossocial e transtorno do espectro autista: orientações para}
\chapter{Deficiência psicossocial e transtorno do espectro autista: orientações para a}
\chapter{Deficiência psicossocial e transtorno do espectro autista: orientações para a realização}
\chapter{Deficiência psicossocial e transtorno do espectro autista: orientações para a realização de}
\chapter{Deficiência psicossocial e transtorno do espectro autista: orientações para a realização de eventos}
\chapter{Deficiência psicossocial e transtorno do espectro autista: orientações para a realização de eventos e}
\chapter{Deficiência psicossocial e transtorno do espectro autista: orientações para a realização de eventos e apresentações}
\chapter{Deficiência psicossocial e transtorno do espectro autista: orientações para a realização de eventos e apresentações acessíveis}
\chapter{Deficiência psicossocial e transtorno do espectro autista: orientações para a realização de eventos e apresentações acessíveis em}
\chapter{Deficiência psicossocial e transtorno do espectro autista: orientações para a realização de eventos e apresentações acessíveis em todo}
\chapter{Comunicação alternativa e}
\chapter{Comunicação alternativa e aumentativa}
\chapter{Comunicação alternativa e aumentativa na}
\chapter{Comunicação alternativa e aumentativa na participação}
\chapter{Comunicação alternativa e aumentativa na participação social}
\chapter{Comunicação alternativa e aumentativa na participação social de}
\chapter{Comunicação alternativa e aumentativa na participação social de pessoas}
\chapter{Comunicação alternativa e aumentativa na participação social de pessoas com}
\chapter{Comunicação alternativa e aumentativa na participação social de pessoas com deficiência}
\chapter{Comunicação alternativa e aumentativa na participação social de pessoas com deficiência intelectual}
\chapter{Comunicação alternativa e aumentativa na participação social de pessoas com deficiência intelectual em}
\chapter{Comunicação alternativa e aumentativa na participação social de pessoas com deficiência intelectual em eventos}
\chapter{Comunicação alternativa e aumentativa na participação social de pessoas com deficiência intelectual em eventos públicos}
\chapter{Comunicação alternativa e aumentativa na participação social de pessoas com deficiência intelectual em eventos públicos de}
\chapter{Comunicação alternativa e aumentativa na participação social de pessoas com deficiência intelectual em eventos públicos de grande}
\chapter{Comunicação alternativa e aumentativa na participação social de pessoas com deficiência intelectual em eventos públicos de grande porte}
\chapter{Psicopedagogicamente inconstitucionalíssimo otorrinolaringologista}
\chapter{Psicopedagogicamente inconstitucionalíssimo otorrinolaringologista socioeducativo}
\chapter{Psicopedagogicamente inconstitucionalíssimo otorrinolaringologista socioeducativo ético-profissional}
\chapter{Psicopedagogicamente inconstitucionalíssimo otorrinolaringologista socioeducativo ético-profissional sociodemográfico}
\chapter{Psicopedagogicamente inconstitucionalíssimo otorrinolaringologista socioeducativo ético-profissional sociodemográfico interdisciplinaridade}
\chapter{Psicopedagogicamente inconstitucionalíssimo otorrinolaringologista socioeducativo ético-profissional sociodemográfico interdisciplinaridade neuropsicológica}
\chapter{Psicopedagogicamente inconstitucionalíssimo otorrinolaringologista socioeducativo ético-profissional sociodemográfico interdisciplinaridade neuropsicológica avaliação}
\chapter{Psicopedagogicamente inconstitucionalíssimo otorrinolaringologista socioeducativo ético-profissional sociodemográfico interdisciplinaridade neuropsicológica avaliação}
\chapter{Psicopedagogicamente inconstitucionalíssimo otorrinolaringologista socioeducativo ético-profissional sociodemográfico interdisciplinaridade neuropsicológica avaliação}
\chapter{Psicopedagogicamente inconstitucionalíssimo otorrinolaringologista socioeducativo ético-profissional sociodemográfico interdisciplinaridade neuropsicológica avaliação}
\chapter{Psicopedagogicamente inconstitucionalíssimo otorrinolaringologista socioeducativo ético-profissional sociodemográfico interdisciplinaridade neuropsicológica avaliação}
\chapter{Psicopedagogicamente inconstitucionalíssimo otorrinolaringologista socioeducativo ético-profissional sociodemográfico interdisciplinaridade neuropsicológica avaliação}
\chapter{Psicopedagogicamente inconstitucionalíssimo otorrinolaringologista socioeducativo ético-profissional sociodemográfico interdisciplinaridade neuropsicológica avaliação}
\chapter{Psicopedagogicamente inconstitucionalíssimo otorrinolaringologista socioeducativo ético-profissional sociodemográfico interdisciplinaridade neuropsicológica avaliação}
\chapter{Psicopedagogicamente inconstitucionalíssimo otorrinolaringologista socioeducativo ético-profissional sociodemográfico interdisciplinaridade neuropsicológica avaliação}
\chapter{Psicopedagogicamente inconstitucionalíssimo otorrinolaringologista socioeducativo ético-profissional sociodemográfico interdisciplinaridade neuropsicológica avaliação}
\end{document}
'   % S: o que a .sty define
% Cada variante só redefine o \l@chapter depois da carga da .sty.
% ===================================================================
\DocumentMetadata{lang=pt-BR, pdfstandard=ua-2, pdfversion=2.0, tagging=on}
\documentclass[11pt,a5paper,twoside,openright]{book}
\usepackage{livros_crp_acessivel_book}
\hypersetup{pdftitle={Sumario longo}, pdfauthor={CRP SP}}
\makeatletter
\providecommand\var{S}
\def\crp@vR{R}\def\crp@vB{B}\def\crp@vC{C}
\def\crp@lchap#1{%  #1 = acréscimo à margem direita (rightskip)
  \renewcommand*\l@chapter[2]{%
    \ifnum \c@tocdepth >\m@ne
        \addpenalty{-\@highpenalty}\addvspace{0.5em \@plus\p@}%
        \setlength\@tempdima{1.5em}%
        \begingroup
            \parindent \z@ \rightskip \@pnumwidth #1
            \parfillskip -\@pnumwidth
            \leavevmode \bfseries
            \advance\leftskip\@tempdima \hskip -\leftskip
            ##1\nobreak\hfil
            \nobreak\hb@xt@\@pnumwidth{\hss ##2\kern-\p@\kern\p@}\par
            \penalty\@highpenalty
        \endgroup
    \fi}}
\ifx\var\crp@vR \crp@lchap{}\fi
\ifx\var\crp@vB \crp@lchap{\@plus 4em}\fi
\ifx\var\crp@vC \crp@lchap{\@plus 4em \tolerance=10000 \hbadness=10000 }\fi
\makeatother
\begin{document}
\frontmatter
\tableofcontents
\mainmatter
\chapter{Deficiência psicossocial e}
\chapter{Deficiência psicossocial e transtorno}
\chapter{Deficiência psicossocial e transtorno do}
\chapter{Deficiência psicossocial e transtorno do espectro}
\chapter{Deficiência psicossocial e transtorno do espectro autista:}
\chapter{Deficiência psicossocial e transtorno do espectro autista: orientações}
\chapter{Deficiência psicossocial e transtorno do espectro autista: orientações para}
\chapter{Deficiência psicossocial e transtorno do espectro autista: orientações para a}
\chapter{Deficiência psicossocial e transtorno do espectro autista: orientações para a realização}
\chapter{Deficiência psicossocial e transtorno do espectro autista: orientações para a realização de}
\chapter{Deficiência psicossocial e transtorno do espectro autista: orientações para a realização de eventos}
\chapter{Deficiência psicossocial e transtorno do espectro autista: orientações para a realização de eventos e}
\chapter{Deficiência psicossocial e transtorno do espectro autista: orientações para a realização de eventos e apresentações}
\chapter{Deficiência psicossocial e transtorno do espectro autista: orientações para a realização de eventos e apresentações acessíveis}
\chapter{Deficiência psicossocial e transtorno do espectro autista: orientações para a realização de eventos e apresentações acessíveis em}
\chapter{Deficiência psicossocial e transtorno do espectro autista: orientações para a realização de eventos e apresentações acessíveis em todo}
\chapter{Comunicação alternativa e}
\chapter{Comunicação alternativa e aumentativa}
\chapter{Comunicação alternativa e aumentativa na}
\chapter{Comunicação alternativa e aumentativa na participação}
\chapter{Comunicação alternativa e aumentativa na participação social}
\chapter{Comunicação alternativa e aumentativa na participação social de}
\chapter{Comunicação alternativa e aumentativa na participação social de pessoas}
\chapter{Comunicação alternativa e aumentativa na participação social de pessoas com}
\chapter{Comunicação alternativa e aumentativa na participação social de pessoas com deficiência}
\chapter{Comunicação alternativa e aumentativa na participação social de pessoas com deficiência intelectual}
\chapter{Comunicação alternativa e aumentativa na participação social de pessoas com deficiência intelectual em}
\chapter{Comunicação alternativa e aumentativa na participação social de pessoas com deficiência intelectual em eventos}
\chapter{Comunicação alternativa e aumentativa na participação social de pessoas com deficiência intelectual em eventos públicos}
\chapter{Comunicação alternativa e aumentativa na participação social de pessoas com deficiência intelectual em eventos públicos de}
\chapter{Comunicação alternativa e aumentativa na participação social de pessoas com deficiência intelectual em eventos públicos de grande}
\chapter{Comunicação alternativa e aumentativa na participação social de pessoas com deficiência intelectual em eventos públicos de grande porte}
\chapter{Psicopedagogicamente inconstitucionalíssimo otorrinolaringologista}
\chapter{Psicopedagogicamente inconstitucionalíssimo otorrinolaringologista socioeducativo}
\chapter{Psicopedagogicamente inconstitucionalíssimo otorrinolaringologista socioeducativo ético-profissional}
\chapter{Psicopedagogicamente inconstitucionalíssimo otorrinolaringologista socioeducativo ético-profissional sociodemográfico}
\chapter{Psicopedagogicamente inconstitucionalíssimo otorrinolaringologista socioeducativo ético-profissional sociodemográfico interdisciplinaridade}
\chapter{Psicopedagogicamente inconstitucionalíssimo otorrinolaringologista socioeducativo ético-profissional sociodemográfico interdisciplinaridade neuropsicológica}
\chapter{Psicopedagogicamente inconstitucionalíssimo otorrinolaringologista socioeducativo ético-profissional sociodemográfico interdisciplinaridade neuropsicológica avaliação}
\chapter{Psicopedagogicamente inconstitucionalíssimo otorrinolaringologista socioeducativo ético-profissional sociodemográfico interdisciplinaridade neuropsicológica avaliação}
\chapter{Psicopedagogicamente inconstitucionalíssimo otorrinolaringologista socioeducativo ético-profissional sociodemográfico interdisciplinaridade neuropsicológica avaliação}
\chapter{Psicopedagogicamente inconstitucionalíssimo otorrinolaringologista socioeducativo ético-profissional sociodemográfico interdisciplinaridade neuropsicológica avaliação}
\chapter{Psicopedagogicamente inconstitucionalíssimo otorrinolaringologista socioeducativo ético-profissional sociodemográfico interdisciplinaridade neuropsicológica avaliação}
\chapter{Psicopedagogicamente inconstitucionalíssimo otorrinolaringologista socioeducativo ético-profissional sociodemográfico interdisciplinaridade neuropsicológica avaliação}
\chapter{Psicopedagogicamente inconstitucionalíssimo otorrinolaringologista socioeducativo ético-profissional sociodemográfico interdisciplinaridade neuropsicológica avaliação}
\chapter{Psicopedagogicamente inconstitucionalíssimo otorrinolaringologista socioeducativo ético-profissional sociodemográfico interdisciplinaridade neuropsicológica avaliação}
\chapter{Psicopedagogicamente inconstitucionalíssimo otorrinolaringologista socioeducativo ético-profissional sociodemográfico interdisciplinaridade neuropsicológica avaliação}
\chapter{Psicopedagogicamente inconstitucionalíssimo otorrinolaringologista socioeducativo ético-profissional sociodemográfico interdisciplinaridade neuropsicológica avaliação}
\end{document}
'   % S: o que a .sty define
% Cada variante só redefine o \l@chapter depois da carga da .sty.
% ===================================================================
\DocumentMetadata{lang=pt-BR, pdfstandard=ua-2, pdfversion=2.0, tagging=on}
\documentclass[11pt,a5paper,twoside,openright]{book}
\usepackage{livros_crp_acessivel_book}
\hypersetup{pdftitle={Sumario longo}, pdfauthor={CRP SP}}
\makeatletter
\providecommand\var{S}
\def\crp@vR{R}\def\crp@vB{B}\def\crp@vC{C}
\def\crp@lchap#1{%  #1 = acréscimo à margem direita (rightskip)
  \renewcommand*\l@chapter[2]{%
    \ifnum \c@tocdepth >\m@ne
        \addpenalty{-\@highpenalty}\addvspace{0.5em \@plus\p@}%
        \setlength\@tempdima{1.5em}%
        \begingroup
            \parindent \z@ \rightskip \@pnumwidth #1
            \parfillskip -\@pnumwidth
            \leavevmode \bfseries
            \advance\leftskip\@tempdima \hskip -\leftskip
            ##1\nobreak\hfil
            \nobreak\hb@xt@\@pnumwidth{\hss ##2\kern-\p@\kern\p@}\par
            \penalty\@highpenalty
        \endgroup
    \fi}}
\ifx\var\crp@vR \crp@lchap{}\fi
\ifx\var\crp@vB \crp@lchap{\@plus 4em}\fi
\ifx\var\crp@vC \crp@lchap{\@plus 4em \tolerance=10000 \hbadness=10000 }\fi
\makeatother
\begin{document}
\frontmatter
\tableofcontents
\mainmatter
\chapter{Deficiência psicossocial e}
\chapter{Deficiência psicossocial e transtorno}
\chapter{Deficiência psicossocial e transtorno do}
\chapter{Deficiência psicossocial e transtorno do espectro}
\chapter{Deficiência psicossocial e transtorno do espectro autista:}
\chapter{Deficiência psicossocial e transtorno do espectro autista: orientações}
\chapter{Deficiência psicossocial e transtorno do espectro autista: orientações para}
\chapter{Deficiência psicossocial e transtorno do espectro autista: orientações para a}
\chapter{Deficiência psicossocial e transtorno do espectro autista: orientações para a realização}
\chapter{Deficiência psicossocial e transtorno do espectro autista: orientações para a realização de}
\chapter{Deficiência psicossocial e transtorno do espectro autista: orientações para a realização de eventos}
\chapter{Deficiência psicossocial e transtorno do espectro autista: orientações para a realização de eventos e}
\chapter{Deficiência psicossocial e transtorno do espectro autista: orientações para a realização de eventos e apresentações}
\chapter{Deficiência psicossocial e transtorno do espectro autista: orientações para a realização de eventos e apresentações acessíveis}
\chapter{Deficiência psicossocial e transtorno do espectro autista: orientações para a realização de eventos e apresentações acessíveis em}
\chapter{Deficiência psicossocial e transtorno do espectro autista: orientações para a realização de eventos e apresentações acessíveis em todo}
\chapter{Comunicação alternativa e}
\chapter{Comunicação alternativa e aumentativa}
\chapter{Comunicação alternativa e aumentativa na}
\chapter{Comunicação alternativa e aumentativa na participação}
\chapter{Comunicação alternativa e aumentativa na participação social}
\chapter{Comunicação alternativa e aumentativa na participação social de}
\chapter{Comunicação alternativa e aumentativa na participação social de pessoas}
\chapter{Comunicação alternativa e aumentativa na participação social de pessoas com}
\chapter{Comunicação alternativa e aumentativa na participação social de pessoas com deficiência}
\chapter{Comunicação alternativa e aumentativa na participação social de pessoas com deficiência intelectual}
\chapter{Comunicação alternativa e aumentativa na participação social de pessoas com deficiência intelectual em}
\chapter{Comunicação alternativa e aumentativa na participação social de pessoas com deficiência intelectual em eventos}
\chapter{Comunicação alternativa e aumentativa na participação social de pessoas com deficiência intelectual em eventos públicos}
\chapter{Comunicação alternativa e aumentativa na participação social de pessoas com deficiência intelectual em eventos públicos de}
\chapter{Comunicação alternativa e aumentativa na participação social de pessoas com deficiência intelectual em eventos públicos de grande}
\chapter{Comunicação alternativa e aumentativa na participação social de pessoas com deficiência intelectual em eventos públicos de grande porte}
\chapter{Psicopedagogicamente inconstitucionalíssimo otorrinolaringologista}
\chapter{Psicopedagogicamente inconstitucionalíssimo otorrinolaringologista socioeducativo}
\chapter{Psicopedagogicamente inconstitucionalíssimo otorrinolaringologista socioeducativo ético-profissional}
\chapter{Psicopedagogicamente inconstitucionalíssimo otorrinolaringologista socioeducativo ético-profissional sociodemográfico}
\chapter{Psicopedagogicamente inconstitucionalíssimo otorrinolaringologista socioeducativo ético-profissional sociodemográfico interdisciplinaridade}
\chapter{Psicopedagogicamente inconstitucionalíssimo otorrinolaringologista socioeducativo ético-profissional sociodemográfico interdisciplinaridade neuropsicológica}
\chapter{Psicopedagogicamente inconstitucionalíssimo otorrinolaringologista socioeducativo ético-profissional sociodemográfico interdisciplinaridade neuropsicológica avaliação}
\chapter{Psicopedagogicamente inconstitucionalíssimo otorrinolaringologista socioeducativo ético-profissional sociodemográfico interdisciplinaridade neuropsicológica avaliação}
\chapter{Psicopedagogicamente inconstitucionalíssimo otorrinolaringologista socioeducativo ético-profissional sociodemográfico interdisciplinaridade neuropsicológica avaliação}
\chapter{Psicopedagogicamente inconstitucionalíssimo otorrinolaringologista socioeducativo ético-profissional sociodemográfico interdisciplinaridade neuropsicológica avaliação}
\chapter{Psicopedagogicamente inconstitucionalíssimo otorrinolaringologista socioeducativo ético-profissional sociodemográfico interdisciplinaridade neuropsicológica avaliação}
\chapter{Psicopedagogicamente inconstitucionalíssimo otorrinolaringologista socioeducativo ético-profissional sociodemográfico interdisciplinaridade neuropsicológica avaliação}
\chapter{Psicopedagogicamente inconstitucionalíssimo otorrinolaringologista socioeducativo ético-profissional sociodemográfico interdisciplinaridade neuropsicológica avaliação}
\chapter{Psicopedagogicamente inconstitucionalíssimo otorrinolaringologista socioeducativo ético-profissional sociodemográfico interdisciplinaridade neuropsicológica avaliação}
\chapter{Psicopedagogicamente inconstitucionalíssimo otorrinolaringologista socioeducativo ético-profissional sociodemográfico interdisciplinaridade neuropsicológica avaliação}
\chapter{Psicopedagogicamente inconstitucionalíssimo otorrinolaringologista socioeducativo ético-profissional sociodemográfico interdisciplinaridade neuropsicológica avaliação}
\end{document}
'   % R: rígido (até 0.5.x)
%   lualatex '\def\var{B}% !TeX program = lualatex
% ===================================================================
% sumario-longo.tex — mede o sumário sem hifenização (issue #35, "Fora do
% escopo"): com a hifenização desligada, a entrada longa de capítulo pode
% bater no número da página. Gera 3 x 16 capítulos, cada um com um prefixo
% (de 3 a 18 palavras) de uma frase longa: frase 1 de palavras médias,
% frase 2 com números, frase 3 só de palavras muito longas. Conta no log
% os "Overfull \hbox" do sumário: o sumário é a única parte do documento
% sem parágrafo de corpo, então todo Overfull "in paragraph" antes do
% primeiro \chapter é entrada que passou da margem.
%
% Uso (de uma pasta com livros_crp_acessivel_book.sty e crpsp-leg.sty no
% caminho; as fontes NEWJUNE por OSFONTDIR):
%   lualatex '\def\var{R}% !TeX program = lualatex
% ===================================================================
% sumario-longo.tex — mede o sumário sem hifenização (issue #35, "Fora do
% escopo"): com a hifenização desligada, a entrada longa de capítulo pode
% bater no número da página. Gera 3 x 16 capítulos, cada um com um prefixo
% (de 3 a 18 palavras) de uma frase longa: frase 1 de palavras médias,
% frase 2 com números, frase 3 só de palavras muito longas. Conta no log
% os "Overfull \hbox" do sumário: o sumário é a única parte do documento
% sem parágrafo de corpo, então todo Overfull "in paragraph" antes do
% primeiro \chapter é entrada que passou da margem.
%
% Uso (de uma pasta com livros_crp_acessivel_book.sty e crpsp-leg.sty no
% caminho; as fontes NEWJUNE por OSFONTDIR):
%   lualatex '\def\var{R}% !TeX program = lualatex
% ===================================================================
% sumario-longo.tex — mede o sumário sem hifenização (issue #35, "Fora do
% escopo"): com a hifenização desligada, a entrada longa de capítulo pode
% bater no número da página. Gera 3 x 16 capítulos, cada um com um prefixo
% (de 3 a 18 palavras) de uma frase longa: frase 1 de palavras médias,
% frase 2 com números, frase 3 só de palavras muito longas. Conta no log
% os "Overfull \hbox" do sumário: o sumário é a única parte do documento
% sem parágrafo de corpo, então todo Overfull "in paragraph" antes do
% primeiro \chapter é entrada que passou da margem.
%
% Uso (de uma pasta com livros_crp_acessivel_book.sty e crpsp-leg.sty no
% caminho; as fontes NEWJUNE por OSFONTDIR):
%   lualatex '\def\var{R}\input{sumario-longo}'   % R: rígido (até 0.5.x)
%   lualatex '\def\var{B}\input{sumario-longo}'   % B: plus 4em (o do guia)
%   lualatex '\def\var{C}\input{sumario-longo}'   % C: plus 4em + \tolerance 10000
%   lualatex '\def\var{S}\input{sumario-longo}'   % S: o que a .sty define
% Cada variante só redefine o \l@chapter depois da carga da .sty.
% ===================================================================
\DocumentMetadata{lang=pt-BR, pdfstandard=ua-2, pdfversion=2.0, tagging=on}
\documentclass[11pt,a5paper,twoside,openright]{book}
\usepackage{livros_crp_acessivel_book}
\hypersetup{pdftitle={Sumario longo}, pdfauthor={CRP SP}}
\makeatletter
\providecommand\var{S}
\def\crp@vR{R}\def\crp@vB{B}\def\crp@vC{C}
\def\crp@lchap#1{%  #1 = acréscimo à margem direita (rightskip)
  \renewcommand*\l@chapter[2]{%
    \ifnum \c@tocdepth >\m@ne
        \addpenalty{-\@highpenalty}\addvspace{0.5em \@plus\p@}%
        \setlength\@tempdima{1.5em}%
        \begingroup
            \parindent \z@ \rightskip \@pnumwidth #1
            \parfillskip -\@pnumwidth
            \leavevmode \bfseries
            \advance\leftskip\@tempdima \hskip -\leftskip
            ##1\nobreak\hfil
            \nobreak\hb@xt@\@pnumwidth{\hss ##2\kern-\p@\kern\p@}\par
            \penalty\@highpenalty
        \endgroup
    \fi}}
\ifx\var\crp@vR \crp@lchap{}\fi
\ifx\var\crp@vB \crp@lchap{\@plus 4em}\fi
\ifx\var\crp@vC \crp@lchap{\@plus 4em \tolerance=10000 \hbadness=10000 }\fi
\makeatother
\begin{document}
\frontmatter
\tableofcontents
\mainmatter
\chapter{Deficiência psicossocial e}
\chapter{Deficiência psicossocial e transtorno}
\chapter{Deficiência psicossocial e transtorno do}
\chapter{Deficiência psicossocial e transtorno do espectro}
\chapter{Deficiência psicossocial e transtorno do espectro autista:}
\chapter{Deficiência psicossocial e transtorno do espectro autista: orientações}
\chapter{Deficiência psicossocial e transtorno do espectro autista: orientações para}
\chapter{Deficiência psicossocial e transtorno do espectro autista: orientações para a}
\chapter{Deficiência psicossocial e transtorno do espectro autista: orientações para a realização}
\chapter{Deficiência psicossocial e transtorno do espectro autista: orientações para a realização de}
\chapter{Deficiência psicossocial e transtorno do espectro autista: orientações para a realização de eventos}
\chapter{Deficiência psicossocial e transtorno do espectro autista: orientações para a realização de eventos e}
\chapter{Deficiência psicossocial e transtorno do espectro autista: orientações para a realização de eventos e apresentações}
\chapter{Deficiência psicossocial e transtorno do espectro autista: orientações para a realização de eventos e apresentações acessíveis}
\chapter{Deficiência psicossocial e transtorno do espectro autista: orientações para a realização de eventos e apresentações acessíveis em}
\chapter{Deficiência psicossocial e transtorno do espectro autista: orientações para a realização de eventos e apresentações acessíveis em todo}
\chapter{Comunicação alternativa e}
\chapter{Comunicação alternativa e aumentativa}
\chapter{Comunicação alternativa e aumentativa na}
\chapter{Comunicação alternativa e aumentativa na participação}
\chapter{Comunicação alternativa e aumentativa na participação social}
\chapter{Comunicação alternativa e aumentativa na participação social de}
\chapter{Comunicação alternativa e aumentativa na participação social de pessoas}
\chapter{Comunicação alternativa e aumentativa na participação social de pessoas com}
\chapter{Comunicação alternativa e aumentativa na participação social de pessoas com deficiência}
\chapter{Comunicação alternativa e aumentativa na participação social de pessoas com deficiência intelectual}
\chapter{Comunicação alternativa e aumentativa na participação social de pessoas com deficiência intelectual em}
\chapter{Comunicação alternativa e aumentativa na participação social de pessoas com deficiência intelectual em eventos}
\chapter{Comunicação alternativa e aumentativa na participação social de pessoas com deficiência intelectual em eventos públicos}
\chapter{Comunicação alternativa e aumentativa na participação social de pessoas com deficiência intelectual em eventos públicos de}
\chapter{Comunicação alternativa e aumentativa na participação social de pessoas com deficiência intelectual em eventos públicos de grande}
\chapter{Comunicação alternativa e aumentativa na participação social de pessoas com deficiência intelectual em eventos públicos de grande porte}
\chapter{Psicopedagogicamente inconstitucionalíssimo otorrinolaringologista}
\chapter{Psicopedagogicamente inconstitucionalíssimo otorrinolaringologista socioeducativo}
\chapter{Psicopedagogicamente inconstitucionalíssimo otorrinolaringologista socioeducativo ético-profissional}
\chapter{Psicopedagogicamente inconstitucionalíssimo otorrinolaringologista socioeducativo ético-profissional sociodemográfico}
\chapter{Psicopedagogicamente inconstitucionalíssimo otorrinolaringologista socioeducativo ético-profissional sociodemográfico interdisciplinaridade}
\chapter{Psicopedagogicamente inconstitucionalíssimo otorrinolaringologista socioeducativo ético-profissional sociodemográfico interdisciplinaridade neuropsicológica}
\chapter{Psicopedagogicamente inconstitucionalíssimo otorrinolaringologista socioeducativo ético-profissional sociodemográfico interdisciplinaridade neuropsicológica avaliação}
\chapter{Psicopedagogicamente inconstitucionalíssimo otorrinolaringologista socioeducativo ético-profissional sociodemográfico interdisciplinaridade neuropsicológica avaliação}
\chapter{Psicopedagogicamente inconstitucionalíssimo otorrinolaringologista socioeducativo ético-profissional sociodemográfico interdisciplinaridade neuropsicológica avaliação}
\chapter{Psicopedagogicamente inconstitucionalíssimo otorrinolaringologista socioeducativo ético-profissional sociodemográfico interdisciplinaridade neuropsicológica avaliação}
\chapter{Psicopedagogicamente inconstitucionalíssimo otorrinolaringologista socioeducativo ético-profissional sociodemográfico interdisciplinaridade neuropsicológica avaliação}
\chapter{Psicopedagogicamente inconstitucionalíssimo otorrinolaringologista socioeducativo ético-profissional sociodemográfico interdisciplinaridade neuropsicológica avaliação}
\chapter{Psicopedagogicamente inconstitucionalíssimo otorrinolaringologista socioeducativo ético-profissional sociodemográfico interdisciplinaridade neuropsicológica avaliação}
\chapter{Psicopedagogicamente inconstitucionalíssimo otorrinolaringologista socioeducativo ético-profissional sociodemográfico interdisciplinaridade neuropsicológica avaliação}
\chapter{Psicopedagogicamente inconstitucionalíssimo otorrinolaringologista socioeducativo ético-profissional sociodemográfico interdisciplinaridade neuropsicológica avaliação}
\chapter{Psicopedagogicamente inconstitucionalíssimo otorrinolaringologista socioeducativo ético-profissional sociodemográfico interdisciplinaridade neuropsicológica avaliação}
\end{document}
'   % R: rígido (até 0.5.x)
%   lualatex '\def\var{B}% !TeX program = lualatex
% ===================================================================
% sumario-longo.tex — mede o sumário sem hifenização (issue #35, "Fora do
% escopo"): com a hifenização desligada, a entrada longa de capítulo pode
% bater no número da página. Gera 3 x 16 capítulos, cada um com um prefixo
% (de 3 a 18 palavras) de uma frase longa: frase 1 de palavras médias,
% frase 2 com números, frase 3 só de palavras muito longas. Conta no log
% os "Overfull \hbox" do sumário: o sumário é a única parte do documento
% sem parágrafo de corpo, então todo Overfull "in paragraph" antes do
% primeiro \chapter é entrada que passou da margem.
%
% Uso (de uma pasta com livros_crp_acessivel_book.sty e crpsp-leg.sty no
% caminho; as fontes NEWJUNE por OSFONTDIR):
%   lualatex '\def\var{R}\input{sumario-longo}'   % R: rígido (até 0.5.x)
%   lualatex '\def\var{B}\input{sumario-longo}'   % B: plus 4em (o do guia)
%   lualatex '\def\var{C}\input{sumario-longo}'   % C: plus 4em + \tolerance 10000
%   lualatex '\def\var{S}\input{sumario-longo}'   % S: o que a .sty define
% Cada variante só redefine o \l@chapter depois da carga da .sty.
% ===================================================================
\DocumentMetadata{lang=pt-BR, pdfstandard=ua-2, pdfversion=2.0, tagging=on}
\documentclass[11pt,a5paper,twoside,openright]{book}
\usepackage{livros_crp_acessivel_book}
\hypersetup{pdftitle={Sumario longo}, pdfauthor={CRP SP}}
\makeatletter
\providecommand\var{S}
\def\crp@vR{R}\def\crp@vB{B}\def\crp@vC{C}
\def\crp@lchap#1{%  #1 = acréscimo à margem direita (rightskip)
  \renewcommand*\l@chapter[2]{%
    \ifnum \c@tocdepth >\m@ne
        \addpenalty{-\@highpenalty}\addvspace{0.5em \@plus\p@}%
        \setlength\@tempdima{1.5em}%
        \begingroup
            \parindent \z@ \rightskip \@pnumwidth #1
            \parfillskip -\@pnumwidth
            \leavevmode \bfseries
            \advance\leftskip\@tempdima \hskip -\leftskip
            ##1\nobreak\hfil
            \nobreak\hb@xt@\@pnumwidth{\hss ##2\kern-\p@\kern\p@}\par
            \penalty\@highpenalty
        \endgroup
    \fi}}
\ifx\var\crp@vR \crp@lchap{}\fi
\ifx\var\crp@vB \crp@lchap{\@plus 4em}\fi
\ifx\var\crp@vC \crp@lchap{\@plus 4em \tolerance=10000 \hbadness=10000 }\fi
\makeatother
\begin{document}
\frontmatter
\tableofcontents
\mainmatter
\chapter{Deficiência psicossocial e}
\chapter{Deficiência psicossocial e transtorno}
\chapter{Deficiência psicossocial e transtorno do}
\chapter{Deficiência psicossocial e transtorno do espectro}
\chapter{Deficiência psicossocial e transtorno do espectro autista:}
\chapter{Deficiência psicossocial e transtorno do espectro autista: orientações}
\chapter{Deficiência psicossocial e transtorno do espectro autista: orientações para}
\chapter{Deficiência psicossocial e transtorno do espectro autista: orientações para a}
\chapter{Deficiência psicossocial e transtorno do espectro autista: orientações para a realização}
\chapter{Deficiência psicossocial e transtorno do espectro autista: orientações para a realização de}
\chapter{Deficiência psicossocial e transtorno do espectro autista: orientações para a realização de eventos}
\chapter{Deficiência psicossocial e transtorno do espectro autista: orientações para a realização de eventos e}
\chapter{Deficiência psicossocial e transtorno do espectro autista: orientações para a realização de eventos e apresentações}
\chapter{Deficiência psicossocial e transtorno do espectro autista: orientações para a realização de eventos e apresentações acessíveis}
\chapter{Deficiência psicossocial e transtorno do espectro autista: orientações para a realização de eventos e apresentações acessíveis em}
\chapter{Deficiência psicossocial e transtorno do espectro autista: orientações para a realização de eventos e apresentações acessíveis em todo}
\chapter{Comunicação alternativa e}
\chapter{Comunicação alternativa e aumentativa}
\chapter{Comunicação alternativa e aumentativa na}
\chapter{Comunicação alternativa e aumentativa na participação}
\chapter{Comunicação alternativa e aumentativa na participação social}
\chapter{Comunicação alternativa e aumentativa na participação social de}
\chapter{Comunicação alternativa e aumentativa na participação social de pessoas}
\chapter{Comunicação alternativa e aumentativa na participação social de pessoas com}
\chapter{Comunicação alternativa e aumentativa na participação social de pessoas com deficiência}
\chapter{Comunicação alternativa e aumentativa na participação social de pessoas com deficiência intelectual}
\chapter{Comunicação alternativa e aumentativa na participação social de pessoas com deficiência intelectual em}
\chapter{Comunicação alternativa e aumentativa na participação social de pessoas com deficiência intelectual em eventos}
\chapter{Comunicação alternativa e aumentativa na participação social de pessoas com deficiência intelectual em eventos públicos}
\chapter{Comunicação alternativa e aumentativa na participação social de pessoas com deficiência intelectual em eventos públicos de}
\chapter{Comunicação alternativa e aumentativa na participação social de pessoas com deficiência intelectual em eventos públicos de grande}
\chapter{Comunicação alternativa e aumentativa na participação social de pessoas com deficiência intelectual em eventos públicos de grande porte}
\chapter{Psicopedagogicamente inconstitucionalíssimo otorrinolaringologista}
\chapter{Psicopedagogicamente inconstitucionalíssimo otorrinolaringologista socioeducativo}
\chapter{Psicopedagogicamente inconstitucionalíssimo otorrinolaringologista socioeducativo ético-profissional}
\chapter{Psicopedagogicamente inconstitucionalíssimo otorrinolaringologista socioeducativo ético-profissional sociodemográfico}
\chapter{Psicopedagogicamente inconstitucionalíssimo otorrinolaringologista socioeducativo ético-profissional sociodemográfico interdisciplinaridade}
\chapter{Psicopedagogicamente inconstitucionalíssimo otorrinolaringologista socioeducativo ético-profissional sociodemográfico interdisciplinaridade neuropsicológica}
\chapter{Psicopedagogicamente inconstitucionalíssimo otorrinolaringologista socioeducativo ético-profissional sociodemográfico interdisciplinaridade neuropsicológica avaliação}
\chapter{Psicopedagogicamente inconstitucionalíssimo otorrinolaringologista socioeducativo ético-profissional sociodemográfico interdisciplinaridade neuropsicológica avaliação}
\chapter{Psicopedagogicamente inconstitucionalíssimo otorrinolaringologista socioeducativo ético-profissional sociodemográfico interdisciplinaridade neuropsicológica avaliação}
\chapter{Psicopedagogicamente inconstitucionalíssimo otorrinolaringologista socioeducativo ético-profissional sociodemográfico interdisciplinaridade neuropsicológica avaliação}
\chapter{Psicopedagogicamente inconstitucionalíssimo otorrinolaringologista socioeducativo ético-profissional sociodemográfico interdisciplinaridade neuropsicológica avaliação}
\chapter{Psicopedagogicamente inconstitucionalíssimo otorrinolaringologista socioeducativo ético-profissional sociodemográfico interdisciplinaridade neuropsicológica avaliação}
\chapter{Psicopedagogicamente inconstitucionalíssimo otorrinolaringologista socioeducativo ético-profissional sociodemográfico interdisciplinaridade neuropsicológica avaliação}
\chapter{Psicopedagogicamente inconstitucionalíssimo otorrinolaringologista socioeducativo ético-profissional sociodemográfico interdisciplinaridade neuropsicológica avaliação}
\chapter{Psicopedagogicamente inconstitucionalíssimo otorrinolaringologista socioeducativo ético-profissional sociodemográfico interdisciplinaridade neuropsicológica avaliação}
\chapter{Psicopedagogicamente inconstitucionalíssimo otorrinolaringologista socioeducativo ético-profissional sociodemográfico interdisciplinaridade neuropsicológica avaliação}
\end{document}
'   % B: plus 4em (o do guia)
%   lualatex '\def\var{C}% !TeX program = lualatex
% ===================================================================
% sumario-longo.tex — mede o sumário sem hifenização (issue #35, "Fora do
% escopo"): com a hifenização desligada, a entrada longa de capítulo pode
% bater no número da página. Gera 3 x 16 capítulos, cada um com um prefixo
% (de 3 a 18 palavras) de uma frase longa: frase 1 de palavras médias,
% frase 2 com números, frase 3 só de palavras muito longas. Conta no log
% os "Overfull \hbox" do sumário: o sumário é a única parte do documento
% sem parágrafo de corpo, então todo Overfull "in paragraph" antes do
% primeiro \chapter é entrada que passou da margem.
%
% Uso (de uma pasta com livros_crp_acessivel_book.sty e crpsp-leg.sty no
% caminho; as fontes NEWJUNE por OSFONTDIR):
%   lualatex '\def\var{R}\input{sumario-longo}'   % R: rígido (até 0.5.x)
%   lualatex '\def\var{B}\input{sumario-longo}'   % B: plus 4em (o do guia)
%   lualatex '\def\var{C}\input{sumario-longo}'   % C: plus 4em + \tolerance 10000
%   lualatex '\def\var{S}\input{sumario-longo}'   % S: o que a .sty define
% Cada variante só redefine o \l@chapter depois da carga da .sty.
% ===================================================================
\DocumentMetadata{lang=pt-BR, pdfstandard=ua-2, pdfversion=2.0, tagging=on}
\documentclass[11pt,a5paper,twoside,openright]{book}
\usepackage{livros_crp_acessivel_book}
\hypersetup{pdftitle={Sumario longo}, pdfauthor={CRP SP}}
\makeatletter
\providecommand\var{S}
\def\crp@vR{R}\def\crp@vB{B}\def\crp@vC{C}
\def\crp@lchap#1{%  #1 = acréscimo à margem direita (rightskip)
  \renewcommand*\l@chapter[2]{%
    \ifnum \c@tocdepth >\m@ne
        \addpenalty{-\@highpenalty}\addvspace{0.5em \@plus\p@}%
        \setlength\@tempdima{1.5em}%
        \begingroup
            \parindent \z@ \rightskip \@pnumwidth #1
            \parfillskip -\@pnumwidth
            \leavevmode \bfseries
            \advance\leftskip\@tempdima \hskip -\leftskip
            ##1\nobreak\hfil
            \nobreak\hb@xt@\@pnumwidth{\hss ##2\kern-\p@\kern\p@}\par
            \penalty\@highpenalty
        \endgroup
    \fi}}
\ifx\var\crp@vR \crp@lchap{}\fi
\ifx\var\crp@vB \crp@lchap{\@plus 4em}\fi
\ifx\var\crp@vC \crp@lchap{\@plus 4em \tolerance=10000 \hbadness=10000 }\fi
\makeatother
\begin{document}
\frontmatter
\tableofcontents
\mainmatter
\chapter{Deficiência psicossocial e}
\chapter{Deficiência psicossocial e transtorno}
\chapter{Deficiência psicossocial e transtorno do}
\chapter{Deficiência psicossocial e transtorno do espectro}
\chapter{Deficiência psicossocial e transtorno do espectro autista:}
\chapter{Deficiência psicossocial e transtorno do espectro autista: orientações}
\chapter{Deficiência psicossocial e transtorno do espectro autista: orientações para}
\chapter{Deficiência psicossocial e transtorno do espectro autista: orientações para a}
\chapter{Deficiência psicossocial e transtorno do espectro autista: orientações para a realização}
\chapter{Deficiência psicossocial e transtorno do espectro autista: orientações para a realização de}
\chapter{Deficiência psicossocial e transtorno do espectro autista: orientações para a realização de eventos}
\chapter{Deficiência psicossocial e transtorno do espectro autista: orientações para a realização de eventos e}
\chapter{Deficiência psicossocial e transtorno do espectro autista: orientações para a realização de eventos e apresentações}
\chapter{Deficiência psicossocial e transtorno do espectro autista: orientações para a realização de eventos e apresentações acessíveis}
\chapter{Deficiência psicossocial e transtorno do espectro autista: orientações para a realização de eventos e apresentações acessíveis em}
\chapter{Deficiência psicossocial e transtorno do espectro autista: orientações para a realização de eventos e apresentações acessíveis em todo}
\chapter{Comunicação alternativa e}
\chapter{Comunicação alternativa e aumentativa}
\chapter{Comunicação alternativa e aumentativa na}
\chapter{Comunicação alternativa e aumentativa na participação}
\chapter{Comunicação alternativa e aumentativa na participação social}
\chapter{Comunicação alternativa e aumentativa na participação social de}
\chapter{Comunicação alternativa e aumentativa na participação social de pessoas}
\chapter{Comunicação alternativa e aumentativa na participação social de pessoas com}
\chapter{Comunicação alternativa e aumentativa na participação social de pessoas com deficiência}
\chapter{Comunicação alternativa e aumentativa na participação social de pessoas com deficiência intelectual}
\chapter{Comunicação alternativa e aumentativa na participação social de pessoas com deficiência intelectual em}
\chapter{Comunicação alternativa e aumentativa na participação social de pessoas com deficiência intelectual em eventos}
\chapter{Comunicação alternativa e aumentativa na participação social de pessoas com deficiência intelectual em eventos públicos}
\chapter{Comunicação alternativa e aumentativa na participação social de pessoas com deficiência intelectual em eventos públicos de}
\chapter{Comunicação alternativa e aumentativa na participação social de pessoas com deficiência intelectual em eventos públicos de grande}
\chapter{Comunicação alternativa e aumentativa na participação social de pessoas com deficiência intelectual em eventos públicos de grande porte}
\chapter{Psicopedagogicamente inconstitucionalíssimo otorrinolaringologista}
\chapter{Psicopedagogicamente inconstitucionalíssimo otorrinolaringologista socioeducativo}
\chapter{Psicopedagogicamente inconstitucionalíssimo otorrinolaringologista socioeducativo ético-profissional}
\chapter{Psicopedagogicamente inconstitucionalíssimo otorrinolaringologista socioeducativo ético-profissional sociodemográfico}
\chapter{Psicopedagogicamente inconstitucionalíssimo otorrinolaringologista socioeducativo ético-profissional sociodemográfico interdisciplinaridade}
\chapter{Psicopedagogicamente inconstitucionalíssimo otorrinolaringologista socioeducativo ético-profissional sociodemográfico interdisciplinaridade neuropsicológica}
\chapter{Psicopedagogicamente inconstitucionalíssimo otorrinolaringologista socioeducativo ético-profissional sociodemográfico interdisciplinaridade neuropsicológica avaliação}
\chapter{Psicopedagogicamente inconstitucionalíssimo otorrinolaringologista socioeducativo ético-profissional sociodemográfico interdisciplinaridade neuropsicológica avaliação}
\chapter{Psicopedagogicamente inconstitucionalíssimo otorrinolaringologista socioeducativo ético-profissional sociodemográfico interdisciplinaridade neuropsicológica avaliação}
\chapter{Psicopedagogicamente inconstitucionalíssimo otorrinolaringologista socioeducativo ético-profissional sociodemográfico interdisciplinaridade neuropsicológica avaliação}
\chapter{Psicopedagogicamente inconstitucionalíssimo otorrinolaringologista socioeducativo ético-profissional sociodemográfico interdisciplinaridade neuropsicológica avaliação}
\chapter{Psicopedagogicamente inconstitucionalíssimo otorrinolaringologista socioeducativo ético-profissional sociodemográfico interdisciplinaridade neuropsicológica avaliação}
\chapter{Psicopedagogicamente inconstitucionalíssimo otorrinolaringologista socioeducativo ético-profissional sociodemográfico interdisciplinaridade neuropsicológica avaliação}
\chapter{Psicopedagogicamente inconstitucionalíssimo otorrinolaringologista socioeducativo ético-profissional sociodemográfico interdisciplinaridade neuropsicológica avaliação}
\chapter{Psicopedagogicamente inconstitucionalíssimo otorrinolaringologista socioeducativo ético-profissional sociodemográfico interdisciplinaridade neuropsicológica avaliação}
\chapter{Psicopedagogicamente inconstitucionalíssimo otorrinolaringologista socioeducativo ético-profissional sociodemográfico interdisciplinaridade neuropsicológica avaliação}
\end{document}
'   % C: plus 4em + \tolerance 10000
%   lualatex '\def\var{S}% !TeX program = lualatex
% ===================================================================
% sumario-longo.tex — mede o sumário sem hifenização (issue #35, "Fora do
% escopo"): com a hifenização desligada, a entrada longa de capítulo pode
% bater no número da página. Gera 3 x 16 capítulos, cada um com um prefixo
% (de 3 a 18 palavras) de uma frase longa: frase 1 de palavras médias,
% frase 2 com números, frase 3 só de palavras muito longas. Conta no log
% os "Overfull \hbox" do sumário: o sumário é a única parte do documento
% sem parágrafo de corpo, então todo Overfull "in paragraph" antes do
% primeiro \chapter é entrada que passou da margem.
%
% Uso (de uma pasta com livros_crp_acessivel_book.sty e crpsp-leg.sty no
% caminho; as fontes NEWJUNE por OSFONTDIR):
%   lualatex '\def\var{R}\input{sumario-longo}'   % R: rígido (até 0.5.x)
%   lualatex '\def\var{B}\input{sumario-longo}'   % B: plus 4em (o do guia)
%   lualatex '\def\var{C}\input{sumario-longo}'   % C: plus 4em + \tolerance 10000
%   lualatex '\def\var{S}\input{sumario-longo}'   % S: o que a .sty define
% Cada variante só redefine o \l@chapter depois da carga da .sty.
% ===================================================================
\DocumentMetadata{lang=pt-BR, pdfstandard=ua-2, pdfversion=2.0, tagging=on}
\documentclass[11pt,a5paper,twoside,openright]{book}
\usepackage{livros_crp_acessivel_book}
\hypersetup{pdftitle={Sumario longo}, pdfauthor={CRP SP}}
\makeatletter
\providecommand\var{S}
\def\crp@vR{R}\def\crp@vB{B}\def\crp@vC{C}
\def\crp@lchap#1{%  #1 = acréscimo à margem direita (rightskip)
  \renewcommand*\l@chapter[2]{%
    \ifnum \c@tocdepth >\m@ne
        \addpenalty{-\@highpenalty}\addvspace{0.5em \@plus\p@}%
        \setlength\@tempdima{1.5em}%
        \begingroup
            \parindent \z@ \rightskip \@pnumwidth #1
            \parfillskip -\@pnumwidth
            \leavevmode \bfseries
            \advance\leftskip\@tempdima \hskip -\leftskip
            ##1\nobreak\hfil
            \nobreak\hb@xt@\@pnumwidth{\hss ##2\kern-\p@\kern\p@}\par
            \penalty\@highpenalty
        \endgroup
    \fi}}
\ifx\var\crp@vR \crp@lchap{}\fi
\ifx\var\crp@vB \crp@lchap{\@plus 4em}\fi
\ifx\var\crp@vC \crp@lchap{\@plus 4em \tolerance=10000 \hbadness=10000 }\fi
\makeatother
\begin{document}
\frontmatter
\tableofcontents
\mainmatter
\chapter{Deficiência psicossocial e}
\chapter{Deficiência psicossocial e transtorno}
\chapter{Deficiência psicossocial e transtorno do}
\chapter{Deficiência psicossocial e transtorno do espectro}
\chapter{Deficiência psicossocial e transtorno do espectro autista:}
\chapter{Deficiência psicossocial e transtorno do espectro autista: orientações}
\chapter{Deficiência psicossocial e transtorno do espectro autista: orientações para}
\chapter{Deficiência psicossocial e transtorno do espectro autista: orientações para a}
\chapter{Deficiência psicossocial e transtorno do espectro autista: orientações para a realização}
\chapter{Deficiência psicossocial e transtorno do espectro autista: orientações para a realização de}
\chapter{Deficiência psicossocial e transtorno do espectro autista: orientações para a realização de eventos}
\chapter{Deficiência psicossocial e transtorno do espectro autista: orientações para a realização de eventos e}
\chapter{Deficiência psicossocial e transtorno do espectro autista: orientações para a realização de eventos e apresentações}
\chapter{Deficiência psicossocial e transtorno do espectro autista: orientações para a realização de eventos e apresentações acessíveis}
\chapter{Deficiência psicossocial e transtorno do espectro autista: orientações para a realização de eventos e apresentações acessíveis em}
\chapter{Deficiência psicossocial e transtorno do espectro autista: orientações para a realização de eventos e apresentações acessíveis em todo}
\chapter{Comunicação alternativa e}
\chapter{Comunicação alternativa e aumentativa}
\chapter{Comunicação alternativa e aumentativa na}
\chapter{Comunicação alternativa e aumentativa na participação}
\chapter{Comunicação alternativa e aumentativa na participação social}
\chapter{Comunicação alternativa e aumentativa na participação social de}
\chapter{Comunicação alternativa e aumentativa na participação social de pessoas}
\chapter{Comunicação alternativa e aumentativa na participação social de pessoas com}
\chapter{Comunicação alternativa e aumentativa na participação social de pessoas com deficiência}
\chapter{Comunicação alternativa e aumentativa na participação social de pessoas com deficiência intelectual}
\chapter{Comunicação alternativa e aumentativa na participação social de pessoas com deficiência intelectual em}
\chapter{Comunicação alternativa e aumentativa na participação social de pessoas com deficiência intelectual em eventos}
\chapter{Comunicação alternativa e aumentativa na participação social de pessoas com deficiência intelectual em eventos públicos}
\chapter{Comunicação alternativa e aumentativa na participação social de pessoas com deficiência intelectual em eventos públicos de}
\chapter{Comunicação alternativa e aumentativa na participação social de pessoas com deficiência intelectual em eventos públicos de grande}
\chapter{Comunicação alternativa e aumentativa na participação social de pessoas com deficiência intelectual em eventos públicos de grande porte}
\chapter{Psicopedagogicamente inconstitucionalíssimo otorrinolaringologista}
\chapter{Psicopedagogicamente inconstitucionalíssimo otorrinolaringologista socioeducativo}
\chapter{Psicopedagogicamente inconstitucionalíssimo otorrinolaringologista socioeducativo ético-profissional}
\chapter{Psicopedagogicamente inconstitucionalíssimo otorrinolaringologista socioeducativo ético-profissional sociodemográfico}
\chapter{Psicopedagogicamente inconstitucionalíssimo otorrinolaringologista socioeducativo ético-profissional sociodemográfico interdisciplinaridade}
\chapter{Psicopedagogicamente inconstitucionalíssimo otorrinolaringologista socioeducativo ético-profissional sociodemográfico interdisciplinaridade neuropsicológica}
\chapter{Psicopedagogicamente inconstitucionalíssimo otorrinolaringologista socioeducativo ético-profissional sociodemográfico interdisciplinaridade neuropsicológica avaliação}
\chapter{Psicopedagogicamente inconstitucionalíssimo otorrinolaringologista socioeducativo ético-profissional sociodemográfico interdisciplinaridade neuropsicológica avaliação}
\chapter{Psicopedagogicamente inconstitucionalíssimo otorrinolaringologista socioeducativo ético-profissional sociodemográfico interdisciplinaridade neuropsicológica avaliação}
\chapter{Psicopedagogicamente inconstitucionalíssimo otorrinolaringologista socioeducativo ético-profissional sociodemográfico interdisciplinaridade neuropsicológica avaliação}
\chapter{Psicopedagogicamente inconstitucionalíssimo otorrinolaringologista socioeducativo ético-profissional sociodemográfico interdisciplinaridade neuropsicológica avaliação}
\chapter{Psicopedagogicamente inconstitucionalíssimo otorrinolaringologista socioeducativo ético-profissional sociodemográfico interdisciplinaridade neuropsicológica avaliação}
\chapter{Psicopedagogicamente inconstitucionalíssimo otorrinolaringologista socioeducativo ético-profissional sociodemográfico interdisciplinaridade neuropsicológica avaliação}
\chapter{Psicopedagogicamente inconstitucionalíssimo otorrinolaringologista socioeducativo ético-profissional sociodemográfico interdisciplinaridade neuropsicológica avaliação}
\chapter{Psicopedagogicamente inconstitucionalíssimo otorrinolaringologista socioeducativo ético-profissional sociodemográfico interdisciplinaridade neuropsicológica avaliação}
\chapter{Psicopedagogicamente inconstitucionalíssimo otorrinolaringologista socioeducativo ético-profissional sociodemográfico interdisciplinaridade neuropsicológica avaliação}
\end{document}
'   % S: o que a .sty define
% Cada variante só redefine o \l@chapter depois da carga da .sty.
% ===================================================================
\DocumentMetadata{lang=pt-BR, pdfstandard=ua-2, pdfversion=2.0, tagging=on}
\documentclass[11pt,a5paper,twoside,openright]{book}
\usepackage{livros_crp_acessivel_book}
\hypersetup{pdftitle={Sumario longo}, pdfauthor={CRP SP}}
\makeatletter
\providecommand\var{S}
\def\crp@vR{R}\def\crp@vB{B}\def\crp@vC{C}
\def\crp@lchap#1{%  #1 = acréscimo à margem direita (rightskip)
  \renewcommand*\l@chapter[2]{%
    \ifnum \c@tocdepth >\m@ne
        \addpenalty{-\@highpenalty}\addvspace{0.5em \@plus\p@}%
        \setlength\@tempdima{1.5em}%
        \begingroup
            \parindent \z@ \rightskip \@pnumwidth #1
            \parfillskip -\@pnumwidth
            \leavevmode \bfseries
            \advance\leftskip\@tempdima \hskip -\leftskip
            ##1\nobreak\hfil
            \nobreak\hb@xt@\@pnumwidth{\hss ##2\kern-\p@\kern\p@}\par
            \penalty\@highpenalty
        \endgroup
    \fi}}
\ifx\var\crp@vR \crp@lchap{}\fi
\ifx\var\crp@vB \crp@lchap{\@plus 4em}\fi
\ifx\var\crp@vC \crp@lchap{\@plus 4em \tolerance=10000 \hbadness=10000 }\fi
\makeatother
\begin{document}
\frontmatter
\tableofcontents
\mainmatter
\chapter{Deficiência psicossocial e}
\chapter{Deficiência psicossocial e transtorno}
\chapter{Deficiência psicossocial e transtorno do}
\chapter{Deficiência psicossocial e transtorno do espectro}
\chapter{Deficiência psicossocial e transtorno do espectro autista:}
\chapter{Deficiência psicossocial e transtorno do espectro autista: orientações}
\chapter{Deficiência psicossocial e transtorno do espectro autista: orientações para}
\chapter{Deficiência psicossocial e transtorno do espectro autista: orientações para a}
\chapter{Deficiência psicossocial e transtorno do espectro autista: orientações para a realização}
\chapter{Deficiência psicossocial e transtorno do espectro autista: orientações para a realização de}
\chapter{Deficiência psicossocial e transtorno do espectro autista: orientações para a realização de eventos}
\chapter{Deficiência psicossocial e transtorno do espectro autista: orientações para a realização de eventos e}
\chapter{Deficiência psicossocial e transtorno do espectro autista: orientações para a realização de eventos e apresentações}
\chapter{Deficiência psicossocial e transtorno do espectro autista: orientações para a realização de eventos e apresentações acessíveis}
\chapter{Deficiência psicossocial e transtorno do espectro autista: orientações para a realização de eventos e apresentações acessíveis em}
\chapter{Deficiência psicossocial e transtorno do espectro autista: orientações para a realização de eventos e apresentações acessíveis em todo}
\chapter{Comunicação alternativa e}
\chapter{Comunicação alternativa e aumentativa}
\chapter{Comunicação alternativa e aumentativa na}
\chapter{Comunicação alternativa e aumentativa na participação}
\chapter{Comunicação alternativa e aumentativa na participação social}
\chapter{Comunicação alternativa e aumentativa na participação social de}
\chapter{Comunicação alternativa e aumentativa na participação social de pessoas}
\chapter{Comunicação alternativa e aumentativa na participação social de pessoas com}
\chapter{Comunicação alternativa e aumentativa na participação social de pessoas com deficiência}
\chapter{Comunicação alternativa e aumentativa na participação social de pessoas com deficiência intelectual}
\chapter{Comunicação alternativa e aumentativa na participação social de pessoas com deficiência intelectual em}
\chapter{Comunicação alternativa e aumentativa na participação social de pessoas com deficiência intelectual em eventos}
\chapter{Comunicação alternativa e aumentativa na participação social de pessoas com deficiência intelectual em eventos públicos}
\chapter{Comunicação alternativa e aumentativa na participação social de pessoas com deficiência intelectual em eventos públicos de}
\chapter{Comunicação alternativa e aumentativa na participação social de pessoas com deficiência intelectual em eventos públicos de grande}
\chapter{Comunicação alternativa e aumentativa na participação social de pessoas com deficiência intelectual em eventos públicos de grande porte}
\chapter{Psicopedagogicamente inconstitucionalíssimo otorrinolaringologista}
\chapter{Psicopedagogicamente inconstitucionalíssimo otorrinolaringologista socioeducativo}
\chapter{Psicopedagogicamente inconstitucionalíssimo otorrinolaringologista socioeducativo ético-profissional}
\chapter{Psicopedagogicamente inconstitucionalíssimo otorrinolaringologista socioeducativo ético-profissional sociodemográfico}
\chapter{Psicopedagogicamente inconstitucionalíssimo otorrinolaringologista socioeducativo ético-profissional sociodemográfico interdisciplinaridade}
\chapter{Psicopedagogicamente inconstitucionalíssimo otorrinolaringologista socioeducativo ético-profissional sociodemográfico interdisciplinaridade neuropsicológica}
\chapter{Psicopedagogicamente inconstitucionalíssimo otorrinolaringologista socioeducativo ético-profissional sociodemográfico interdisciplinaridade neuropsicológica avaliação}
\chapter{Psicopedagogicamente inconstitucionalíssimo otorrinolaringologista socioeducativo ético-profissional sociodemográfico interdisciplinaridade neuropsicológica avaliação}
\chapter{Psicopedagogicamente inconstitucionalíssimo otorrinolaringologista socioeducativo ético-profissional sociodemográfico interdisciplinaridade neuropsicológica avaliação}
\chapter{Psicopedagogicamente inconstitucionalíssimo otorrinolaringologista socioeducativo ético-profissional sociodemográfico interdisciplinaridade neuropsicológica avaliação}
\chapter{Psicopedagogicamente inconstitucionalíssimo otorrinolaringologista socioeducativo ético-profissional sociodemográfico interdisciplinaridade neuropsicológica avaliação}
\chapter{Psicopedagogicamente inconstitucionalíssimo otorrinolaringologista socioeducativo ético-profissional sociodemográfico interdisciplinaridade neuropsicológica avaliação}
\chapter{Psicopedagogicamente inconstitucionalíssimo otorrinolaringologista socioeducativo ético-profissional sociodemográfico interdisciplinaridade neuropsicológica avaliação}
\chapter{Psicopedagogicamente inconstitucionalíssimo otorrinolaringologista socioeducativo ético-profissional sociodemográfico interdisciplinaridade neuropsicológica avaliação}
\chapter{Psicopedagogicamente inconstitucionalíssimo otorrinolaringologista socioeducativo ético-profissional sociodemográfico interdisciplinaridade neuropsicológica avaliação}
\chapter{Psicopedagogicamente inconstitucionalíssimo otorrinolaringologista socioeducativo ético-profissional sociodemográfico interdisciplinaridade neuropsicológica avaliação}
\end{document}
'   % R: rígido (até 0.5.x)
%   lualatex '\def\var{B}% !TeX program = lualatex
% ===================================================================
% sumario-longo.tex — mede o sumário sem hifenização (issue #35, "Fora do
% escopo"): com a hifenização desligada, a entrada longa de capítulo pode
% bater no número da página. Gera 3 x 16 capítulos, cada um com um prefixo
% (de 3 a 18 palavras) de uma frase longa: frase 1 de palavras médias,
% frase 2 com números, frase 3 só de palavras muito longas. Conta no log
% os "Overfull \hbox" do sumário: o sumário é a única parte do documento
% sem parágrafo de corpo, então todo Overfull "in paragraph" antes do
% primeiro \chapter é entrada que passou da margem.
%
% Uso (de uma pasta com livros_crp_acessivel_book.sty e crpsp-leg.sty no
% caminho; as fontes NEWJUNE por OSFONTDIR):
%   lualatex '\def\var{R}% !TeX program = lualatex
% ===================================================================
% sumario-longo.tex — mede o sumário sem hifenização (issue #35, "Fora do
% escopo"): com a hifenização desligada, a entrada longa de capítulo pode
% bater no número da página. Gera 3 x 16 capítulos, cada um com um prefixo
% (de 3 a 18 palavras) de uma frase longa: frase 1 de palavras médias,
% frase 2 com números, frase 3 só de palavras muito longas. Conta no log
% os "Overfull \hbox" do sumário: o sumário é a única parte do documento
% sem parágrafo de corpo, então todo Overfull "in paragraph" antes do
% primeiro \chapter é entrada que passou da margem.
%
% Uso (de uma pasta com livros_crp_acessivel_book.sty e crpsp-leg.sty no
% caminho; as fontes NEWJUNE por OSFONTDIR):
%   lualatex '\def\var{R}\input{sumario-longo}'   % R: rígido (até 0.5.x)
%   lualatex '\def\var{B}\input{sumario-longo}'   % B: plus 4em (o do guia)
%   lualatex '\def\var{C}\input{sumario-longo}'   % C: plus 4em + \tolerance 10000
%   lualatex '\def\var{S}\input{sumario-longo}'   % S: o que a .sty define
% Cada variante só redefine o \l@chapter depois da carga da .sty.
% ===================================================================
\DocumentMetadata{lang=pt-BR, pdfstandard=ua-2, pdfversion=2.0, tagging=on}
\documentclass[11pt,a5paper,twoside,openright]{book}
\usepackage{livros_crp_acessivel_book}
\hypersetup{pdftitle={Sumario longo}, pdfauthor={CRP SP}}
\makeatletter
\providecommand\var{S}
\def\crp@vR{R}\def\crp@vB{B}\def\crp@vC{C}
\def\crp@lchap#1{%  #1 = acréscimo à margem direita (rightskip)
  \renewcommand*\l@chapter[2]{%
    \ifnum \c@tocdepth >\m@ne
        \addpenalty{-\@highpenalty}\addvspace{0.5em \@plus\p@}%
        \setlength\@tempdima{1.5em}%
        \begingroup
            \parindent \z@ \rightskip \@pnumwidth #1
            \parfillskip -\@pnumwidth
            \leavevmode \bfseries
            \advance\leftskip\@tempdima \hskip -\leftskip
            ##1\nobreak\hfil
            \nobreak\hb@xt@\@pnumwidth{\hss ##2\kern-\p@\kern\p@}\par
            \penalty\@highpenalty
        \endgroup
    \fi}}
\ifx\var\crp@vR \crp@lchap{}\fi
\ifx\var\crp@vB \crp@lchap{\@plus 4em}\fi
\ifx\var\crp@vC \crp@lchap{\@plus 4em \tolerance=10000 \hbadness=10000 }\fi
\makeatother
\begin{document}
\frontmatter
\tableofcontents
\mainmatter
\chapter{Deficiência psicossocial e}
\chapter{Deficiência psicossocial e transtorno}
\chapter{Deficiência psicossocial e transtorno do}
\chapter{Deficiência psicossocial e transtorno do espectro}
\chapter{Deficiência psicossocial e transtorno do espectro autista:}
\chapter{Deficiência psicossocial e transtorno do espectro autista: orientações}
\chapter{Deficiência psicossocial e transtorno do espectro autista: orientações para}
\chapter{Deficiência psicossocial e transtorno do espectro autista: orientações para a}
\chapter{Deficiência psicossocial e transtorno do espectro autista: orientações para a realização}
\chapter{Deficiência psicossocial e transtorno do espectro autista: orientações para a realização de}
\chapter{Deficiência psicossocial e transtorno do espectro autista: orientações para a realização de eventos}
\chapter{Deficiência psicossocial e transtorno do espectro autista: orientações para a realização de eventos e}
\chapter{Deficiência psicossocial e transtorno do espectro autista: orientações para a realização de eventos e apresentações}
\chapter{Deficiência psicossocial e transtorno do espectro autista: orientações para a realização de eventos e apresentações acessíveis}
\chapter{Deficiência psicossocial e transtorno do espectro autista: orientações para a realização de eventos e apresentações acessíveis em}
\chapter{Deficiência psicossocial e transtorno do espectro autista: orientações para a realização de eventos e apresentações acessíveis em todo}
\chapter{Comunicação alternativa e}
\chapter{Comunicação alternativa e aumentativa}
\chapter{Comunicação alternativa e aumentativa na}
\chapter{Comunicação alternativa e aumentativa na participação}
\chapter{Comunicação alternativa e aumentativa na participação social}
\chapter{Comunicação alternativa e aumentativa na participação social de}
\chapter{Comunicação alternativa e aumentativa na participação social de pessoas}
\chapter{Comunicação alternativa e aumentativa na participação social de pessoas com}
\chapter{Comunicação alternativa e aumentativa na participação social de pessoas com deficiência}
\chapter{Comunicação alternativa e aumentativa na participação social de pessoas com deficiência intelectual}
\chapter{Comunicação alternativa e aumentativa na participação social de pessoas com deficiência intelectual em}
\chapter{Comunicação alternativa e aumentativa na participação social de pessoas com deficiência intelectual em eventos}
\chapter{Comunicação alternativa e aumentativa na participação social de pessoas com deficiência intelectual em eventos públicos}
\chapter{Comunicação alternativa e aumentativa na participação social de pessoas com deficiência intelectual em eventos públicos de}
\chapter{Comunicação alternativa e aumentativa na participação social de pessoas com deficiência intelectual em eventos públicos de grande}
\chapter{Comunicação alternativa e aumentativa na participação social de pessoas com deficiência intelectual em eventos públicos de grande porte}
\chapter{Psicopedagogicamente inconstitucionalíssimo otorrinolaringologista}
\chapter{Psicopedagogicamente inconstitucionalíssimo otorrinolaringologista socioeducativo}
\chapter{Psicopedagogicamente inconstitucionalíssimo otorrinolaringologista socioeducativo ético-profissional}
\chapter{Psicopedagogicamente inconstitucionalíssimo otorrinolaringologista socioeducativo ético-profissional sociodemográfico}
\chapter{Psicopedagogicamente inconstitucionalíssimo otorrinolaringologista socioeducativo ético-profissional sociodemográfico interdisciplinaridade}
\chapter{Psicopedagogicamente inconstitucionalíssimo otorrinolaringologista socioeducativo ético-profissional sociodemográfico interdisciplinaridade neuropsicológica}
\chapter{Psicopedagogicamente inconstitucionalíssimo otorrinolaringologista socioeducativo ético-profissional sociodemográfico interdisciplinaridade neuropsicológica avaliação}
\chapter{Psicopedagogicamente inconstitucionalíssimo otorrinolaringologista socioeducativo ético-profissional sociodemográfico interdisciplinaridade neuropsicológica avaliação}
\chapter{Psicopedagogicamente inconstitucionalíssimo otorrinolaringologista socioeducativo ético-profissional sociodemográfico interdisciplinaridade neuropsicológica avaliação}
\chapter{Psicopedagogicamente inconstitucionalíssimo otorrinolaringologista socioeducativo ético-profissional sociodemográfico interdisciplinaridade neuropsicológica avaliação}
\chapter{Psicopedagogicamente inconstitucionalíssimo otorrinolaringologista socioeducativo ético-profissional sociodemográfico interdisciplinaridade neuropsicológica avaliação}
\chapter{Psicopedagogicamente inconstitucionalíssimo otorrinolaringologista socioeducativo ético-profissional sociodemográfico interdisciplinaridade neuropsicológica avaliação}
\chapter{Psicopedagogicamente inconstitucionalíssimo otorrinolaringologista socioeducativo ético-profissional sociodemográfico interdisciplinaridade neuropsicológica avaliação}
\chapter{Psicopedagogicamente inconstitucionalíssimo otorrinolaringologista socioeducativo ético-profissional sociodemográfico interdisciplinaridade neuropsicológica avaliação}
\chapter{Psicopedagogicamente inconstitucionalíssimo otorrinolaringologista socioeducativo ético-profissional sociodemográfico interdisciplinaridade neuropsicológica avaliação}
\chapter{Psicopedagogicamente inconstitucionalíssimo otorrinolaringologista socioeducativo ético-profissional sociodemográfico interdisciplinaridade neuropsicológica avaliação}
\end{document}
'   % R: rígido (até 0.5.x)
%   lualatex '\def\var{B}% !TeX program = lualatex
% ===================================================================
% sumario-longo.tex — mede o sumário sem hifenização (issue #35, "Fora do
% escopo"): com a hifenização desligada, a entrada longa de capítulo pode
% bater no número da página. Gera 3 x 16 capítulos, cada um com um prefixo
% (de 3 a 18 palavras) de uma frase longa: frase 1 de palavras médias,
% frase 2 com números, frase 3 só de palavras muito longas. Conta no log
% os "Overfull \hbox" do sumário: o sumário é a única parte do documento
% sem parágrafo de corpo, então todo Overfull "in paragraph" antes do
% primeiro \chapter é entrada que passou da margem.
%
% Uso (de uma pasta com livros_crp_acessivel_book.sty e crpsp-leg.sty no
% caminho; as fontes NEWJUNE por OSFONTDIR):
%   lualatex '\def\var{R}\input{sumario-longo}'   % R: rígido (até 0.5.x)
%   lualatex '\def\var{B}\input{sumario-longo}'   % B: plus 4em (o do guia)
%   lualatex '\def\var{C}\input{sumario-longo}'   % C: plus 4em + \tolerance 10000
%   lualatex '\def\var{S}\input{sumario-longo}'   % S: o que a .sty define
% Cada variante só redefine o \l@chapter depois da carga da .sty.
% ===================================================================
\DocumentMetadata{lang=pt-BR, pdfstandard=ua-2, pdfversion=2.0, tagging=on}
\documentclass[11pt,a5paper,twoside,openright]{book}
\usepackage{livros_crp_acessivel_book}
\hypersetup{pdftitle={Sumario longo}, pdfauthor={CRP SP}}
\makeatletter
\providecommand\var{S}
\def\crp@vR{R}\def\crp@vB{B}\def\crp@vC{C}
\def\crp@lchap#1{%  #1 = acréscimo à margem direita (rightskip)
  \renewcommand*\l@chapter[2]{%
    \ifnum \c@tocdepth >\m@ne
        \addpenalty{-\@highpenalty}\addvspace{0.5em \@plus\p@}%
        \setlength\@tempdima{1.5em}%
        \begingroup
            \parindent \z@ \rightskip \@pnumwidth #1
            \parfillskip -\@pnumwidth
            \leavevmode \bfseries
            \advance\leftskip\@tempdima \hskip -\leftskip
            ##1\nobreak\hfil
            \nobreak\hb@xt@\@pnumwidth{\hss ##2\kern-\p@\kern\p@}\par
            \penalty\@highpenalty
        \endgroup
    \fi}}
\ifx\var\crp@vR \crp@lchap{}\fi
\ifx\var\crp@vB \crp@lchap{\@plus 4em}\fi
\ifx\var\crp@vC \crp@lchap{\@plus 4em \tolerance=10000 \hbadness=10000 }\fi
\makeatother
\begin{document}
\frontmatter
\tableofcontents
\mainmatter
\chapter{Deficiência psicossocial e}
\chapter{Deficiência psicossocial e transtorno}
\chapter{Deficiência psicossocial e transtorno do}
\chapter{Deficiência psicossocial e transtorno do espectro}
\chapter{Deficiência psicossocial e transtorno do espectro autista:}
\chapter{Deficiência psicossocial e transtorno do espectro autista: orientações}
\chapter{Deficiência psicossocial e transtorno do espectro autista: orientações para}
\chapter{Deficiência psicossocial e transtorno do espectro autista: orientações para a}
\chapter{Deficiência psicossocial e transtorno do espectro autista: orientações para a realização}
\chapter{Deficiência psicossocial e transtorno do espectro autista: orientações para a realização de}
\chapter{Deficiência psicossocial e transtorno do espectro autista: orientações para a realização de eventos}
\chapter{Deficiência psicossocial e transtorno do espectro autista: orientações para a realização de eventos e}
\chapter{Deficiência psicossocial e transtorno do espectro autista: orientações para a realização de eventos e apresentações}
\chapter{Deficiência psicossocial e transtorno do espectro autista: orientações para a realização de eventos e apresentações acessíveis}
\chapter{Deficiência psicossocial e transtorno do espectro autista: orientações para a realização de eventos e apresentações acessíveis em}
\chapter{Deficiência psicossocial e transtorno do espectro autista: orientações para a realização de eventos e apresentações acessíveis em todo}
\chapter{Comunicação alternativa e}
\chapter{Comunicação alternativa e aumentativa}
\chapter{Comunicação alternativa e aumentativa na}
\chapter{Comunicação alternativa e aumentativa na participação}
\chapter{Comunicação alternativa e aumentativa na participação social}
\chapter{Comunicação alternativa e aumentativa na participação social de}
\chapter{Comunicação alternativa e aumentativa na participação social de pessoas}
\chapter{Comunicação alternativa e aumentativa na participação social de pessoas com}
\chapter{Comunicação alternativa e aumentativa na participação social de pessoas com deficiência}
\chapter{Comunicação alternativa e aumentativa na participação social de pessoas com deficiência intelectual}
\chapter{Comunicação alternativa e aumentativa na participação social de pessoas com deficiência intelectual em}
\chapter{Comunicação alternativa e aumentativa na participação social de pessoas com deficiência intelectual em eventos}
\chapter{Comunicação alternativa e aumentativa na participação social de pessoas com deficiência intelectual em eventos públicos}
\chapter{Comunicação alternativa e aumentativa na participação social de pessoas com deficiência intelectual em eventos públicos de}
\chapter{Comunicação alternativa e aumentativa na participação social de pessoas com deficiência intelectual em eventos públicos de grande}
\chapter{Comunicação alternativa e aumentativa na participação social de pessoas com deficiência intelectual em eventos públicos de grande porte}
\chapter{Psicopedagogicamente inconstitucionalíssimo otorrinolaringologista}
\chapter{Psicopedagogicamente inconstitucionalíssimo otorrinolaringologista socioeducativo}
\chapter{Psicopedagogicamente inconstitucionalíssimo otorrinolaringologista socioeducativo ético-profissional}
\chapter{Psicopedagogicamente inconstitucionalíssimo otorrinolaringologista socioeducativo ético-profissional sociodemográfico}
\chapter{Psicopedagogicamente inconstitucionalíssimo otorrinolaringologista socioeducativo ético-profissional sociodemográfico interdisciplinaridade}
\chapter{Psicopedagogicamente inconstitucionalíssimo otorrinolaringologista socioeducativo ético-profissional sociodemográfico interdisciplinaridade neuropsicológica}
\chapter{Psicopedagogicamente inconstitucionalíssimo otorrinolaringologista socioeducativo ético-profissional sociodemográfico interdisciplinaridade neuropsicológica avaliação}
\chapter{Psicopedagogicamente inconstitucionalíssimo otorrinolaringologista socioeducativo ético-profissional sociodemográfico interdisciplinaridade neuropsicológica avaliação}
\chapter{Psicopedagogicamente inconstitucionalíssimo otorrinolaringologista socioeducativo ético-profissional sociodemográfico interdisciplinaridade neuropsicológica avaliação}
\chapter{Psicopedagogicamente inconstitucionalíssimo otorrinolaringologista socioeducativo ético-profissional sociodemográfico interdisciplinaridade neuropsicológica avaliação}
\chapter{Psicopedagogicamente inconstitucionalíssimo otorrinolaringologista socioeducativo ético-profissional sociodemográfico interdisciplinaridade neuropsicológica avaliação}
\chapter{Psicopedagogicamente inconstitucionalíssimo otorrinolaringologista socioeducativo ético-profissional sociodemográfico interdisciplinaridade neuropsicológica avaliação}
\chapter{Psicopedagogicamente inconstitucionalíssimo otorrinolaringologista socioeducativo ético-profissional sociodemográfico interdisciplinaridade neuropsicológica avaliação}
\chapter{Psicopedagogicamente inconstitucionalíssimo otorrinolaringologista socioeducativo ético-profissional sociodemográfico interdisciplinaridade neuropsicológica avaliação}
\chapter{Psicopedagogicamente inconstitucionalíssimo otorrinolaringologista socioeducativo ético-profissional sociodemográfico interdisciplinaridade neuropsicológica avaliação}
\chapter{Psicopedagogicamente inconstitucionalíssimo otorrinolaringologista socioeducativo ético-profissional sociodemográfico interdisciplinaridade neuropsicológica avaliação}
\end{document}
'   % B: plus 4em (o do guia)
%   lualatex '\def\var{C}% !TeX program = lualatex
% ===================================================================
% sumario-longo.tex — mede o sumário sem hifenização (issue #35, "Fora do
% escopo"): com a hifenização desligada, a entrada longa de capítulo pode
% bater no número da página. Gera 3 x 16 capítulos, cada um com um prefixo
% (de 3 a 18 palavras) de uma frase longa: frase 1 de palavras médias,
% frase 2 com números, frase 3 só de palavras muito longas. Conta no log
% os "Overfull \hbox" do sumário: o sumário é a única parte do documento
% sem parágrafo de corpo, então todo Overfull "in paragraph" antes do
% primeiro \chapter é entrada que passou da margem.
%
% Uso (de uma pasta com livros_crp_acessivel_book.sty e crpsp-leg.sty no
% caminho; as fontes NEWJUNE por OSFONTDIR):
%   lualatex '\def\var{R}\input{sumario-longo}'   % R: rígido (até 0.5.x)
%   lualatex '\def\var{B}\input{sumario-longo}'   % B: plus 4em (o do guia)
%   lualatex '\def\var{C}\input{sumario-longo}'   % C: plus 4em + \tolerance 10000
%   lualatex '\def\var{S}\input{sumario-longo}'   % S: o que a .sty define
% Cada variante só redefine o \l@chapter depois da carga da .sty.
% ===================================================================
\DocumentMetadata{lang=pt-BR, pdfstandard=ua-2, pdfversion=2.0, tagging=on}
\documentclass[11pt,a5paper,twoside,openright]{book}
\usepackage{livros_crp_acessivel_book}
\hypersetup{pdftitle={Sumario longo}, pdfauthor={CRP SP}}
\makeatletter
\providecommand\var{S}
\def\crp@vR{R}\def\crp@vB{B}\def\crp@vC{C}
\def\crp@lchap#1{%  #1 = acréscimo à margem direita (rightskip)
  \renewcommand*\l@chapter[2]{%
    \ifnum \c@tocdepth >\m@ne
        \addpenalty{-\@highpenalty}\addvspace{0.5em \@plus\p@}%
        \setlength\@tempdima{1.5em}%
        \begingroup
            \parindent \z@ \rightskip \@pnumwidth #1
            \parfillskip -\@pnumwidth
            \leavevmode \bfseries
            \advance\leftskip\@tempdima \hskip -\leftskip
            ##1\nobreak\hfil
            \nobreak\hb@xt@\@pnumwidth{\hss ##2\kern-\p@\kern\p@}\par
            \penalty\@highpenalty
        \endgroup
    \fi}}
\ifx\var\crp@vR \crp@lchap{}\fi
\ifx\var\crp@vB \crp@lchap{\@plus 4em}\fi
\ifx\var\crp@vC \crp@lchap{\@plus 4em \tolerance=10000 \hbadness=10000 }\fi
\makeatother
\begin{document}
\frontmatter
\tableofcontents
\mainmatter
\chapter{Deficiência psicossocial e}
\chapter{Deficiência psicossocial e transtorno}
\chapter{Deficiência psicossocial e transtorno do}
\chapter{Deficiência psicossocial e transtorno do espectro}
\chapter{Deficiência psicossocial e transtorno do espectro autista:}
\chapter{Deficiência psicossocial e transtorno do espectro autista: orientações}
\chapter{Deficiência psicossocial e transtorno do espectro autista: orientações para}
\chapter{Deficiência psicossocial e transtorno do espectro autista: orientações para a}
\chapter{Deficiência psicossocial e transtorno do espectro autista: orientações para a realização}
\chapter{Deficiência psicossocial e transtorno do espectro autista: orientações para a realização de}
\chapter{Deficiência psicossocial e transtorno do espectro autista: orientações para a realização de eventos}
\chapter{Deficiência psicossocial e transtorno do espectro autista: orientações para a realização de eventos e}
\chapter{Deficiência psicossocial e transtorno do espectro autista: orientações para a realização de eventos e apresentações}
\chapter{Deficiência psicossocial e transtorno do espectro autista: orientações para a realização de eventos e apresentações acessíveis}
\chapter{Deficiência psicossocial e transtorno do espectro autista: orientações para a realização de eventos e apresentações acessíveis em}
\chapter{Deficiência psicossocial e transtorno do espectro autista: orientações para a realização de eventos e apresentações acessíveis em todo}
\chapter{Comunicação alternativa e}
\chapter{Comunicação alternativa e aumentativa}
\chapter{Comunicação alternativa e aumentativa na}
\chapter{Comunicação alternativa e aumentativa na participação}
\chapter{Comunicação alternativa e aumentativa na participação social}
\chapter{Comunicação alternativa e aumentativa na participação social de}
\chapter{Comunicação alternativa e aumentativa na participação social de pessoas}
\chapter{Comunicação alternativa e aumentativa na participação social de pessoas com}
\chapter{Comunicação alternativa e aumentativa na participação social de pessoas com deficiência}
\chapter{Comunicação alternativa e aumentativa na participação social de pessoas com deficiência intelectual}
\chapter{Comunicação alternativa e aumentativa na participação social de pessoas com deficiência intelectual em}
\chapter{Comunicação alternativa e aumentativa na participação social de pessoas com deficiência intelectual em eventos}
\chapter{Comunicação alternativa e aumentativa na participação social de pessoas com deficiência intelectual em eventos públicos}
\chapter{Comunicação alternativa e aumentativa na participação social de pessoas com deficiência intelectual em eventos públicos de}
\chapter{Comunicação alternativa e aumentativa na participação social de pessoas com deficiência intelectual em eventos públicos de grande}
\chapter{Comunicação alternativa e aumentativa na participação social de pessoas com deficiência intelectual em eventos públicos de grande porte}
\chapter{Psicopedagogicamente inconstitucionalíssimo otorrinolaringologista}
\chapter{Psicopedagogicamente inconstitucionalíssimo otorrinolaringologista socioeducativo}
\chapter{Psicopedagogicamente inconstitucionalíssimo otorrinolaringologista socioeducativo ético-profissional}
\chapter{Psicopedagogicamente inconstitucionalíssimo otorrinolaringologista socioeducativo ético-profissional sociodemográfico}
\chapter{Psicopedagogicamente inconstitucionalíssimo otorrinolaringologista socioeducativo ético-profissional sociodemográfico interdisciplinaridade}
\chapter{Psicopedagogicamente inconstitucionalíssimo otorrinolaringologista socioeducativo ético-profissional sociodemográfico interdisciplinaridade neuropsicológica}
\chapter{Psicopedagogicamente inconstitucionalíssimo otorrinolaringologista socioeducativo ético-profissional sociodemográfico interdisciplinaridade neuropsicológica avaliação}
\chapter{Psicopedagogicamente inconstitucionalíssimo otorrinolaringologista socioeducativo ético-profissional sociodemográfico interdisciplinaridade neuropsicológica avaliação}
\chapter{Psicopedagogicamente inconstitucionalíssimo otorrinolaringologista socioeducativo ético-profissional sociodemográfico interdisciplinaridade neuropsicológica avaliação}
\chapter{Psicopedagogicamente inconstitucionalíssimo otorrinolaringologista socioeducativo ético-profissional sociodemográfico interdisciplinaridade neuropsicológica avaliação}
\chapter{Psicopedagogicamente inconstitucionalíssimo otorrinolaringologista socioeducativo ético-profissional sociodemográfico interdisciplinaridade neuropsicológica avaliação}
\chapter{Psicopedagogicamente inconstitucionalíssimo otorrinolaringologista socioeducativo ético-profissional sociodemográfico interdisciplinaridade neuropsicológica avaliação}
\chapter{Psicopedagogicamente inconstitucionalíssimo otorrinolaringologista socioeducativo ético-profissional sociodemográfico interdisciplinaridade neuropsicológica avaliação}
\chapter{Psicopedagogicamente inconstitucionalíssimo otorrinolaringologista socioeducativo ético-profissional sociodemográfico interdisciplinaridade neuropsicológica avaliação}
\chapter{Psicopedagogicamente inconstitucionalíssimo otorrinolaringologista socioeducativo ético-profissional sociodemográfico interdisciplinaridade neuropsicológica avaliação}
\chapter{Psicopedagogicamente inconstitucionalíssimo otorrinolaringologista socioeducativo ético-profissional sociodemográfico interdisciplinaridade neuropsicológica avaliação}
\end{document}
'   % C: plus 4em + \tolerance 10000
%   lualatex '\def\var{S}% !TeX program = lualatex
% ===================================================================
% sumario-longo.tex — mede o sumário sem hifenização (issue #35, "Fora do
% escopo"): com a hifenização desligada, a entrada longa de capítulo pode
% bater no número da página. Gera 3 x 16 capítulos, cada um com um prefixo
% (de 3 a 18 palavras) de uma frase longa: frase 1 de palavras médias,
% frase 2 com números, frase 3 só de palavras muito longas. Conta no log
% os "Overfull \hbox" do sumário: o sumário é a única parte do documento
% sem parágrafo de corpo, então todo Overfull "in paragraph" antes do
% primeiro \chapter é entrada que passou da margem.
%
% Uso (de uma pasta com livros_crp_acessivel_book.sty e crpsp-leg.sty no
% caminho; as fontes NEWJUNE por OSFONTDIR):
%   lualatex '\def\var{R}\input{sumario-longo}'   % R: rígido (até 0.5.x)
%   lualatex '\def\var{B}\input{sumario-longo}'   % B: plus 4em (o do guia)
%   lualatex '\def\var{C}\input{sumario-longo}'   % C: plus 4em + \tolerance 10000
%   lualatex '\def\var{S}\input{sumario-longo}'   % S: o que a .sty define
% Cada variante só redefine o \l@chapter depois da carga da .sty.
% ===================================================================
\DocumentMetadata{lang=pt-BR, pdfstandard=ua-2, pdfversion=2.0, tagging=on}
\documentclass[11pt,a5paper,twoside,openright]{book}
\usepackage{livros_crp_acessivel_book}
\hypersetup{pdftitle={Sumario longo}, pdfauthor={CRP SP}}
\makeatletter
\providecommand\var{S}
\def\crp@vR{R}\def\crp@vB{B}\def\crp@vC{C}
\def\crp@lchap#1{%  #1 = acréscimo à margem direita (rightskip)
  \renewcommand*\l@chapter[2]{%
    \ifnum \c@tocdepth >\m@ne
        \addpenalty{-\@highpenalty}\addvspace{0.5em \@plus\p@}%
        \setlength\@tempdima{1.5em}%
        \begingroup
            \parindent \z@ \rightskip \@pnumwidth #1
            \parfillskip -\@pnumwidth
            \leavevmode \bfseries
            \advance\leftskip\@tempdima \hskip -\leftskip
            ##1\nobreak\hfil
            \nobreak\hb@xt@\@pnumwidth{\hss ##2\kern-\p@\kern\p@}\par
            \penalty\@highpenalty
        \endgroup
    \fi}}
\ifx\var\crp@vR \crp@lchap{}\fi
\ifx\var\crp@vB \crp@lchap{\@plus 4em}\fi
\ifx\var\crp@vC \crp@lchap{\@plus 4em \tolerance=10000 \hbadness=10000 }\fi
\makeatother
\begin{document}
\frontmatter
\tableofcontents
\mainmatter
\chapter{Deficiência psicossocial e}
\chapter{Deficiência psicossocial e transtorno}
\chapter{Deficiência psicossocial e transtorno do}
\chapter{Deficiência psicossocial e transtorno do espectro}
\chapter{Deficiência psicossocial e transtorno do espectro autista:}
\chapter{Deficiência psicossocial e transtorno do espectro autista: orientações}
\chapter{Deficiência psicossocial e transtorno do espectro autista: orientações para}
\chapter{Deficiência psicossocial e transtorno do espectro autista: orientações para a}
\chapter{Deficiência psicossocial e transtorno do espectro autista: orientações para a realização}
\chapter{Deficiência psicossocial e transtorno do espectro autista: orientações para a realização de}
\chapter{Deficiência psicossocial e transtorno do espectro autista: orientações para a realização de eventos}
\chapter{Deficiência psicossocial e transtorno do espectro autista: orientações para a realização de eventos e}
\chapter{Deficiência psicossocial e transtorno do espectro autista: orientações para a realização de eventos e apresentações}
\chapter{Deficiência psicossocial e transtorno do espectro autista: orientações para a realização de eventos e apresentações acessíveis}
\chapter{Deficiência psicossocial e transtorno do espectro autista: orientações para a realização de eventos e apresentações acessíveis em}
\chapter{Deficiência psicossocial e transtorno do espectro autista: orientações para a realização de eventos e apresentações acessíveis em todo}
\chapter{Comunicação alternativa e}
\chapter{Comunicação alternativa e aumentativa}
\chapter{Comunicação alternativa e aumentativa na}
\chapter{Comunicação alternativa e aumentativa na participação}
\chapter{Comunicação alternativa e aumentativa na participação social}
\chapter{Comunicação alternativa e aumentativa na participação social de}
\chapter{Comunicação alternativa e aumentativa na participação social de pessoas}
\chapter{Comunicação alternativa e aumentativa na participação social de pessoas com}
\chapter{Comunicação alternativa e aumentativa na participação social de pessoas com deficiência}
\chapter{Comunicação alternativa e aumentativa na participação social de pessoas com deficiência intelectual}
\chapter{Comunicação alternativa e aumentativa na participação social de pessoas com deficiência intelectual em}
\chapter{Comunicação alternativa e aumentativa na participação social de pessoas com deficiência intelectual em eventos}
\chapter{Comunicação alternativa e aumentativa na participação social de pessoas com deficiência intelectual em eventos públicos}
\chapter{Comunicação alternativa e aumentativa na participação social de pessoas com deficiência intelectual em eventos públicos de}
\chapter{Comunicação alternativa e aumentativa na participação social de pessoas com deficiência intelectual em eventos públicos de grande}
\chapter{Comunicação alternativa e aumentativa na participação social de pessoas com deficiência intelectual em eventos públicos de grande porte}
\chapter{Psicopedagogicamente inconstitucionalíssimo otorrinolaringologista}
\chapter{Psicopedagogicamente inconstitucionalíssimo otorrinolaringologista socioeducativo}
\chapter{Psicopedagogicamente inconstitucionalíssimo otorrinolaringologista socioeducativo ético-profissional}
\chapter{Psicopedagogicamente inconstitucionalíssimo otorrinolaringologista socioeducativo ético-profissional sociodemográfico}
\chapter{Psicopedagogicamente inconstitucionalíssimo otorrinolaringologista socioeducativo ético-profissional sociodemográfico interdisciplinaridade}
\chapter{Psicopedagogicamente inconstitucionalíssimo otorrinolaringologista socioeducativo ético-profissional sociodemográfico interdisciplinaridade neuropsicológica}
\chapter{Psicopedagogicamente inconstitucionalíssimo otorrinolaringologista socioeducativo ético-profissional sociodemográfico interdisciplinaridade neuropsicológica avaliação}
\chapter{Psicopedagogicamente inconstitucionalíssimo otorrinolaringologista socioeducativo ético-profissional sociodemográfico interdisciplinaridade neuropsicológica avaliação}
\chapter{Psicopedagogicamente inconstitucionalíssimo otorrinolaringologista socioeducativo ético-profissional sociodemográfico interdisciplinaridade neuropsicológica avaliação}
\chapter{Psicopedagogicamente inconstitucionalíssimo otorrinolaringologista socioeducativo ético-profissional sociodemográfico interdisciplinaridade neuropsicológica avaliação}
\chapter{Psicopedagogicamente inconstitucionalíssimo otorrinolaringologista socioeducativo ético-profissional sociodemográfico interdisciplinaridade neuropsicológica avaliação}
\chapter{Psicopedagogicamente inconstitucionalíssimo otorrinolaringologista socioeducativo ético-profissional sociodemográfico interdisciplinaridade neuropsicológica avaliação}
\chapter{Psicopedagogicamente inconstitucionalíssimo otorrinolaringologista socioeducativo ético-profissional sociodemográfico interdisciplinaridade neuropsicológica avaliação}
\chapter{Psicopedagogicamente inconstitucionalíssimo otorrinolaringologista socioeducativo ético-profissional sociodemográfico interdisciplinaridade neuropsicológica avaliação}
\chapter{Psicopedagogicamente inconstitucionalíssimo otorrinolaringologista socioeducativo ético-profissional sociodemográfico interdisciplinaridade neuropsicológica avaliação}
\chapter{Psicopedagogicamente inconstitucionalíssimo otorrinolaringologista socioeducativo ético-profissional sociodemográfico interdisciplinaridade neuropsicológica avaliação}
\end{document}
'   % S: o que a .sty define
% Cada variante só redefine o \l@chapter depois da carga da .sty.
% ===================================================================
\DocumentMetadata{lang=pt-BR, pdfstandard=ua-2, pdfversion=2.0, tagging=on}
\documentclass[11pt,a5paper,twoside,openright]{book}
\usepackage{livros_crp_acessivel_book}
\hypersetup{pdftitle={Sumario longo}, pdfauthor={CRP SP}}
\makeatletter
\providecommand\var{S}
\def\crp@vR{R}\def\crp@vB{B}\def\crp@vC{C}
\def\crp@lchap#1{%  #1 = acréscimo à margem direita (rightskip)
  \renewcommand*\l@chapter[2]{%
    \ifnum \c@tocdepth >\m@ne
        \addpenalty{-\@highpenalty}\addvspace{0.5em \@plus\p@}%
        \setlength\@tempdima{1.5em}%
        \begingroup
            \parindent \z@ \rightskip \@pnumwidth #1
            \parfillskip -\@pnumwidth
            \leavevmode \bfseries
            \advance\leftskip\@tempdima \hskip -\leftskip
            ##1\nobreak\hfil
            \nobreak\hb@xt@\@pnumwidth{\hss ##2\kern-\p@\kern\p@}\par
            \penalty\@highpenalty
        \endgroup
    \fi}}
\ifx\var\crp@vR \crp@lchap{}\fi
\ifx\var\crp@vB \crp@lchap{\@plus 4em}\fi
\ifx\var\crp@vC \crp@lchap{\@plus 4em \tolerance=10000 \hbadness=10000 }\fi
\makeatother
\begin{document}
\frontmatter
\tableofcontents
\mainmatter
\chapter{Deficiência psicossocial e}
\chapter{Deficiência psicossocial e transtorno}
\chapter{Deficiência psicossocial e transtorno do}
\chapter{Deficiência psicossocial e transtorno do espectro}
\chapter{Deficiência psicossocial e transtorno do espectro autista:}
\chapter{Deficiência psicossocial e transtorno do espectro autista: orientações}
\chapter{Deficiência psicossocial e transtorno do espectro autista: orientações para}
\chapter{Deficiência psicossocial e transtorno do espectro autista: orientações para a}
\chapter{Deficiência psicossocial e transtorno do espectro autista: orientações para a realização}
\chapter{Deficiência psicossocial e transtorno do espectro autista: orientações para a realização de}
\chapter{Deficiência psicossocial e transtorno do espectro autista: orientações para a realização de eventos}
\chapter{Deficiência psicossocial e transtorno do espectro autista: orientações para a realização de eventos e}
\chapter{Deficiência psicossocial e transtorno do espectro autista: orientações para a realização de eventos e apresentações}
\chapter{Deficiência psicossocial e transtorno do espectro autista: orientações para a realização de eventos e apresentações acessíveis}
\chapter{Deficiência psicossocial e transtorno do espectro autista: orientações para a realização de eventos e apresentações acessíveis em}
\chapter{Deficiência psicossocial e transtorno do espectro autista: orientações para a realização de eventos e apresentações acessíveis em todo}
\chapter{Comunicação alternativa e}
\chapter{Comunicação alternativa e aumentativa}
\chapter{Comunicação alternativa e aumentativa na}
\chapter{Comunicação alternativa e aumentativa na participação}
\chapter{Comunicação alternativa e aumentativa na participação social}
\chapter{Comunicação alternativa e aumentativa na participação social de}
\chapter{Comunicação alternativa e aumentativa na participação social de pessoas}
\chapter{Comunicação alternativa e aumentativa na participação social de pessoas com}
\chapter{Comunicação alternativa e aumentativa na participação social de pessoas com deficiência}
\chapter{Comunicação alternativa e aumentativa na participação social de pessoas com deficiência intelectual}
\chapter{Comunicação alternativa e aumentativa na participação social de pessoas com deficiência intelectual em}
\chapter{Comunicação alternativa e aumentativa na participação social de pessoas com deficiência intelectual em eventos}
\chapter{Comunicação alternativa e aumentativa na participação social de pessoas com deficiência intelectual em eventos públicos}
\chapter{Comunicação alternativa e aumentativa na participação social de pessoas com deficiência intelectual em eventos públicos de}
\chapter{Comunicação alternativa e aumentativa na participação social de pessoas com deficiência intelectual em eventos públicos de grande}
\chapter{Comunicação alternativa e aumentativa na participação social de pessoas com deficiência intelectual em eventos públicos de grande porte}
\chapter{Psicopedagogicamente inconstitucionalíssimo otorrinolaringologista}
\chapter{Psicopedagogicamente inconstitucionalíssimo otorrinolaringologista socioeducativo}
\chapter{Psicopedagogicamente inconstitucionalíssimo otorrinolaringologista socioeducativo ético-profissional}
\chapter{Psicopedagogicamente inconstitucionalíssimo otorrinolaringologista socioeducativo ético-profissional sociodemográfico}
\chapter{Psicopedagogicamente inconstitucionalíssimo otorrinolaringologista socioeducativo ético-profissional sociodemográfico interdisciplinaridade}
\chapter{Psicopedagogicamente inconstitucionalíssimo otorrinolaringologista socioeducativo ético-profissional sociodemográfico interdisciplinaridade neuropsicológica}
\chapter{Psicopedagogicamente inconstitucionalíssimo otorrinolaringologista socioeducativo ético-profissional sociodemográfico interdisciplinaridade neuropsicológica avaliação}
\chapter{Psicopedagogicamente inconstitucionalíssimo otorrinolaringologista socioeducativo ético-profissional sociodemográfico interdisciplinaridade neuropsicológica avaliação}
\chapter{Psicopedagogicamente inconstitucionalíssimo otorrinolaringologista socioeducativo ético-profissional sociodemográfico interdisciplinaridade neuropsicológica avaliação}
\chapter{Psicopedagogicamente inconstitucionalíssimo otorrinolaringologista socioeducativo ético-profissional sociodemográfico interdisciplinaridade neuropsicológica avaliação}
\chapter{Psicopedagogicamente inconstitucionalíssimo otorrinolaringologista socioeducativo ético-profissional sociodemográfico interdisciplinaridade neuropsicológica avaliação}
\chapter{Psicopedagogicamente inconstitucionalíssimo otorrinolaringologista socioeducativo ético-profissional sociodemográfico interdisciplinaridade neuropsicológica avaliação}
\chapter{Psicopedagogicamente inconstitucionalíssimo otorrinolaringologista socioeducativo ético-profissional sociodemográfico interdisciplinaridade neuropsicológica avaliação}
\chapter{Psicopedagogicamente inconstitucionalíssimo otorrinolaringologista socioeducativo ético-profissional sociodemográfico interdisciplinaridade neuropsicológica avaliação}
\chapter{Psicopedagogicamente inconstitucionalíssimo otorrinolaringologista socioeducativo ético-profissional sociodemográfico interdisciplinaridade neuropsicológica avaliação}
\chapter{Psicopedagogicamente inconstitucionalíssimo otorrinolaringologista socioeducativo ético-profissional sociodemográfico interdisciplinaridade neuropsicológica avaliação}
\end{document}
'   % B: plus 4em (o do guia)
%   lualatex '\def\var{C}% !TeX program = lualatex
% ===================================================================
% sumario-longo.tex — mede o sumário sem hifenização (issue #35, "Fora do
% escopo"): com a hifenização desligada, a entrada longa de capítulo pode
% bater no número da página. Gera 3 x 16 capítulos, cada um com um prefixo
% (de 3 a 18 palavras) de uma frase longa: frase 1 de palavras médias,
% frase 2 com números, frase 3 só de palavras muito longas. Conta no log
% os "Overfull \hbox" do sumário: o sumário é a única parte do documento
% sem parágrafo de corpo, então todo Overfull "in paragraph" antes do
% primeiro \chapter é entrada que passou da margem.
%
% Uso (de uma pasta com livros_crp_acessivel_book.sty e crpsp-leg.sty no
% caminho; as fontes NEWJUNE por OSFONTDIR):
%   lualatex '\def\var{R}% !TeX program = lualatex
% ===================================================================
% sumario-longo.tex — mede o sumário sem hifenização (issue #35, "Fora do
% escopo"): com a hifenização desligada, a entrada longa de capítulo pode
% bater no número da página. Gera 3 x 16 capítulos, cada um com um prefixo
% (de 3 a 18 palavras) de uma frase longa: frase 1 de palavras médias,
% frase 2 com números, frase 3 só de palavras muito longas. Conta no log
% os "Overfull \hbox" do sumário: o sumário é a única parte do documento
% sem parágrafo de corpo, então todo Overfull "in paragraph" antes do
% primeiro \chapter é entrada que passou da margem.
%
% Uso (de uma pasta com livros_crp_acessivel_book.sty e crpsp-leg.sty no
% caminho; as fontes NEWJUNE por OSFONTDIR):
%   lualatex '\def\var{R}\input{sumario-longo}'   % R: rígido (até 0.5.x)
%   lualatex '\def\var{B}\input{sumario-longo}'   % B: plus 4em (o do guia)
%   lualatex '\def\var{C}\input{sumario-longo}'   % C: plus 4em + \tolerance 10000
%   lualatex '\def\var{S}\input{sumario-longo}'   % S: o que a .sty define
% Cada variante só redefine o \l@chapter depois da carga da .sty.
% ===================================================================
\DocumentMetadata{lang=pt-BR, pdfstandard=ua-2, pdfversion=2.0, tagging=on}
\documentclass[11pt,a5paper,twoside,openright]{book}
\usepackage{livros_crp_acessivel_book}
\hypersetup{pdftitle={Sumario longo}, pdfauthor={CRP SP}}
\makeatletter
\providecommand\var{S}
\def\crp@vR{R}\def\crp@vB{B}\def\crp@vC{C}
\def\crp@lchap#1{%  #1 = acréscimo à margem direita (rightskip)
  \renewcommand*\l@chapter[2]{%
    \ifnum \c@tocdepth >\m@ne
        \addpenalty{-\@highpenalty}\addvspace{0.5em \@plus\p@}%
        \setlength\@tempdima{1.5em}%
        \begingroup
            \parindent \z@ \rightskip \@pnumwidth #1
            \parfillskip -\@pnumwidth
            \leavevmode \bfseries
            \advance\leftskip\@tempdima \hskip -\leftskip
            ##1\nobreak\hfil
            \nobreak\hb@xt@\@pnumwidth{\hss ##2\kern-\p@\kern\p@}\par
            \penalty\@highpenalty
        \endgroup
    \fi}}
\ifx\var\crp@vR \crp@lchap{}\fi
\ifx\var\crp@vB \crp@lchap{\@plus 4em}\fi
\ifx\var\crp@vC \crp@lchap{\@plus 4em \tolerance=10000 \hbadness=10000 }\fi
\makeatother
\begin{document}
\frontmatter
\tableofcontents
\mainmatter
\chapter{Deficiência psicossocial e}
\chapter{Deficiência psicossocial e transtorno}
\chapter{Deficiência psicossocial e transtorno do}
\chapter{Deficiência psicossocial e transtorno do espectro}
\chapter{Deficiência psicossocial e transtorno do espectro autista:}
\chapter{Deficiência psicossocial e transtorno do espectro autista: orientações}
\chapter{Deficiência psicossocial e transtorno do espectro autista: orientações para}
\chapter{Deficiência psicossocial e transtorno do espectro autista: orientações para a}
\chapter{Deficiência psicossocial e transtorno do espectro autista: orientações para a realização}
\chapter{Deficiência psicossocial e transtorno do espectro autista: orientações para a realização de}
\chapter{Deficiência psicossocial e transtorno do espectro autista: orientações para a realização de eventos}
\chapter{Deficiência psicossocial e transtorno do espectro autista: orientações para a realização de eventos e}
\chapter{Deficiência psicossocial e transtorno do espectro autista: orientações para a realização de eventos e apresentações}
\chapter{Deficiência psicossocial e transtorno do espectro autista: orientações para a realização de eventos e apresentações acessíveis}
\chapter{Deficiência psicossocial e transtorno do espectro autista: orientações para a realização de eventos e apresentações acessíveis em}
\chapter{Deficiência psicossocial e transtorno do espectro autista: orientações para a realização de eventos e apresentações acessíveis em todo}
\chapter{Comunicação alternativa e}
\chapter{Comunicação alternativa e aumentativa}
\chapter{Comunicação alternativa e aumentativa na}
\chapter{Comunicação alternativa e aumentativa na participação}
\chapter{Comunicação alternativa e aumentativa na participação social}
\chapter{Comunicação alternativa e aumentativa na participação social de}
\chapter{Comunicação alternativa e aumentativa na participação social de pessoas}
\chapter{Comunicação alternativa e aumentativa na participação social de pessoas com}
\chapter{Comunicação alternativa e aumentativa na participação social de pessoas com deficiência}
\chapter{Comunicação alternativa e aumentativa na participação social de pessoas com deficiência intelectual}
\chapter{Comunicação alternativa e aumentativa na participação social de pessoas com deficiência intelectual em}
\chapter{Comunicação alternativa e aumentativa na participação social de pessoas com deficiência intelectual em eventos}
\chapter{Comunicação alternativa e aumentativa na participação social de pessoas com deficiência intelectual em eventos públicos}
\chapter{Comunicação alternativa e aumentativa na participação social de pessoas com deficiência intelectual em eventos públicos de}
\chapter{Comunicação alternativa e aumentativa na participação social de pessoas com deficiência intelectual em eventos públicos de grande}
\chapter{Comunicação alternativa e aumentativa na participação social de pessoas com deficiência intelectual em eventos públicos de grande porte}
\chapter{Psicopedagogicamente inconstitucionalíssimo otorrinolaringologista}
\chapter{Psicopedagogicamente inconstitucionalíssimo otorrinolaringologista socioeducativo}
\chapter{Psicopedagogicamente inconstitucionalíssimo otorrinolaringologista socioeducativo ético-profissional}
\chapter{Psicopedagogicamente inconstitucionalíssimo otorrinolaringologista socioeducativo ético-profissional sociodemográfico}
\chapter{Psicopedagogicamente inconstitucionalíssimo otorrinolaringologista socioeducativo ético-profissional sociodemográfico interdisciplinaridade}
\chapter{Psicopedagogicamente inconstitucionalíssimo otorrinolaringologista socioeducativo ético-profissional sociodemográfico interdisciplinaridade neuropsicológica}
\chapter{Psicopedagogicamente inconstitucionalíssimo otorrinolaringologista socioeducativo ético-profissional sociodemográfico interdisciplinaridade neuropsicológica avaliação}
\chapter{Psicopedagogicamente inconstitucionalíssimo otorrinolaringologista socioeducativo ético-profissional sociodemográfico interdisciplinaridade neuropsicológica avaliação}
\chapter{Psicopedagogicamente inconstitucionalíssimo otorrinolaringologista socioeducativo ético-profissional sociodemográfico interdisciplinaridade neuropsicológica avaliação}
\chapter{Psicopedagogicamente inconstitucionalíssimo otorrinolaringologista socioeducativo ético-profissional sociodemográfico interdisciplinaridade neuropsicológica avaliação}
\chapter{Psicopedagogicamente inconstitucionalíssimo otorrinolaringologista socioeducativo ético-profissional sociodemográfico interdisciplinaridade neuropsicológica avaliação}
\chapter{Psicopedagogicamente inconstitucionalíssimo otorrinolaringologista socioeducativo ético-profissional sociodemográfico interdisciplinaridade neuropsicológica avaliação}
\chapter{Psicopedagogicamente inconstitucionalíssimo otorrinolaringologista socioeducativo ético-profissional sociodemográfico interdisciplinaridade neuropsicológica avaliação}
\chapter{Psicopedagogicamente inconstitucionalíssimo otorrinolaringologista socioeducativo ético-profissional sociodemográfico interdisciplinaridade neuropsicológica avaliação}
\chapter{Psicopedagogicamente inconstitucionalíssimo otorrinolaringologista socioeducativo ético-profissional sociodemográfico interdisciplinaridade neuropsicológica avaliação}
\chapter{Psicopedagogicamente inconstitucionalíssimo otorrinolaringologista socioeducativo ético-profissional sociodemográfico interdisciplinaridade neuropsicológica avaliação}
\end{document}
'   % R: rígido (até 0.5.x)
%   lualatex '\def\var{B}% !TeX program = lualatex
% ===================================================================
% sumario-longo.tex — mede o sumário sem hifenização (issue #35, "Fora do
% escopo"): com a hifenização desligada, a entrada longa de capítulo pode
% bater no número da página. Gera 3 x 16 capítulos, cada um com um prefixo
% (de 3 a 18 palavras) de uma frase longa: frase 1 de palavras médias,
% frase 2 com números, frase 3 só de palavras muito longas. Conta no log
% os "Overfull \hbox" do sumário: o sumário é a única parte do documento
% sem parágrafo de corpo, então todo Overfull "in paragraph" antes do
% primeiro \chapter é entrada que passou da margem.
%
% Uso (de uma pasta com livros_crp_acessivel_book.sty e crpsp-leg.sty no
% caminho; as fontes NEWJUNE por OSFONTDIR):
%   lualatex '\def\var{R}\input{sumario-longo}'   % R: rígido (até 0.5.x)
%   lualatex '\def\var{B}\input{sumario-longo}'   % B: plus 4em (o do guia)
%   lualatex '\def\var{C}\input{sumario-longo}'   % C: plus 4em + \tolerance 10000
%   lualatex '\def\var{S}\input{sumario-longo}'   % S: o que a .sty define
% Cada variante só redefine o \l@chapter depois da carga da .sty.
% ===================================================================
\DocumentMetadata{lang=pt-BR, pdfstandard=ua-2, pdfversion=2.0, tagging=on}
\documentclass[11pt,a5paper,twoside,openright]{book}
\usepackage{livros_crp_acessivel_book}
\hypersetup{pdftitle={Sumario longo}, pdfauthor={CRP SP}}
\makeatletter
\providecommand\var{S}
\def\crp@vR{R}\def\crp@vB{B}\def\crp@vC{C}
\def\crp@lchap#1{%  #1 = acréscimo à margem direita (rightskip)
  \renewcommand*\l@chapter[2]{%
    \ifnum \c@tocdepth >\m@ne
        \addpenalty{-\@highpenalty}\addvspace{0.5em \@plus\p@}%
        \setlength\@tempdima{1.5em}%
        \begingroup
            \parindent \z@ \rightskip \@pnumwidth #1
            \parfillskip -\@pnumwidth
            \leavevmode \bfseries
            \advance\leftskip\@tempdima \hskip -\leftskip
            ##1\nobreak\hfil
            \nobreak\hb@xt@\@pnumwidth{\hss ##2\kern-\p@\kern\p@}\par
            \penalty\@highpenalty
        \endgroup
    \fi}}
\ifx\var\crp@vR \crp@lchap{}\fi
\ifx\var\crp@vB \crp@lchap{\@plus 4em}\fi
\ifx\var\crp@vC \crp@lchap{\@plus 4em \tolerance=10000 \hbadness=10000 }\fi
\makeatother
\begin{document}
\frontmatter
\tableofcontents
\mainmatter
\chapter{Deficiência psicossocial e}
\chapter{Deficiência psicossocial e transtorno}
\chapter{Deficiência psicossocial e transtorno do}
\chapter{Deficiência psicossocial e transtorno do espectro}
\chapter{Deficiência psicossocial e transtorno do espectro autista:}
\chapter{Deficiência psicossocial e transtorno do espectro autista: orientações}
\chapter{Deficiência psicossocial e transtorno do espectro autista: orientações para}
\chapter{Deficiência psicossocial e transtorno do espectro autista: orientações para a}
\chapter{Deficiência psicossocial e transtorno do espectro autista: orientações para a realização}
\chapter{Deficiência psicossocial e transtorno do espectro autista: orientações para a realização de}
\chapter{Deficiência psicossocial e transtorno do espectro autista: orientações para a realização de eventos}
\chapter{Deficiência psicossocial e transtorno do espectro autista: orientações para a realização de eventos e}
\chapter{Deficiência psicossocial e transtorno do espectro autista: orientações para a realização de eventos e apresentações}
\chapter{Deficiência psicossocial e transtorno do espectro autista: orientações para a realização de eventos e apresentações acessíveis}
\chapter{Deficiência psicossocial e transtorno do espectro autista: orientações para a realização de eventos e apresentações acessíveis em}
\chapter{Deficiência psicossocial e transtorno do espectro autista: orientações para a realização de eventos e apresentações acessíveis em todo}
\chapter{Comunicação alternativa e}
\chapter{Comunicação alternativa e aumentativa}
\chapter{Comunicação alternativa e aumentativa na}
\chapter{Comunicação alternativa e aumentativa na participação}
\chapter{Comunicação alternativa e aumentativa na participação social}
\chapter{Comunicação alternativa e aumentativa na participação social de}
\chapter{Comunicação alternativa e aumentativa na participação social de pessoas}
\chapter{Comunicação alternativa e aumentativa na participação social de pessoas com}
\chapter{Comunicação alternativa e aumentativa na participação social de pessoas com deficiência}
\chapter{Comunicação alternativa e aumentativa na participação social de pessoas com deficiência intelectual}
\chapter{Comunicação alternativa e aumentativa na participação social de pessoas com deficiência intelectual em}
\chapter{Comunicação alternativa e aumentativa na participação social de pessoas com deficiência intelectual em eventos}
\chapter{Comunicação alternativa e aumentativa na participação social de pessoas com deficiência intelectual em eventos públicos}
\chapter{Comunicação alternativa e aumentativa na participação social de pessoas com deficiência intelectual em eventos públicos de}
\chapter{Comunicação alternativa e aumentativa na participação social de pessoas com deficiência intelectual em eventos públicos de grande}
\chapter{Comunicação alternativa e aumentativa na participação social de pessoas com deficiência intelectual em eventos públicos de grande porte}
\chapter{Psicopedagogicamente inconstitucionalíssimo otorrinolaringologista}
\chapter{Psicopedagogicamente inconstitucionalíssimo otorrinolaringologista socioeducativo}
\chapter{Psicopedagogicamente inconstitucionalíssimo otorrinolaringologista socioeducativo ético-profissional}
\chapter{Psicopedagogicamente inconstitucionalíssimo otorrinolaringologista socioeducativo ético-profissional sociodemográfico}
\chapter{Psicopedagogicamente inconstitucionalíssimo otorrinolaringologista socioeducativo ético-profissional sociodemográfico interdisciplinaridade}
\chapter{Psicopedagogicamente inconstitucionalíssimo otorrinolaringologista socioeducativo ético-profissional sociodemográfico interdisciplinaridade neuropsicológica}
\chapter{Psicopedagogicamente inconstitucionalíssimo otorrinolaringologista socioeducativo ético-profissional sociodemográfico interdisciplinaridade neuropsicológica avaliação}
\chapter{Psicopedagogicamente inconstitucionalíssimo otorrinolaringologista socioeducativo ético-profissional sociodemográfico interdisciplinaridade neuropsicológica avaliação}
\chapter{Psicopedagogicamente inconstitucionalíssimo otorrinolaringologista socioeducativo ético-profissional sociodemográfico interdisciplinaridade neuropsicológica avaliação}
\chapter{Psicopedagogicamente inconstitucionalíssimo otorrinolaringologista socioeducativo ético-profissional sociodemográfico interdisciplinaridade neuropsicológica avaliação}
\chapter{Psicopedagogicamente inconstitucionalíssimo otorrinolaringologista socioeducativo ético-profissional sociodemográfico interdisciplinaridade neuropsicológica avaliação}
\chapter{Psicopedagogicamente inconstitucionalíssimo otorrinolaringologista socioeducativo ético-profissional sociodemográfico interdisciplinaridade neuropsicológica avaliação}
\chapter{Psicopedagogicamente inconstitucionalíssimo otorrinolaringologista socioeducativo ético-profissional sociodemográfico interdisciplinaridade neuropsicológica avaliação}
\chapter{Psicopedagogicamente inconstitucionalíssimo otorrinolaringologista socioeducativo ético-profissional sociodemográfico interdisciplinaridade neuropsicológica avaliação}
\chapter{Psicopedagogicamente inconstitucionalíssimo otorrinolaringologista socioeducativo ético-profissional sociodemográfico interdisciplinaridade neuropsicológica avaliação}
\chapter{Psicopedagogicamente inconstitucionalíssimo otorrinolaringologista socioeducativo ético-profissional sociodemográfico interdisciplinaridade neuropsicológica avaliação}
\end{document}
'   % B: plus 4em (o do guia)
%   lualatex '\def\var{C}% !TeX program = lualatex
% ===================================================================
% sumario-longo.tex — mede o sumário sem hifenização (issue #35, "Fora do
% escopo"): com a hifenização desligada, a entrada longa de capítulo pode
% bater no número da página. Gera 3 x 16 capítulos, cada um com um prefixo
% (de 3 a 18 palavras) de uma frase longa: frase 1 de palavras médias,
% frase 2 com números, frase 3 só de palavras muito longas. Conta no log
% os "Overfull \hbox" do sumário: o sumário é a única parte do documento
% sem parágrafo de corpo, então todo Overfull "in paragraph" antes do
% primeiro \chapter é entrada que passou da margem.
%
% Uso (de uma pasta com livros_crp_acessivel_book.sty e crpsp-leg.sty no
% caminho; as fontes NEWJUNE por OSFONTDIR):
%   lualatex '\def\var{R}\input{sumario-longo}'   % R: rígido (até 0.5.x)
%   lualatex '\def\var{B}\input{sumario-longo}'   % B: plus 4em (o do guia)
%   lualatex '\def\var{C}\input{sumario-longo}'   % C: plus 4em + \tolerance 10000
%   lualatex '\def\var{S}\input{sumario-longo}'   % S: o que a .sty define
% Cada variante só redefine o \l@chapter depois da carga da .sty.
% ===================================================================
\DocumentMetadata{lang=pt-BR, pdfstandard=ua-2, pdfversion=2.0, tagging=on}
\documentclass[11pt,a5paper,twoside,openright]{book}
\usepackage{livros_crp_acessivel_book}
\hypersetup{pdftitle={Sumario longo}, pdfauthor={CRP SP}}
\makeatletter
\providecommand\var{S}
\def\crp@vR{R}\def\crp@vB{B}\def\crp@vC{C}
\def\crp@lchap#1{%  #1 = acréscimo à margem direita (rightskip)
  \renewcommand*\l@chapter[2]{%
    \ifnum \c@tocdepth >\m@ne
        \addpenalty{-\@highpenalty}\addvspace{0.5em \@plus\p@}%
        \setlength\@tempdima{1.5em}%
        \begingroup
            \parindent \z@ \rightskip \@pnumwidth #1
            \parfillskip -\@pnumwidth
            \leavevmode \bfseries
            \advance\leftskip\@tempdima \hskip -\leftskip
            ##1\nobreak\hfil
            \nobreak\hb@xt@\@pnumwidth{\hss ##2\kern-\p@\kern\p@}\par
            \penalty\@highpenalty
        \endgroup
    \fi}}
\ifx\var\crp@vR \crp@lchap{}\fi
\ifx\var\crp@vB \crp@lchap{\@plus 4em}\fi
\ifx\var\crp@vC \crp@lchap{\@plus 4em \tolerance=10000 \hbadness=10000 }\fi
\makeatother
\begin{document}
\frontmatter
\tableofcontents
\mainmatter
\chapter{Deficiência psicossocial e}
\chapter{Deficiência psicossocial e transtorno}
\chapter{Deficiência psicossocial e transtorno do}
\chapter{Deficiência psicossocial e transtorno do espectro}
\chapter{Deficiência psicossocial e transtorno do espectro autista:}
\chapter{Deficiência psicossocial e transtorno do espectro autista: orientações}
\chapter{Deficiência psicossocial e transtorno do espectro autista: orientações para}
\chapter{Deficiência psicossocial e transtorno do espectro autista: orientações para a}
\chapter{Deficiência psicossocial e transtorno do espectro autista: orientações para a realização}
\chapter{Deficiência psicossocial e transtorno do espectro autista: orientações para a realização de}
\chapter{Deficiência psicossocial e transtorno do espectro autista: orientações para a realização de eventos}
\chapter{Deficiência psicossocial e transtorno do espectro autista: orientações para a realização de eventos e}
\chapter{Deficiência psicossocial e transtorno do espectro autista: orientações para a realização de eventos e apresentações}
\chapter{Deficiência psicossocial e transtorno do espectro autista: orientações para a realização de eventos e apresentações acessíveis}
\chapter{Deficiência psicossocial e transtorno do espectro autista: orientações para a realização de eventos e apresentações acessíveis em}
\chapter{Deficiência psicossocial e transtorno do espectro autista: orientações para a realização de eventos e apresentações acessíveis em todo}
\chapter{Comunicação alternativa e}
\chapter{Comunicação alternativa e aumentativa}
\chapter{Comunicação alternativa e aumentativa na}
\chapter{Comunicação alternativa e aumentativa na participação}
\chapter{Comunicação alternativa e aumentativa na participação social}
\chapter{Comunicação alternativa e aumentativa na participação social de}
\chapter{Comunicação alternativa e aumentativa na participação social de pessoas}
\chapter{Comunicação alternativa e aumentativa na participação social de pessoas com}
\chapter{Comunicação alternativa e aumentativa na participação social de pessoas com deficiência}
\chapter{Comunicação alternativa e aumentativa na participação social de pessoas com deficiência intelectual}
\chapter{Comunicação alternativa e aumentativa na participação social de pessoas com deficiência intelectual em}
\chapter{Comunicação alternativa e aumentativa na participação social de pessoas com deficiência intelectual em eventos}
\chapter{Comunicação alternativa e aumentativa na participação social de pessoas com deficiência intelectual em eventos públicos}
\chapter{Comunicação alternativa e aumentativa na participação social de pessoas com deficiência intelectual em eventos públicos de}
\chapter{Comunicação alternativa e aumentativa na participação social de pessoas com deficiência intelectual em eventos públicos de grande}
\chapter{Comunicação alternativa e aumentativa na participação social de pessoas com deficiência intelectual em eventos públicos de grande porte}
\chapter{Psicopedagogicamente inconstitucionalíssimo otorrinolaringologista}
\chapter{Psicopedagogicamente inconstitucionalíssimo otorrinolaringologista socioeducativo}
\chapter{Psicopedagogicamente inconstitucionalíssimo otorrinolaringologista socioeducativo ético-profissional}
\chapter{Psicopedagogicamente inconstitucionalíssimo otorrinolaringologista socioeducativo ético-profissional sociodemográfico}
\chapter{Psicopedagogicamente inconstitucionalíssimo otorrinolaringologista socioeducativo ético-profissional sociodemográfico interdisciplinaridade}
\chapter{Psicopedagogicamente inconstitucionalíssimo otorrinolaringologista socioeducativo ético-profissional sociodemográfico interdisciplinaridade neuropsicológica}
\chapter{Psicopedagogicamente inconstitucionalíssimo otorrinolaringologista socioeducativo ético-profissional sociodemográfico interdisciplinaridade neuropsicológica avaliação}
\chapter{Psicopedagogicamente inconstitucionalíssimo otorrinolaringologista socioeducativo ético-profissional sociodemográfico interdisciplinaridade neuropsicológica avaliação}
\chapter{Psicopedagogicamente inconstitucionalíssimo otorrinolaringologista socioeducativo ético-profissional sociodemográfico interdisciplinaridade neuropsicológica avaliação}
\chapter{Psicopedagogicamente inconstitucionalíssimo otorrinolaringologista socioeducativo ético-profissional sociodemográfico interdisciplinaridade neuropsicológica avaliação}
\chapter{Psicopedagogicamente inconstitucionalíssimo otorrinolaringologista socioeducativo ético-profissional sociodemográfico interdisciplinaridade neuropsicológica avaliação}
\chapter{Psicopedagogicamente inconstitucionalíssimo otorrinolaringologista socioeducativo ético-profissional sociodemográfico interdisciplinaridade neuropsicológica avaliação}
\chapter{Psicopedagogicamente inconstitucionalíssimo otorrinolaringologista socioeducativo ético-profissional sociodemográfico interdisciplinaridade neuropsicológica avaliação}
\chapter{Psicopedagogicamente inconstitucionalíssimo otorrinolaringologista socioeducativo ético-profissional sociodemográfico interdisciplinaridade neuropsicológica avaliação}
\chapter{Psicopedagogicamente inconstitucionalíssimo otorrinolaringologista socioeducativo ético-profissional sociodemográfico interdisciplinaridade neuropsicológica avaliação}
\chapter{Psicopedagogicamente inconstitucionalíssimo otorrinolaringologista socioeducativo ético-profissional sociodemográfico interdisciplinaridade neuropsicológica avaliação}
\end{document}
'   % C: plus 4em + \tolerance 10000
%   lualatex '\def\var{S}% !TeX program = lualatex
% ===================================================================
% sumario-longo.tex — mede o sumário sem hifenização (issue #35, "Fora do
% escopo"): com a hifenização desligada, a entrada longa de capítulo pode
% bater no número da página. Gera 3 x 16 capítulos, cada um com um prefixo
% (de 3 a 18 palavras) de uma frase longa: frase 1 de palavras médias,
% frase 2 com números, frase 3 só de palavras muito longas. Conta no log
% os "Overfull \hbox" do sumário: o sumário é a única parte do documento
% sem parágrafo de corpo, então todo Overfull "in paragraph" antes do
% primeiro \chapter é entrada que passou da margem.
%
% Uso (de uma pasta com livros_crp_acessivel_book.sty e crpsp-leg.sty no
% caminho; as fontes NEWJUNE por OSFONTDIR):
%   lualatex '\def\var{R}\input{sumario-longo}'   % R: rígido (até 0.5.x)
%   lualatex '\def\var{B}\input{sumario-longo}'   % B: plus 4em (o do guia)
%   lualatex '\def\var{C}\input{sumario-longo}'   % C: plus 4em + \tolerance 10000
%   lualatex '\def\var{S}\input{sumario-longo}'   % S: o que a .sty define
% Cada variante só redefine o \l@chapter depois da carga da .sty.
% ===================================================================
\DocumentMetadata{lang=pt-BR, pdfstandard=ua-2, pdfversion=2.0, tagging=on}
\documentclass[11pt,a5paper,twoside,openright]{book}
\usepackage{livros_crp_acessivel_book}
\hypersetup{pdftitle={Sumario longo}, pdfauthor={CRP SP}}
\makeatletter
\providecommand\var{S}
\def\crp@vR{R}\def\crp@vB{B}\def\crp@vC{C}
\def\crp@lchap#1{%  #1 = acréscimo à margem direita (rightskip)
  \renewcommand*\l@chapter[2]{%
    \ifnum \c@tocdepth >\m@ne
        \addpenalty{-\@highpenalty}\addvspace{0.5em \@plus\p@}%
        \setlength\@tempdima{1.5em}%
        \begingroup
            \parindent \z@ \rightskip \@pnumwidth #1
            \parfillskip -\@pnumwidth
            \leavevmode \bfseries
            \advance\leftskip\@tempdima \hskip -\leftskip
            ##1\nobreak\hfil
            \nobreak\hb@xt@\@pnumwidth{\hss ##2\kern-\p@\kern\p@}\par
            \penalty\@highpenalty
        \endgroup
    \fi}}
\ifx\var\crp@vR \crp@lchap{}\fi
\ifx\var\crp@vB \crp@lchap{\@plus 4em}\fi
\ifx\var\crp@vC \crp@lchap{\@plus 4em \tolerance=10000 \hbadness=10000 }\fi
\makeatother
\begin{document}
\frontmatter
\tableofcontents
\mainmatter
\chapter{Deficiência psicossocial e}
\chapter{Deficiência psicossocial e transtorno}
\chapter{Deficiência psicossocial e transtorno do}
\chapter{Deficiência psicossocial e transtorno do espectro}
\chapter{Deficiência psicossocial e transtorno do espectro autista:}
\chapter{Deficiência psicossocial e transtorno do espectro autista: orientações}
\chapter{Deficiência psicossocial e transtorno do espectro autista: orientações para}
\chapter{Deficiência psicossocial e transtorno do espectro autista: orientações para a}
\chapter{Deficiência psicossocial e transtorno do espectro autista: orientações para a realização}
\chapter{Deficiência psicossocial e transtorno do espectro autista: orientações para a realização de}
\chapter{Deficiência psicossocial e transtorno do espectro autista: orientações para a realização de eventos}
\chapter{Deficiência psicossocial e transtorno do espectro autista: orientações para a realização de eventos e}
\chapter{Deficiência psicossocial e transtorno do espectro autista: orientações para a realização de eventos e apresentações}
\chapter{Deficiência psicossocial e transtorno do espectro autista: orientações para a realização de eventos e apresentações acessíveis}
\chapter{Deficiência psicossocial e transtorno do espectro autista: orientações para a realização de eventos e apresentações acessíveis em}
\chapter{Deficiência psicossocial e transtorno do espectro autista: orientações para a realização de eventos e apresentações acessíveis em todo}
\chapter{Comunicação alternativa e}
\chapter{Comunicação alternativa e aumentativa}
\chapter{Comunicação alternativa e aumentativa na}
\chapter{Comunicação alternativa e aumentativa na participação}
\chapter{Comunicação alternativa e aumentativa na participação social}
\chapter{Comunicação alternativa e aumentativa na participação social de}
\chapter{Comunicação alternativa e aumentativa na participação social de pessoas}
\chapter{Comunicação alternativa e aumentativa na participação social de pessoas com}
\chapter{Comunicação alternativa e aumentativa na participação social de pessoas com deficiência}
\chapter{Comunicação alternativa e aumentativa na participação social de pessoas com deficiência intelectual}
\chapter{Comunicação alternativa e aumentativa na participação social de pessoas com deficiência intelectual em}
\chapter{Comunicação alternativa e aumentativa na participação social de pessoas com deficiência intelectual em eventos}
\chapter{Comunicação alternativa e aumentativa na participação social de pessoas com deficiência intelectual em eventos públicos}
\chapter{Comunicação alternativa e aumentativa na participação social de pessoas com deficiência intelectual em eventos públicos de}
\chapter{Comunicação alternativa e aumentativa na participação social de pessoas com deficiência intelectual em eventos públicos de grande}
\chapter{Comunicação alternativa e aumentativa na participação social de pessoas com deficiência intelectual em eventos públicos de grande porte}
\chapter{Psicopedagogicamente inconstitucionalíssimo otorrinolaringologista}
\chapter{Psicopedagogicamente inconstitucionalíssimo otorrinolaringologista socioeducativo}
\chapter{Psicopedagogicamente inconstitucionalíssimo otorrinolaringologista socioeducativo ético-profissional}
\chapter{Psicopedagogicamente inconstitucionalíssimo otorrinolaringologista socioeducativo ético-profissional sociodemográfico}
\chapter{Psicopedagogicamente inconstitucionalíssimo otorrinolaringologista socioeducativo ético-profissional sociodemográfico interdisciplinaridade}
\chapter{Psicopedagogicamente inconstitucionalíssimo otorrinolaringologista socioeducativo ético-profissional sociodemográfico interdisciplinaridade neuropsicológica}
\chapter{Psicopedagogicamente inconstitucionalíssimo otorrinolaringologista socioeducativo ético-profissional sociodemográfico interdisciplinaridade neuropsicológica avaliação}
\chapter{Psicopedagogicamente inconstitucionalíssimo otorrinolaringologista socioeducativo ético-profissional sociodemográfico interdisciplinaridade neuropsicológica avaliação}
\chapter{Psicopedagogicamente inconstitucionalíssimo otorrinolaringologista socioeducativo ético-profissional sociodemográfico interdisciplinaridade neuropsicológica avaliação}
\chapter{Psicopedagogicamente inconstitucionalíssimo otorrinolaringologista socioeducativo ético-profissional sociodemográfico interdisciplinaridade neuropsicológica avaliação}
\chapter{Psicopedagogicamente inconstitucionalíssimo otorrinolaringologista socioeducativo ético-profissional sociodemográfico interdisciplinaridade neuropsicológica avaliação}
\chapter{Psicopedagogicamente inconstitucionalíssimo otorrinolaringologista socioeducativo ético-profissional sociodemográfico interdisciplinaridade neuropsicológica avaliação}
\chapter{Psicopedagogicamente inconstitucionalíssimo otorrinolaringologista socioeducativo ético-profissional sociodemográfico interdisciplinaridade neuropsicológica avaliação}
\chapter{Psicopedagogicamente inconstitucionalíssimo otorrinolaringologista socioeducativo ético-profissional sociodemográfico interdisciplinaridade neuropsicológica avaliação}
\chapter{Psicopedagogicamente inconstitucionalíssimo otorrinolaringologista socioeducativo ético-profissional sociodemográfico interdisciplinaridade neuropsicológica avaliação}
\chapter{Psicopedagogicamente inconstitucionalíssimo otorrinolaringologista socioeducativo ético-profissional sociodemográfico interdisciplinaridade neuropsicológica avaliação}
\end{document}
'   % S: o que a .sty define
% Cada variante só redefine o \l@chapter depois da carga da .sty.
% ===================================================================
\DocumentMetadata{lang=pt-BR, pdfstandard=ua-2, pdfversion=2.0, tagging=on}
\documentclass[11pt,a5paper,twoside,openright]{book}
\usepackage{livros_crp_acessivel_book}
\hypersetup{pdftitle={Sumario longo}, pdfauthor={CRP SP}}
\makeatletter
\providecommand\var{S}
\def\crp@vR{R}\def\crp@vB{B}\def\crp@vC{C}
\def\crp@lchap#1{%  #1 = acréscimo à margem direita (rightskip)
  \renewcommand*\l@chapter[2]{%
    \ifnum \c@tocdepth >\m@ne
        \addpenalty{-\@highpenalty}\addvspace{0.5em \@plus\p@}%
        \setlength\@tempdima{1.5em}%
        \begingroup
            \parindent \z@ \rightskip \@pnumwidth #1
            \parfillskip -\@pnumwidth
            \leavevmode \bfseries
            \advance\leftskip\@tempdima \hskip -\leftskip
            ##1\nobreak\hfil
            \nobreak\hb@xt@\@pnumwidth{\hss ##2\kern-\p@\kern\p@}\par
            \penalty\@highpenalty
        \endgroup
    \fi}}
\ifx\var\crp@vR \crp@lchap{}\fi
\ifx\var\crp@vB \crp@lchap{\@plus 4em}\fi
\ifx\var\crp@vC \crp@lchap{\@plus 4em \tolerance=10000 \hbadness=10000 }\fi
\makeatother
\begin{document}
\frontmatter
\tableofcontents
\mainmatter
\chapter{Deficiência psicossocial e}
\chapter{Deficiência psicossocial e transtorno}
\chapter{Deficiência psicossocial e transtorno do}
\chapter{Deficiência psicossocial e transtorno do espectro}
\chapter{Deficiência psicossocial e transtorno do espectro autista:}
\chapter{Deficiência psicossocial e transtorno do espectro autista: orientações}
\chapter{Deficiência psicossocial e transtorno do espectro autista: orientações para}
\chapter{Deficiência psicossocial e transtorno do espectro autista: orientações para a}
\chapter{Deficiência psicossocial e transtorno do espectro autista: orientações para a realização}
\chapter{Deficiência psicossocial e transtorno do espectro autista: orientações para a realização de}
\chapter{Deficiência psicossocial e transtorno do espectro autista: orientações para a realização de eventos}
\chapter{Deficiência psicossocial e transtorno do espectro autista: orientações para a realização de eventos e}
\chapter{Deficiência psicossocial e transtorno do espectro autista: orientações para a realização de eventos e apresentações}
\chapter{Deficiência psicossocial e transtorno do espectro autista: orientações para a realização de eventos e apresentações acessíveis}
\chapter{Deficiência psicossocial e transtorno do espectro autista: orientações para a realização de eventos e apresentações acessíveis em}
\chapter{Deficiência psicossocial e transtorno do espectro autista: orientações para a realização de eventos e apresentações acessíveis em todo}
\chapter{Comunicação alternativa e}
\chapter{Comunicação alternativa e aumentativa}
\chapter{Comunicação alternativa e aumentativa na}
\chapter{Comunicação alternativa e aumentativa na participação}
\chapter{Comunicação alternativa e aumentativa na participação social}
\chapter{Comunicação alternativa e aumentativa na participação social de}
\chapter{Comunicação alternativa e aumentativa na participação social de pessoas}
\chapter{Comunicação alternativa e aumentativa na participação social de pessoas com}
\chapter{Comunicação alternativa e aumentativa na participação social de pessoas com deficiência}
\chapter{Comunicação alternativa e aumentativa na participação social de pessoas com deficiência intelectual}
\chapter{Comunicação alternativa e aumentativa na participação social de pessoas com deficiência intelectual em}
\chapter{Comunicação alternativa e aumentativa na participação social de pessoas com deficiência intelectual em eventos}
\chapter{Comunicação alternativa e aumentativa na participação social de pessoas com deficiência intelectual em eventos públicos}
\chapter{Comunicação alternativa e aumentativa na participação social de pessoas com deficiência intelectual em eventos públicos de}
\chapter{Comunicação alternativa e aumentativa na participação social de pessoas com deficiência intelectual em eventos públicos de grande}
\chapter{Comunicação alternativa e aumentativa na participação social de pessoas com deficiência intelectual em eventos públicos de grande porte}
\chapter{Psicopedagogicamente inconstitucionalíssimo otorrinolaringologista}
\chapter{Psicopedagogicamente inconstitucionalíssimo otorrinolaringologista socioeducativo}
\chapter{Psicopedagogicamente inconstitucionalíssimo otorrinolaringologista socioeducativo ético-profissional}
\chapter{Psicopedagogicamente inconstitucionalíssimo otorrinolaringologista socioeducativo ético-profissional sociodemográfico}
\chapter{Psicopedagogicamente inconstitucionalíssimo otorrinolaringologista socioeducativo ético-profissional sociodemográfico interdisciplinaridade}
\chapter{Psicopedagogicamente inconstitucionalíssimo otorrinolaringologista socioeducativo ético-profissional sociodemográfico interdisciplinaridade neuropsicológica}
\chapter{Psicopedagogicamente inconstitucionalíssimo otorrinolaringologista socioeducativo ético-profissional sociodemográfico interdisciplinaridade neuropsicológica avaliação}
\chapter{Psicopedagogicamente inconstitucionalíssimo otorrinolaringologista socioeducativo ético-profissional sociodemográfico interdisciplinaridade neuropsicológica avaliação}
\chapter{Psicopedagogicamente inconstitucionalíssimo otorrinolaringologista socioeducativo ético-profissional sociodemográfico interdisciplinaridade neuropsicológica avaliação}
\chapter{Psicopedagogicamente inconstitucionalíssimo otorrinolaringologista socioeducativo ético-profissional sociodemográfico interdisciplinaridade neuropsicológica avaliação}
\chapter{Psicopedagogicamente inconstitucionalíssimo otorrinolaringologista socioeducativo ético-profissional sociodemográfico interdisciplinaridade neuropsicológica avaliação}
\chapter{Psicopedagogicamente inconstitucionalíssimo otorrinolaringologista socioeducativo ético-profissional sociodemográfico interdisciplinaridade neuropsicológica avaliação}
\chapter{Psicopedagogicamente inconstitucionalíssimo otorrinolaringologista socioeducativo ético-profissional sociodemográfico interdisciplinaridade neuropsicológica avaliação}
\chapter{Psicopedagogicamente inconstitucionalíssimo otorrinolaringologista socioeducativo ético-profissional sociodemográfico interdisciplinaridade neuropsicológica avaliação}
\chapter{Psicopedagogicamente inconstitucionalíssimo otorrinolaringologista socioeducativo ético-profissional sociodemográfico interdisciplinaridade neuropsicológica avaliação}
\chapter{Psicopedagogicamente inconstitucionalíssimo otorrinolaringologista socioeducativo ético-profissional sociodemográfico interdisciplinaridade neuropsicológica avaliação}
\end{document}
'   % C: plus 4em + \tolerance 10000
%   lualatex '\def\var{S}% !TeX program = lualatex
% ===================================================================
% sumario-longo.tex — mede o sumário sem hifenização (issue #35, "Fora do
% escopo"): com a hifenização desligada, a entrada longa de capítulo pode
% bater no número da página. Gera 3 x 16 capítulos, cada um com um prefixo
% (de 3 a 18 palavras) de uma frase longa: frase 1 de palavras médias,
% frase 2 com números, frase 3 só de palavras muito longas. Conta no log
% os "Overfull \hbox" do sumário: o sumário é a única parte do documento
% sem parágrafo de corpo, então todo Overfull "in paragraph" antes do
% primeiro \chapter é entrada que passou da margem.
%
% Uso (de uma pasta com livros_crp_acessivel_book.sty e crpsp-leg.sty no
% caminho; as fontes NEWJUNE por OSFONTDIR):
%   lualatex '\def\var{R}% !TeX program = lualatex
% ===================================================================
% sumario-longo.tex — mede o sumário sem hifenização (issue #35, "Fora do
% escopo"): com a hifenização desligada, a entrada longa de capítulo pode
% bater no número da página. Gera 3 x 16 capítulos, cada um com um prefixo
% (de 3 a 18 palavras) de uma frase longa: frase 1 de palavras médias,
% frase 2 com números, frase 3 só de palavras muito longas. Conta no log
% os "Overfull \hbox" do sumário: o sumário é a única parte do documento
% sem parágrafo de corpo, então todo Overfull "in paragraph" antes do
% primeiro \chapter é entrada que passou da margem.
%
% Uso (de uma pasta com livros_crp_acessivel_book.sty e crpsp-leg.sty no
% caminho; as fontes NEWJUNE por OSFONTDIR):
%   lualatex '\def\var{R}\input{sumario-longo}'   % R: rígido (até 0.5.x)
%   lualatex '\def\var{B}\input{sumario-longo}'   % B: plus 4em (o do guia)
%   lualatex '\def\var{C}\input{sumario-longo}'   % C: plus 4em + \tolerance 10000
%   lualatex '\def\var{S}\input{sumario-longo}'   % S: o que a .sty define
% Cada variante só redefine o \l@chapter depois da carga da .sty.
% ===================================================================
\DocumentMetadata{lang=pt-BR, pdfstandard=ua-2, pdfversion=2.0, tagging=on}
\documentclass[11pt,a5paper,twoside,openright]{book}
\usepackage{livros_crp_acessivel_book}
\hypersetup{pdftitle={Sumario longo}, pdfauthor={CRP SP}}
\makeatletter
\providecommand\var{S}
\def\crp@vR{R}\def\crp@vB{B}\def\crp@vC{C}
\def\crp@lchap#1{%  #1 = acréscimo à margem direita (rightskip)
  \renewcommand*\l@chapter[2]{%
    \ifnum \c@tocdepth >\m@ne
        \addpenalty{-\@highpenalty}\addvspace{0.5em \@plus\p@}%
        \setlength\@tempdima{1.5em}%
        \begingroup
            \parindent \z@ \rightskip \@pnumwidth #1
            \parfillskip -\@pnumwidth
            \leavevmode \bfseries
            \advance\leftskip\@tempdima \hskip -\leftskip
            ##1\nobreak\hfil
            \nobreak\hb@xt@\@pnumwidth{\hss ##2\kern-\p@\kern\p@}\par
            \penalty\@highpenalty
        \endgroup
    \fi}}
\ifx\var\crp@vR \crp@lchap{}\fi
\ifx\var\crp@vB \crp@lchap{\@plus 4em}\fi
\ifx\var\crp@vC \crp@lchap{\@plus 4em \tolerance=10000 \hbadness=10000 }\fi
\makeatother
\begin{document}
\frontmatter
\tableofcontents
\mainmatter
\chapter{Deficiência psicossocial e}
\chapter{Deficiência psicossocial e transtorno}
\chapter{Deficiência psicossocial e transtorno do}
\chapter{Deficiência psicossocial e transtorno do espectro}
\chapter{Deficiência psicossocial e transtorno do espectro autista:}
\chapter{Deficiência psicossocial e transtorno do espectro autista: orientações}
\chapter{Deficiência psicossocial e transtorno do espectro autista: orientações para}
\chapter{Deficiência psicossocial e transtorno do espectro autista: orientações para a}
\chapter{Deficiência psicossocial e transtorno do espectro autista: orientações para a realização}
\chapter{Deficiência psicossocial e transtorno do espectro autista: orientações para a realização de}
\chapter{Deficiência psicossocial e transtorno do espectro autista: orientações para a realização de eventos}
\chapter{Deficiência psicossocial e transtorno do espectro autista: orientações para a realização de eventos e}
\chapter{Deficiência psicossocial e transtorno do espectro autista: orientações para a realização de eventos e apresentações}
\chapter{Deficiência psicossocial e transtorno do espectro autista: orientações para a realização de eventos e apresentações acessíveis}
\chapter{Deficiência psicossocial e transtorno do espectro autista: orientações para a realização de eventos e apresentações acessíveis em}
\chapter{Deficiência psicossocial e transtorno do espectro autista: orientações para a realização de eventos e apresentações acessíveis em todo}
\chapter{Comunicação alternativa e}
\chapter{Comunicação alternativa e aumentativa}
\chapter{Comunicação alternativa e aumentativa na}
\chapter{Comunicação alternativa e aumentativa na participação}
\chapter{Comunicação alternativa e aumentativa na participação social}
\chapter{Comunicação alternativa e aumentativa na participação social de}
\chapter{Comunicação alternativa e aumentativa na participação social de pessoas}
\chapter{Comunicação alternativa e aumentativa na participação social de pessoas com}
\chapter{Comunicação alternativa e aumentativa na participação social de pessoas com deficiência}
\chapter{Comunicação alternativa e aumentativa na participação social de pessoas com deficiência intelectual}
\chapter{Comunicação alternativa e aumentativa na participação social de pessoas com deficiência intelectual em}
\chapter{Comunicação alternativa e aumentativa na participação social de pessoas com deficiência intelectual em eventos}
\chapter{Comunicação alternativa e aumentativa na participação social de pessoas com deficiência intelectual em eventos públicos}
\chapter{Comunicação alternativa e aumentativa na participação social de pessoas com deficiência intelectual em eventos públicos de}
\chapter{Comunicação alternativa e aumentativa na participação social de pessoas com deficiência intelectual em eventos públicos de grande}
\chapter{Comunicação alternativa e aumentativa na participação social de pessoas com deficiência intelectual em eventos públicos de grande porte}
\chapter{Psicopedagogicamente inconstitucionalíssimo otorrinolaringologista}
\chapter{Psicopedagogicamente inconstitucionalíssimo otorrinolaringologista socioeducativo}
\chapter{Psicopedagogicamente inconstitucionalíssimo otorrinolaringologista socioeducativo ético-profissional}
\chapter{Psicopedagogicamente inconstitucionalíssimo otorrinolaringologista socioeducativo ético-profissional sociodemográfico}
\chapter{Psicopedagogicamente inconstitucionalíssimo otorrinolaringologista socioeducativo ético-profissional sociodemográfico interdisciplinaridade}
\chapter{Psicopedagogicamente inconstitucionalíssimo otorrinolaringologista socioeducativo ético-profissional sociodemográfico interdisciplinaridade neuropsicológica}
\chapter{Psicopedagogicamente inconstitucionalíssimo otorrinolaringologista socioeducativo ético-profissional sociodemográfico interdisciplinaridade neuropsicológica avaliação}
\chapter{Psicopedagogicamente inconstitucionalíssimo otorrinolaringologista socioeducativo ético-profissional sociodemográfico interdisciplinaridade neuropsicológica avaliação}
\chapter{Psicopedagogicamente inconstitucionalíssimo otorrinolaringologista socioeducativo ético-profissional sociodemográfico interdisciplinaridade neuropsicológica avaliação}
\chapter{Psicopedagogicamente inconstitucionalíssimo otorrinolaringologista socioeducativo ético-profissional sociodemográfico interdisciplinaridade neuropsicológica avaliação}
\chapter{Psicopedagogicamente inconstitucionalíssimo otorrinolaringologista socioeducativo ético-profissional sociodemográfico interdisciplinaridade neuropsicológica avaliação}
\chapter{Psicopedagogicamente inconstitucionalíssimo otorrinolaringologista socioeducativo ético-profissional sociodemográfico interdisciplinaridade neuropsicológica avaliação}
\chapter{Psicopedagogicamente inconstitucionalíssimo otorrinolaringologista socioeducativo ético-profissional sociodemográfico interdisciplinaridade neuropsicológica avaliação}
\chapter{Psicopedagogicamente inconstitucionalíssimo otorrinolaringologista socioeducativo ético-profissional sociodemográfico interdisciplinaridade neuropsicológica avaliação}
\chapter{Psicopedagogicamente inconstitucionalíssimo otorrinolaringologista socioeducativo ético-profissional sociodemográfico interdisciplinaridade neuropsicológica avaliação}
\chapter{Psicopedagogicamente inconstitucionalíssimo otorrinolaringologista socioeducativo ético-profissional sociodemográfico interdisciplinaridade neuropsicológica avaliação}
\end{document}
'   % R: rígido (até 0.5.x)
%   lualatex '\def\var{B}% !TeX program = lualatex
% ===================================================================
% sumario-longo.tex — mede o sumário sem hifenização (issue #35, "Fora do
% escopo"): com a hifenização desligada, a entrada longa de capítulo pode
% bater no número da página. Gera 3 x 16 capítulos, cada um com um prefixo
% (de 3 a 18 palavras) de uma frase longa: frase 1 de palavras médias,
% frase 2 com números, frase 3 só de palavras muito longas. Conta no log
% os "Overfull \hbox" do sumário: o sumário é a única parte do documento
% sem parágrafo de corpo, então todo Overfull "in paragraph" antes do
% primeiro \chapter é entrada que passou da margem.
%
% Uso (de uma pasta com livros_crp_acessivel_book.sty e crpsp-leg.sty no
% caminho; as fontes NEWJUNE por OSFONTDIR):
%   lualatex '\def\var{R}\input{sumario-longo}'   % R: rígido (até 0.5.x)
%   lualatex '\def\var{B}\input{sumario-longo}'   % B: plus 4em (o do guia)
%   lualatex '\def\var{C}\input{sumario-longo}'   % C: plus 4em + \tolerance 10000
%   lualatex '\def\var{S}\input{sumario-longo}'   % S: o que a .sty define
% Cada variante só redefine o \l@chapter depois da carga da .sty.
% ===================================================================
\DocumentMetadata{lang=pt-BR, pdfstandard=ua-2, pdfversion=2.0, tagging=on}
\documentclass[11pt,a5paper,twoside,openright]{book}
\usepackage{livros_crp_acessivel_book}
\hypersetup{pdftitle={Sumario longo}, pdfauthor={CRP SP}}
\makeatletter
\providecommand\var{S}
\def\crp@vR{R}\def\crp@vB{B}\def\crp@vC{C}
\def\crp@lchap#1{%  #1 = acréscimo à margem direita (rightskip)
  \renewcommand*\l@chapter[2]{%
    \ifnum \c@tocdepth >\m@ne
        \addpenalty{-\@highpenalty}\addvspace{0.5em \@plus\p@}%
        \setlength\@tempdima{1.5em}%
        \begingroup
            \parindent \z@ \rightskip \@pnumwidth #1
            \parfillskip -\@pnumwidth
            \leavevmode \bfseries
            \advance\leftskip\@tempdima \hskip -\leftskip
            ##1\nobreak\hfil
            \nobreak\hb@xt@\@pnumwidth{\hss ##2\kern-\p@\kern\p@}\par
            \penalty\@highpenalty
        \endgroup
    \fi}}
\ifx\var\crp@vR \crp@lchap{}\fi
\ifx\var\crp@vB \crp@lchap{\@plus 4em}\fi
\ifx\var\crp@vC \crp@lchap{\@plus 4em \tolerance=10000 \hbadness=10000 }\fi
\makeatother
\begin{document}
\frontmatter
\tableofcontents
\mainmatter
\chapter{Deficiência psicossocial e}
\chapter{Deficiência psicossocial e transtorno}
\chapter{Deficiência psicossocial e transtorno do}
\chapter{Deficiência psicossocial e transtorno do espectro}
\chapter{Deficiência psicossocial e transtorno do espectro autista:}
\chapter{Deficiência psicossocial e transtorno do espectro autista: orientações}
\chapter{Deficiência psicossocial e transtorno do espectro autista: orientações para}
\chapter{Deficiência psicossocial e transtorno do espectro autista: orientações para a}
\chapter{Deficiência psicossocial e transtorno do espectro autista: orientações para a realização}
\chapter{Deficiência psicossocial e transtorno do espectro autista: orientações para a realização de}
\chapter{Deficiência psicossocial e transtorno do espectro autista: orientações para a realização de eventos}
\chapter{Deficiência psicossocial e transtorno do espectro autista: orientações para a realização de eventos e}
\chapter{Deficiência psicossocial e transtorno do espectro autista: orientações para a realização de eventos e apresentações}
\chapter{Deficiência psicossocial e transtorno do espectro autista: orientações para a realização de eventos e apresentações acessíveis}
\chapter{Deficiência psicossocial e transtorno do espectro autista: orientações para a realização de eventos e apresentações acessíveis em}
\chapter{Deficiência psicossocial e transtorno do espectro autista: orientações para a realização de eventos e apresentações acessíveis em todo}
\chapter{Comunicação alternativa e}
\chapter{Comunicação alternativa e aumentativa}
\chapter{Comunicação alternativa e aumentativa na}
\chapter{Comunicação alternativa e aumentativa na participação}
\chapter{Comunicação alternativa e aumentativa na participação social}
\chapter{Comunicação alternativa e aumentativa na participação social de}
\chapter{Comunicação alternativa e aumentativa na participação social de pessoas}
\chapter{Comunicação alternativa e aumentativa na participação social de pessoas com}
\chapter{Comunicação alternativa e aumentativa na participação social de pessoas com deficiência}
\chapter{Comunicação alternativa e aumentativa na participação social de pessoas com deficiência intelectual}
\chapter{Comunicação alternativa e aumentativa na participação social de pessoas com deficiência intelectual em}
\chapter{Comunicação alternativa e aumentativa na participação social de pessoas com deficiência intelectual em eventos}
\chapter{Comunicação alternativa e aumentativa na participação social de pessoas com deficiência intelectual em eventos públicos}
\chapter{Comunicação alternativa e aumentativa na participação social de pessoas com deficiência intelectual em eventos públicos de}
\chapter{Comunicação alternativa e aumentativa na participação social de pessoas com deficiência intelectual em eventos públicos de grande}
\chapter{Comunicação alternativa e aumentativa na participação social de pessoas com deficiência intelectual em eventos públicos de grande porte}
\chapter{Psicopedagogicamente inconstitucionalíssimo otorrinolaringologista}
\chapter{Psicopedagogicamente inconstitucionalíssimo otorrinolaringologista socioeducativo}
\chapter{Psicopedagogicamente inconstitucionalíssimo otorrinolaringologista socioeducativo ético-profissional}
\chapter{Psicopedagogicamente inconstitucionalíssimo otorrinolaringologista socioeducativo ético-profissional sociodemográfico}
\chapter{Psicopedagogicamente inconstitucionalíssimo otorrinolaringologista socioeducativo ético-profissional sociodemográfico interdisciplinaridade}
\chapter{Psicopedagogicamente inconstitucionalíssimo otorrinolaringologista socioeducativo ético-profissional sociodemográfico interdisciplinaridade neuropsicológica}
\chapter{Psicopedagogicamente inconstitucionalíssimo otorrinolaringologista socioeducativo ético-profissional sociodemográfico interdisciplinaridade neuropsicológica avaliação}
\chapter{Psicopedagogicamente inconstitucionalíssimo otorrinolaringologista socioeducativo ético-profissional sociodemográfico interdisciplinaridade neuropsicológica avaliação}
\chapter{Psicopedagogicamente inconstitucionalíssimo otorrinolaringologista socioeducativo ético-profissional sociodemográfico interdisciplinaridade neuropsicológica avaliação}
\chapter{Psicopedagogicamente inconstitucionalíssimo otorrinolaringologista socioeducativo ético-profissional sociodemográfico interdisciplinaridade neuropsicológica avaliação}
\chapter{Psicopedagogicamente inconstitucionalíssimo otorrinolaringologista socioeducativo ético-profissional sociodemográfico interdisciplinaridade neuropsicológica avaliação}
\chapter{Psicopedagogicamente inconstitucionalíssimo otorrinolaringologista socioeducativo ético-profissional sociodemográfico interdisciplinaridade neuropsicológica avaliação}
\chapter{Psicopedagogicamente inconstitucionalíssimo otorrinolaringologista socioeducativo ético-profissional sociodemográfico interdisciplinaridade neuropsicológica avaliação}
\chapter{Psicopedagogicamente inconstitucionalíssimo otorrinolaringologista socioeducativo ético-profissional sociodemográfico interdisciplinaridade neuropsicológica avaliação}
\chapter{Psicopedagogicamente inconstitucionalíssimo otorrinolaringologista socioeducativo ético-profissional sociodemográfico interdisciplinaridade neuropsicológica avaliação}
\chapter{Psicopedagogicamente inconstitucionalíssimo otorrinolaringologista socioeducativo ético-profissional sociodemográfico interdisciplinaridade neuropsicológica avaliação}
\end{document}
'   % B: plus 4em (o do guia)
%   lualatex '\def\var{C}% !TeX program = lualatex
% ===================================================================
% sumario-longo.tex — mede o sumário sem hifenização (issue #35, "Fora do
% escopo"): com a hifenização desligada, a entrada longa de capítulo pode
% bater no número da página. Gera 3 x 16 capítulos, cada um com um prefixo
% (de 3 a 18 palavras) de uma frase longa: frase 1 de palavras médias,
% frase 2 com números, frase 3 só de palavras muito longas. Conta no log
% os "Overfull \hbox" do sumário: o sumário é a única parte do documento
% sem parágrafo de corpo, então todo Overfull "in paragraph" antes do
% primeiro \chapter é entrada que passou da margem.
%
% Uso (de uma pasta com livros_crp_acessivel_book.sty e crpsp-leg.sty no
% caminho; as fontes NEWJUNE por OSFONTDIR):
%   lualatex '\def\var{R}\input{sumario-longo}'   % R: rígido (até 0.5.x)
%   lualatex '\def\var{B}\input{sumario-longo}'   % B: plus 4em (o do guia)
%   lualatex '\def\var{C}\input{sumario-longo}'   % C: plus 4em + \tolerance 10000
%   lualatex '\def\var{S}\input{sumario-longo}'   % S: o que a .sty define
% Cada variante só redefine o \l@chapter depois da carga da .sty.
% ===================================================================
\DocumentMetadata{lang=pt-BR, pdfstandard=ua-2, pdfversion=2.0, tagging=on}
\documentclass[11pt,a5paper,twoside,openright]{book}
\usepackage{livros_crp_acessivel_book}
\hypersetup{pdftitle={Sumario longo}, pdfauthor={CRP SP}}
\makeatletter
\providecommand\var{S}
\def\crp@vR{R}\def\crp@vB{B}\def\crp@vC{C}
\def\crp@lchap#1{%  #1 = acréscimo à margem direita (rightskip)
  \renewcommand*\l@chapter[2]{%
    \ifnum \c@tocdepth >\m@ne
        \addpenalty{-\@highpenalty}\addvspace{0.5em \@plus\p@}%
        \setlength\@tempdima{1.5em}%
        \begingroup
            \parindent \z@ \rightskip \@pnumwidth #1
            \parfillskip -\@pnumwidth
            \leavevmode \bfseries
            \advance\leftskip\@tempdima \hskip -\leftskip
            ##1\nobreak\hfil
            \nobreak\hb@xt@\@pnumwidth{\hss ##2\kern-\p@\kern\p@}\par
            \penalty\@highpenalty
        \endgroup
    \fi}}
\ifx\var\crp@vR \crp@lchap{}\fi
\ifx\var\crp@vB \crp@lchap{\@plus 4em}\fi
\ifx\var\crp@vC \crp@lchap{\@plus 4em \tolerance=10000 \hbadness=10000 }\fi
\makeatother
\begin{document}
\frontmatter
\tableofcontents
\mainmatter
\chapter{Deficiência psicossocial e}
\chapter{Deficiência psicossocial e transtorno}
\chapter{Deficiência psicossocial e transtorno do}
\chapter{Deficiência psicossocial e transtorno do espectro}
\chapter{Deficiência psicossocial e transtorno do espectro autista:}
\chapter{Deficiência psicossocial e transtorno do espectro autista: orientações}
\chapter{Deficiência psicossocial e transtorno do espectro autista: orientações para}
\chapter{Deficiência psicossocial e transtorno do espectro autista: orientações para a}
\chapter{Deficiência psicossocial e transtorno do espectro autista: orientações para a realização}
\chapter{Deficiência psicossocial e transtorno do espectro autista: orientações para a realização de}
\chapter{Deficiência psicossocial e transtorno do espectro autista: orientações para a realização de eventos}
\chapter{Deficiência psicossocial e transtorno do espectro autista: orientações para a realização de eventos e}
\chapter{Deficiência psicossocial e transtorno do espectro autista: orientações para a realização de eventos e apresentações}
\chapter{Deficiência psicossocial e transtorno do espectro autista: orientações para a realização de eventos e apresentações acessíveis}
\chapter{Deficiência psicossocial e transtorno do espectro autista: orientações para a realização de eventos e apresentações acessíveis em}
\chapter{Deficiência psicossocial e transtorno do espectro autista: orientações para a realização de eventos e apresentações acessíveis em todo}
\chapter{Comunicação alternativa e}
\chapter{Comunicação alternativa e aumentativa}
\chapter{Comunicação alternativa e aumentativa na}
\chapter{Comunicação alternativa e aumentativa na participação}
\chapter{Comunicação alternativa e aumentativa na participação social}
\chapter{Comunicação alternativa e aumentativa na participação social de}
\chapter{Comunicação alternativa e aumentativa na participação social de pessoas}
\chapter{Comunicação alternativa e aumentativa na participação social de pessoas com}
\chapter{Comunicação alternativa e aumentativa na participação social de pessoas com deficiência}
\chapter{Comunicação alternativa e aumentativa na participação social de pessoas com deficiência intelectual}
\chapter{Comunicação alternativa e aumentativa na participação social de pessoas com deficiência intelectual em}
\chapter{Comunicação alternativa e aumentativa na participação social de pessoas com deficiência intelectual em eventos}
\chapter{Comunicação alternativa e aumentativa na participação social de pessoas com deficiência intelectual em eventos públicos}
\chapter{Comunicação alternativa e aumentativa na participação social de pessoas com deficiência intelectual em eventos públicos de}
\chapter{Comunicação alternativa e aumentativa na participação social de pessoas com deficiência intelectual em eventos públicos de grande}
\chapter{Comunicação alternativa e aumentativa na participação social de pessoas com deficiência intelectual em eventos públicos de grande porte}
\chapter{Psicopedagogicamente inconstitucionalíssimo otorrinolaringologista}
\chapter{Psicopedagogicamente inconstitucionalíssimo otorrinolaringologista socioeducativo}
\chapter{Psicopedagogicamente inconstitucionalíssimo otorrinolaringologista socioeducativo ético-profissional}
\chapter{Psicopedagogicamente inconstitucionalíssimo otorrinolaringologista socioeducativo ético-profissional sociodemográfico}
\chapter{Psicopedagogicamente inconstitucionalíssimo otorrinolaringologista socioeducativo ético-profissional sociodemográfico interdisciplinaridade}
\chapter{Psicopedagogicamente inconstitucionalíssimo otorrinolaringologista socioeducativo ético-profissional sociodemográfico interdisciplinaridade neuropsicológica}
\chapter{Psicopedagogicamente inconstitucionalíssimo otorrinolaringologista socioeducativo ético-profissional sociodemográfico interdisciplinaridade neuropsicológica avaliação}
\chapter{Psicopedagogicamente inconstitucionalíssimo otorrinolaringologista socioeducativo ético-profissional sociodemográfico interdisciplinaridade neuropsicológica avaliação}
\chapter{Psicopedagogicamente inconstitucionalíssimo otorrinolaringologista socioeducativo ético-profissional sociodemográfico interdisciplinaridade neuropsicológica avaliação}
\chapter{Psicopedagogicamente inconstitucionalíssimo otorrinolaringologista socioeducativo ético-profissional sociodemográfico interdisciplinaridade neuropsicológica avaliação}
\chapter{Psicopedagogicamente inconstitucionalíssimo otorrinolaringologista socioeducativo ético-profissional sociodemográfico interdisciplinaridade neuropsicológica avaliação}
\chapter{Psicopedagogicamente inconstitucionalíssimo otorrinolaringologista socioeducativo ético-profissional sociodemográfico interdisciplinaridade neuropsicológica avaliação}
\chapter{Psicopedagogicamente inconstitucionalíssimo otorrinolaringologista socioeducativo ético-profissional sociodemográfico interdisciplinaridade neuropsicológica avaliação}
\chapter{Psicopedagogicamente inconstitucionalíssimo otorrinolaringologista socioeducativo ético-profissional sociodemográfico interdisciplinaridade neuropsicológica avaliação}
\chapter{Psicopedagogicamente inconstitucionalíssimo otorrinolaringologista socioeducativo ético-profissional sociodemográfico interdisciplinaridade neuropsicológica avaliação}
\chapter{Psicopedagogicamente inconstitucionalíssimo otorrinolaringologista socioeducativo ético-profissional sociodemográfico interdisciplinaridade neuropsicológica avaliação}
\end{document}
'   % C: plus 4em + \tolerance 10000
%   lualatex '\def\var{S}% !TeX program = lualatex
% ===================================================================
% sumario-longo.tex — mede o sumário sem hifenização (issue #35, "Fora do
% escopo"): com a hifenização desligada, a entrada longa de capítulo pode
% bater no número da página. Gera 3 x 16 capítulos, cada um com um prefixo
% (de 3 a 18 palavras) de uma frase longa: frase 1 de palavras médias,
% frase 2 com números, frase 3 só de palavras muito longas. Conta no log
% os "Overfull \hbox" do sumário: o sumário é a única parte do documento
% sem parágrafo de corpo, então todo Overfull "in paragraph" antes do
% primeiro \chapter é entrada que passou da margem.
%
% Uso (de uma pasta com livros_crp_acessivel_book.sty e crpsp-leg.sty no
% caminho; as fontes NEWJUNE por OSFONTDIR):
%   lualatex '\def\var{R}\input{sumario-longo}'   % R: rígido (até 0.5.x)
%   lualatex '\def\var{B}\input{sumario-longo}'   % B: plus 4em (o do guia)
%   lualatex '\def\var{C}\input{sumario-longo}'   % C: plus 4em + \tolerance 10000
%   lualatex '\def\var{S}\input{sumario-longo}'   % S: o que a .sty define
% Cada variante só redefine o \l@chapter depois da carga da .sty.
% ===================================================================
\DocumentMetadata{lang=pt-BR, pdfstandard=ua-2, pdfversion=2.0, tagging=on}
\documentclass[11pt,a5paper,twoside,openright]{book}
\usepackage{livros_crp_acessivel_book}
\hypersetup{pdftitle={Sumario longo}, pdfauthor={CRP SP}}
\makeatletter
\providecommand\var{S}
\def\crp@vR{R}\def\crp@vB{B}\def\crp@vC{C}
\def\crp@lchap#1{%  #1 = acréscimo à margem direita (rightskip)
  \renewcommand*\l@chapter[2]{%
    \ifnum \c@tocdepth >\m@ne
        \addpenalty{-\@highpenalty}\addvspace{0.5em \@plus\p@}%
        \setlength\@tempdima{1.5em}%
        \begingroup
            \parindent \z@ \rightskip \@pnumwidth #1
            \parfillskip -\@pnumwidth
            \leavevmode \bfseries
            \advance\leftskip\@tempdima \hskip -\leftskip
            ##1\nobreak\hfil
            \nobreak\hb@xt@\@pnumwidth{\hss ##2\kern-\p@\kern\p@}\par
            \penalty\@highpenalty
        \endgroup
    \fi}}
\ifx\var\crp@vR \crp@lchap{}\fi
\ifx\var\crp@vB \crp@lchap{\@plus 4em}\fi
\ifx\var\crp@vC \crp@lchap{\@plus 4em \tolerance=10000 \hbadness=10000 }\fi
\makeatother
\begin{document}
\frontmatter
\tableofcontents
\mainmatter
\chapter{Deficiência psicossocial e}
\chapter{Deficiência psicossocial e transtorno}
\chapter{Deficiência psicossocial e transtorno do}
\chapter{Deficiência psicossocial e transtorno do espectro}
\chapter{Deficiência psicossocial e transtorno do espectro autista:}
\chapter{Deficiência psicossocial e transtorno do espectro autista: orientações}
\chapter{Deficiência psicossocial e transtorno do espectro autista: orientações para}
\chapter{Deficiência psicossocial e transtorno do espectro autista: orientações para a}
\chapter{Deficiência psicossocial e transtorno do espectro autista: orientações para a realização}
\chapter{Deficiência psicossocial e transtorno do espectro autista: orientações para a realização de}
\chapter{Deficiência psicossocial e transtorno do espectro autista: orientações para a realização de eventos}
\chapter{Deficiência psicossocial e transtorno do espectro autista: orientações para a realização de eventos e}
\chapter{Deficiência psicossocial e transtorno do espectro autista: orientações para a realização de eventos e apresentações}
\chapter{Deficiência psicossocial e transtorno do espectro autista: orientações para a realização de eventos e apresentações acessíveis}
\chapter{Deficiência psicossocial e transtorno do espectro autista: orientações para a realização de eventos e apresentações acessíveis em}
\chapter{Deficiência psicossocial e transtorno do espectro autista: orientações para a realização de eventos e apresentações acessíveis em todo}
\chapter{Comunicação alternativa e}
\chapter{Comunicação alternativa e aumentativa}
\chapter{Comunicação alternativa e aumentativa na}
\chapter{Comunicação alternativa e aumentativa na participação}
\chapter{Comunicação alternativa e aumentativa na participação social}
\chapter{Comunicação alternativa e aumentativa na participação social de}
\chapter{Comunicação alternativa e aumentativa na participação social de pessoas}
\chapter{Comunicação alternativa e aumentativa na participação social de pessoas com}
\chapter{Comunicação alternativa e aumentativa na participação social de pessoas com deficiência}
\chapter{Comunicação alternativa e aumentativa na participação social de pessoas com deficiência intelectual}
\chapter{Comunicação alternativa e aumentativa na participação social de pessoas com deficiência intelectual em}
\chapter{Comunicação alternativa e aumentativa na participação social de pessoas com deficiência intelectual em eventos}
\chapter{Comunicação alternativa e aumentativa na participação social de pessoas com deficiência intelectual em eventos públicos}
\chapter{Comunicação alternativa e aumentativa na participação social de pessoas com deficiência intelectual em eventos públicos de}
\chapter{Comunicação alternativa e aumentativa na participação social de pessoas com deficiência intelectual em eventos públicos de grande}
\chapter{Comunicação alternativa e aumentativa na participação social de pessoas com deficiência intelectual em eventos públicos de grande porte}
\chapter{Psicopedagogicamente inconstitucionalíssimo otorrinolaringologista}
\chapter{Psicopedagogicamente inconstitucionalíssimo otorrinolaringologista socioeducativo}
\chapter{Psicopedagogicamente inconstitucionalíssimo otorrinolaringologista socioeducativo ético-profissional}
\chapter{Psicopedagogicamente inconstitucionalíssimo otorrinolaringologista socioeducativo ético-profissional sociodemográfico}
\chapter{Psicopedagogicamente inconstitucionalíssimo otorrinolaringologista socioeducativo ético-profissional sociodemográfico interdisciplinaridade}
\chapter{Psicopedagogicamente inconstitucionalíssimo otorrinolaringologista socioeducativo ético-profissional sociodemográfico interdisciplinaridade neuropsicológica}
\chapter{Psicopedagogicamente inconstitucionalíssimo otorrinolaringologista socioeducativo ético-profissional sociodemográfico interdisciplinaridade neuropsicológica avaliação}
\chapter{Psicopedagogicamente inconstitucionalíssimo otorrinolaringologista socioeducativo ético-profissional sociodemográfico interdisciplinaridade neuropsicológica avaliação}
\chapter{Psicopedagogicamente inconstitucionalíssimo otorrinolaringologista socioeducativo ético-profissional sociodemográfico interdisciplinaridade neuropsicológica avaliação}
\chapter{Psicopedagogicamente inconstitucionalíssimo otorrinolaringologista socioeducativo ético-profissional sociodemográfico interdisciplinaridade neuropsicológica avaliação}
\chapter{Psicopedagogicamente inconstitucionalíssimo otorrinolaringologista socioeducativo ético-profissional sociodemográfico interdisciplinaridade neuropsicológica avaliação}
\chapter{Psicopedagogicamente inconstitucionalíssimo otorrinolaringologista socioeducativo ético-profissional sociodemográfico interdisciplinaridade neuropsicológica avaliação}
\chapter{Psicopedagogicamente inconstitucionalíssimo otorrinolaringologista socioeducativo ético-profissional sociodemográfico interdisciplinaridade neuropsicológica avaliação}
\chapter{Psicopedagogicamente inconstitucionalíssimo otorrinolaringologista socioeducativo ético-profissional sociodemográfico interdisciplinaridade neuropsicológica avaliação}
\chapter{Psicopedagogicamente inconstitucionalíssimo otorrinolaringologista socioeducativo ético-profissional sociodemográfico interdisciplinaridade neuropsicológica avaliação}
\chapter{Psicopedagogicamente inconstitucionalíssimo otorrinolaringologista socioeducativo ético-profissional sociodemográfico interdisciplinaridade neuropsicológica avaliação}
\end{document}
'   % S: o que a .sty define
% Cada variante só redefine o \l@chapter depois da carga da .sty.
% ===================================================================
\DocumentMetadata{lang=pt-BR, pdfstandard=ua-2, pdfversion=2.0, tagging=on}
\documentclass[11pt,a5paper,twoside,openright]{book}
\usepackage{livros_crp_acessivel_book}
\hypersetup{pdftitle={Sumario longo}, pdfauthor={CRP SP}}
\makeatletter
\providecommand\var{S}
\def\crp@vR{R}\def\crp@vB{B}\def\crp@vC{C}
\def\crp@lchap#1{%  #1 = acréscimo à margem direita (rightskip)
  \renewcommand*\l@chapter[2]{%
    \ifnum \c@tocdepth >\m@ne
        \addpenalty{-\@highpenalty}\addvspace{0.5em \@plus\p@}%
        \setlength\@tempdima{1.5em}%
        \begingroup
            \parindent \z@ \rightskip \@pnumwidth #1
            \parfillskip -\@pnumwidth
            \leavevmode \bfseries
            \advance\leftskip\@tempdima \hskip -\leftskip
            ##1\nobreak\hfil
            \nobreak\hb@xt@\@pnumwidth{\hss ##2\kern-\p@\kern\p@}\par
            \penalty\@highpenalty
        \endgroup
    \fi}}
\ifx\var\crp@vR \crp@lchap{}\fi
\ifx\var\crp@vB \crp@lchap{\@plus 4em}\fi
\ifx\var\crp@vC \crp@lchap{\@plus 4em \tolerance=10000 \hbadness=10000 }\fi
\makeatother
\begin{document}
\frontmatter
\tableofcontents
\mainmatter
\chapter{Deficiência psicossocial e}
\chapter{Deficiência psicossocial e transtorno}
\chapter{Deficiência psicossocial e transtorno do}
\chapter{Deficiência psicossocial e transtorno do espectro}
\chapter{Deficiência psicossocial e transtorno do espectro autista:}
\chapter{Deficiência psicossocial e transtorno do espectro autista: orientações}
\chapter{Deficiência psicossocial e transtorno do espectro autista: orientações para}
\chapter{Deficiência psicossocial e transtorno do espectro autista: orientações para a}
\chapter{Deficiência psicossocial e transtorno do espectro autista: orientações para a realização}
\chapter{Deficiência psicossocial e transtorno do espectro autista: orientações para a realização de}
\chapter{Deficiência psicossocial e transtorno do espectro autista: orientações para a realização de eventos}
\chapter{Deficiência psicossocial e transtorno do espectro autista: orientações para a realização de eventos e}
\chapter{Deficiência psicossocial e transtorno do espectro autista: orientações para a realização de eventos e apresentações}
\chapter{Deficiência psicossocial e transtorno do espectro autista: orientações para a realização de eventos e apresentações acessíveis}
\chapter{Deficiência psicossocial e transtorno do espectro autista: orientações para a realização de eventos e apresentações acessíveis em}
\chapter{Deficiência psicossocial e transtorno do espectro autista: orientações para a realização de eventos e apresentações acessíveis em todo}
\chapter{Comunicação alternativa e}
\chapter{Comunicação alternativa e aumentativa}
\chapter{Comunicação alternativa e aumentativa na}
\chapter{Comunicação alternativa e aumentativa na participação}
\chapter{Comunicação alternativa e aumentativa na participação social}
\chapter{Comunicação alternativa e aumentativa na participação social de}
\chapter{Comunicação alternativa e aumentativa na participação social de pessoas}
\chapter{Comunicação alternativa e aumentativa na participação social de pessoas com}
\chapter{Comunicação alternativa e aumentativa na participação social de pessoas com deficiência}
\chapter{Comunicação alternativa e aumentativa na participação social de pessoas com deficiência intelectual}
\chapter{Comunicação alternativa e aumentativa na participação social de pessoas com deficiência intelectual em}
\chapter{Comunicação alternativa e aumentativa na participação social de pessoas com deficiência intelectual em eventos}
\chapter{Comunicação alternativa e aumentativa na participação social de pessoas com deficiência intelectual em eventos públicos}
\chapter{Comunicação alternativa e aumentativa na participação social de pessoas com deficiência intelectual em eventos públicos de}
\chapter{Comunicação alternativa e aumentativa na participação social de pessoas com deficiência intelectual em eventos públicos de grande}
\chapter{Comunicação alternativa e aumentativa na participação social de pessoas com deficiência intelectual em eventos públicos de grande porte}
\chapter{Psicopedagogicamente inconstitucionalíssimo otorrinolaringologista}
\chapter{Psicopedagogicamente inconstitucionalíssimo otorrinolaringologista socioeducativo}
\chapter{Psicopedagogicamente inconstitucionalíssimo otorrinolaringologista socioeducativo ético-profissional}
\chapter{Psicopedagogicamente inconstitucionalíssimo otorrinolaringologista socioeducativo ético-profissional sociodemográfico}
\chapter{Psicopedagogicamente inconstitucionalíssimo otorrinolaringologista socioeducativo ético-profissional sociodemográfico interdisciplinaridade}
\chapter{Psicopedagogicamente inconstitucionalíssimo otorrinolaringologista socioeducativo ético-profissional sociodemográfico interdisciplinaridade neuropsicológica}
\chapter{Psicopedagogicamente inconstitucionalíssimo otorrinolaringologista socioeducativo ético-profissional sociodemográfico interdisciplinaridade neuropsicológica avaliação}
\chapter{Psicopedagogicamente inconstitucionalíssimo otorrinolaringologista socioeducativo ético-profissional sociodemográfico interdisciplinaridade neuropsicológica avaliação}
\chapter{Psicopedagogicamente inconstitucionalíssimo otorrinolaringologista socioeducativo ético-profissional sociodemográfico interdisciplinaridade neuropsicológica avaliação}
\chapter{Psicopedagogicamente inconstitucionalíssimo otorrinolaringologista socioeducativo ético-profissional sociodemográfico interdisciplinaridade neuropsicológica avaliação}
\chapter{Psicopedagogicamente inconstitucionalíssimo otorrinolaringologista socioeducativo ético-profissional sociodemográfico interdisciplinaridade neuropsicológica avaliação}
\chapter{Psicopedagogicamente inconstitucionalíssimo otorrinolaringologista socioeducativo ético-profissional sociodemográfico interdisciplinaridade neuropsicológica avaliação}
\chapter{Psicopedagogicamente inconstitucionalíssimo otorrinolaringologista socioeducativo ético-profissional sociodemográfico interdisciplinaridade neuropsicológica avaliação}
\chapter{Psicopedagogicamente inconstitucionalíssimo otorrinolaringologista socioeducativo ético-profissional sociodemográfico interdisciplinaridade neuropsicológica avaliação}
\chapter{Psicopedagogicamente inconstitucionalíssimo otorrinolaringologista socioeducativo ético-profissional sociodemográfico interdisciplinaridade neuropsicológica avaliação}
\chapter{Psicopedagogicamente inconstitucionalíssimo otorrinolaringologista socioeducativo ético-profissional sociodemográfico interdisciplinaridade neuropsicológica avaliação}
\end{document}
'   % S: o que a .sty define
% Cada variante só redefine o \l@chapter depois da carga da .sty.
% ===================================================================
\DocumentMetadata{lang=pt-BR, pdfstandard=ua-2, pdfversion=2.0, tagging=on}
\documentclass[11pt,a5paper,twoside,openright]{book}
\usepackage{livros_crp_acessivel_book}
\hypersetup{pdftitle={Sumario longo}, pdfauthor={CRP SP}}
\makeatletter
\providecommand\var{S}
\def\crp@vR{R}\def\crp@vB{B}\def\crp@vC{C}
\def\crp@lchap#1{%  #1 = acréscimo à margem direita (rightskip)
  \renewcommand*\l@chapter[2]{%
    \ifnum \c@tocdepth >\m@ne
        \addpenalty{-\@highpenalty}\addvspace{0.5em \@plus\p@}%
        \setlength\@tempdima{1.5em}%
        \begingroup
            \parindent \z@ \rightskip \@pnumwidth #1
            \parfillskip -\@pnumwidth
            \leavevmode \bfseries
            \advance\leftskip\@tempdima \hskip -\leftskip
            ##1\nobreak\hfil
            \nobreak\hb@xt@\@pnumwidth{\hss ##2\kern-\p@\kern\p@}\par
            \penalty\@highpenalty
        \endgroup
    \fi}}
\ifx\var\crp@vR \crp@lchap{}\fi
\ifx\var\crp@vB \crp@lchap{\@plus 4em}\fi
\ifx\var\crp@vC \crp@lchap{\@plus 4em \tolerance=10000 \hbadness=10000 }\fi
\makeatother
\begin{document}
\frontmatter
\tableofcontents
\mainmatter
\chapter{Deficiência psicossocial e}
\chapter{Deficiência psicossocial e transtorno}
\chapter{Deficiência psicossocial e transtorno do}
\chapter{Deficiência psicossocial e transtorno do espectro}
\chapter{Deficiência psicossocial e transtorno do espectro autista:}
\chapter{Deficiência psicossocial e transtorno do espectro autista: orientações}
\chapter{Deficiência psicossocial e transtorno do espectro autista: orientações para}
\chapter{Deficiência psicossocial e transtorno do espectro autista: orientações para a}
\chapter{Deficiência psicossocial e transtorno do espectro autista: orientações para a realização}
\chapter{Deficiência psicossocial e transtorno do espectro autista: orientações para a realização de}
\chapter{Deficiência psicossocial e transtorno do espectro autista: orientações para a realização de eventos}
\chapter{Deficiência psicossocial e transtorno do espectro autista: orientações para a realização de eventos e}
\chapter{Deficiência psicossocial e transtorno do espectro autista: orientações para a realização de eventos e apresentações}
\chapter{Deficiência psicossocial e transtorno do espectro autista: orientações para a realização de eventos e apresentações acessíveis}
\chapter{Deficiência psicossocial e transtorno do espectro autista: orientações para a realização de eventos e apresentações acessíveis em}
\chapter{Deficiência psicossocial e transtorno do espectro autista: orientações para a realização de eventos e apresentações acessíveis em todo}
\chapter{Comunicação alternativa e}
\chapter{Comunicação alternativa e aumentativa}
\chapter{Comunicação alternativa e aumentativa na}
\chapter{Comunicação alternativa e aumentativa na participação}
\chapter{Comunicação alternativa e aumentativa na participação social}
\chapter{Comunicação alternativa e aumentativa na participação social de}
\chapter{Comunicação alternativa e aumentativa na participação social de pessoas}
\chapter{Comunicação alternativa e aumentativa na participação social de pessoas com}
\chapter{Comunicação alternativa e aumentativa na participação social de pessoas com deficiência}
\chapter{Comunicação alternativa e aumentativa na participação social de pessoas com deficiência intelectual}
\chapter{Comunicação alternativa e aumentativa na participação social de pessoas com deficiência intelectual em}
\chapter{Comunicação alternativa e aumentativa na participação social de pessoas com deficiência intelectual em eventos}
\chapter{Comunicação alternativa e aumentativa na participação social de pessoas com deficiência intelectual em eventos públicos}
\chapter{Comunicação alternativa e aumentativa na participação social de pessoas com deficiência intelectual em eventos públicos de}
\chapter{Comunicação alternativa e aumentativa na participação social de pessoas com deficiência intelectual em eventos públicos de grande}
\chapter{Comunicação alternativa e aumentativa na participação social de pessoas com deficiência intelectual em eventos públicos de grande porte}
\chapter{Psicopedagogicamente inconstitucionalíssimo otorrinolaringologista}
\chapter{Psicopedagogicamente inconstitucionalíssimo otorrinolaringologista socioeducativo}
\chapter{Psicopedagogicamente inconstitucionalíssimo otorrinolaringologista socioeducativo ético-profissional}
\chapter{Psicopedagogicamente inconstitucionalíssimo otorrinolaringologista socioeducativo ético-profissional sociodemográfico}
\chapter{Psicopedagogicamente inconstitucionalíssimo otorrinolaringologista socioeducativo ético-profissional sociodemográfico interdisciplinaridade}
\chapter{Psicopedagogicamente inconstitucionalíssimo otorrinolaringologista socioeducativo ético-profissional sociodemográfico interdisciplinaridade neuropsicológica}
\chapter{Psicopedagogicamente inconstitucionalíssimo otorrinolaringologista socioeducativo ético-profissional sociodemográfico interdisciplinaridade neuropsicológica avaliação}
\chapter{Psicopedagogicamente inconstitucionalíssimo otorrinolaringologista socioeducativo ético-profissional sociodemográfico interdisciplinaridade neuropsicológica avaliação}
\chapter{Psicopedagogicamente inconstitucionalíssimo otorrinolaringologista socioeducativo ético-profissional sociodemográfico interdisciplinaridade neuropsicológica avaliação}
\chapter{Psicopedagogicamente inconstitucionalíssimo otorrinolaringologista socioeducativo ético-profissional sociodemográfico interdisciplinaridade neuropsicológica avaliação}
\chapter{Psicopedagogicamente inconstitucionalíssimo otorrinolaringologista socioeducativo ético-profissional sociodemográfico interdisciplinaridade neuropsicológica avaliação}
\chapter{Psicopedagogicamente inconstitucionalíssimo otorrinolaringologista socioeducativo ético-profissional sociodemográfico interdisciplinaridade neuropsicológica avaliação}
\chapter{Psicopedagogicamente inconstitucionalíssimo otorrinolaringologista socioeducativo ético-profissional sociodemográfico interdisciplinaridade neuropsicológica avaliação}
\chapter{Psicopedagogicamente inconstitucionalíssimo otorrinolaringologista socioeducativo ético-profissional sociodemográfico interdisciplinaridade neuropsicológica avaliação}
\chapter{Psicopedagogicamente inconstitucionalíssimo otorrinolaringologista socioeducativo ético-profissional sociodemográfico interdisciplinaridade neuropsicológica avaliação}
\chapter{Psicopedagogicamente inconstitucionalíssimo otorrinolaringologista socioeducativo ético-profissional sociodemográfico interdisciplinaridade neuropsicológica avaliação}
\end{document}
'   % B: plus 4em (o do guia)
%   lualatex '\def\var{C}% !TeX program = lualatex
% ===================================================================
% sumario-longo.tex — mede o sumário sem hifenização (issue #35, "Fora do
% escopo"): com a hifenização desligada, a entrada longa de capítulo pode
% bater no número da página. Gera 3 x 16 capítulos, cada um com um prefixo
% (de 3 a 18 palavras) de uma frase longa: frase 1 de palavras médias,
% frase 2 com números, frase 3 só de palavras muito longas. Conta no log
% os "Overfull \hbox" do sumário: o sumário é a única parte do documento
% sem parágrafo de corpo, então todo Overfull "in paragraph" antes do
% primeiro \chapter é entrada que passou da margem.
%
% Uso (de uma pasta com livros_crp_acessivel_book.sty e crpsp-leg.sty no
% caminho; as fontes NEWJUNE por OSFONTDIR):
%   lualatex '\def\var{R}% !TeX program = lualatex
% ===================================================================
% sumario-longo.tex — mede o sumário sem hifenização (issue #35, "Fora do
% escopo"): com a hifenização desligada, a entrada longa de capítulo pode
% bater no número da página. Gera 3 x 16 capítulos, cada um com um prefixo
% (de 3 a 18 palavras) de uma frase longa: frase 1 de palavras médias,
% frase 2 com números, frase 3 só de palavras muito longas. Conta no log
% os "Overfull \hbox" do sumário: o sumário é a única parte do documento
% sem parágrafo de corpo, então todo Overfull "in paragraph" antes do
% primeiro \chapter é entrada que passou da margem.
%
% Uso (de uma pasta com livros_crp_acessivel_book.sty e crpsp-leg.sty no
% caminho; as fontes NEWJUNE por OSFONTDIR):
%   lualatex '\def\var{R}% !TeX program = lualatex
% ===================================================================
% sumario-longo.tex — mede o sumário sem hifenização (issue #35, "Fora do
% escopo"): com a hifenização desligada, a entrada longa de capítulo pode
% bater no número da página. Gera 3 x 16 capítulos, cada um com um prefixo
% (de 3 a 18 palavras) de uma frase longa: frase 1 de palavras médias,
% frase 2 com números, frase 3 só de palavras muito longas. Conta no log
% os "Overfull \hbox" do sumário: o sumário é a única parte do documento
% sem parágrafo de corpo, então todo Overfull "in paragraph" antes do
% primeiro \chapter é entrada que passou da margem.
%
% Uso (de uma pasta com livros_crp_acessivel_book.sty e crpsp-leg.sty no
% caminho; as fontes NEWJUNE por OSFONTDIR):
%   lualatex '\def\var{R}\input{sumario-longo}'   % R: rígido (até 0.5.x)
%   lualatex '\def\var{B}\input{sumario-longo}'   % B: plus 4em (o do guia)
%   lualatex '\def\var{C}\input{sumario-longo}'   % C: plus 4em + \tolerance 10000
%   lualatex '\def\var{S}\input{sumario-longo}'   % S: o que a .sty define
% Cada variante só redefine o \l@chapter depois da carga da .sty.
% ===================================================================
\DocumentMetadata{lang=pt-BR, pdfstandard=ua-2, pdfversion=2.0, tagging=on}
\documentclass[11pt,a5paper,twoside,openright]{book}
\usepackage{livros_crp_acessivel_book}
\hypersetup{pdftitle={Sumario longo}, pdfauthor={CRP SP}}
\makeatletter
\providecommand\var{S}
\def\crp@vR{R}\def\crp@vB{B}\def\crp@vC{C}
\def\crp@lchap#1{%  #1 = acréscimo à margem direita (rightskip)
  \renewcommand*\l@chapter[2]{%
    \ifnum \c@tocdepth >\m@ne
        \addpenalty{-\@highpenalty}\addvspace{0.5em \@plus\p@}%
        \setlength\@tempdima{1.5em}%
        \begingroup
            \parindent \z@ \rightskip \@pnumwidth #1
            \parfillskip -\@pnumwidth
            \leavevmode \bfseries
            \advance\leftskip\@tempdima \hskip -\leftskip
            ##1\nobreak\hfil
            \nobreak\hb@xt@\@pnumwidth{\hss ##2\kern-\p@\kern\p@}\par
            \penalty\@highpenalty
        \endgroup
    \fi}}
\ifx\var\crp@vR \crp@lchap{}\fi
\ifx\var\crp@vB \crp@lchap{\@plus 4em}\fi
\ifx\var\crp@vC \crp@lchap{\@plus 4em \tolerance=10000 \hbadness=10000 }\fi
\makeatother
\begin{document}
\frontmatter
\tableofcontents
\mainmatter
\chapter{Deficiência psicossocial e}
\chapter{Deficiência psicossocial e transtorno}
\chapter{Deficiência psicossocial e transtorno do}
\chapter{Deficiência psicossocial e transtorno do espectro}
\chapter{Deficiência psicossocial e transtorno do espectro autista:}
\chapter{Deficiência psicossocial e transtorno do espectro autista: orientações}
\chapter{Deficiência psicossocial e transtorno do espectro autista: orientações para}
\chapter{Deficiência psicossocial e transtorno do espectro autista: orientações para a}
\chapter{Deficiência psicossocial e transtorno do espectro autista: orientações para a realização}
\chapter{Deficiência psicossocial e transtorno do espectro autista: orientações para a realização de}
\chapter{Deficiência psicossocial e transtorno do espectro autista: orientações para a realização de eventos}
\chapter{Deficiência psicossocial e transtorno do espectro autista: orientações para a realização de eventos e}
\chapter{Deficiência psicossocial e transtorno do espectro autista: orientações para a realização de eventos e apresentações}
\chapter{Deficiência psicossocial e transtorno do espectro autista: orientações para a realização de eventos e apresentações acessíveis}
\chapter{Deficiência psicossocial e transtorno do espectro autista: orientações para a realização de eventos e apresentações acessíveis em}
\chapter{Deficiência psicossocial e transtorno do espectro autista: orientações para a realização de eventos e apresentações acessíveis em todo}
\chapter{Comunicação alternativa e}
\chapter{Comunicação alternativa e aumentativa}
\chapter{Comunicação alternativa e aumentativa na}
\chapter{Comunicação alternativa e aumentativa na participação}
\chapter{Comunicação alternativa e aumentativa na participação social}
\chapter{Comunicação alternativa e aumentativa na participação social de}
\chapter{Comunicação alternativa e aumentativa na participação social de pessoas}
\chapter{Comunicação alternativa e aumentativa na participação social de pessoas com}
\chapter{Comunicação alternativa e aumentativa na participação social de pessoas com deficiência}
\chapter{Comunicação alternativa e aumentativa na participação social de pessoas com deficiência intelectual}
\chapter{Comunicação alternativa e aumentativa na participação social de pessoas com deficiência intelectual em}
\chapter{Comunicação alternativa e aumentativa na participação social de pessoas com deficiência intelectual em eventos}
\chapter{Comunicação alternativa e aumentativa na participação social de pessoas com deficiência intelectual em eventos públicos}
\chapter{Comunicação alternativa e aumentativa na participação social de pessoas com deficiência intelectual em eventos públicos de}
\chapter{Comunicação alternativa e aumentativa na participação social de pessoas com deficiência intelectual em eventos públicos de grande}
\chapter{Comunicação alternativa e aumentativa na participação social de pessoas com deficiência intelectual em eventos públicos de grande porte}
\chapter{Psicopedagogicamente inconstitucionalíssimo otorrinolaringologista}
\chapter{Psicopedagogicamente inconstitucionalíssimo otorrinolaringologista socioeducativo}
\chapter{Psicopedagogicamente inconstitucionalíssimo otorrinolaringologista socioeducativo ético-profissional}
\chapter{Psicopedagogicamente inconstitucionalíssimo otorrinolaringologista socioeducativo ético-profissional sociodemográfico}
\chapter{Psicopedagogicamente inconstitucionalíssimo otorrinolaringologista socioeducativo ético-profissional sociodemográfico interdisciplinaridade}
\chapter{Psicopedagogicamente inconstitucionalíssimo otorrinolaringologista socioeducativo ético-profissional sociodemográfico interdisciplinaridade neuropsicológica}
\chapter{Psicopedagogicamente inconstitucionalíssimo otorrinolaringologista socioeducativo ético-profissional sociodemográfico interdisciplinaridade neuropsicológica avaliação}
\chapter{Psicopedagogicamente inconstitucionalíssimo otorrinolaringologista socioeducativo ético-profissional sociodemográfico interdisciplinaridade neuropsicológica avaliação}
\chapter{Psicopedagogicamente inconstitucionalíssimo otorrinolaringologista socioeducativo ético-profissional sociodemográfico interdisciplinaridade neuropsicológica avaliação}
\chapter{Psicopedagogicamente inconstitucionalíssimo otorrinolaringologista socioeducativo ético-profissional sociodemográfico interdisciplinaridade neuropsicológica avaliação}
\chapter{Psicopedagogicamente inconstitucionalíssimo otorrinolaringologista socioeducativo ético-profissional sociodemográfico interdisciplinaridade neuropsicológica avaliação}
\chapter{Psicopedagogicamente inconstitucionalíssimo otorrinolaringologista socioeducativo ético-profissional sociodemográfico interdisciplinaridade neuropsicológica avaliação}
\chapter{Psicopedagogicamente inconstitucionalíssimo otorrinolaringologista socioeducativo ético-profissional sociodemográfico interdisciplinaridade neuropsicológica avaliação}
\chapter{Psicopedagogicamente inconstitucionalíssimo otorrinolaringologista socioeducativo ético-profissional sociodemográfico interdisciplinaridade neuropsicológica avaliação}
\chapter{Psicopedagogicamente inconstitucionalíssimo otorrinolaringologista socioeducativo ético-profissional sociodemográfico interdisciplinaridade neuropsicológica avaliação}
\chapter{Psicopedagogicamente inconstitucionalíssimo otorrinolaringologista socioeducativo ético-profissional sociodemográfico interdisciplinaridade neuropsicológica avaliação}
\end{document}
'   % R: rígido (até 0.5.x)
%   lualatex '\def\var{B}% !TeX program = lualatex
% ===================================================================
% sumario-longo.tex — mede o sumário sem hifenização (issue #35, "Fora do
% escopo"): com a hifenização desligada, a entrada longa de capítulo pode
% bater no número da página. Gera 3 x 16 capítulos, cada um com um prefixo
% (de 3 a 18 palavras) de uma frase longa: frase 1 de palavras médias,
% frase 2 com números, frase 3 só de palavras muito longas. Conta no log
% os "Overfull \hbox" do sumário: o sumário é a única parte do documento
% sem parágrafo de corpo, então todo Overfull "in paragraph" antes do
% primeiro \chapter é entrada que passou da margem.
%
% Uso (de uma pasta com livros_crp_acessivel_book.sty e crpsp-leg.sty no
% caminho; as fontes NEWJUNE por OSFONTDIR):
%   lualatex '\def\var{R}\input{sumario-longo}'   % R: rígido (até 0.5.x)
%   lualatex '\def\var{B}\input{sumario-longo}'   % B: plus 4em (o do guia)
%   lualatex '\def\var{C}\input{sumario-longo}'   % C: plus 4em + \tolerance 10000
%   lualatex '\def\var{S}\input{sumario-longo}'   % S: o que a .sty define
% Cada variante só redefine o \l@chapter depois da carga da .sty.
% ===================================================================
\DocumentMetadata{lang=pt-BR, pdfstandard=ua-2, pdfversion=2.0, tagging=on}
\documentclass[11pt,a5paper,twoside,openright]{book}
\usepackage{livros_crp_acessivel_book}
\hypersetup{pdftitle={Sumario longo}, pdfauthor={CRP SP}}
\makeatletter
\providecommand\var{S}
\def\crp@vR{R}\def\crp@vB{B}\def\crp@vC{C}
\def\crp@lchap#1{%  #1 = acréscimo à margem direita (rightskip)
  \renewcommand*\l@chapter[2]{%
    \ifnum \c@tocdepth >\m@ne
        \addpenalty{-\@highpenalty}\addvspace{0.5em \@plus\p@}%
        \setlength\@tempdima{1.5em}%
        \begingroup
            \parindent \z@ \rightskip \@pnumwidth #1
            \parfillskip -\@pnumwidth
            \leavevmode \bfseries
            \advance\leftskip\@tempdima \hskip -\leftskip
            ##1\nobreak\hfil
            \nobreak\hb@xt@\@pnumwidth{\hss ##2\kern-\p@\kern\p@}\par
            \penalty\@highpenalty
        \endgroup
    \fi}}
\ifx\var\crp@vR \crp@lchap{}\fi
\ifx\var\crp@vB \crp@lchap{\@plus 4em}\fi
\ifx\var\crp@vC \crp@lchap{\@plus 4em \tolerance=10000 \hbadness=10000 }\fi
\makeatother
\begin{document}
\frontmatter
\tableofcontents
\mainmatter
\chapter{Deficiência psicossocial e}
\chapter{Deficiência psicossocial e transtorno}
\chapter{Deficiência psicossocial e transtorno do}
\chapter{Deficiência psicossocial e transtorno do espectro}
\chapter{Deficiência psicossocial e transtorno do espectro autista:}
\chapter{Deficiência psicossocial e transtorno do espectro autista: orientações}
\chapter{Deficiência psicossocial e transtorno do espectro autista: orientações para}
\chapter{Deficiência psicossocial e transtorno do espectro autista: orientações para a}
\chapter{Deficiência psicossocial e transtorno do espectro autista: orientações para a realização}
\chapter{Deficiência psicossocial e transtorno do espectro autista: orientações para a realização de}
\chapter{Deficiência psicossocial e transtorno do espectro autista: orientações para a realização de eventos}
\chapter{Deficiência psicossocial e transtorno do espectro autista: orientações para a realização de eventos e}
\chapter{Deficiência psicossocial e transtorno do espectro autista: orientações para a realização de eventos e apresentações}
\chapter{Deficiência psicossocial e transtorno do espectro autista: orientações para a realização de eventos e apresentações acessíveis}
\chapter{Deficiência psicossocial e transtorno do espectro autista: orientações para a realização de eventos e apresentações acessíveis em}
\chapter{Deficiência psicossocial e transtorno do espectro autista: orientações para a realização de eventos e apresentações acessíveis em todo}
\chapter{Comunicação alternativa e}
\chapter{Comunicação alternativa e aumentativa}
\chapter{Comunicação alternativa e aumentativa na}
\chapter{Comunicação alternativa e aumentativa na participação}
\chapter{Comunicação alternativa e aumentativa na participação social}
\chapter{Comunicação alternativa e aumentativa na participação social de}
\chapter{Comunicação alternativa e aumentativa na participação social de pessoas}
\chapter{Comunicação alternativa e aumentativa na participação social de pessoas com}
\chapter{Comunicação alternativa e aumentativa na participação social de pessoas com deficiência}
\chapter{Comunicação alternativa e aumentativa na participação social de pessoas com deficiência intelectual}
\chapter{Comunicação alternativa e aumentativa na participação social de pessoas com deficiência intelectual em}
\chapter{Comunicação alternativa e aumentativa na participação social de pessoas com deficiência intelectual em eventos}
\chapter{Comunicação alternativa e aumentativa na participação social de pessoas com deficiência intelectual em eventos públicos}
\chapter{Comunicação alternativa e aumentativa na participação social de pessoas com deficiência intelectual em eventos públicos de}
\chapter{Comunicação alternativa e aumentativa na participação social de pessoas com deficiência intelectual em eventos públicos de grande}
\chapter{Comunicação alternativa e aumentativa na participação social de pessoas com deficiência intelectual em eventos públicos de grande porte}
\chapter{Psicopedagogicamente inconstitucionalíssimo otorrinolaringologista}
\chapter{Psicopedagogicamente inconstitucionalíssimo otorrinolaringologista socioeducativo}
\chapter{Psicopedagogicamente inconstitucionalíssimo otorrinolaringologista socioeducativo ético-profissional}
\chapter{Psicopedagogicamente inconstitucionalíssimo otorrinolaringologista socioeducativo ético-profissional sociodemográfico}
\chapter{Psicopedagogicamente inconstitucionalíssimo otorrinolaringologista socioeducativo ético-profissional sociodemográfico interdisciplinaridade}
\chapter{Psicopedagogicamente inconstitucionalíssimo otorrinolaringologista socioeducativo ético-profissional sociodemográfico interdisciplinaridade neuropsicológica}
\chapter{Psicopedagogicamente inconstitucionalíssimo otorrinolaringologista socioeducativo ético-profissional sociodemográfico interdisciplinaridade neuropsicológica avaliação}
\chapter{Psicopedagogicamente inconstitucionalíssimo otorrinolaringologista socioeducativo ético-profissional sociodemográfico interdisciplinaridade neuropsicológica avaliação}
\chapter{Psicopedagogicamente inconstitucionalíssimo otorrinolaringologista socioeducativo ético-profissional sociodemográfico interdisciplinaridade neuropsicológica avaliação}
\chapter{Psicopedagogicamente inconstitucionalíssimo otorrinolaringologista socioeducativo ético-profissional sociodemográfico interdisciplinaridade neuropsicológica avaliação}
\chapter{Psicopedagogicamente inconstitucionalíssimo otorrinolaringologista socioeducativo ético-profissional sociodemográfico interdisciplinaridade neuropsicológica avaliação}
\chapter{Psicopedagogicamente inconstitucionalíssimo otorrinolaringologista socioeducativo ético-profissional sociodemográfico interdisciplinaridade neuropsicológica avaliação}
\chapter{Psicopedagogicamente inconstitucionalíssimo otorrinolaringologista socioeducativo ético-profissional sociodemográfico interdisciplinaridade neuropsicológica avaliação}
\chapter{Psicopedagogicamente inconstitucionalíssimo otorrinolaringologista socioeducativo ético-profissional sociodemográfico interdisciplinaridade neuropsicológica avaliação}
\chapter{Psicopedagogicamente inconstitucionalíssimo otorrinolaringologista socioeducativo ético-profissional sociodemográfico interdisciplinaridade neuropsicológica avaliação}
\chapter{Psicopedagogicamente inconstitucionalíssimo otorrinolaringologista socioeducativo ético-profissional sociodemográfico interdisciplinaridade neuropsicológica avaliação}
\end{document}
'   % B: plus 4em (o do guia)
%   lualatex '\def\var{C}% !TeX program = lualatex
% ===================================================================
% sumario-longo.tex — mede o sumário sem hifenização (issue #35, "Fora do
% escopo"): com a hifenização desligada, a entrada longa de capítulo pode
% bater no número da página. Gera 3 x 16 capítulos, cada um com um prefixo
% (de 3 a 18 palavras) de uma frase longa: frase 1 de palavras médias,
% frase 2 com números, frase 3 só de palavras muito longas. Conta no log
% os "Overfull \hbox" do sumário: o sumário é a única parte do documento
% sem parágrafo de corpo, então todo Overfull "in paragraph" antes do
% primeiro \chapter é entrada que passou da margem.
%
% Uso (de uma pasta com livros_crp_acessivel_book.sty e crpsp-leg.sty no
% caminho; as fontes NEWJUNE por OSFONTDIR):
%   lualatex '\def\var{R}\input{sumario-longo}'   % R: rígido (até 0.5.x)
%   lualatex '\def\var{B}\input{sumario-longo}'   % B: plus 4em (o do guia)
%   lualatex '\def\var{C}\input{sumario-longo}'   % C: plus 4em + \tolerance 10000
%   lualatex '\def\var{S}\input{sumario-longo}'   % S: o que a .sty define
% Cada variante só redefine o \l@chapter depois da carga da .sty.
% ===================================================================
\DocumentMetadata{lang=pt-BR, pdfstandard=ua-2, pdfversion=2.0, tagging=on}
\documentclass[11pt,a5paper,twoside,openright]{book}
\usepackage{livros_crp_acessivel_book}
\hypersetup{pdftitle={Sumario longo}, pdfauthor={CRP SP}}
\makeatletter
\providecommand\var{S}
\def\crp@vR{R}\def\crp@vB{B}\def\crp@vC{C}
\def\crp@lchap#1{%  #1 = acréscimo à margem direita (rightskip)
  \renewcommand*\l@chapter[2]{%
    \ifnum \c@tocdepth >\m@ne
        \addpenalty{-\@highpenalty}\addvspace{0.5em \@plus\p@}%
        \setlength\@tempdima{1.5em}%
        \begingroup
            \parindent \z@ \rightskip \@pnumwidth #1
            \parfillskip -\@pnumwidth
            \leavevmode \bfseries
            \advance\leftskip\@tempdima \hskip -\leftskip
            ##1\nobreak\hfil
            \nobreak\hb@xt@\@pnumwidth{\hss ##2\kern-\p@\kern\p@}\par
            \penalty\@highpenalty
        \endgroup
    \fi}}
\ifx\var\crp@vR \crp@lchap{}\fi
\ifx\var\crp@vB \crp@lchap{\@plus 4em}\fi
\ifx\var\crp@vC \crp@lchap{\@plus 4em \tolerance=10000 \hbadness=10000 }\fi
\makeatother
\begin{document}
\frontmatter
\tableofcontents
\mainmatter
\chapter{Deficiência psicossocial e}
\chapter{Deficiência psicossocial e transtorno}
\chapter{Deficiência psicossocial e transtorno do}
\chapter{Deficiência psicossocial e transtorno do espectro}
\chapter{Deficiência psicossocial e transtorno do espectro autista:}
\chapter{Deficiência psicossocial e transtorno do espectro autista: orientações}
\chapter{Deficiência psicossocial e transtorno do espectro autista: orientações para}
\chapter{Deficiência psicossocial e transtorno do espectro autista: orientações para a}
\chapter{Deficiência psicossocial e transtorno do espectro autista: orientações para a realização}
\chapter{Deficiência psicossocial e transtorno do espectro autista: orientações para a realização de}
\chapter{Deficiência psicossocial e transtorno do espectro autista: orientações para a realização de eventos}
\chapter{Deficiência psicossocial e transtorno do espectro autista: orientações para a realização de eventos e}
\chapter{Deficiência psicossocial e transtorno do espectro autista: orientações para a realização de eventos e apresentações}
\chapter{Deficiência psicossocial e transtorno do espectro autista: orientações para a realização de eventos e apresentações acessíveis}
\chapter{Deficiência psicossocial e transtorno do espectro autista: orientações para a realização de eventos e apresentações acessíveis em}
\chapter{Deficiência psicossocial e transtorno do espectro autista: orientações para a realização de eventos e apresentações acessíveis em todo}
\chapter{Comunicação alternativa e}
\chapter{Comunicação alternativa e aumentativa}
\chapter{Comunicação alternativa e aumentativa na}
\chapter{Comunicação alternativa e aumentativa na participação}
\chapter{Comunicação alternativa e aumentativa na participação social}
\chapter{Comunicação alternativa e aumentativa na participação social de}
\chapter{Comunicação alternativa e aumentativa na participação social de pessoas}
\chapter{Comunicação alternativa e aumentativa na participação social de pessoas com}
\chapter{Comunicação alternativa e aumentativa na participação social de pessoas com deficiência}
\chapter{Comunicação alternativa e aumentativa na participação social de pessoas com deficiência intelectual}
\chapter{Comunicação alternativa e aumentativa na participação social de pessoas com deficiência intelectual em}
\chapter{Comunicação alternativa e aumentativa na participação social de pessoas com deficiência intelectual em eventos}
\chapter{Comunicação alternativa e aumentativa na participação social de pessoas com deficiência intelectual em eventos públicos}
\chapter{Comunicação alternativa e aumentativa na participação social de pessoas com deficiência intelectual em eventos públicos de}
\chapter{Comunicação alternativa e aumentativa na participação social de pessoas com deficiência intelectual em eventos públicos de grande}
\chapter{Comunicação alternativa e aumentativa na participação social de pessoas com deficiência intelectual em eventos públicos de grande porte}
\chapter{Psicopedagogicamente inconstitucionalíssimo otorrinolaringologista}
\chapter{Psicopedagogicamente inconstitucionalíssimo otorrinolaringologista socioeducativo}
\chapter{Psicopedagogicamente inconstitucionalíssimo otorrinolaringologista socioeducativo ético-profissional}
\chapter{Psicopedagogicamente inconstitucionalíssimo otorrinolaringologista socioeducativo ético-profissional sociodemográfico}
\chapter{Psicopedagogicamente inconstitucionalíssimo otorrinolaringologista socioeducativo ético-profissional sociodemográfico interdisciplinaridade}
\chapter{Psicopedagogicamente inconstitucionalíssimo otorrinolaringologista socioeducativo ético-profissional sociodemográfico interdisciplinaridade neuropsicológica}
\chapter{Psicopedagogicamente inconstitucionalíssimo otorrinolaringologista socioeducativo ético-profissional sociodemográfico interdisciplinaridade neuropsicológica avaliação}
\chapter{Psicopedagogicamente inconstitucionalíssimo otorrinolaringologista socioeducativo ético-profissional sociodemográfico interdisciplinaridade neuropsicológica avaliação}
\chapter{Psicopedagogicamente inconstitucionalíssimo otorrinolaringologista socioeducativo ético-profissional sociodemográfico interdisciplinaridade neuropsicológica avaliação}
\chapter{Psicopedagogicamente inconstitucionalíssimo otorrinolaringologista socioeducativo ético-profissional sociodemográfico interdisciplinaridade neuropsicológica avaliação}
\chapter{Psicopedagogicamente inconstitucionalíssimo otorrinolaringologista socioeducativo ético-profissional sociodemográfico interdisciplinaridade neuropsicológica avaliação}
\chapter{Psicopedagogicamente inconstitucionalíssimo otorrinolaringologista socioeducativo ético-profissional sociodemográfico interdisciplinaridade neuropsicológica avaliação}
\chapter{Psicopedagogicamente inconstitucionalíssimo otorrinolaringologista socioeducativo ético-profissional sociodemográfico interdisciplinaridade neuropsicológica avaliação}
\chapter{Psicopedagogicamente inconstitucionalíssimo otorrinolaringologista socioeducativo ético-profissional sociodemográfico interdisciplinaridade neuropsicológica avaliação}
\chapter{Psicopedagogicamente inconstitucionalíssimo otorrinolaringologista socioeducativo ético-profissional sociodemográfico interdisciplinaridade neuropsicológica avaliação}
\chapter{Psicopedagogicamente inconstitucionalíssimo otorrinolaringologista socioeducativo ético-profissional sociodemográfico interdisciplinaridade neuropsicológica avaliação}
\end{document}
'   % C: plus 4em + \tolerance 10000
%   lualatex '\def\var{S}% !TeX program = lualatex
% ===================================================================
% sumario-longo.tex — mede o sumário sem hifenização (issue #35, "Fora do
% escopo"): com a hifenização desligada, a entrada longa de capítulo pode
% bater no número da página. Gera 3 x 16 capítulos, cada um com um prefixo
% (de 3 a 18 palavras) de uma frase longa: frase 1 de palavras médias,
% frase 2 com números, frase 3 só de palavras muito longas. Conta no log
% os "Overfull \hbox" do sumário: o sumário é a única parte do documento
% sem parágrafo de corpo, então todo Overfull "in paragraph" antes do
% primeiro \chapter é entrada que passou da margem.
%
% Uso (de uma pasta com livros_crp_acessivel_book.sty e crpsp-leg.sty no
% caminho; as fontes NEWJUNE por OSFONTDIR):
%   lualatex '\def\var{R}\input{sumario-longo}'   % R: rígido (até 0.5.x)
%   lualatex '\def\var{B}\input{sumario-longo}'   % B: plus 4em (o do guia)
%   lualatex '\def\var{C}\input{sumario-longo}'   % C: plus 4em + \tolerance 10000
%   lualatex '\def\var{S}\input{sumario-longo}'   % S: o que a .sty define
% Cada variante só redefine o \l@chapter depois da carga da .sty.
% ===================================================================
\DocumentMetadata{lang=pt-BR, pdfstandard=ua-2, pdfversion=2.0, tagging=on}
\documentclass[11pt,a5paper,twoside,openright]{book}
\usepackage{livros_crp_acessivel_book}
\hypersetup{pdftitle={Sumario longo}, pdfauthor={CRP SP}}
\makeatletter
\providecommand\var{S}
\def\crp@vR{R}\def\crp@vB{B}\def\crp@vC{C}
\def\crp@lchap#1{%  #1 = acréscimo à margem direita (rightskip)
  \renewcommand*\l@chapter[2]{%
    \ifnum \c@tocdepth >\m@ne
        \addpenalty{-\@highpenalty}\addvspace{0.5em \@plus\p@}%
        \setlength\@tempdima{1.5em}%
        \begingroup
            \parindent \z@ \rightskip \@pnumwidth #1
            \parfillskip -\@pnumwidth
            \leavevmode \bfseries
            \advance\leftskip\@tempdima \hskip -\leftskip
            ##1\nobreak\hfil
            \nobreak\hb@xt@\@pnumwidth{\hss ##2\kern-\p@\kern\p@}\par
            \penalty\@highpenalty
        \endgroup
    \fi}}
\ifx\var\crp@vR \crp@lchap{}\fi
\ifx\var\crp@vB \crp@lchap{\@plus 4em}\fi
\ifx\var\crp@vC \crp@lchap{\@plus 4em \tolerance=10000 \hbadness=10000 }\fi
\makeatother
\begin{document}
\frontmatter
\tableofcontents
\mainmatter
\chapter{Deficiência psicossocial e}
\chapter{Deficiência psicossocial e transtorno}
\chapter{Deficiência psicossocial e transtorno do}
\chapter{Deficiência psicossocial e transtorno do espectro}
\chapter{Deficiência psicossocial e transtorno do espectro autista:}
\chapter{Deficiência psicossocial e transtorno do espectro autista: orientações}
\chapter{Deficiência psicossocial e transtorno do espectro autista: orientações para}
\chapter{Deficiência psicossocial e transtorno do espectro autista: orientações para a}
\chapter{Deficiência psicossocial e transtorno do espectro autista: orientações para a realização}
\chapter{Deficiência psicossocial e transtorno do espectro autista: orientações para a realização de}
\chapter{Deficiência psicossocial e transtorno do espectro autista: orientações para a realização de eventos}
\chapter{Deficiência psicossocial e transtorno do espectro autista: orientações para a realização de eventos e}
\chapter{Deficiência psicossocial e transtorno do espectro autista: orientações para a realização de eventos e apresentações}
\chapter{Deficiência psicossocial e transtorno do espectro autista: orientações para a realização de eventos e apresentações acessíveis}
\chapter{Deficiência psicossocial e transtorno do espectro autista: orientações para a realização de eventos e apresentações acessíveis em}
\chapter{Deficiência psicossocial e transtorno do espectro autista: orientações para a realização de eventos e apresentações acessíveis em todo}
\chapter{Comunicação alternativa e}
\chapter{Comunicação alternativa e aumentativa}
\chapter{Comunicação alternativa e aumentativa na}
\chapter{Comunicação alternativa e aumentativa na participação}
\chapter{Comunicação alternativa e aumentativa na participação social}
\chapter{Comunicação alternativa e aumentativa na participação social de}
\chapter{Comunicação alternativa e aumentativa na participação social de pessoas}
\chapter{Comunicação alternativa e aumentativa na participação social de pessoas com}
\chapter{Comunicação alternativa e aumentativa na participação social de pessoas com deficiência}
\chapter{Comunicação alternativa e aumentativa na participação social de pessoas com deficiência intelectual}
\chapter{Comunicação alternativa e aumentativa na participação social de pessoas com deficiência intelectual em}
\chapter{Comunicação alternativa e aumentativa na participação social de pessoas com deficiência intelectual em eventos}
\chapter{Comunicação alternativa e aumentativa na participação social de pessoas com deficiência intelectual em eventos públicos}
\chapter{Comunicação alternativa e aumentativa na participação social de pessoas com deficiência intelectual em eventos públicos de}
\chapter{Comunicação alternativa e aumentativa na participação social de pessoas com deficiência intelectual em eventos públicos de grande}
\chapter{Comunicação alternativa e aumentativa na participação social de pessoas com deficiência intelectual em eventos públicos de grande porte}
\chapter{Psicopedagogicamente inconstitucionalíssimo otorrinolaringologista}
\chapter{Psicopedagogicamente inconstitucionalíssimo otorrinolaringologista socioeducativo}
\chapter{Psicopedagogicamente inconstitucionalíssimo otorrinolaringologista socioeducativo ético-profissional}
\chapter{Psicopedagogicamente inconstitucionalíssimo otorrinolaringologista socioeducativo ético-profissional sociodemográfico}
\chapter{Psicopedagogicamente inconstitucionalíssimo otorrinolaringologista socioeducativo ético-profissional sociodemográfico interdisciplinaridade}
\chapter{Psicopedagogicamente inconstitucionalíssimo otorrinolaringologista socioeducativo ético-profissional sociodemográfico interdisciplinaridade neuropsicológica}
\chapter{Psicopedagogicamente inconstitucionalíssimo otorrinolaringologista socioeducativo ético-profissional sociodemográfico interdisciplinaridade neuropsicológica avaliação}
\chapter{Psicopedagogicamente inconstitucionalíssimo otorrinolaringologista socioeducativo ético-profissional sociodemográfico interdisciplinaridade neuropsicológica avaliação}
\chapter{Psicopedagogicamente inconstitucionalíssimo otorrinolaringologista socioeducativo ético-profissional sociodemográfico interdisciplinaridade neuropsicológica avaliação}
\chapter{Psicopedagogicamente inconstitucionalíssimo otorrinolaringologista socioeducativo ético-profissional sociodemográfico interdisciplinaridade neuropsicológica avaliação}
\chapter{Psicopedagogicamente inconstitucionalíssimo otorrinolaringologista socioeducativo ético-profissional sociodemográfico interdisciplinaridade neuropsicológica avaliação}
\chapter{Psicopedagogicamente inconstitucionalíssimo otorrinolaringologista socioeducativo ético-profissional sociodemográfico interdisciplinaridade neuropsicológica avaliação}
\chapter{Psicopedagogicamente inconstitucionalíssimo otorrinolaringologista socioeducativo ético-profissional sociodemográfico interdisciplinaridade neuropsicológica avaliação}
\chapter{Psicopedagogicamente inconstitucionalíssimo otorrinolaringologista socioeducativo ético-profissional sociodemográfico interdisciplinaridade neuropsicológica avaliação}
\chapter{Psicopedagogicamente inconstitucionalíssimo otorrinolaringologista socioeducativo ético-profissional sociodemográfico interdisciplinaridade neuropsicológica avaliação}
\chapter{Psicopedagogicamente inconstitucionalíssimo otorrinolaringologista socioeducativo ético-profissional sociodemográfico interdisciplinaridade neuropsicológica avaliação}
\end{document}
'   % S: o que a .sty define
% Cada variante só redefine o \l@chapter depois da carga da .sty.
% ===================================================================
\DocumentMetadata{lang=pt-BR, pdfstandard=ua-2, pdfversion=2.0, tagging=on}
\documentclass[11pt,a5paper,twoside,openright]{book}
\usepackage{livros_crp_acessivel_book}
\hypersetup{pdftitle={Sumario longo}, pdfauthor={CRP SP}}
\makeatletter
\providecommand\var{S}
\def\crp@vR{R}\def\crp@vB{B}\def\crp@vC{C}
\def\crp@lchap#1{%  #1 = acréscimo à margem direita (rightskip)
  \renewcommand*\l@chapter[2]{%
    \ifnum \c@tocdepth >\m@ne
        \addpenalty{-\@highpenalty}\addvspace{0.5em \@plus\p@}%
        \setlength\@tempdima{1.5em}%
        \begingroup
            \parindent \z@ \rightskip \@pnumwidth #1
            \parfillskip -\@pnumwidth
            \leavevmode \bfseries
            \advance\leftskip\@tempdima \hskip -\leftskip
            ##1\nobreak\hfil
            \nobreak\hb@xt@\@pnumwidth{\hss ##2\kern-\p@\kern\p@}\par
            \penalty\@highpenalty
        \endgroup
    \fi}}
\ifx\var\crp@vR \crp@lchap{}\fi
\ifx\var\crp@vB \crp@lchap{\@plus 4em}\fi
\ifx\var\crp@vC \crp@lchap{\@plus 4em \tolerance=10000 \hbadness=10000 }\fi
\makeatother
\begin{document}
\frontmatter
\tableofcontents
\mainmatter
\chapter{Deficiência psicossocial e}
\chapter{Deficiência psicossocial e transtorno}
\chapter{Deficiência psicossocial e transtorno do}
\chapter{Deficiência psicossocial e transtorno do espectro}
\chapter{Deficiência psicossocial e transtorno do espectro autista:}
\chapter{Deficiência psicossocial e transtorno do espectro autista: orientações}
\chapter{Deficiência psicossocial e transtorno do espectro autista: orientações para}
\chapter{Deficiência psicossocial e transtorno do espectro autista: orientações para a}
\chapter{Deficiência psicossocial e transtorno do espectro autista: orientações para a realização}
\chapter{Deficiência psicossocial e transtorno do espectro autista: orientações para a realização de}
\chapter{Deficiência psicossocial e transtorno do espectro autista: orientações para a realização de eventos}
\chapter{Deficiência psicossocial e transtorno do espectro autista: orientações para a realização de eventos e}
\chapter{Deficiência psicossocial e transtorno do espectro autista: orientações para a realização de eventos e apresentações}
\chapter{Deficiência psicossocial e transtorno do espectro autista: orientações para a realização de eventos e apresentações acessíveis}
\chapter{Deficiência psicossocial e transtorno do espectro autista: orientações para a realização de eventos e apresentações acessíveis em}
\chapter{Deficiência psicossocial e transtorno do espectro autista: orientações para a realização de eventos e apresentações acessíveis em todo}
\chapter{Comunicação alternativa e}
\chapter{Comunicação alternativa e aumentativa}
\chapter{Comunicação alternativa e aumentativa na}
\chapter{Comunicação alternativa e aumentativa na participação}
\chapter{Comunicação alternativa e aumentativa na participação social}
\chapter{Comunicação alternativa e aumentativa na participação social de}
\chapter{Comunicação alternativa e aumentativa na participação social de pessoas}
\chapter{Comunicação alternativa e aumentativa na participação social de pessoas com}
\chapter{Comunicação alternativa e aumentativa na participação social de pessoas com deficiência}
\chapter{Comunicação alternativa e aumentativa na participação social de pessoas com deficiência intelectual}
\chapter{Comunicação alternativa e aumentativa na participação social de pessoas com deficiência intelectual em}
\chapter{Comunicação alternativa e aumentativa na participação social de pessoas com deficiência intelectual em eventos}
\chapter{Comunicação alternativa e aumentativa na participação social de pessoas com deficiência intelectual em eventos públicos}
\chapter{Comunicação alternativa e aumentativa na participação social de pessoas com deficiência intelectual em eventos públicos de}
\chapter{Comunicação alternativa e aumentativa na participação social de pessoas com deficiência intelectual em eventos públicos de grande}
\chapter{Comunicação alternativa e aumentativa na participação social de pessoas com deficiência intelectual em eventos públicos de grande porte}
\chapter{Psicopedagogicamente inconstitucionalíssimo otorrinolaringologista}
\chapter{Psicopedagogicamente inconstitucionalíssimo otorrinolaringologista socioeducativo}
\chapter{Psicopedagogicamente inconstitucionalíssimo otorrinolaringologista socioeducativo ético-profissional}
\chapter{Psicopedagogicamente inconstitucionalíssimo otorrinolaringologista socioeducativo ético-profissional sociodemográfico}
\chapter{Psicopedagogicamente inconstitucionalíssimo otorrinolaringologista socioeducativo ético-profissional sociodemográfico interdisciplinaridade}
\chapter{Psicopedagogicamente inconstitucionalíssimo otorrinolaringologista socioeducativo ético-profissional sociodemográfico interdisciplinaridade neuropsicológica}
\chapter{Psicopedagogicamente inconstitucionalíssimo otorrinolaringologista socioeducativo ético-profissional sociodemográfico interdisciplinaridade neuropsicológica avaliação}
\chapter{Psicopedagogicamente inconstitucionalíssimo otorrinolaringologista socioeducativo ético-profissional sociodemográfico interdisciplinaridade neuropsicológica avaliação}
\chapter{Psicopedagogicamente inconstitucionalíssimo otorrinolaringologista socioeducativo ético-profissional sociodemográfico interdisciplinaridade neuropsicológica avaliação}
\chapter{Psicopedagogicamente inconstitucionalíssimo otorrinolaringologista socioeducativo ético-profissional sociodemográfico interdisciplinaridade neuropsicológica avaliação}
\chapter{Psicopedagogicamente inconstitucionalíssimo otorrinolaringologista socioeducativo ético-profissional sociodemográfico interdisciplinaridade neuropsicológica avaliação}
\chapter{Psicopedagogicamente inconstitucionalíssimo otorrinolaringologista socioeducativo ético-profissional sociodemográfico interdisciplinaridade neuropsicológica avaliação}
\chapter{Psicopedagogicamente inconstitucionalíssimo otorrinolaringologista socioeducativo ético-profissional sociodemográfico interdisciplinaridade neuropsicológica avaliação}
\chapter{Psicopedagogicamente inconstitucionalíssimo otorrinolaringologista socioeducativo ético-profissional sociodemográfico interdisciplinaridade neuropsicológica avaliação}
\chapter{Psicopedagogicamente inconstitucionalíssimo otorrinolaringologista socioeducativo ético-profissional sociodemográfico interdisciplinaridade neuropsicológica avaliação}
\chapter{Psicopedagogicamente inconstitucionalíssimo otorrinolaringologista socioeducativo ético-profissional sociodemográfico interdisciplinaridade neuropsicológica avaliação}
\end{document}
'   % R: rígido (até 0.5.x)
%   lualatex '\def\var{B}% !TeX program = lualatex
% ===================================================================
% sumario-longo.tex — mede o sumário sem hifenização (issue #35, "Fora do
% escopo"): com a hifenização desligada, a entrada longa de capítulo pode
% bater no número da página. Gera 3 x 16 capítulos, cada um com um prefixo
% (de 3 a 18 palavras) de uma frase longa: frase 1 de palavras médias,
% frase 2 com números, frase 3 só de palavras muito longas. Conta no log
% os "Overfull \hbox" do sumário: o sumário é a única parte do documento
% sem parágrafo de corpo, então todo Overfull "in paragraph" antes do
% primeiro \chapter é entrada que passou da margem.
%
% Uso (de uma pasta com livros_crp_acessivel_book.sty e crpsp-leg.sty no
% caminho; as fontes NEWJUNE por OSFONTDIR):
%   lualatex '\def\var{R}% !TeX program = lualatex
% ===================================================================
% sumario-longo.tex — mede o sumário sem hifenização (issue #35, "Fora do
% escopo"): com a hifenização desligada, a entrada longa de capítulo pode
% bater no número da página. Gera 3 x 16 capítulos, cada um com um prefixo
% (de 3 a 18 palavras) de uma frase longa: frase 1 de palavras médias,
% frase 2 com números, frase 3 só de palavras muito longas. Conta no log
% os "Overfull \hbox" do sumário: o sumário é a única parte do documento
% sem parágrafo de corpo, então todo Overfull "in paragraph" antes do
% primeiro \chapter é entrada que passou da margem.
%
% Uso (de uma pasta com livros_crp_acessivel_book.sty e crpsp-leg.sty no
% caminho; as fontes NEWJUNE por OSFONTDIR):
%   lualatex '\def\var{R}\input{sumario-longo}'   % R: rígido (até 0.5.x)
%   lualatex '\def\var{B}\input{sumario-longo}'   % B: plus 4em (o do guia)
%   lualatex '\def\var{C}\input{sumario-longo}'   % C: plus 4em + \tolerance 10000
%   lualatex '\def\var{S}\input{sumario-longo}'   % S: o que a .sty define
% Cada variante só redefine o \l@chapter depois da carga da .sty.
% ===================================================================
\DocumentMetadata{lang=pt-BR, pdfstandard=ua-2, pdfversion=2.0, tagging=on}
\documentclass[11pt,a5paper,twoside,openright]{book}
\usepackage{livros_crp_acessivel_book}
\hypersetup{pdftitle={Sumario longo}, pdfauthor={CRP SP}}
\makeatletter
\providecommand\var{S}
\def\crp@vR{R}\def\crp@vB{B}\def\crp@vC{C}
\def\crp@lchap#1{%  #1 = acréscimo à margem direita (rightskip)
  \renewcommand*\l@chapter[2]{%
    \ifnum \c@tocdepth >\m@ne
        \addpenalty{-\@highpenalty}\addvspace{0.5em \@plus\p@}%
        \setlength\@tempdima{1.5em}%
        \begingroup
            \parindent \z@ \rightskip \@pnumwidth #1
            \parfillskip -\@pnumwidth
            \leavevmode \bfseries
            \advance\leftskip\@tempdima \hskip -\leftskip
            ##1\nobreak\hfil
            \nobreak\hb@xt@\@pnumwidth{\hss ##2\kern-\p@\kern\p@}\par
            \penalty\@highpenalty
        \endgroup
    \fi}}
\ifx\var\crp@vR \crp@lchap{}\fi
\ifx\var\crp@vB \crp@lchap{\@plus 4em}\fi
\ifx\var\crp@vC \crp@lchap{\@plus 4em \tolerance=10000 \hbadness=10000 }\fi
\makeatother
\begin{document}
\frontmatter
\tableofcontents
\mainmatter
\chapter{Deficiência psicossocial e}
\chapter{Deficiência psicossocial e transtorno}
\chapter{Deficiência psicossocial e transtorno do}
\chapter{Deficiência psicossocial e transtorno do espectro}
\chapter{Deficiência psicossocial e transtorno do espectro autista:}
\chapter{Deficiência psicossocial e transtorno do espectro autista: orientações}
\chapter{Deficiência psicossocial e transtorno do espectro autista: orientações para}
\chapter{Deficiência psicossocial e transtorno do espectro autista: orientações para a}
\chapter{Deficiência psicossocial e transtorno do espectro autista: orientações para a realização}
\chapter{Deficiência psicossocial e transtorno do espectro autista: orientações para a realização de}
\chapter{Deficiência psicossocial e transtorno do espectro autista: orientações para a realização de eventos}
\chapter{Deficiência psicossocial e transtorno do espectro autista: orientações para a realização de eventos e}
\chapter{Deficiência psicossocial e transtorno do espectro autista: orientações para a realização de eventos e apresentações}
\chapter{Deficiência psicossocial e transtorno do espectro autista: orientações para a realização de eventos e apresentações acessíveis}
\chapter{Deficiência psicossocial e transtorno do espectro autista: orientações para a realização de eventos e apresentações acessíveis em}
\chapter{Deficiência psicossocial e transtorno do espectro autista: orientações para a realização de eventos e apresentações acessíveis em todo}
\chapter{Comunicação alternativa e}
\chapter{Comunicação alternativa e aumentativa}
\chapter{Comunicação alternativa e aumentativa na}
\chapter{Comunicação alternativa e aumentativa na participação}
\chapter{Comunicação alternativa e aumentativa na participação social}
\chapter{Comunicação alternativa e aumentativa na participação social de}
\chapter{Comunicação alternativa e aumentativa na participação social de pessoas}
\chapter{Comunicação alternativa e aumentativa na participação social de pessoas com}
\chapter{Comunicação alternativa e aumentativa na participação social de pessoas com deficiência}
\chapter{Comunicação alternativa e aumentativa na participação social de pessoas com deficiência intelectual}
\chapter{Comunicação alternativa e aumentativa na participação social de pessoas com deficiência intelectual em}
\chapter{Comunicação alternativa e aumentativa na participação social de pessoas com deficiência intelectual em eventos}
\chapter{Comunicação alternativa e aumentativa na participação social de pessoas com deficiência intelectual em eventos públicos}
\chapter{Comunicação alternativa e aumentativa na participação social de pessoas com deficiência intelectual em eventos públicos de}
\chapter{Comunicação alternativa e aumentativa na participação social de pessoas com deficiência intelectual em eventos públicos de grande}
\chapter{Comunicação alternativa e aumentativa na participação social de pessoas com deficiência intelectual em eventos públicos de grande porte}
\chapter{Psicopedagogicamente inconstitucionalíssimo otorrinolaringologista}
\chapter{Psicopedagogicamente inconstitucionalíssimo otorrinolaringologista socioeducativo}
\chapter{Psicopedagogicamente inconstitucionalíssimo otorrinolaringologista socioeducativo ético-profissional}
\chapter{Psicopedagogicamente inconstitucionalíssimo otorrinolaringologista socioeducativo ético-profissional sociodemográfico}
\chapter{Psicopedagogicamente inconstitucionalíssimo otorrinolaringologista socioeducativo ético-profissional sociodemográfico interdisciplinaridade}
\chapter{Psicopedagogicamente inconstitucionalíssimo otorrinolaringologista socioeducativo ético-profissional sociodemográfico interdisciplinaridade neuropsicológica}
\chapter{Psicopedagogicamente inconstitucionalíssimo otorrinolaringologista socioeducativo ético-profissional sociodemográfico interdisciplinaridade neuropsicológica avaliação}
\chapter{Psicopedagogicamente inconstitucionalíssimo otorrinolaringologista socioeducativo ético-profissional sociodemográfico interdisciplinaridade neuropsicológica avaliação}
\chapter{Psicopedagogicamente inconstitucionalíssimo otorrinolaringologista socioeducativo ético-profissional sociodemográfico interdisciplinaridade neuropsicológica avaliação}
\chapter{Psicopedagogicamente inconstitucionalíssimo otorrinolaringologista socioeducativo ético-profissional sociodemográfico interdisciplinaridade neuropsicológica avaliação}
\chapter{Psicopedagogicamente inconstitucionalíssimo otorrinolaringologista socioeducativo ético-profissional sociodemográfico interdisciplinaridade neuropsicológica avaliação}
\chapter{Psicopedagogicamente inconstitucionalíssimo otorrinolaringologista socioeducativo ético-profissional sociodemográfico interdisciplinaridade neuropsicológica avaliação}
\chapter{Psicopedagogicamente inconstitucionalíssimo otorrinolaringologista socioeducativo ético-profissional sociodemográfico interdisciplinaridade neuropsicológica avaliação}
\chapter{Psicopedagogicamente inconstitucionalíssimo otorrinolaringologista socioeducativo ético-profissional sociodemográfico interdisciplinaridade neuropsicológica avaliação}
\chapter{Psicopedagogicamente inconstitucionalíssimo otorrinolaringologista socioeducativo ético-profissional sociodemográfico interdisciplinaridade neuropsicológica avaliação}
\chapter{Psicopedagogicamente inconstitucionalíssimo otorrinolaringologista socioeducativo ético-profissional sociodemográfico interdisciplinaridade neuropsicológica avaliação}
\end{document}
'   % R: rígido (até 0.5.x)
%   lualatex '\def\var{B}% !TeX program = lualatex
% ===================================================================
% sumario-longo.tex — mede o sumário sem hifenização (issue #35, "Fora do
% escopo"): com a hifenização desligada, a entrada longa de capítulo pode
% bater no número da página. Gera 3 x 16 capítulos, cada um com um prefixo
% (de 3 a 18 palavras) de uma frase longa: frase 1 de palavras médias,
% frase 2 com números, frase 3 só de palavras muito longas. Conta no log
% os "Overfull \hbox" do sumário: o sumário é a única parte do documento
% sem parágrafo de corpo, então todo Overfull "in paragraph" antes do
% primeiro \chapter é entrada que passou da margem.
%
% Uso (de uma pasta com livros_crp_acessivel_book.sty e crpsp-leg.sty no
% caminho; as fontes NEWJUNE por OSFONTDIR):
%   lualatex '\def\var{R}\input{sumario-longo}'   % R: rígido (até 0.5.x)
%   lualatex '\def\var{B}\input{sumario-longo}'   % B: plus 4em (o do guia)
%   lualatex '\def\var{C}\input{sumario-longo}'   % C: plus 4em + \tolerance 10000
%   lualatex '\def\var{S}\input{sumario-longo}'   % S: o que a .sty define
% Cada variante só redefine o \l@chapter depois da carga da .sty.
% ===================================================================
\DocumentMetadata{lang=pt-BR, pdfstandard=ua-2, pdfversion=2.0, tagging=on}
\documentclass[11pt,a5paper,twoside,openright]{book}
\usepackage{livros_crp_acessivel_book}
\hypersetup{pdftitle={Sumario longo}, pdfauthor={CRP SP}}
\makeatletter
\providecommand\var{S}
\def\crp@vR{R}\def\crp@vB{B}\def\crp@vC{C}
\def\crp@lchap#1{%  #1 = acréscimo à margem direita (rightskip)
  \renewcommand*\l@chapter[2]{%
    \ifnum \c@tocdepth >\m@ne
        \addpenalty{-\@highpenalty}\addvspace{0.5em \@plus\p@}%
        \setlength\@tempdima{1.5em}%
        \begingroup
            \parindent \z@ \rightskip \@pnumwidth #1
            \parfillskip -\@pnumwidth
            \leavevmode \bfseries
            \advance\leftskip\@tempdima \hskip -\leftskip
            ##1\nobreak\hfil
            \nobreak\hb@xt@\@pnumwidth{\hss ##2\kern-\p@\kern\p@}\par
            \penalty\@highpenalty
        \endgroup
    \fi}}
\ifx\var\crp@vR \crp@lchap{}\fi
\ifx\var\crp@vB \crp@lchap{\@plus 4em}\fi
\ifx\var\crp@vC \crp@lchap{\@plus 4em \tolerance=10000 \hbadness=10000 }\fi
\makeatother
\begin{document}
\frontmatter
\tableofcontents
\mainmatter
\chapter{Deficiência psicossocial e}
\chapter{Deficiência psicossocial e transtorno}
\chapter{Deficiência psicossocial e transtorno do}
\chapter{Deficiência psicossocial e transtorno do espectro}
\chapter{Deficiência psicossocial e transtorno do espectro autista:}
\chapter{Deficiência psicossocial e transtorno do espectro autista: orientações}
\chapter{Deficiência psicossocial e transtorno do espectro autista: orientações para}
\chapter{Deficiência psicossocial e transtorno do espectro autista: orientações para a}
\chapter{Deficiência psicossocial e transtorno do espectro autista: orientações para a realização}
\chapter{Deficiência psicossocial e transtorno do espectro autista: orientações para a realização de}
\chapter{Deficiência psicossocial e transtorno do espectro autista: orientações para a realização de eventos}
\chapter{Deficiência psicossocial e transtorno do espectro autista: orientações para a realização de eventos e}
\chapter{Deficiência psicossocial e transtorno do espectro autista: orientações para a realização de eventos e apresentações}
\chapter{Deficiência psicossocial e transtorno do espectro autista: orientações para a realização de eventos e apresentações acessíveis}
\chapter{Deficiência psicossocial e transtorno do espectro autista: orientações para a realização de eventos e apresentações acessíveis em}
\chapter{Deficiência psicossocial e transtorno do espectro autista: orientações para a realização de eventos e apresentações acessíveis em todo}
\chapter{Comunicação alternativa e}
\chapter{Comunicação alternativa e aumentativa}
\chapter{Comunicação alternativa e aumentativa na}
\chapter{Comunicação alternativa e aumentativa na participação}
\chapter{Comunicação alternativa e aumentativa na participação social}
\chapter{Comunicação alternativa e aumentativa na participação social de}
\chapter{Comunicação alternativa e aumentativa na participação social de pessoas}
\chapter{Comunicação alternativa e aumentativa na participação social de pessoas com}
\chapter{Comunicação alternativa e aumentativa na participação social de pessoas com deficiência}
\chapter{Comunicação alternativa e aumentativa na participação social de pessoas com deficiência intelectual}
\chapter{Comunicação alternativa e aumentativa na participação social de pessoas com deficiência intelectual em}
\chapter{Comunicação alternativa e aumentativa na participação social de pessoas com deficiência intelectual em eventos}
\chapter{Comunicação alternativa e aumentativa na participação social de pessoas com deficiência intelectual em eventos públicos}
\chapter{Comunicação alternativa e aumentativa na participação social de pessoas com deficiência intelectual em eventos públicos de}
\chapter{Comunicação alternativa e aumentativa na participação social de pessoas com deficiência intelectual em eventos públicos de grande}
\chapter{Comunicação alternativa e aumentativa na participação social de pessoas com deficiência intelectual em eventos públicos de grande porte}
\chapter{Psicopedagogicamente inconstitucionalíssimo otorrinolaringologista}
\chapter{Psicopedagogicamente inconstitucionalíssimo otorrinolaringologista socioeducativo}
\chapter{Psicopedagogicamente inconstitucionalíssimo otorrinolaringologista socioeducativo ético-profissional}
\chapter{Psicopedagogicamente inconstitucionalíssimo otorrinolaringologista socioeducativo ético-profissional sociodemográfico}
\chapter{Psicopedagogicamente inconstitucionalíssimo otorrinolaringologista socioeducativo ético-profissional sociodemográfico interdisciplinaridade}
\chapter{Psicopedagogicamente inconstitucionalíssimo otorrinolaringologista socioeducativo ético-profissional sociodemográfico interdisciplinaridade neuropsicológica}
\chapter{Psicopedagogicamente inconstitucionalíssimo otorrinolaringologista socioeducativo ético-profissional sociodemográfico interdisciplinaridade neuropsicológica avaliação}
\chapter{Psicopedagogicamente inconstitucionalíssimo otorrinolaringologista socioeducativo ético-profissional sociodemográfico interdisciplinaridade neuropsicológica avaliação}
\chapter{Psicopedagogicamente inconstitucionalíssimo otorrinolaringologista socioeducativo ético-profissional sociodemográfico interdisciplinaridade neuropsicológica avaliação}
\chapter{Psicopedagogicamente inconstitucionalíssimo otorrinolaringologista socioeducativo ético-profissional sociodemográfico interdisciplinaridade neuropsicológica avaliação}
\chapter{Psicopedagogicamente inconstitucionalíssimo otorrinolaringologista socioeducativo ético-profissional sociodemográfico interdisciplinaridade neuropsicológica avaliação}
\chapter{Psicopedagogicamente inconstitucionalíssimo otorrinolaringologista socioeducativo ético-profissional sociodemográfico interdisciplinaridade neuropsicológica avaliação}
\chapter{Psicopedagogicamente inconstitucionalíssimo otorrinolaringologista socioeducativo ético-profissional sociodemográfico interdisciplinaridade neuropsicológica avaliação}
\chapter{Psicopedagogicamente inconstitucionalíssimo otorrinolaringologista socioeducativo ético-profissional sociodemográfico interdisciplinaridade neuropsicológica avaliação}
\chapter{Psicopedagogicamente inconstitucionalíssimo otorrinolaringologista socioeducativo ético-profissional sociodemográfico interdisciplinaridade neuropsicológica avaliação}
\chapter{Psicopedagogicamente inconstitucionalíssimo otorrinolaringologista socioeducativo ético-profissional sociodemográfico interdisciplinaridade neuropsicológica avaliação}
\end{document}
'   % B: plus 4em (o do guia)
%   lualatex '\def\var{C}% !TeX program = lualatex
% ===================================================================
% sumario-longo.tex — mede o sumário sem hifenização (issue #35, "Fora do
% escopo"): com a hifenização desligada, a entrada longa de capítulo pode
% bater no número da página. Gera 3 x 16 capítulos, cada um com um prefixo
% (de 3 a 18 palavras) de uma frase longa: frase 1 de palavras médias,
% frase 2 com números, frase 3 só de palavras muito longas. Conta no log
% os "Overfull \hbox" do sumário: o sumário é a única parte do documento
% sem parágrafo de corpo, então todo Overfull "in paragraph" antes do
% primeiro \chapter é entrada que passou da margem.
%
% Uso (de uma pasta com livros_crp_acessivel_book.sty e crpsp-leg.sty no
% caminho; as fontes NEWJUNE por OSFONTDIR):
%   lualatex '\def\var{R}\input{sumario-longo}'   % R: rígido (até 0.5.x)
%   lualatex '\def\var{B}\input{sumario-longo}'   % B: plus 4em (o do guia)
%   lualatex '\def\var{C}\input{sumario-longo}'   % C: plus 4em + \tolerance 10000
%   lualatex '\def\var{S}\input{sumario-longo}'   % S: o que a .sty define
% Cada variante só redefine o \l@chapter depois da carga da .sty.
% ===================================================================
\DocumentMetadata{lang=pt-BR, pdfstandard=ua-2, pdfversion=2.0, tagging=on}
\documentclass[11pt,a5paper,twoside,openright]{book}
\usepackage{livros_crp_acessivel_book}
\hypersetup{pdftitle={Sumario longo}, pdfauthor={CRP SP}}
\makeatletter
\providecommand\var{S}
\def\crp@vR{R}\def\crp@vB{B}\def\crp@vC{C}
\def\crp@lchap#1{%  #1 = acréscimo à margem direita (rightskip)
  \renewcommand*\l@chapter[2]{%
    \ifnum \c@tocdepth >\m@ne
        \addpenalty{-\@highpenalty}\addvspace{0.5em \@plus\p@}%
        \setlength\@tempdima{1.5em}%
        \begingroup
            \parindent \z@ \rightskip \@pnumwidth #1
            \parfillskip -\@pnumwidth
            \leavevmode \bfseries
            \advance\leftskip\@tempdima \hskip -\leftskip
            ##1\nobreak\hfil
            \nobreak\hb@xt@\@pnumwidth{\hss ##2\kern-\p@\kern\p@}\par
            \penalty\@highpenalty
        \endgroup
    \fi}}
\ifx\var\crp@vR \crp@lchap{}\fi
\ifx\var\crp@vB \crp@lchap{\@plus 4em}\fi
\ifx\var\crp@vC \crp@lchap{\@plus 4em \tolerance=10000 \hbadness=10000 }\fi
\makeatother
\begin{document}
\frontmatter
\tableofcontents
\mainmatter
\chapter{Deficiência psicossocial e}
\chapter{Deficiência psicossocial e transtorno}
\chapter{Deficiência psicossocial e transtorno do}
\chapter{Deficiência psicossocial e transtorno do espectro}
\chapter{Deficiência psicossocial e transtorno do espectro autista:}
\chapter{Deficiência psicossocial e transtorno do espectro autista: orientações}
\chapter{Deficiência psicossocial e transtorno do espectro autista: orientações para}
\chapter{Deficiência psicossocial e transtorno do espectro autista: orientações para a}
\chapter{Deficiência psicossocial e transtorno do espectro autista: orientações para a realização}
\chapter{Deficiência psicossocial e transtorno do espectro autista: orientações para a realização de}
\chapter{Deficiência psicossocial e transtorno do espectro autista: orientações para a realização de eventos}
\chapter{Deficiência psicossocial e transtorno do espectro autista: orientações para a realização de eventos e}
\chapter{Deficiência psicossocial e transtorno do espectro autista: orientações para a realização de eventos e apresentações}
\chapter{Deficiência psicossocial e transtorno do espectro autista: orientações para a realização de eventos e apresentações acessíveis}
\chapter{Deficiência psicossocial e transtorno do espectro autista: orientações para a realização de eventos e apresentações acessíveis em}
\chapter{Deficiência psicossocial e transtorno do espectro autista: orientações para a realização de eventos e apresentações acessíveis em todo}
\chapter{Comunicação alternativa e}
\chapter{Comunicação alternativa e aumentativa}
\chapter{Comunicação alternativa e aumentativa na}
\chapter{Comunicação alternativa e aumentativa na participação}
\chapter{Comunicação alternativa e aumentativa na participação social}
\chapter{Comunicação alternativa e aumentativa na participação social de}
\chapter{Comunicação alternativa e aumentativa na participação social de pessoas}
\chapter{Comunicação alternativa e aumentativa na participação social de pessoas com}
\chapter{Comunicação alternativa e aumentativa na participação social de pessoas com deficiência}
\chapter{Comunicação alternativa e aumentativa na participação social de pessoas com deficiência intelectual}
\chapter{Comunicação alternativa e aumentativa na participação social de pessoas com deficiência intelectual em}
\chapter{Comunicação alternativa e aumentativa na participação social de pessoas com deficiência intelectual em eventos}
\chapter{Comunicação alternativa e aumentativa na participação social de pessoas com deficiência intelectual em eventos públicos}
\chapter{Comunicação alternativa e aumentativa na participação social de pessoas com deficiência intelectual em eventos públicos de}
\chapter{Comunicação alternativa e aumentativa na participação social de pessoas com deficiência intelectual em eventos públicos de grande}
\chapter{Comunicação alternativa e aumentativa na participação social de pessoas com deficiência intelectual em eventos públicos de grande porte}
\chapter{Psicopedagogicamente inconstitucionalíssimo otorrinolaringologista}
\chapter{Psicopedagogicamente inconstitucionalíssimo otorrinolaringologista socioeducativo}
\chapter{Psicopedagogicamente inconstitucionalíssimo otorrinolaringologista socioeducativo ético-profissional}
\chapter{Psicopedagogicamente inconstitucionalíssimo otorrinolaringologista socioeducativo ético-profissional sociodemográfico}
\chapter{Psicopedagogicamente inconstitucionalíssimo otorrinolaringologista socioeducativo ético-profissional sociodemográfico interdisciplinaridade}
\chapter{Psicopedagogicamente inconstitucionalíssimo otorrinolaringologista socioeducativo ético-profissional sociodemográfico interdisciplinaridade neuropsicológica}
\chapter{Psicopedagogicamente inconstitucionalíssimo otorrinolaringologista socioeducativo ético-profissional sociodemográfico interdisciplinaridade neuropsicológica avaliação}
\chapter{Psicopedagogicamente inconstitucionalíssimo otorrinolaringologista socioeducativo ético-profissional sociodemográfico interdisciplinaridade neuropsicológica avaliação}
\chapter{Psicopedagogicamente inconstitucionalíssimo otorrinolaringologista socioeducativo ético-profissional sociodemográfico interdisciplinaridade neuropsicológica avaliação}
\chapter{Psicopedagogicamente inconstitucionalíssimo otorrinolaringologista socioeducativo ético-profissional sociodemográfico interdisciplinaridade neuropsicológica avaliação}
\chapter{Psicopedagogicamente inconstitucionalíssimo otorrinolaringologista socioeducativo ético-profissional sociodemográfico interdisciplinaridade neuropsicológica avaliação}
\chapter{Psicopedagogicamente inconstitucionalíssimo otorrinolaringologista socioeducativo ético-profissional sociodemográfico interdisciplinaridade neuropsicológica avaliação}
\chapter{Psicopedagogicamente inconstitucionalíssimo otorrinolaringologista socioeducativo ético-profissional sociodemográfico interdisciplinaridade neuropsicológica avaliação}
\chapter{Psicopedagogicamente inconstitucionalíssimo otorrinolaringologista socioeducativo ético-profissional sociodemográfico interdisciplinaridade neuropsicológica avaliação}
\chapter{Psicopedagogicamente inconstitucionalíssimo otorrinolaringologista socioeducativo ético-profissional sociodemográfico interdisciplinaridade neuropsicológica avaliação}
\chapter{Psicopedagogicamente inconstitucionalíssimo otorrinolaringologista socioeducativo ético-profissional sociodemográfico interdisciplinaridade neuropsicológica avaliação}
\end{document}
'   % C: plus 4em + \tolerance 10000
%   lualatex '\def\var{S}% !TeX program = lualatex
% ===================================================================
% sumario-longo.tex — mede o sumário sem hifenização (issue #35, "Fora do
% escopo"): com a hifenização desligada, a entrada longa de capítulo pode
% bater no número da página. Gera 3 x 16 capítulos, cada um com um prefixo
% (de 3 a 18 palavras) de uma frase longa: frase 1 de palavras médias,
% frase 2 com números, frase 3 só de palavras muito longas. Conta no log
% os "Overfull \hbox" do sumário: o sumário é a única parte do documento
% sem parágrafo de corpo, então todo Overfull "in paragraph" antes do
% primeiro \chapter é entrada que passou da margem.
%
% Uso (de uma pasta com livros_crp_acessivel_book.sty e crpsp-leg.sty no
% caminho; as fontes NEWJUNE por OSFONTDIR):
%   lualatex '\def\var{R}\input{sumario-longo}'   % R: rígido (até 0.5.x)
%   lualatex '\def\var{B}\input{sumario-longo}'   % B: plus 4em (o do guia)
%   lualatex '\def\var{C}\input{sumario-longo}'   % C: plus 4em + \tolerance 10000
%   lualatex '\def\var{S}\input{sumario-longo}'   % S: o que a .sty define
% Cada variante só redefine o \l@chapter depois da carga da .sty.
% ===================================================================
\DocumentMetadata{lang=pt-BR, pdfstandard=ua-2, pdfversion=2.0, tagging=on}
\documentclass[11pt,a5paper,twoside,openright]{book}
\usepackage{livros_crp_acessivel_book}
\hypersetup{pdftitle={Sumario longo}, pdfauthor={CRP SP}}
\makeatletter
\providecommand\var{S}
\def\crp@vR{R}\def\crp@vB{B}\def\crp@vC{C}
\def\crp@lchap#1{%  #1 = acréscimo à margem direita (rightskip)
  \renewcommand*\l@chapter[2]{%
    \ifnum \c@tocdepth >\m@ne
        \addpenalty{-\@highpenalty}\addvspace{0.5em \@plus\p@}%
        \setlength\@tempdima{1.5em}%
        \begingroup
            \parindent \z@ \rightskip \@pnumwidth #1
            \parfillskip -\@pnumwidth
            \leavevmode \bfseries
            \advance\leftskip\@tempdima \hskip -\leftskip
            ##1\nobreak\hfil
            \nobreak\hb@xt@\@pnumwidth{\hss ##2\kern-\p@\kern\p@}\par
            \penalty\@highpenalty
        \endgroup
    \fi}}
\ifx\var\crp@vR \crp@lchap{}\fi
\ifx\var\crp@vB \crp@lchap{\@plus 4em}\fi
\ifx\var\crp@vC \crp@lchap{\@plus 4em \tolerance=10000 \hbadness=10000 }\fi
\makeatother
\begin{document}
\frontmatter
\tableofcontents
\mainmatter
\chapter{Deficiência psicossocial e}
\chapter{Deficiência psicossocial e transtorno}
\chapter{Deficiência psicossocial e transtorno do}
\chapter{Deficiência psicossocial e transtorno do espectro}
\chapter{Deficiência psicossocial e transtorno do espectro autista:}
\chapter{Deficiência psicossocial e transtorno do espectro autista: orientações}
\chapter{Deficiência psicossocial e transtorno do espectro autista: orientações para}
\chapter{Deficiência psicossocial e transtorno do espectro autista: orientações para a}
\chapter{Deficiência psicossocial e transtorno do espectro autista: orientações para a realização}
\chapter{Deficiência psicossocial e transtorno do espectro autista: orientações para a realização de}
\chapter{Deficiência psicossocial e transtorno do espectro autista: orientações para a realização de eventos}
\chapter{Deficiência psicossocial e transtorno do espectro autista: orientações para a realização de eventos e}
\chapter{Deficiência psicossocial e transtorno do espectro autista: orientações para a realização de eventos e apresentações}
\chapter{Deficiência psicossocial e transtorno do espectro autista: orientações para a realização de eventos e apresentações acessíveis}
\chapter{Deficiência psicossocial e transtorno do espectro autista: orientações para a realização de eventos e apresentações acessíveis em}
\chapter{Deficiência psicossocial e transtorno do espectro autista: orientações para a realização de eventos e apresentações acessíveis em todo}
\chapter{Comunicação alternativa e}
\chapter{Comunicação alternativa e aumentativa}
\chapter{Comunicação alternativa e aumentativa na}
\chapter{Comunicação alternativa e aumentativa na participação}
\chapter{Comunicação alternativa e aumentativa na participação social}
\chapter{Comunicação alternativa e aumentativa na participação social de}
\chapter{Comunicação alternativa e aumentativa na participação social de pessoas}
\chapter{Comunicação alternativa e aumentativa na participação social de pessoas com}
\chapter{Comunicação alternativa e aumentativa na participação social de pessoas com deficiência}
\chapter{Comunicação alternativa e aumentativa na participação social de pessoas com deficiência intelectual}
\chapter{Comunicação alternativa e aumentativa na participação social de pessoas com deficiência intelectual em}
\chapter{Comunicação alternativa e aumentativa na participação social de pessoas com deficiência intelectual em eventos}
\chapter{Comunicação alternativa e aumentativa na participação social de pessoas com deficiência intelectual em eventos públicos}
\chapter{Comunicação alternativa e aumentativa na participação social de pessoas com deficiência intelectual em eventos públicos de}
\chapter{Comunicação alternativa e aumentativa na participação social de pessoas com deficiência intelectual em eventos públicos de grande}
\chapter{Comunicação alternativa e aumentativa na participação social de pessoas com deficiência intelectual em eventos públicos de grande porte}
\chapter{Psicopedagogicamente inconstitucionalíssimo otorrinolaringologista}
\chapter{Psicopedagogicamente inconstitucionalíssimo otorrinolaringologista socioeducativo}
\chapter{Psicopedagogicamente inconstitucionalíssimo otorrinolaringologista socioeducativo ético-profissional}
\chapter{Psicopedagogicamente inconstitucionalíssimo otorrinolaringologista socioeducativo ético-profissional sociodemográfico}
\chapter{Psicopedagogicamente inconstitucionalíssimo otorrinolaringologista socioeducativo ético-profissional sociodemográfico interdisciplinaridade}
\chapter{Psicopedagogicamente inconstitucionalíssimo otorrinolaringologista socioeducativo ético-profissional sociodemográfico interdisciplinaridade neuropsicológica}
\chapter{Psicopedagogicamente inconstitucionalíssimo otorrinolaringologista socioeducativo ético-profissional sociodemográfico interdisciplinaridade neuropsicológica avaliação}
\chapter{Psicopedagogicamente inconstitucionalíssimo otorrinolaringologista socioeducativo ético-profissional sociodemográfico interdisciplinaridade neuropsicológica avaliação}
\chapter{Psicopedagogicamente inconstitucionalíssimo otorrinolaringologista socioeducativo ético-profissional sociodemográfico interdisciplinaridade neuropsicológica avaliação}
\chapter{Psicopedagogicamente inconstitucionalíssimo otorrinolaringologista socioeducativo ético-profissional sociodemográfico interdisciplinaridade neuropsicológica avaliação}
\chapter{Psicopedagogicamente inconstitucionalíssimo otorrinolaringologista socioeducativo ético-profissional sociodemográfico interdisciplinaridade neuropsicológica avaliação}
\chapter{Psicopedagogicamente inconstitucionalíssimo otorrinolaringologista socioeducativo ético-profissional sociodemográfico interdisciplinaridade neuropsicológica avaliação}
\chapter{Psicopedagogicamente inconstitucionalíssimo otorrinolaringologista socioeducativo ético-profissional sociodemográfico interdisciplinaridade neuropsicológica avaliação}
\chapter{Psicopedagogicamente inconstitucionalíssimo otorrinolaringologista socioeducativo ético-profissional sociodemográfico interdisciplinaridade neuropsicológica avaliação}
\chapter{Psicopedagogicamente inconstitucionalíssimo otorrinolaringologista socioeducativo ético-profissional sociodemográfico interdisciplinaridade neuropsicológica avaliação}
\chapter{Psicopedagogicamente inconstitucionalíssimo otorrinolaringologista socioeducativo ético-profissional sociodemográfico interdisciplinaridade neuropsicológica avaliação}
\end{document}
'   % S: o que a .sty define
% Cada variante só redefine o \l@chapter depois da carga da .sty.
% ===================================================================
\DocumentMetadata{lang=pt-BR, pdfstandard=ua-2, pdfversion=2.0, tagging=on}
\documentclass[11pt,a5paper,twoside,openright]{book}
\usepackage{livros_crp_acessivel_book}
\hypersetup{pdftitle={Sumario longo}, pdfauthor={CRP SP}}
\makeatletter
\providecommand\var{S}
\def\crp@vR{R}\def\crp@vB{B}\def\crp@vC{C}
\def\crp@lchap#1{%  #1 = acréscimo à margem direita (rightskip)
  \renewcommand*\l@chapter[2]{%
    \ifnum \c@tocdepth >\m@ne
        \addpenalty{-\@highpenalty}\addvspace{0.5em \@plus\p@}%
        \setlength\@tempdima{1.5em}%
        \begingroup
            \parindent \z@ \rightskip \@pnumwidth #1
            \parfillskip -\@pnumwidth
            \leavevmode \bfseries
            \advance\leftskip\@tempdima \hskip -\leftskip
            ##1\nobreak\hfil
            \nobreak\hb@xt@\@pnumwidth{\hss ##2\kern-\p@\kern\p@}\par
            \penalty\@highpenalty
        \endgroup
    \fi}}
\ifx\var\crp@vR \crp@lchap{}\fi
\ifx\var\crp@vB \crp@lchap{\@plus 4em}\fi
\ifx\var\crp@vC \crp@lchap{\@plus 4em \tolerance=10000 \hbadness=10000 }\fi
\makeatother
\begin{document}
\frontmatter
\tableofcontents
\mainmatter
\chapter{Deficiência psicossocial e}
\chapter{Deficiência psicossocial e transtorno}
\chapter{Deficiência psicossocial e transtorno do}
\chapter{Deficiência psicossocial e transtorno do espectro}
\chapter{Deficiência psicossocial e transtorno do espectro autista:}
\chapter{Deficiência psicossocial e transtorno do espectro autista: orientações}
\chapter{Deficiência psicossocial e transtorno do espectro autista: orientações para}
\chapter{Deficiência psicossocial e transtorno do espectro autista: orientações para a}
\chapter{Deficiência psicossocial e transtorno do espectro autista: orientações para a realização}
\chapter{Deficiência psicossocial e transtorno do espectro autista: orientações para a realização de}
\chapter{Deficiência psicossocial e transtorno do espectro autista: orientações para a realização de eventos}
\chapter{Deficiência psicossocial e transtorno do espectro autista: orientações para a realização de eventos e}
\chapter{Deficiência psicossocial e transtorno do espectro autista: orientações para a realização de eventos e apresentações}
\chapter{Deficiência psicossocial e transtorno do espectro autista: orientações para a realização de eventos e apresentações acessíveis}
\chapter{Deficiência psicossocial e transtorno do espectro autista: orientações para a realização de eventos e apresentações acessíveis em}
\chapter{Deficiência psicossocial e transtorno do espectro autista: orientações para a realização de eventos e apresentações acessíveis em todo}
\chapter{Comunicação alternativa e}
\chapter{Comunicação alternativa e aumentativa}
\chapter{Comunicação alternativa e aumentativa na}
\chapter{Comunicação alternativa e aumentativa na participação}
\chapter{Comunicação alternativa e aumentativa na participação social}
\chapter{Comunicação alternativa e aumentativa na participação social de}
\chapter{Comunicação alternativa e aumentativa na participação social de pessoas}
\chapter{Comunicação alternativa e aumentativa na participação social de pessoas com}
\chapter{Comunicação alternativa e aumentativa na participação social de pessoas com deficiência}
\chapter{Comunicação alternativa e aumentativa na participação social de pessoas com deficiência intelectual}
\chapter{Comunicação alternativa e aumentativa na participação social de pessoas com deficiência intelectual em}
\chapter{Comunicação alternativa e aumentativa na participação social de pessoas com deficiência intelectual em eventos}
\chapter{Comunicação alternativa e aumentativa na participação social de pessoas com deficiência intelectual em eventos públicos}
\chapter{Comunicação alternativa e aumentativa na participação social de pessoas com deficiência intelectual em eventos públicos de}
\chapter{Comunicação alternativa e aumentativa na participação social de pessoas com deficiência intelectual em eventos públicos de grande}
\chapter{Comunicação alternativa e aumentativa na participação social de pessoas com deficiência intelectual em eventos públicos de grande porte}
\chapter{Psicopedagogicamente inconstitucionalíssimo otorrinolaringologista}
\chapter{Psicopedagogicamente inconstitucionalíssimo otorrinolaringologista socioeducativo}
\chapter{Psicopedagogicamente inconstitucionalíssimo otorrinolaringologista socioeducativo ético-profissional}
\chapter{Psicopedagogicamente inconstitucionalíssimo otorrinolaringologista socioeducativo ético-profissional sociodemográfico}
\chapter{Psicopedagogicamente inconstitucionalíssimo otorrinolaringologista socioeducativo ético-profissional sociodemográfico interdisciplinaridade}
\chapter{Psicopedagogicamente inconstitucionalíssimo otorrinolaringologista socioeducativo ético-profissional sociodemográfico interdisciplinaridade neuropsicológica}
\chapter{Psicopedagogicamente inconstitucionalíssimo otorrinolaringologista socioeducativo ético-profissional sociodemográfico interdisciplinaridade neuropsicológica avaliação}
\chapter{Psicopedagogicamente inconstitucionalíssimo otorrinolaringologista socioeducativo ético-profissional sociodemográfico interdisciplinaridade neuropsicológica avaliação}
\chapter{Psicopedagogicamente inconstitucionalíssimo otorrinolaringologista socioeducativo ético-profissional sociodemográfico interdisciplinaridade neuropsicológica avaliação}
\chapter{Psicopedagogicamente inconstitucionalíssimo otorrinolaringologista socioeducativo ético-profissional sociodemográfico interdisciplinaridade neuropsicológica avaliação}
\chapter{Psicopedagogicamente inconstitucionalíssimo otorrinolaringologista socioeducativo ético-profissional sociodemográfico interdisciplinaridade neuropsicológica avaliação}
\chapter{Psicopedagogicamente inconstitucionalíssimo otorrinolaringologista socioeducativo ético-profissional sociodemográfico interdisciplinaridade neuropsicológica avaliação}
\chapter{Psicopedagogicamente inconstitucionalíssimo otorrinolaringologista socioeducativo ético-profissional sociodemográfico interdisciplinaridade neuropsicológica avaliação}
\chapter{Psicopedagogicamente inconstitucionalíssimo otorrinolaringologista socioeducativo ético-profissional sociodemográfico interdisciplinaridade neuropsicológica avaliação}
\chapter{Psicopedagogicamente inconstitucionalíssimo otorrinolaringologista socioeducativo ético-profissional sociodemográfico interdisciplinaridade neuropsicológica avaliação}
\chapter{Psicopedagogicamente inconstitucionalíssimo otorrinolaringologista socioeducativo ético-profissional sociodemográfico interdisciplinaridade neuropsicológica avaliação}
\end{document}
'   % B: plus 4em (o do guia)
%   lualatex '\def\var{C}% !TeX program = lualatex
% ===================================================================
% sumario-longo.tex — mede o sumário sem hifenização (issue #35, "Fora do
% escopo"): com a hifenização desligada, a entrada longa de capítulo pode
% bater no número da página. Gera 3 x 16 capítulos, cada um com um prefixo
% (de 3 a 18 palavras) de uma frase longa: frase 1 de palavras médias,
% frase 2 com números, frase 3 só de palavras muito longas. Conta no log
% os "Overfull \hbox" do sumário: o sumário é a única parte do documento
% sem parágrafo de corpo, então todo Overfull "in paragraph" antes do
% primeiro \chapter é entrada que passou da margem.
%
% Uso (de uma pasta com livros_crp_acessivel_book.sty e crpsp-leg.sty no
% caminho; as fontes NEWJUNE por OSFONTDIR):
%   lualatex '\def\var{R}% !TeX program = lualatex
% ===================================================================
% sumario-longo.tex — mede o sumário sem hifenização (issue #35, "Fora do
% escopo"): com a hifenização desligada, a entrada longa de capítulo pode
% bater no número da página. Gera 3 x 16 capítulos, cada um com um prefixo
% (de 3 a 18 palavras) de uma frase longa: frase 1 de palavras médias,
% frase 2 com números, frase 3 só de palavras muito longas. Conta no log
% os "Overfull \hbox" do sumário: o sumário é a única parte do documento
% sem parágrafo de corpo, então todo Overfull "in paragraph" antes do
% primeiro \chapter é entrada que passou da margem.
%
% Uso (de uma pasta com livros_crp_acessivel_book.sty e crpsp-leg.sty no
% caminho; as fontes NEWJUNE por OSFONTDIR):
%   lualatex '\def\var{R}\input{sumario-longo}'   % R: rígido (até 0.5.x)
%   lualatex '\def\var{B}\input{sumario-longo}'   % B: plus 4em (o do guia)
%   lualatex '\def\var{C}\input{sumario-longo}'   % C: plus 4em + \tolerance 10000
%   lualatex '\def\var{S}\input{sumario-longo}'   % S: o que a .sty define
% Cada variante só redefine o \l@chapter depois da carga da .sty.
% ===================================================================
\DocumentMetadata{lang=pt-BR, pdfstandard=ua-2, pdfversion=2.0, tagging=on}
\documentclass[11pt,a5paper,twoside,openright]{book}
\usepackage{livros_crp_acessivel_book}
\hypersetup{pdftitle={Sumario longo}, pdfauthor={CRP SP}}
\makeatletter
\providecommand\var{S}
\def\crp@vR{R}\def\crp@vB{B}\def\crp@vC{C}
\def\crp@lchap#1{%  #1 = acréscimo à margem direita (rightskip)
  \renewcommand*\l@chapter[2]{%
    \ifnum \c@tocdepth >\m@ne
        \addpenalty{-\@highpenalty}\addvspace{0.5em \@plus\p@}%
        \setlength\@tempdima{1.5em}%
        \begingroup
            \parindent \z@ \rightskip \@pnumwidth #1
            \parfillskip -\@pnumwidth
            \leavevmode \bfseries
            \advance\leftskip\@tempdima \hskip -\leftskip
            ##1\nobreak\hfil
            \nobreak\hb@xt@\@pnumwidth{\hss ##2\kern-\p@\kern\p@}\par
            \penalty\@highpenalty
        \endgroup
    \fi}}
\ifx\var\crp@vR \crp@lchap{}\fi
\ifx\var\crp@vB \crp@lchap{\@plus 4em}\fi
\ifx\var\crp@vC \crp@lchap{\@plus 4em \tolerance=10000 \hbadness=10000 }\fi
\makeatother
\begin{document}
\frontmatter
\tableofcontents
\mainmatter
\chapter{Deficiência psicossocial e}
\chapter{Deficiência psicossocial e transtorno}
\chapter{Deficiência psicossocial e transtorno do}
\chapter{Deficiência psicossocial e transtorno do espectro}
\chapter{Deficiência psicossocial e transtorno do espectro autista:}
\chapter{Deficiência psicossocial e transtorno do espectro autista: orientações}
\chapter{Deficiência psicossocial e transtorno do espectro autista: orientações para}
\chapter{Deficiência psicossocial e transtorno do espectro autista: orientações para a}
\chapter{Deficiência psicossocial e transtorno do espectro autista: orientações para a realização}
\chapter{Deficiência psicossocial e transtorno do espectro autista: orientações para a realização de}
\chapter{Deficiência psicossocial e transtorno do espectro autista: orientações para a realização de eventos}
\chapter{Deficiência psicossocial e transtorno do espectro autista: orientações para a realização de eventos e}
\chapter{Deficiência psicossocial e transtorno do espectro autista: orientações para a realização de eventos e apresentações}
\chapter{Deficiência psicossocial e transtorno do espectro autista: orientações para a realização de eventos e apresentações acessíveis}
\chapter{Deficiência psicossocial e transtorno do espectro autista: orientações para a realização de eventos e apresentações acessíveis em}
\chapter{Deficiência psicossocial e transtorno do espectro autista: orientações para a realização de eventos e apresentações acessíveis em todo}
\chapter{Comunicação alternativa e}
\chapter{Comunicação alternativa e aumentativa}
\chapter{Comunicação alternativa e aumentativa na}
\chapter{Comunicação alternativa e aumentativa na participação}
\chapter{Comunicação alternativa e aumentativa na participação social}
\chapter{Comunicação alternativa e aumentativa na participação social de}
\chapter{Comunicação alternativa e aumentativa na participação social de pessoas}
\chapter{Comunicação alternativa e aumentativa na participação social de pessoas com}
\chapter{Comunicação alternativa e aumentativa na participação social de pessoas com deficiência}
\chapter{Comunicação alternativa e aumentativa na participação social de pessoas com deficiência intelectual}
\chapter{Comunicação alternativa e aumentativa na participação social de pessoas com deficiência intelectual em}
\chapter{Comunicação alternativa e aumentativa na participação social de pessoas com deficiência intelectual em eventos}
\chapter{Comunicação alternativa e aumentativa na participação social de pessoas com deficiência intelectual em eventos públicos}
\chapter{Comunicação alternativa e aumentativa na participação social de pessoas com deficiência intelectual em eventos públicos de}
\chapter{Comunicação alternativa e aumentativa na participação social de pessoas com deficiência intelectual em eventos públicos de grande}
\chapter{Comunicação alternativa e aumentativa na participação social de pessoas com deficiência intelectual em eventos públicos de grande porte}
\chapter{Psicopedagogicamente inconstitucionalíssimo otorrinolaringologista}
\chapter{Psicopedagogicamente inconstitucionalíssimo otorrinolaringologista socioeducativo}
\chapter{Psicopedagogicamente inconstitucionalíssimo otorrinolaringologista socioeducativo ético-profissional}
\chapter{Psicopedagogicamente inconstitucionalíssimo otorrinolaringologista socioeducativo ético-profissional sociodemográfico}
\chapter{Psicopedagogicamente inconstitucionalíssimo otorrinolaringologista socioeducativo ético-profissional sociodemográfico interdisciplinaridade}
\chapter{Psicopedagogicamente inconstitucionalíssimo otorrinolaringologista socioeducativo ético-profissional sociodemográfico interdisciplinaridade neuropsicológica}
\chapter{Psicopedagogicamente inconstitucionalíssimo otorrinolaringologista socioeducativo ético-profissional sociodemográfico interdisciplinaridade neuropsicológica avaliação}
\chapter{Psicopedagogicamente inconstitucionalíssimo otorrinolaringologista socioeducativo ético-profissional sociodemográfico interdisciplinaridade neuropsicológica avaliação}
\chapter{Psicopedagogicamente inconstitucionalíssimo otorrinolaringologista socioeducativo ético-profissional sociodemográfico interdisciplinaridade neuropsicológica avaliação}
\chapter{Psicopedagogicamente inconstitucionalíssimo otorrinolaringologista socioeducativo ético-profissional sociodemográfico interdisciplinaridade neuropsicológica avaliação}
\chapter{Psicopedagogicamente inconstitucionalíssimo otorrinolaringologista socioeducativo ético-profissional sociodemográfico interdisciplinaridade neuropsicológica avaliação}
\chapter{Psicopedagogicamente inconstitucionalíssimo otorrinolaringologista socioeducativo ético-profissional sociodemográfico interdisciplinaridade neuropsicológica avaliação}
\chapter{Psicopedagogicamente inconstitucionalíssimo otorrinolaringologista socioeducativo ético-profissional sociodemográfico interdisciplinaridade neuropsicológica avaliação}
\chapter{Psicopedagogicamente inconstitucionalíssimo otorrinolaringologista socioeducativo ético-profissional sociodemográfico interdisciplinaridade neuropsicológica avaliação}
\chapter{Psicopedagogicamente inconstitucionalíssimo otorrinolaringologista socioeducativo ético-profissional sociodemográfico interdisciplinaridade neuropsicológica avaliação}
\chapter{Psicopedagogicamente inconstitucionalíssimo otorrinolaringologista socioeducativo ético-profissional sociodemográfico interdisciplinaridade neuropsicológica avaliação}
\end{document}
'   % R: rígido (até 0.5.x)
%   lualatex '\def\var{B}% !TeX program = lualatex
% ===================================================================
% sumario-longo.tex — mede o sumário sem hifenização (issue #35, "Fora do
% escopo"): com a hifenização desligada, a entrada longa de capítulo pode
% bater no número da página. Gera 3 x 16 capítulos, cada um com um prefixo
% (de 3 a 18 palavras) de uma frase longa: frase 1 de palavras médias,
% frase 2 com números, frase 3 só de palavras muito longas. Conta no log
% os "Overfull \hbox" do sumário: o sumário é a única parte do documento
% sem parágrafo de corpo, então todo Overfull "in paragraph" antes do
% primeiro \chapter é entrada que passou da margem.
%
% Uso (de uma pasta com livros_crp_acessivel_book.sty e crpsp-leg.sty no
% caminho; as fontes NEWJUNE por OSFONTDIR):
%   lualatex '\def\var{R}\input{sumario-longo}'   % R: rígido (até 0.5.x)
%   lualatex '\def\var{B}\input{sumario-longo}'   % B: plus 4em (o do guia)
%   lualatex '\def\var{C}\input{sumario-longo}'   % C: plus 4em + \tolerance 10000
%   lualatex '\def\var{S}\input{sumario-longo}'   % S: o que a .sty define
% Cada variante só redefine o \l@chapter depois da carga da .sty.
% ===================================================================
\DocumentMetadata{lang=pt-BR, pdfstandard=ua-2, pdfversion=2.0, tagging=on}
\documentclass[11pt,a5paper,twoside,openright]{book}
\usepackage{livros_crp_acessivel_book}
\hypersetup{pdftitle={Sumario longo}, pdfauthor={CRP SP}}
\makeatletter
\providecommand\var{S}
\def\crp@vR{R}\def\crp@vB{B}\def\crp@vC{C}
\def\crp@lchap#1{%  #1 = acréscimo à margem direita (rightskip)
  \renewcommand*\l@chapter[2]{%
    \ifnum \c@tocdepth >\m@ne
        \addpenalty{-\@highpenalty}\addvspace{0.5em \@plus\p@}%
        \setlength\@tempdima{1.5em}%
        \begingroup
            \parindent \z@ \rightskip \@pnumwidth #1
            \parfillskip -\@pnumwidth
            \leavevmode \bfseries
            \advance\leftskip\@tempdima \hskip -\leftskip
            ##1\nobreak\hfil
            \nobreak\hb@xt@\@pnumwidth{\hss ##2\kern-\p@\kern\p@}\par
            \penalty\@highpenalty
        \endgroup
    \fi}}
\ifx\var\crp@vR \crp@lchap{}\fi
\ifx\var\crp@vB \crp@lchap{\@plus 4em}\fi
\ifx\var\crp@vC \crp@lchap{\@plus 4em \tolerance=10000 \hbadness=10000 }\fi
\makeatother
\begin{document}
\frontmatter
\tableofcontents
\mainmatter
\chapter{Deficiência psicossocial e}
\chapter{Deficiência psicossocial e transtorno}
\chapter{Deficiência psicossocial e transtorno do}
\chapter{Deficiência psicossocial e transtorno do espectro}
\chapter{Deficiência psicossocial e transtorno do espectro autista:}
\chapter{Deficiência psicossocial e transtorno do espectro autista: orientações}
\chapter{Deficiência psicossocial e transtorno do espectro autista: orientações para}
\chapter{Deficiência psicossocial e transtorno do espectro autista: orientações para a}
\chapter{Deficiência psicossocial e transtorno do espectro autista: orientações para a realização}
\chapter{Deficiência psicossocial e transtorno do espectro autista: orientações para a realização de}
\chapter{Deficiência psicossocial e transtorno do espectro autista: orientações para a realização de eventos}
\chapter{Deficiência psicossocial e transtorno do espectro autista: orientações para a realização de eventos e}
\chapter{Deficiência psicossocial e transtorno do espectro autista: orientações para a realização de eventos e apresentações}
\chapter{Deficiência psicossocial e transtorno do espectro autista: orientações para a realização de eventos e apresentações acessíveis}
\chapter{Deficiência psicossocial e transtorno do espectro autista: orientações para a realização de eventos e apresentações acessíveis em}
\chapter{Deficiência psicossocial e transtorno do espectro autista: orientações para a realização de eventos e apresentações acessíveis em todo}
\chapter{Comunicação alternativa e}
\chapter{Comunicação alternativa e aumentativa}
\chapter{Comunicação alternativa e aumentativa na}
\chapter{Comunicação alternativa e aumentativa na participação}
\chapter{Comunicação alternativa e aumentativa na participação social}
\chapter{Comunicação alternativa e aumentativa na participação social de}
\chapter{Comunicação alternativa e aumentativa na participação social de pessoas}
\chapter{Comunicação alternativa e aumentativa na participação social de pessoas com}
\chapter{Comunicação alternativa e aumentativa na participação social de pessoas com deficiência}
\chapter{Comunicação alternativa e aumentativa na participação social de pessoas com deficiência intelectual}
\chapter{Comunicação alternativa e aumentativa na participação social de pessoas com deficiência intelectual em}
\chapter{Comunicação alternativa e aumentativa na participação social de pessoas com deficiência intelectual em eventos}
\chapter{Comunicação alternativa e aumentativa na participação social de pessoas com deficiência intelectual em eventos públicos}
\chapter{Comunicação alternativa e aumentativa na participação social de pessoas com deficiência intelectual em eventos públicos de}
\chapter{Comunicação alternativa e aumentativa na participação social de pessoas com deficiência intelectual em eventos públicos de grande}
\chapter{Comunicação alternativa e aumentativa na participação social de pessoas com deficiência intelectual em eventos públicos de grande porte}
\chapter{Psicopedagogicamente inconstitucionalíssimo otorrinolaringologista}
\chapter{Psicopedagogicamente inconstitucionalíssimo otorrinolaringologista socioeducativo}
\chapter{Psicopedagogicamente inconstitucionalíssimo otorrinolaringologista socioeducativo ético-profissional}
\chapter{Psicopedagogicamente inconstitucionalíssimo otorrinolaringologista socioeducativo ético-profissional sociodemográfico}
\chapter{Psicopedagogicamente inconstitucionalíssimo otorrinolaringologista socioeducativo ético-profissional sociodemográfico interdisciplinaridade}
\chapter{Psicopedagogicamente inconstitucionalíssimo otorrinolaringologista socioeducativo ético-profissional sociodemográfico interdisciplinaridade neuropsicológica}
\chapter{Psicopedagogicamente inconstitucionalíssimo otorrinolaringologista socioeducativo ético-profissional sociodemográfico interdisciplinaridade neuropsicológica avaliação}
\chapter{Psicopedagogicamente inconstitucionalíssimo otorrinolaringologista socioeducativo ético-profissional sociodemográfico interdisciplinaridade neuropsicológica avaliação}
\chapter{Psicopedagogicamente inconstitucionalíssimo otorrinolaringologista socioeducativo ético-profissional sociodemográfico interdisciplinaridade neuropsicológica avaliação}
\chapter{Psicopedagogicamente inconstitucionalíssimo otorrinolaringologista socioeducativo ético-profissional sociodemográfico interdisciplinaridade neuropsicológica avaliação}
\chapter{Psicopedagogicamente inconstitucionalíssimo otorrinolaringologista socioeducativo ético-profissional sociodemográfico interdisciplinaridade neuropsicológica avaliação}
\chapter{Psicopedagogicamente inconstitucionalíssimo otorrinolaringologista socioeducativo ético-profissional sociodemográfico interdisciplinaridade neuropsicológica avaliação}
\chapter{Psicopedagogicamente inconstitucionalíssimo otorrinolaringologista socioeducativo ético-profissional sociodemográfico interdisciplinaridade neuropsicológica avaliação}
\chapter{Psicopedagogicamente inconstitucionalíssimo otorrinolaringologista socioeducativo ético-profissional sociodemográfico interdisciplinaridade neuropsicológica avaliação}
\chapter{Psicopedagogicamente inconstitucionalíssimo otorrinolaringologista socioeducativo ético-profissional sociodemográfico interdisciplinaridade neuropsicológica avaliação}
\chapter{Psicopedagogicamente inconstitucionalíssimo otorrinolaringologista socioeducativo ético-profissional sociodemográfico interdisciplinaridade neuropsicológica avaliação}
\end{document}
'   % B: plus 4em (o do guia)
%   lualatex '\def\var{C}% !TeX program = lualatex
% ===================================================================
% sumario-longo.tex — mede o sumário sem hifenização (issue #35, "Fora do
% escopo"): com a hifenização desligada, a entrada longa de capítulo pode
% bater no número da página. Gera 3 x 16 capítulos, cada um com um prefixo
% (de 3 a 18 palavras) de uma frase longa: frase 1 de palavras médias,
% frase 2 com números, frase 3 só de palavras muito longas. Conta no log
% os "Overfull \hbox" do sumário: o sumário é a única parte do documento
% sem parágrafo de corpo, então todo Overfull "in paragraph" antes do
% primeiro \chapter é entrada que passou da margem.
%
% Uso (de uma pasta com livros_crp_acessivel_book.sty e crpsp-leg.sty no
% caminho; as fontes NEWJUNE por OSFONTDIR):
%   lualatex '\def\var{R}\input{sumario-longo}'   % R: rígido (até 0.5.x)
%   lualatex '\def\var{B}\input{sumario-longo}'   % B: plus 4em (o do guia)
%   lualatex '\def\var{C}\input{sumario-longo}'   % C: plus 4em + \tolerance 10000
%   lualatex '\def\var{S}\input{sumario-longo}'   % S: o que a .sty define
% Cada variante só redefine o \l@chapter depois da carga da .sty.
% ===================================================================
\DocumentMetadata{lang=pt-BR, pdfstandard=ua-2, pdfversion=2.0, tagging=on}
\documentclass[11pt,a5paper,twoside,openright]{book}
\usepackage{livros_crp_acessivel_book}
\hypersetup{pdftitle={Sumario longo}, pdfauthor={CRP SP}}
\makeatletter
\providecommand\var{S}
\def\crp@vR{R}\def\crp@vB{B}\def\crp@vC{C}
\def\crp@lchap#1{%  #1 = acréscimo à margem direita (rightskip)
  \renewcommand*\l@chapter[2]{%
    \ifnum \c@tocdepth >\m@ne
        \addpenalty{-\@highpenalty}\addvspace{0.5em \@plus\p@}%
        \setlength\@tempdima{1.5em}%
        \begingroup
            \parindent \z@ \rightskip \@pnumwidth #1
            \parfillskip -\@pnumwidth
            \leavevmode \bfseries
            \advance\leftskip\@tempdima \hskip -\leftskip
            ##1\nobreak\hfil
            \nobreak\hb@xt@\@pnumwidth{\hss ##2\kern-\p@\kern\p@}\par
            \penalty\@highpenalty
        \endgroup
    \fi}}
\ifx\var\crp@vR \crp@lchap{}\fi
\ifx\var\crp@vB \crp@lchap{\@plus 4em}\fi
\ifx\var\crp@vC \crp@lchap{\@plus 4em \tolerance=10000 \hbadness=10000 }\fi
\makeatother
\begin{document}
\frontmatter
\tableofcontents
\mainmatter
\chapter{Deficiência psicossocial e}
\chapter{Deficiência psicossocial e transtorno}
\chapter{Deficiência psicossocial e transtorno do}
\chapter{Deficiência psicossocial e transtorno do espectro}
\chapter{Deficiência psicossocial e transtorno do espectro autista:}
\chapter{Deficiência psicossocial e transtorno do espectro autista: orientações}
\chapter{Deficiência psicossocial e transtorno do espectro autista: orientações para}
\chapter{Deficiência psicossocial e transtorno do espectro autista: orientações para a}
\chapter{Deficiência psicossocial e transtorno do espectro autista: orientações para a realização}
\chapter{Deficiência psicossocial e transtorno do espectro autista: orientações para a realização de}
\chapter{Deficiência psicossocial e transtorno do espectro autista: orientações para a realização de eventos}
\chapter{Deficiência psicossocial e transtorno do espectro autista: orientações para a realização de eventos e}
\chapter{Deficiência psicossocial e transtorno do espectro autista: orientações para a realização de eventos e apresentações}
\chapter{Deficiência psicossocial e transtorno do espectro autista: orientações para a realização de eventos e apresentações acessíveis}
\chapter{Deficiência psicossocial e transtorno do espectro autista: orientações para a realização de eventos e apresentações acessíveis em}
\chapter{Deficiência psicossocial e transtorno do espectro autista: orientações para a realização de eventos e apresentações acessíveis em todo}
\chapter{Comunicação alternativa e}
\chapter{Comunicação alternativa e aumentativa}
\chapter{Comunicação alternativa e aumentativa na}
\chapter{Comunicação alternativa e aumentativa na participação}
\chapter{Comunicação alternativa e aumentativa na participação social}
\chapter{Comunicação alternativa e aumentativa na participação social de}
\chapter{Comunicação alternativa e aumentativa na participação social de pessoas}
\chapter{Comunicação alternativa e aumentativa na participação social de pessoas com}
\chapter{Comunicação alternativa e aumentativa na participação social de pessoas com deficiência}
\chapter{Comunicação alternativa e aumentativa na participação social de pessoas com deficiência intelectual}
\chapter{Comunicação alternativa e aumentativa na participação social de pessoas com deficiência intelectual em}
\chapter{Comunicação alternativa e aumentativa na participação social de pessoas com deficiência intelectual em eventos}
\chapter{Comunicação alternativa e aumentativa na participação social de pessoas com deficiência intelectual em eventos públicos}
\chapter{Comunicação alternativa e aumentativa na participação social de pessoas com deficiência intelectual em eventos públicos de}
\chapter{Comunicação alternativa e aumentativa na participação social de pessoas com deficiência intelectual em eventos públicos de grande}
\chapter{Comunicação alternativa e aumentativa na participação social de pessoas com deficiência intelectual em eventos públicos de grande porte}
\chapter{Psicopedagogicamente inconstitucionalíssimo otorrinolaringologista}
\chapter{Psicopedagogicamente inconstitucionalíssimo otorrinolaringologista socioeducativo}
\chapter{Psicopedagogicamente inconstitucionalíssimo otorrinolaringologista socioeducativo ético-profissional}
\chapter{Psicopedagogicamente inconstitucionalíssimo otorrinolaringologista socioeducativo ético-profissional sociodemográfico}
\chapter{Psicopedagogicamente inconstitucionalíssimo otorrinolaringologista socioeducativo ético-profissional sociodemográfico interdisciplinaridade}
\chapter{Psicopedagogicamente inconstitucionalíssimo otorrinolaringologista socioeducativo ético-profissional sociodemográfico interdisciplinaridade neuropsicológica}
\chapter{Psicopedagogicamente inconstitucionalíssimo otorrinolaringologista socioeducativo ético-profissional sociodemográfico interdisciplinaridade neuropsicológica avaliação}
\chapter{Psicopedagogicamente inconstitucionalíssimo otorrinolaringologista socioeducativo ético-profissional sociodemográfico interdisciplinaridade neuropsicológica avaliação}
\chapter{Psicopedagogicamente inconstitucionalíssimo otorrinolaringologista socioeducativo ético-profissional sociodemográfico interdisciplinaridade neuropsicológica avaliação}
\chapter{Psicopedagogicamente inconstitucionalíssimo otorrinolaringologista socioeducativo ético-profissional sociodemográfico interdisciplinaridade neuropsicológica avaliação}
\chapter{Psicopedagogicamente inconstitucionalíssimo otorrinolaringologista socioeducativo ético-profissional sociodemográfico interdisciplinaridade neuropsicológica avaliação}
\chapter{Psicopedagogicamente inconstitucionalíssimo otorrinolaringologista socioeducativo ético-profissional sociodemográfico interdisciplinaridade neuropsicológica avaliação}
\chapter{Psicopedagogicamente inconstitucionalíssimo otorrinolaringologista socioeducativo ético-profissional sociodemográfico interdisciplinaridade neuropsicológica avaliação}
\chapter{Psicopedagogicamente inconstitucionalíssimo otorrinolaringologista socioeducativo ético-profissional sociodemográfico interdisciplinaridade neuropsicológica avaliação}
\chapter{Psicopedagogicamente inconstitucionalíssimo otorrinolaringologista socioeducativo ético-profissional sociodemográfico interdisciplinaridade neuropsicológica avaliação}
\chapter{Psicopedagogicamente inconstitucionalíssimo otorrinolaringologista socioeducativo ético-profissional sociodemográfico interdisciplinaridade neuropsicológica avaliação}
\end{document}
'   % C: plus 4em + \tolerance 10000
%   lualatex '\def\var{S}% !TeX program = lualatex
% ===================================================================
% sumario-longo.tex — mede o sumário sem hifenização (issue #35, "Fora do
% escopo"): com a hifenização desligada, a entrada longa de capítulo pode
% bater no número da página. Gera 3 x 16 capítulos, cada um com um prefixo
% (de 3 a 18 palavras) de uma frase longa: frase 1 de palavras médias,
% frase 2 com números, frase 3 só de palavras muito longas. Conta no log
% os "Overfull \hbox" do sumário: o sumário é a única parte do documento
% sem parágrafo de corpo, então todo Overfull "in paragraph" antes do
% primeiro \chapter é entrada que passou da margem.
%
% Uso (de uma pasta com livros_crp_acessivel_book.sty e crpsp-leg.sty no
% caminho; as fontes NEWJUNE por OSFONTDIR):
%   lualatex '\def\var{R}\input{sumario-longo}'   % R: rígido (até 0.5.x)
%   lualatex '\def\var{B}\input{sumario-longo}'   % B: plus 4em (o do guia)
%   lualatex '\def\var{C}\input{sumario-longo}'   % C: plus 4em + \tolerance 10000
%   lualatex '\def\var{S}\input{sumario-longo}'   % S: o que a .sty define
% Cada variante só redefine o \l@chapter depois da carga da .sty.
% ===================================================================
\DocumentMetadata{lang=pt-BR, pdfstandard=ua-2, pdfversion=2.0, tagging=on}
\documentclass[11pt,a5paper,twoside,openright]{book}
\usepackage{livros_crp_acessivel_book}
\hypersetup{pdftitle={Sumario longo}, pdfauthor={CRP SP}}
\makeatletter
\providecommand\var{S}
\def\crp@vR{R}\def\crp@vB{B}\def\crp@vC{C}
\def\crp@lchap#1{%  #1 = acréscimo à margem direita (rightskip)
  \renewcommand*\l@chapter[2]{%
    \ifnum \c@tocdepth >\m@ne
        \addpenalty{-\@highpenalty}\addvspace{0.5em \@plus\p@}%
        \setlength\@tempdima{1.5em}%
        \begingroup
            \parindent \z@ \rightskip \@pnumwidth #1
            \parfillskip -\@pnumwidth
            \leavevmode \bfseries
            \advance\leftskip\@tempdima \hskip -\leftskip
            ##1\nobreak\hfil
            \nobreak\hb@xt@\@pnumwidth{\hss ##2\kern-\p@\kern\p@}\par
            \penalty\@highpenalty
        \endgroup
    \fi}}
\ifx\var\crp@vR \crp@lchap{}\fi
\ifx\var\crp@vB \crp@lchap{\@plus 4em}\fi
\ifx\var\crp@vC \crp@lchap{\@plus 4em \tolerance=10000 \hbadness=10000 }\fi
\makeatother
\begin{document}
\frontmatter
\tableofcontents
\mainmatter
\chapter{Deficiência psicossocial e}
\chapter{Deficiência psicossocial e transtorno}
\chapter{Deficiência psicossocial e transtorno do}
\chapter{Deficiência psicossocial e transtorno do espectro}
\chapter{Deficiência psicossocial e transtorno do espectro autista:}
\chapter{Deficiência psicossocial e transtorno do espectro autista: orientações}
\chapter{Deficiência psicossocial e transtorno do espectro autista: orientações para}
\chapter{Deficiência psicossocial e transtorno do espectro autista: orientações para a}
\chapter{Deficiência psicossocial e transtorno do espectro autista: orientações para a realização}
\chapter{Deficiência psicossocial e transtorno do espectro autista: orientações para a realização de}
\chapter{Deficiência psicossocial e transtorno do espectro autista: orientações para a realização de eventos}
\chapter{Deficiência psicossocial e transtorno do espectro autista: orientações para a realização de eventos e}
\chapter{Deficiência psicossocial e transtorno do espectro autista: orientações para a realização de eventos e apresentações}
\chapter{Deficiência psicossocial e transtorno do espectro autista: orientações para a realização de eventos e apresentações acessíveis}
\chapter{Deficiência psicossocial e transtorno do espectro autista: orientações para a realização de eventos e apresentações acessíveis em}
\chapter{Deficiência psicossocial e transtorno do espectro autista: orientações para a realização de eventos e apresentações acessíveis em todo}
\chapter{Comunicação alternativa e}
\chapter{Comunicação alternativa e aumentativa}
\chapter{Comunicação alternativa e aumentativa na}
\chapter{Comunicação alternativa e aumentativa na participação}
\chapter{Comunicação alternativa e aumentativa na participação social}
\chapter{Comunicação alternativa e aumentativa na participação social de}
\chapter{Comunicação alternativa e aumentativa na participação social de pessoas}
\chapter{Comunicação alternativa e aumentativa na participação social de pessoas com}
\chapter{Comunicação alternativa e aumentativa na participação social de pessoas com deficiência}
\chapter{Comunicação alternativa e aumentativa na participação social de pessoas com deficiência intelectual}
\chapter{Comunicação alternativa e aumentativa na participação social de pessoas com deficiência intelectual em}
\chapter{Comunicação alternativa e aumentativa na participação social de pessoas com deficiência intelectual em eventos}
\chapter{Comunicação alternativa e aumentativa na participação social de pessoas com deficiência intelectual em eventos públicos}
\chapter{Comunicação alternativa e aumentativa na participação social de pessoas com deficiência intelectual em eventos públicos de}
\chapter{Comunicação alternativa e aumentativa na participação social de pessoas com deficiência intelectual em eventos públicos de grande}
\chapter{Comunicação alternativa e aumentativa na participação social de pessoas com deficiência intelectual em eventos públicos de grande porte}
\chapter{Psicopedagogicamente inconstitucionalíssimo otorrinolaringologista}
\chapter{Psicopedagogicamente inconstitucionalíssimo otorrinolaringologista socioeducativo}
\chapter{Psicopedagogicamente inconstitucionalíssimo otorrinolaringologista socioeducativo ético-profissional}
\chapter{Psicopedagogicamente inconstitucionalíssimo otorrinolaringologista socioeducativo ético-profissional sociodemográfico}
\chapter{Psicopedagogicamente inconstitucionalíssimo otorrinolaringologista socioeducativo ético-profissional sociodemográfico interdisciplinaridade}
\chapter{Psicopedagogicamente inconstitucionalíssimo otorrinolaringologista socioeducativo ético-profissional sociodemográfico interdisciplinaridade neuropsicológica}
\chapter{Psicopedagogicamente inconstitucionalíssimo otorrinolaringologista socioeducativo ético-profissional sociodemográfico interdisciplinaridade neuropsicológica avaliação}
\chapter{Psicopedagogicamente inconstitucionalíssimo otorrinolaringologista socioeducativo ético-profissional sociodemográfico interdisciplinaridade neuropsicológica avaliação}
\chapter{Psicopedagogicamente inconstitucionalíssimo otorrinolaringologista socioeducativo ético-profissional sociodemográfico interdisciplinaridade neuropsicológica avaliação}
\chapter{Psicopedagogicamente inconstitucionalíssimo otorrinolaringologista socioeducativo ético-profissional sociodemográfico interdisciplinaridade neuropsicológica avaliação}
\chapter{Psicopedagogicamente inconstitucionalíssimo otorrinolaringologista socioeducativo ético-profissional sociodemográfico interdisciplinaridade neuropsicológica avaliação}
\chapter{Psicopedagogicamente inconstitucionalíssimo otorrinolaringologista socioeducativo ético-profissional sociodemográfico interdisciplinaridade neuropsicológica avaliação}
\chapter{Psicopedagogicamente inconstitucionalíssimo otorrinolaringologista socioeducativo ético-profissional sociodemográfico interdisciplinaridade neuropsicológica avaliação}
\chapter{Psicopedagogicamente inconstitucionalíssimo otorrinolaringologista socioeducativo ético-profissional sociodemográfico interdisciplinaridade neuropsicológica avaliação}
\chapter{Psicopedagogicamente inconstitucionalíssimo otorrinolaringologista socioeducativo ético-profissional sociodemográfico interdisciplinaridade neuropsicológica avaliação}
\chapter{Psicopedagogicamente inconstitucionalíssimo otorrinolaringologista socioeducativo ético-profissional sociodemográfico interdisciplinaridade neuropsicológica avaliação}
\end{document}
'   % S: o que a .sty define
% Cada variante só redefine o \l@chapter depois da carga da .sty.
% ===================================================================
\DocumentMetadata{lang=pt-BR, pdfstandard=ua-2, pdfversion=2.0, tagging=on}
\documentclass[11pt,a5paper,twoside,openright]{book}
\usepackage{livros_crp_acessivel_book}
\hypersetup{pdftitle={Sumario longo}, pdfauthor={CRP SP}}
\makeatletter
\providecommand\var{S}
\def\crp@vR{R}\def\crp@vB{B}\def\crp@vC{C}
\def\crp@lchap#1{%  #1 = acréscimo à margem direita (rightskip)
  \renewcommand*\l@chapter[2]{%
    \ifnum \c@tocdepth >\m@ne
        \addpenalty{-\@highpenalty}\addvspace{0.5em \@plus\p@}%
        \setlength\@tempdima{1.5em}%
        \begingroup
            \parindent \z@ \rightskip \@pnumwidth #1
            \parfillskip -\@pnumwidth
            \leavevmode \bfseries
            \advance\leftskip\@tempdima \hskip -\leftskip
            ##1\nobreak\hfil
            \nobreak\hb@xt@\@pnumwidth{\hss ##2\kern-\p@\kern\p@}\par
            \penalty\@highpenalty
        \endgroup
    \fi}}
\ifx\var\crp@vR \crp@lchap{}\fi
\ifx\var\crp@vB \crp@lchap{\@plus 4em}\fi
\ifx\var\crp@vC \crp@lchap{\@plus 4em \tolerance=10000 \hbadness=10000 }\fi
\makeatother
\begin{document}
\frontmatter
\tableofcontents
\mainmatter
\chapter{Deficiência psicossocial e}
\chapter{Deficiência psicossocial e transtorno}
\chapter{Deficiência psicossocial e transtorno do}
\chapter{Deficiência psicossocial e transtorno do espectro}
\chapter{Deficiência psicossocial e transtorno do espectro autista:}
\chapter{Deficiência psicossocial e transtorno do espectro autista: orientações}
\chapter{Deficiência psicossocial e transtorno do espectro autista: orientações para}
\chapter{Deficiência psicossocial e transtorno do espectro autista: orientações para a}
\chapter{Deficiência psicossocial e transtorno do espectro autista: orientações para a realização}
\chapter{Deficiência psicossocial e transtorno do espectro autista: orientações para a realização de}
\chapter{Deficiência psicossocial e transtorno do espectro autista: orientações para a realização de eventos}
\chapter{Deficiência psicossocial e transtorno do espectro autista: orientações para a realização de eventos e}
\chapter{Deficiência psicossocial e transtorno do espectro autista: orientações para a realização de eventos e apresentações}
\chapter{Deficiência psicossocial e transtorno do espectro autista: orientações para a realização de eventos e apresentações acessíveis}
\chapter{Deficiência psicossocial e transtorno do espectro autista: orientações para a realização de eventos e apresentações acessíveis em}
\chapter{Deficiência psicossocial e transtorno do espectro autista: orientações para a realização de eventos e apresentações acessíveis em todo}
\chapter{Comunicação alternativa e}
\chapter{Comunicação alternativa e aumentativa}
\chapter{Comunicação alternativa e aumentativa na}
\chapter{Comunicação alternativa e aumentativa na participação}
\chapter{Comunicação alternativa e aumentativa na participação social}
\chapter{Comunicação alternativa e aumentativa na participação social de}
\chapter{Comunicação alternativa e aumentativa na participação social de pessoas}
\chapter{Comunicação alternativa e aumentativa na participação social de pessoas com}
\chapter{Comunicação alternativa e aumentativa na participação social de pessoas com deficiência}
\chapter{Comunicação alternativa e aumentativa na participação social de pessoas com deficiência intelectual}
\chapter{Comunicação alternativa e aumentativa na participação social de pessoas com deficiência intelectual em}
\chapter{Comunicação alternativa e aumentativa na participação social de pessoas com deficiência intelectual em eventos}
\chapter{Comunicação alternativa e aumentativa na participação social de pessoas com deficiência intelectual em eventos públicos}
\chapter{Comunicação alternativa e aumentativa na participação social de pessoas com deficiência intelectual em eventos públicos de}
\chapter{Comunicação alternativa e aumentativa na participação social de pessoas com deficiência intelectual em eventos públicos de grande}
\chapter{Comunicação alternativa e aumentativa na participação social de pessoas com deficiência intelectual em eventos públicos de grande porte}
\chapter{Psicopedagogicamente inconstitucionalíssimo otorrinolaringologista}
\chapter{Psicopedagogicamente inconstitucionalíssimo otorrinolaringologista socioeducativo}
\chapter{Psicopedagogicamente inconstitucionalíssimo otorrinolaringologista socioeducativo ético-profissional}
\chapter{Psicopedagogicamente inconstitucionalíssimo otorrinolaringologista socioeducativo ético-profissional sociodemográfico}
\chapter{Psicopedagogicamente inconstitucionalíssimo otorrinolaringologista socioeducativo ético-profissional sociodemográfico interdisciplinaridade}
\chapter{Psicopedagogicamente inconstitucionalíssimo otorrinolaringologista socioeducativo ético-profissional sociodemográfico interdisciplinaridade neuropsicológica}
\chapter{Psicopedagogicamente inconstitucionalíssimo otorrinolaringologista socioeducativo ético-profissional sociodemográfico interdisciplinaridade neuropsicológica avaliação}
\chapter{Psicopedagogicamente inconstitucionalíssimo otorrinolaringologista socioeducativo ético-profissional sociodemográfico interdisciplinaridade neuropsicológica avaliação}
\chapter{Psicopedagogicamente inconstitucionalíssimo otorrinolaringologista socioeducativo ético-profissional sociodemográfico interdisciplinaridade neuropsicológica avaliação}
\chapter{Psicopedagogicamente inconstitucionalíssimo otorrinolaringologista socioeducativo ético-profissional sociodemográfico interdisciplinaridade neuropsicológica avaliação}
\chapter{Psicopedagogicamente inconstitucionalíssimo otorrinolaringologista socioeducativo ético-profissional sociodemográfico interdisciplinaridade neuropsicológica avaliação}
\chapter{Psicopedagogicamente inconstitucionalíssimo otorrinolaringologista socioeducativo ético-profissional sociodemográfico interdisciplinaridade neuropsicológica avaliação}
\chapter{Psicopedagogicamente inconstitucionalíssimo otorrinolaringologista socioeducativo ético-profissional sociodemográfico interdisciplinaridade neuropsicológica avaliação}
\chapter{Psicopedagogicamente inconstitucionalíssimo otorrinolaringologista socioeducativo ético-profissional sociodemográfico interdisciplinaridade neuropsicológica avaliação}
\chapter{Psicopedagogicamente inconstitucionalíssimo otorrinolaringologista socioeducativo ético-profissional sociodemográfico interdisciplinaridade neuropsicológica avaliação}
\chapter{Psicopedagogicamente inconstitucionalíssimo otorrinolaringologista socioeducativo ético-profissional sociodemográfico interdisciplinaridade neuropsicológica avaliação}
\end{document}
'   % C: plus 4em + \tolerance 10000
%   lualatex '\def\var{S}% !TeX program = lualatex
% ===================================================================
% sumario-longo.tex — mede o sumário sem hifenização (issue #35, "Fora do
% escopo"): com a hifenização desligada, a entrada longa de capítulo pode
% bater no número da página. Gera 3 x 16 capítulos, cada um com um prefixo
% (de 3 a 18 palavras) de uma frase longa: frase 1 de palavras médias,
% frase 2 com números, frase 3 só de palavras muito longas. Conta no log
% os "Overfull \hbox" do sumário: o sumário é a única parte do documento
% sem parágrafo de corpo, então todo Overfull "in paragraph" antes do
% primeiro \chapter é entrada que passou da margem.
%
% Uso (de uma pasta com livros_crp_acessivel_book.sty e crpsp-leg.sty no
% caminho; as fontes NEWJUNE por OSFONTDIR):
%   lualatex '\def\var{R}% !TeX program = lualatex
% ===================================================================
% sumario-longo.tex — mede o sumário sem hifenização (issue #35, "Fora do
% escopo"): com a hifenização desligada, a entrada longa de capítulo pode
% bater no número da página. Gera 3 x 16 capítulos, cada um com um prefixo
% (de 3 a 18 palavras) de uma frase longa: frase 1 de palavras médias,
% frase 2 com números, frase 3 só de palavras muito longas. Conta no log
% os "Overfull \hbox" do sumário: o sumário é a única parte do documento
% sem parágrafo de corpo, então todo Overfull "in paragraph" antes do
% primeiro \chapter é entrada que passou da margem.
%
% Uso (de uma pasta com livros_crp_acessivel_book.sty e crpsp-leg.sty no
% caminho; as fontes NEWJUNE por OSFONTDIR):
%   lualatex '\def\var{R}\input{sumario-longo}'   % R: rígido (até 0.5.x)
%   lualatex '\def\var{B}\input{sumario-longo}'   % B: plus 4em (o do guia)
%   lualatex '\def\var{C}\input{sumario-longo}'   % C: plus 4em + \tolerance 10000
%   lualatex '\def\var{S}\input{sumario-longo}'   % S: o que a .sty define
% Cada variante só redefine o \l@chapter depois da carga da .sty.
% ===================================================================
\DocumentMetadata{lang=pt-BR, pdfstandard=ua-2, pdfversion=2.0, tagging=on}
\documentclass[11pt,a5paper,twoside,openright]{book}
\usepackage{livros_crp_acessivel_book}
\hypersetup{pdftitle={Sumario longo}, pdfauthor={CRP SP}}
\makeatletter
\providecommand\var{S}
\def\crp@vR{R}\def\crp@vB{B}\def\crp@vC{C}
\def\crp@lchap#1{%  #1 = acréscimo à margem direita (rightskip)
  \renewcommand*\l@chapter[2]{%
    \ifnum \c@tocdepth >\m@ne
        \addpenalty{-\@highpenalty}\addvspace{0.5em \@plus\p@}%
        \setlength\@tempdima{1.5em}%
        \begingroup
            \parindent \z@ \rightskip \@pnumwidth #1
            \parfillskip -\@pnumwidth
            \leavevmode \bfseries
            \advance\leftskip\@tempdima \hskip -\leftskip
            ##1\nobreak\hfil
            \nobreak\hb@xt@\@pnumwidth{\hss ##2\kern-\p@\kern\p@}\par
            \penalty\@highpenalty
        \endgroup
    \fi}}
\ifx\var\crp@vR \crp@lchap{}\fi
\ifx\var\crp@vB \crp@lchap{\@plus 4em}\fi
\ifx\var\crp@vC \crp@lchap{\@plus 4em \tolerance=10000 \hbadness=10000 }\fi
\makeatother
\begin{document}
\frontmatter
\tableofcontents
\mainmatter
\chapter{Deficiência psicossocial e}
\chapter{Deficiência psicossocial e transtorno}
\chapter{Deficiência psicossocial e transtorno do}
\chapter{Deficiência psicossocial e transtorno do espectro}
\chapter{Deficiência psicossocial e transtorno do espectro autista:}
\chapter{Deficiência psicossocial e transtorno do espectro autista: orientações}
\chapter{Deficiência psicossocial e transtorno do espectro autista: orientações para}
\chapter{Deficiência psicossocial e transtorno do espectro autista: orientações para a}
\chapter{Deficiência psicossocial e transtorno do espectro autista: orientações para a realização}
\chapter{Deficiência psicossocial e transtorno do espectro autista: orientações para a realização de}
\chapter{Deficiência psicossocial e transtorno do espectro autista: orientações para a realização de eventos}
\chapter{Deficiência psicossocial e transtorno do espectro autista: orientações para a realização de eventos e}
\chapter{Deficiência psicossocial e transtorno do espectro autista: orientações para a realização de eventos e apresentações}
\chapter{Deficiência psicossocial e transtorno do espectro autista: orientações para a realização de eventos e apresentações acessíveis}
\chapter{Deficiência psicossocial e transtorno do espectro autista: orientações para a realização de eventos e apresentações acessíveis em}
\chapter{Deficiência psicossocial e transtorno do espectro autista: orientações para a realização de eventos e apresentações acessíveis em todo}
\chapter{Comunicação alternativa e}
\chapter{Comunicação alternativa e aumentativa}
\chapter{Comunicação alternativa e aumentativa na}
\chapter{Comunicação alternativa e aumentativa na participação}
\chapter{Comunicação alternativa e aumentativa na participação social}
\chapter{Comunicação alternativa e aumentativa na participação social de}
\chapter{Comunicação alternativa e aumentativa na participação social de pessoas}
\chapter{Comunicação alternativa e aumentativa na participação social de pessoas com}
\chapter{Comunicação alternativa e aumentativa na participação social de pessoas com deficiência}
\chapter{Comunicação alternativa e aumentativa na participação social de pessoas com deficiência intelectual}
\chapter{Comunicação alternativa e aumentativa na participação social de pessoas com deficiência intelectual em}
\chapter{Comunicação alternativa e aumentativa na participação social de pessoas com deficiência intelectual em eventos}
\chapter{Comunicação alternativa e aumentativa na participação social de pessoas com deficiência intelectual em eventos públicos}
\chapter{Comunicação alternativa e aumentativa na participação social de pessoas com deficiência intelectual em eventos públicos de}
\chapter{Comunicação alternativa e aumentativa na participação social de pessoas com deficiência intelectual em eventos públicos de grande}
\chapter{Comunicação alternativa e aumentativa na participação social de pessoas com deficiência intelectual em eventos públicos de grande porte}
\chapter{Psicopedagogicamente inconstitucionalíssimo otorrinolaringologista}
\chapter{Psicopedagogicamente inconstitucionalíssimo otorrinolaringologista socioeducativo}
\chapter{Psicopedagogicamente inconstitucionalíssimo otorrinolaringologista socioeducativo ético-profissional}
\chapter{Psicopedagogicamente inconstitucionalíssimo otorrinolaringologista socioeducativo ético-profissional sociodemográfico}
\chapter{Psicopedagogicamente inconstitucionalíssimo otorrinolaringologista socioeducativo ético-profissional sociodemográfico interdisciplinaridade}
\chapter{Psicopedagogicamente inconstitucionalíssimo otorrinolaringologista socioeducativo ético-profissional sociodemográfico interdisciplinaridade neuropsicológica}
\chapter{Psicopedagogicamente inconstitucionalíssimo otorrinolaringologista socioeducativo ético-profissional sociodemográfico interdisciplinaridade neuropsicológica avaliação}
\chapter{Psicopedagogicamente inconstitucionalíssimo otorrinolaringologista socioeducativo ético-profissional sociodemográfico interdisciplinaridade neuropsicológica avaliação}
\chapter{Psicopedagogicamente inconstitucionalíssimo otorrinolaringologista socioeducativo ético-profissional sociodemográfico interdisciplinaridade neuropsicológica avaliação}
\chapter{Psicopedagogicamente inconstitucionalíssimo otorrinolaringologista socioeducativo ético-profissional sociodemográfico interdisciplinaridade neuropsicológica avaliação}
\chapter{Psicopedagogicamente inconstitucionalíssimo otorrinolaringologista socioeducativo ético-profissional sociodemográfico interdisciplinaridade neuropsicológica avaliação}
\chapter{Psicopedagogicamente inconstitucionalíssimo otorrinolaringologista socioeducativo ético-profissional sociodemográfico interdisciplinaridade neuropsicológica avaliação}
\chapter{Psicopedagogicamente inconstitucionalíssimo otorrinolaringologista socioeducativo ético-profissional sociodemográfico interdisciplinaridade neuropsicológica avaliação}
\chapter{Psicopedagogicamente inconstitucionalíssimo otorrinolaringologista socioeducativo ético-profissional sociodemográfico interdisciplinaridade neuropsicológica avaliação}
\chapter{Psicopedagogicamente inconstitucionalíssimo otorrinolaringologista socioeducativo ético-profissional sociodemográfico interdisciplinaridade neuropsicológica avaliação}
\chapter{Psicopedagogicamente inconstitucionalíssimo otorrinolaringologista socioeducativo ético-profissional sociodemográfico interdisciplinaridade neuropsicológica avaliação}
\end{document}
'   % R: rígido (até 0.5.x)
%   lualatex '\def\var{B}% !TeX program = lualatex
% ===================================================================
% sumario-longo.tex — mede o sumário sem hifenização (issue #35, "Fora do
% escopo"): com a hifenização desligada, a entrada longa de capítulo pode
% bater no número da página. Gera 3 x 16 capítulos, cada um com um prefixo
% (de 3 a 18 palavras) de uma frase longa: frase 1 de palavras médias,
% frase 2 com números, frase 3 só de palavras muito longas. Conta no log
% os "Overfull \hbox" do sumário: o sumário é a única parte do documento
% sem parágrafo de corpo, então todo Overfull "in paragraph" antes do
% primeiro \chapter é entrada que passou da margem.
%
% Uso (de uma pasta com livros_crp_acessivel_book.sty e crpsp-leg.sty no
% caminho; as fontes NEWJUNE por OSFONTDIR):
%   lualatex '\def\var{R}\input{sumario-longo}'   % R: rígido (até 0.5.x)
%   lualatex '\def\var{B}\input{sumario-longo}'   % B: plus 4em (o do guia)
%   lualatex '\def\var{C}\input{sumario-longo}'   % C: plus 4em + \tolerance 10000
%   lualatex '\def\var{S}\input{sumario-longo}'   % S: o que a .sty define
% Cada variante só redefine o \l@chapter depois da carga da .sty.
% ===================================================================
\DocumentMetadata{lang=pt-BR, pdfstandard=ua-2, pdfversion=2.0, tagging=on}
\documentclass[11pt,a5paper,twoside,openright]{book}
\usepackage{livros_crp_acessivel_book}
\hypersetup{pdftitle={Sumario longo}, pdfauthor={CRP SP}}
\makeatletter
\providecommand\var{S}
\def\crp@vR{R}\def\crp@vB{B}\def\crp@vC{C}
\def\crp@lchap#1{%  #1 = acréscimo à margem direita (rightskip)
  \renewcommand*\l@chapter[2]{%
    \ifnum \c@tocdepth >\m@ne
        \addpenalty{-\@highpenalty}\addvspace{0.5em \@plus\p@}%
        \setlength\@tempdima{1.5em}%
        \begingroup
            \parindent \z@ \rightskip \@pnumwidth #1
            \parfillskip -\@pnumwidth
            \leavevmode \bfseries
            \advance\leftskip\@tempdima \hskip -\leftskip
            ##1\nobreak\hfil
            \nobreak\hb@xt@\@pnumwidth{\hss ##2\kern-\p@\kern\p@}\par
            \penalty\@highpenalty
        \endgroup
    \fi}}
\ifx\var\crp@vR \crp@lchap{}\fi
\ifx\var\crp@vB \crp@lchap{\@plus 4em}\fi
\ifx\var\crp@vC \crp@lchap{\@plus 4em \tolerance=10000 \hbadness=10000 }\fi
\makeatother
\begin{document}
\frontmatter
\tableofcontents
\mainmatter
\chapter{Deficiência psicossocial e}
\chapter{Deficiência psicossocial e transtorno}
\chapter{Deficiência psicossocial e transtorno do}
\chapter{Deficiência psicossocial e transtorno do espectro}
\chapter{Deficiência psicossocial e transtorno do espectro autista:}
\chapter{Deficiência psicossocial e transtorno do espectro autista: orientações}
\chapter{Deficiência psicossocial e transtorno do espectro autista: orientações para}
\chapter{Deficiência psicossocial e transtorno do espectro autista: orientações para a}
\chapter{Deficiência psicossocial e transtorno do espectro autista: orientações para a realização}
\chapter{Deficiência psicossocial e transtorno do espectro autista: orientações para a realização de}
\chapter{Deficiência psicossocial e transtorno do espectro autista: orientações para a realização de eventos}
\chapter{Deficiência psicossocial e transtorno do espectro autista: orientações para a realização de eventos e}
\chapter{Deficiência psicossocial e transtorno do espectro autista: orientações para a realização de eventos e apresentações}
\chapter{Deficiência psicossocial e transtorno do espectro autista: orientações para a realização de eventos e apresentações acessíveis}
\chapter{Deficiência psicossocial e transtorno do espectro autista: orientações para a realização de eventos e apresentações acessíveis em}
\chapter{Deficiência psicossocial e transtorno do espectro autista: orientações para a realização de eventos e apresentações acessíveis em todo}
\chapter{Comunicação alternativa e}
\chapter{Comunicação alternativa e aumentativa}
\chapter{Comunicação alternativa e aumentativa na}
\chapter{Comunicação alternativa e aumentativa na participação}
\chapter{Comunicação alternativa e aumentativa na participação social}
\chapter{Comunicação alternativa e aumentativa na participação social de}
\chapter{Comunicação alternativa e aumentativa na participação social de pessoas}
\chapter{Comunicação alternativa e aumentativa na participação social de pessoas com}
\chapter{Comunicação alternativa e aumentativa na participação social de pessoas com deficiência}
\chapter{Comunicação alternativa e aumentativa na participação social de pessoas com deficiência intelectual}
\chapter{Comunicação alternativa e aumentativa na participação social de pessoas com deficiência intelectual em}
\chapter{Comunicação alternativa e aumentativa na participação social de pessoas com deficiência intelectual em eventos}
\chapter{Comunicação alternativa e aumentativa na participação social de pessoas com deficiência intelectual em eventos públicos}
\chapter{Comunicação alternativa e aumentativa na participação social de pessoas com deficiência intelectual em eventos públicos de}
\chapter{Comunicação alternativa e aumentativa na participação social de pessoas com deficiência intelectual em eventos públicos de grande}
\chapter{Comunicação alternativa e aumentativa na participação social de pessoas com deficiência intelectual em eventos públicos de grande porte}
\chapter{Psicopedagogicamente inconstitucionalíssimo otorrinolaringologista}
\chapter{Psicopedagogicamente inconstitucionalíssimo otorrinolaringologista socioeducativo}
\chapter{Psicopedagogicamente inconstitucionalíssimo otorrinolaringologista socioeducativo ético-profissional}
\chapter{Psicopedagogicamente inconstitucionalíssimo otorrinolaringologista socioeducativo ético-profissional sociodemográfico}
\chapter{Psicopedagogicamente inconstitucionalíssimo otorrinolaringologista socioeducativo ético-profissional sociodemográfico interdisciplinaridade}
\chapter{Psicopedagogicamente inconstitucionalíssimo otorrinolaringologista socioeducativo ético-profissional sociodemográfico interdisciplinaridade neuropsicológica}
\chapter{Psicopedagogicamente inconstitucionalíssimo otorrinolaringologista socioeducativo ético-profissional sociodemográfico interdisciplinaridade neuropsicológica avaliação}
\chapter{Psicopedagogicamente inconstitucionalíssimo otorrinolaringologista socioeducativo ético-profissional sociodemográfico interdisciplinaridade neuropsicológica avaliação}
\chapter{Psicopedagogicamente inconstitucionalíssimo otorrinolaringologista socioeducativo ético-profissional sociodemográfico interdisciplinaridade neuropsicológica avaliação}
\chapter{Psicopedagogicamente inconstitucionalíssimo otorrinolaringologista socioeducativo ético-profissional sociodemográfico interdisciplinaridade neuropsicológica avaliação}
\chapter{Psicopedagogicamente inconstitucionalíssimo otorrinolaringologista socioeducativo ético-profissional sociodemográfico interdisciplinaridade neuropsicológica avaliação}
\chapter{Psicopedagogicamente inconstitucionalíssimo otorrinolaringologista socioeducativo ético-profissional sociodemográfico interdisciplinaridade neuropsicológica avaliação}
\chapter{Psicopedagogicamente inconstitucionalíssimo otorrinolaringologista socioeducativo ético-profissional sociodemográfico interdisciplinaridade neuropsicológica avaliação}
\chapter{Psicopedagogicamente inconstitucionalíssimo otorrinolaringologista socioeducativo ético-profissional sociodemográfico interdisciplinaridade neuropsicológica avaliação}
\chapter{Psicopedagogicamente inconstitucionalíssimo otorrinolaringologista socioeducativo ético-profissional sociodemográfico interdisciplinaridade neuropsicológica avaliação}
\chapter{Psicopedagogicamente inconstitucionalíssimo otorrinolaringologista socioeducativo ético-profissional sociodemográfico interdisciplinaridade neuropsicológica avaliação}
\end{document}
'   % B: plus 4em (o do guia)
%   lualatex '\def\var{C}% !TeX program = lualatex
% ===================================================================
% sumario-longo.tex — mede o sumário sem hifenização (issue #35, "Fora do
% escopo"): com a hifenização desligada, a entrada longa de capítulo pode
% bater no número da página. Gera 3 x 16 capítulos, cada um com um prefixo
% (de 3 a 18 palavras) de uma frase longa: frase 1 de palavras médias,
% frase 2 com números, frase 3 só de palavras muito longas. Conta no log
% os "Overfull \hbox" do sumário: o sumário é a única parte do documento
% sem parágrafo de corpo, então todo Overfull "in paragraph" antes do
% primeiro \chapter é entrada que passou da margem.
%
% Uso (de uma pasta com livros_crp_acessivel_book.sty e crpsp-leg.sty no
% caminho; as fontes NEWJUNE por OSFONTDIR):
%   lualatex '\def\var{R}\input{sumario-longo}'   % R: rígido (até 0.5.x)
%   lualatex '\def\var{B}\input{sumario-longo}'   % B: plus 4em (o do guia)
%   lualatex '\def\var{C}\input{sumario-longo}'   % C: plus 4em + \tolerance 10000
%   lualatex '\def\var{S}\input{sumario-longo}'   % S: o que a .sty define
% Cada variante só redefine o \l@chapter depois da carga da .sty.
% ===================================================================
\DocumentMetadata{lang=pt-BR, pdfstandard=ua-2, pdfversion=2.0, tagging=on}
\documentclass[11pt,a5paper,twoside,openright]{book}
\usepackage{livros_crp_acessivel_book}
\hypersetup{pdftitle={Sumario longo}, pdfauthor={CRP SP}}
\makeatletter
\providecommand\var{S}
\def\crp@vR{R}\def\crp@vB{B}\def\crp@vC{C}
\def\crp@lchap#1{%  #1 = acréscimo à margem direita (rightskip)
  \renewcommand*\l@chapter[2]{%
    \ifnum \c@tocdepth >\m@ne
        \addpenalty{-\@highpenalty}\addvspace{0.5em \@plus\p@}%
        \setlength\@tempdima{1.5em}%
        \begingroup
            \parindent \z@ \rightskip \@pnumwidth #1
            \parfillskip -\@pnumwidth
            \leavevmode \bfseries
            \advance\leftskip\@tempdima \hskip -\leftskip
            ##1\nobreak\hfil
            \nobreak\hb@xt@\@pnumwidth{\hss ##2\kern-\p@\kern\p@}\par
            \penalty\@highpenalty
        \endgroup
    \fi}}
\ifx\var\crp@vR \crp@lchap{}\fi
\ifx\var\crp@vB \crp@lchap{\@plus 4em}\fi
\ifx\var\crp@vC \crp@lchap{\@plus 4em \tolerance=10000 \hbadness=10000 }\fi
\makeatother
\begin{document}
\frontmatter
\tableofcontents
\mainmatter
\chapter{Deficiência psicossocial e}
\chapter{Deficiência psicossocial e transtorno}
\chapter{Deficiência psicossocial e transtorno do}
\chapter{Deficiência psicossocial e transtorno do espectro}
\chapter{Deficiência psicossocial e transtorno do espectro autista:}
\chapter{Deficiência psicossocial e transtorno do espectro autista: orientações}
\chapter{Deficiência psicossocial e transtorno do espectro autista: orientações para}
\chapter{Deficiência psicossocial e transtorno do espectro autista: orientações para a}
\chapter{Deficiência psicossocial e transtorno do espectro autista: orientações para a realização}
\chapter{Deficiência psicossocial e transtorno do espectro autista: orientações para a realização de}
\chapter{Deficiência psicossocial e transtorno do espectro autista: orientações para a realização de eventos}
\chapter{Deficiência psicossocial e transtorno do espectro autista: orientações para a realização de eventos e}
\chapter{Deficiência psicossocial e transtorno do espectro autista: orientações para a realização de eventos e apresentações}
\chapter{Deficiência psicossocial e transtorno do espectro autista: orientações para a realização de eventos e apresentações acessíveis}
\chapter{Deficiência psicossocial e transtorno do espectro autista: orientações para a realização de eventos e apresentações acessíveis em}
\chapter{Deficiência psicossocial e transtorno do espectro autista: orientações para a realização de eventos e apresentações acessíveis em todo}
\chapter{Comunicação alternativa e}
\chapter{Comunicação alternativa e aumentativa}
\chapter{Comunicação alternativa e aumentativa na}
\chapter{Comunicação alternativa e aumentativa na participação}
\chapter{Comunicação alternativa e aumentativa na participação social}
\chapter{Comunicação alternativa e aumentativa na participação social de}
\chapter{Comunicação alternativa e aumentativa na participação social de pessoas}
\chapter{Comunicação alternativa e aumentativa na participação social de pessoas com}
\chapter{Comunicação alternativa e aumentativa na participação social de pessoas com deficiência}
\chapter{Comunicação alternativa e aumentativa na participação social de pessoas com deficiência intelectual}
\chapter{Comunicação alternativa e aumentativa na participação social de pessoas com deficiência intelectual em}
\chapter{Comunicação alternativa e aumentativa na participação social de pessoas com deficiência intelectual em eventos}
\chapter{Comunicação alternativa e aumentativa na participação social de pessoas com deficiência intelectual em eventos públicos}
\chapter{Comunicação alternativa e aumentativa na participação social de pessoas com deficiência intelectual em eventos públicos de}
\chapter{Comunicação alternativa e aumentativa na participação social de pessoas com deficiência intelectual em eventos públicos de grande}
\chapter{Comunicação alternativa e aumentativa na participação social de pessoas com deficiência intelectual em eventos públicos de grande porte}
\chapter{Psicopedagogicamente inconstitucionalíssimo otorrinolaringologista}
\chapter{Psicopedagogicamente inconstitucionalíssimo otorrinolaringologista socioeducativo}
\chapter{Psicopedagogicamente inconstitucionalíssimo otorrinolaringologista socioeducativo ético-profissional}
\chapter{Psicopedagogicamente inconstitucionalíssimo otorrinolaringologista socioeducativo ético-profissional sociodemográfico}
\chapter{Psicopedagogicamente inconstitucionalíssimo otorrinolaringologista socioeducativo ético-profissional sociodemográfico interdisciplinaridade}
\chapter{Psicopedagogicamente inconstitucionalíssimo otorrinolaringologista socioeducativo ético-profissional sociodemográfico interdisciplinaridade neuropsicológica}
\chapter{Psicopedagogicamente inconstitucionalíssimo otorrinolaringologista socioeducativo ético-profissional sociodemográfico interdisciplinaridade neuropsicológica avaliação}
\chapter{Psicopedagogicamente inconstitucionalíssimo otorrinolaringologista socioeducativo ético-profissional sociodemográfico interdisciplinaridade neuropsicológica avaliação}
\chapter{Psicopedagogicamente inconstitucionalíssimo otorrinolaringologista socioeducativo ético-profissional sociodemográfico interdisciplinaridade neuropsicológica avaliação}
\chapter{Psicopedagogicamente inconstitucionalíssimo otorrinolaringologista socioeducativo ético-profissional sociodemográfico interdisciplinaridade neuropsicológica avaliação}
\chapter{Psicopedagogicamente inconstitucionalíssimo otorrinolaringologista socioeducativo ético-profissional sociodemográfico interdisciplinaridade neuropsicológica avaliação}
\chapter{Psicopedagogicamente inconstitucionalíssimo otorrinolaringologista socioeducativo ético-profissional sociodemográfico interdisciplinaridade neuropsicológica avaliação}
\chapter{Psicopedagogicamente inconstitucionalíssimo otorrinolaringologista socioeducativo ético-profissional sociodemográfico interdisciplinaridade neuropsicológica avaliação}
\chapter{Psicopedagogicamente inconstitucionalíssimo otorrinolaringologista socioeducativo ético-profissional sociodemográfico interdisciplinaridade neuropsicológica avaliação}
\chapter{Psicopedagogicamente inconstitucionalíssimo otorrinolaringologista socioeducativo ético-profissional sociodemográfico interdisciplinaridade neuropsicológica avaliação}
\chapter{Psicopedagogicamente inconstitucionalíssimo otorrinolaringologista socioeducativo ético-profissional sociodemográfico interdisciplinaridade neuropsicológica avaliação}
\end{document}
'   % C: plus 4em + \tolerance 10000
%   lualatex '\def\var{S}% !TeX program = lualatex
% ===================================================================
% sumario-longo.tex — mede o sumário sem hifenização (issue #35, "Fora do
% escopo"): com a hifenização desligada, a entrada longa de capítulo pode
% bater no número da página. Gera 3 x 16 capítulos, cada um com um prefixo
% (de 3 a 18 palavras) de uma frase longa: frase 1 de palavras médias,
% frase 2 com números, frase 3 só de palavras muito longas. Conta no log
% os "Overfull \hbox" do sumário: o sumário é a única parte do documento
% sem parágrafo de corpo, então todo Overfull "in paragraph" antes do
% primeiro \chapter é entrada que passou da margem.
%
% Uso (de uma pasta com livros_crp_acessivel_book.sty e crpsp-leg.sty no
% caminho; as fontes NEWJUNE por OSFONTDIR):
%   lualatex '\def\var{R}\input{sumario-longo}'   % R: rígido (até 0.5.x)
%   lualatex '\def\var{B}\input{sumario-longo}'   % B: plus 4em (o do guia)
%   lualatex '\def\var{C}\input{sumario-longo}'   % C: plus 4em + \tolerance 10000
%   lualatex '\def\var{S}\input{sumario-longo}'   % S: o que a .sty define
% Cada variante só redefine o \l@chapter depois da carga da .sty.
% ===================================================================
\DocumentMetadata{lang=pt-BR, pdfstandard=ua-2, pdfversion=2.0, tagging=on}
\documentclass[11pt,a5paper,twoside,openright]{book}
\usepackage{livros_crp_acessivel_book}
\hypersetup{pdftitle={Sumario longo}, pdfauthor={CRP SP}}
\makeatletter
\providecommand\var{S}
\def\crp@vR{R}\def\crp@vB{B}\def\crp@vC{C}
\def\crp@lchap#1{%  #1 = acréscimo à margem direita (rightskip)
  \renewcommand*\l@chapter[2]{%
    \ifnum \c@tocdepth >\m@ne
        \addpenalty{-\@highpenalty}\addvspace{0.5em \@plus\p@}%
        \setlength\@tempdima{1.5em}%
        \begingroup
            \parindent \z@ \rightskip \@pnumwidth #1
            \parfillskip -\@pnumwidth
            \leavevmode \bfseries
            \advance\leftskip\@tempdima \hskip -\leftskip
            ##1\nobreak\hfil
            \nobreak\hb@xt@\@pnumwidth{\hss ##2\kern-\p@\kern\p@}\par
            \penalty\@highpenalty
        \endgroup
    \fi}}
\ifx\var\crp@vR \crp@lchap{}\fi
\ifx\var\crp@vB \crp@lchap{\@plus 4em}\fi
\ifx\var\crp@vC \crp@lchap{\@plus 4em \tolerance=10000 \hbadness=10000 }\fi
\makeatother
\begin{document}
\frontmatter
\tableofcontents
\mainmatter
\chapter{Deficiência psicossocial e}
\chapter{Deficiência psicossocial e transtorno}
\chapter{Deficiência psicossocial e transtorno do}
\chapter{Deficiência psicossocial e transtorno do espectro}
\chapter{Deficiência psicossocial e transtorno do espectro autista:}
\chapter{Deficiência psicossocial e transtorno do espectro autista: orientações}
\chapter{Deficiência psicossocial e transtorno do espectro autista: orientações para}
\chapter{Deficiência psicossocial e transtorno do espectro autista: orientações para a}
\chapter{Deficiência psicossocial e transtorno do espectro autista: orientações para a realização}
\chapter{Deficiência psicossocial e transtorno do espectro autista: orientações para a realização de}
\chapter{Deficiência psicossocial e transtorno do espectro autista: orientações para a realização de eventos}
\chapter{Deficiência psicossocial e transtorno do espectro autista: orientações para a realização de eventos e}
\chapter{Deficiência psicossocial e transtorno do espectro autista: orientações para a realização de eventos e apresentações}
\chapter{Deficiência psicossocial e transtorno do espectro autista: orientações para a realização de eventos e apresentações acessíveis}
\chapter{Deficiência psicossocial e transtorno do espectro autista: orientações para a realização de eventos e apresentações acessíveis em}
\chapter{Deficiência psicossocial e transtorno do espectro autista: orientações para a realização de eventos e apresentações acessíveis em todo}
\chapter{Comunicação alternativa e}
\chapter{Comunicação alternativa e aumentativa}
\chapter{Comunicação alternativa e aumentativa na}
\chapter{Comunicação alternativa e aumentativa na participação}
\chapter{Comunicação alternativa e aumentativa na participação social}
\chapter{Comunicação alternativa e aumentativa na participação social de}
\chapter{Comunicação alternativa e aumentativa na participação social de pessoas}
\chapter{Comunicação alternativa e aumentativa na participação social de pessoas com}
\chapter{Comunicação alternativa e aumentativa na participação social de pessoas com deficiência}
\chapter{Comunicação alternativa e aumentativa na participação social de pessoas com deficiência intelectual}
\chapter{Comunicação alternativa e aumentativa na participação social de pessoas com deficiência intelectual em}
\chapter{Comunicação alternativa e aumentativa na participação social de pessoas com deficiência intelectual em eventos}
\chapter{Comunicação alternativa e aumentativa na participação social de pessoas com deficiência intelectual em eventos públicos}
\chapter{Comunicação alternativa e aumentativa na participação social de pessoas com deficiência intelectual em eventos públicos de}
\chapter{Comunicação alternativa e aumentativa na participação social de pessoas com deficiência intelectual em eventos públicos de grande}
\chapter{Comunicação alternativa e aumentativa na participação social de pessoas com deficiência intelectual em eventos públicos de grande porte}
\chapter{Psicopedagogicamente inconstitucionalíssimo otorrinolaringologista}
\chapter{Psicopedagogicamente inconstitucionalíssimo otorrinolaringologista socioeducativo}
\chapter{Psicopedagogicamente inconstitucionalíssimo otorrinolaringologista socioeducativo ético-profissional}
\chapter{Psicopedagogicamente inconstitucionalíssimo otorrinolaringologista socioeducativo ético-profissional sociodemográfico}
\chapter{Psicopedagogicamente inconstitucionalíssimo otorrinolaringologista socioeducativo ético-profissional sociodemográfico interdisciplinaridade}
\chapter{Psicopedagogicamente inconstitucionalíssimo otorrinolaringologista socioeducativo ético-profissional sociodemográfico interdisciplinaridade neuropsicológica}
\chapter{Psicopedagogicamente inconstitucionalíssimo otorrinolaringologista socioeducativo ético-profissional sociodemográfico interdisciplinaridade neuropsicológica avaliação}
\chapter{Psicopedagogicamente inconstitucionalíssimo otorrinolaringologista socioeducativo ético-profissional sociodemográfico interdisciplinaridade neuropsicológica avaliação}
\chapter{Psicopedagogicamente inconstitucionalíssimo otorrinolaringologista socioeducativo ético-profissional sociodemográfico interdisciplinaridade neuropsicológica avaliação}
\chapter{Psicopedagogicamente inconstitucionalíssimo otorrinolaringologista socioeducativo ético-profissional sociodemográfico interdisciplinaridade neuropsicológica avaliação}
\chapter{Psicopedagogicamente inconstitucionalíssimo otorrinolaringologista socioeducativo ético-profissional sociodemográfico interdisciplinaridade neuropsicológica avaliação}
\chapter{Psicopedagogicamente inconstitucionalíssimo otorrinolaringologista socioeducativo ético-profissional sociodemográfico interdisciplinaridade neuropsicológica avaliação}
\chapter{Psicopedagogicamente inconstitucionalíssimo otorrinolaringologista socioeducativo ético-profissional sociodemográfico interdisciplinaridade neuropsicológica avaliação}
\chapter{Psicopedagogicamente inconstitucionalíssimo otorrinolaringologista socioeducativo ético-profissional sociodemográfico interdisciplinaridade neuropsicológica avaliação}
\chapter{Psicopedagogicamente inconstitucionalíssimo otorrinolaringologista socioeducativo ético-profissional sociodemográfico interdisciplinaridade neuropsicológica avaliação}
\chapter{Psicopedagogicamente inconstitucionalíssimo otorrinolaringologista socioeducativo ético-profissional sociodemográfico interdisciplinaridade neuropsicológica avaliação}
\end{document}
'   % S: o que a .sty define
% Cada variante só redefine o \l@chapter depois da carga da .sty.
% ===================================================================
\DocumentMetadata{lang=pt-BR, pdfstandard=ua-2, pdfversion=2.0, tagging=on}
\documentclass[11pt,a5paper,twoside,openright]{book}
\usepackage{livros_crp_acessivel_book}
\hypersetup{pdftitle={Sumario longo}, pdfauthor={CRP SP}}
\makeatletter
\providecommand\var{S}
\def\crp@vR{R}\def\crp@vB{B}\def\crp@vC{C}
\def\crp@lchap#1{%  #1 = acréscimo à margem direita (rightskip)
  \renewcommand*\l@chapter[2]{%
    \ifnum \c@tocdepth >\m@ne
        \addpenalty{-\@highpenalty}\addvspace{0.5em \@plus\p@}%
        \setlength\@tempdima{1.5em}%
        \begingroup
            \parindent \z@ \rightskip \@pnumwidth #1
            \parfillskip -\@pnumwidth
            \leavevmode \bfseries
            \advance\leftskip\@tempdima \hskip -\leftskip
            ##1\nobreak\hfil
            \nobreak\hb@xt@\@pnumwidth{\hss ##2\kern-\p@\kern\p@}\par
            \penalty\@highpenalty
        \endgroup
    \fi}}
\ifx\var\crp@vR \crp@lchap{}\fi
\ifx\var\crp@vB \crp@lchap{\@plus 4em}\fi
\ifx\var\crp@vC \crp@lchap{\@plus 4em \tolerance=10000 \hbadness=10000 }\fi
\makeatother
\begin{document}
\frontmatter
\tableofcontents
\mainmatter
\chapter{Deficiência psicossocial e}
\chapter{Deficiência psicossocial e transtorno}
\chapter{Deficiência psicossocial e transtorno do}
\chapter{Deficiência psicossocial e transtorno do espectro}
\chapter{Deficiência psicossocial e transtorno do espectro autista:}
\chapter{Deficiência psicossocial e transtorno do espectro autista: orientações}
\chapter{Deficiência psicossocial e transtorno do espectro autista: orientações para}
\chapter{Deficiência psicossocial e transtorno do espectro autista: orientações para a}
\chapter{Deficiência psicossocial e transtorno do espectro autista: orientações para a realização}
\chapter{Deficiência psicossocial e transtorno do espectro autista: orientações para a realização de}
\chapter{Deficiência psicossocial e transtorno do espectro autista: orientações para a realização de eventos}
\chapter{Deficiência psicossocial e transtorno do espectro autista: orientações para a realização de eventos e}
\chapter{Deficiência psicossocial e transtorno do espectro autista: orientações para a realização de eventos e apresentações}
\chapter{Deficiência psicossocial e transtorno do espectro autista: orientações para a realização de eventos e apresentações acessíveis}
\chapter{Deficiência psicossocial e transtorno do espectro autista: orientações para a realização de eventos e apresentações acessíveis em}
\chapter{Deficiência psicossocial e transtorno do espectro autista: orientações para a realização de eventos e apresentações acessíveis em todo}
\chapter{Comunicação alternativa e}
\chapter{Comunicação alternativa e aumentativa}
\chapter{Comunicação alternativa e aumentativa na}
\chapter{Comunicação alternativa e aumentativa na participação}
\chapter{Comunicação alternativa e aumentativa na participação social}
\chapter{Comunicação alternativa e aumentativa na participação social de}
\chapter{Comunicação alternativa e aumentativa na participação social de pessoas}
\chapter{Comunicação alternativa e aumentativa na participação social de pessoas com}
\chapter{Comunicação alternativa e aumentativa na participação social de pessoas com deficiência}
\chapter{Comunicação alternativa e aumentativa na participação social de pessoas com deficiência intelectual}
\chapter{Comunicação alternativa e aumentativa na participação social de pessoas com deficiência intelectual em}
\chapter{Comunicação alternativa e aumentativa na participação social de pessoas com deficiência intelectual em eventos}
\chapter{Comunicação alternativa e aumentativa na participação social de pessoas com deficiência intelectual em eventos públicos}
\chapter{Comunicação alternativa e aumentativa na participação social de pessoas com deficiência intelectual em eventos públicos de}
\chapter{Comunicação alternativa e aumentativa na participação social de pessoas com deficiência intelectual em eventos públicos de grande}
\chapter{Comunicação alternativa e aumentativa na participação social de pessoas com deficiência intelectual em eventos públicos de grande porte}
\chapter{Psicopedagogicamente inconstitucionalíssimo otorrinolaringologista}
\chapter{Psicopedagogicamente inconstitucionalíssimo otorrinolaringologista socioeducativo}
\chapter{Psicopedagogicamente inconstitucionalíssimo otorrinolaringologista socioeducativo ético-profissional}
\chapter{Psicopedagogicamente inconstitucionalíssimo otorrinolaringologista socioeducativo ético-profissional sociodemográfico}
\chapter{Psicopedagogicamente inconstitucionalíssimo otorrinolaringologista socioeducativo ético-profissional sociodemográfico interdisciplinaridade}
\chapter{Psicopedagogicamente inconstitucionalíssimo otorrinolaringologista socioeducativo ético-profissional sociodemográfico interdisciplinaridade neuropsicológica}
\chapter{Psicopedagogicamente inconstitucionalíssimo otorrinolaringologista socioeducativo ético-profissional sociodemográfico interdisciplinaridade neuropsicológica avaliação}
\chapter{Psicopedagogicamente inconstitucionalíssimo otorrinolaringologista socioeducativo ético-profissional sociodemográfico interdisciplinaridade neuropsicológica avaliação}
\chapter{Psicopedagogicamente inconstitucionalíssimo otorrinolaringologista socioeducativo ético-profissional sociodemográfico interdisciplinaridade neuropsicológica avaliação}
\chapter{Psicopedagogicamente inconstitucionalíssimo otorrinolaringologista socioeducativo ético-profissional sociodemográfico interdisciplinaridade neuropsicológica avaliação}
\chapter{Psicopedagogicamente inconstitucionalíssimo otorrinolaringologista socioeducativo ético-profissional sociodemográfico interdisciplinaridade neuropsicológica avaliação}
\chapter{Psicopedagogicamente inconstitucionalíssimo otorrinolaringologista socioeducativo ético-profissional sociodemográfico interdisciplinaridade neuropsicológica avaliação}
\chapter{Psicopedagogicamente inconstitucionalíssimo otorrinolaringologista socioeducativo ético-profissional sociodemográfico interdisciplinaridade neuropsicológica avaliação}
\chapter{Psicopedagogicamente inconstitucionalíssimo otorrinolaringologista socioeducativo ético-profissional sociodemográfico interdisciplinaridade neuropsicológica avaliação}
\chapter{Psicopedagogicamente inconstitucionalíssimo otorrinolaringologista socioeducativo ético-profissional sociodemográfico interdisciplinaridade neuropsicológica avaliação}
\chapter{Psicopedagogicamente inconstitucionalíssimo otorrinolaringologista socioeducativo ético-profissional sociodemográfico interdisciplinaridade neuropsicológica avaliação}
\end{document}
'   % S: o que a .sty define
% Cada variante só redefine o \l@chapter depois da carga da .sty.
% ===================================================================
\DocumentMetadata{lang=pt-BR, pdfstandard=ua-2, pdfversion=2.0, tagging=on}
\documentclass[11pt,a5paper,twoside,openright]{book}
\usepackage{livros_crp_acessivel_book}
\hypersetup{pdftitle={Sumario longo}, pdfauthor={CRP SP}}
\makeatletter
\providecommand\var{S}
\def\crp@vR{R}\def\crp@vB{B}\def\crp@vC{C}
\def\crp@lchap#1{%  #1 = acréscimo à margem direita (rightskip)
  \renewcommand*\l@chapter[2]{%
    \ifnum \c@tocdepth >\m@ne
        \addpenalty{-\@highpenalty}\addvspace{0.5em \@plus\p@}%
        \setlength\@tempdima{1.5em}%
        \begingroup
            \parindent \z@ \rightskip \@pnumwidth #1
            \parfillskip -\@pnumwidth
            \leavevmode \bfseries
            \advance\leftskip\@tempdima \hskip -\leftskip
            ##1\nobreak\hfil
            \nobreak\hb@xt@\@pnumwidth{\hss ##2\kern-\p@\kern\p@}\par
            \penalty\@highpenalty
        \endgroup
    \fi}}
\ifx\var\crp@vR \crp@lchap{}\fi
\ifx\var\crp@vB \crp@lchap{\@plus 4em}\fi
\ifx\var\crp@vC \crp@lchap{\@plus 4em \tolerance=10000 \hbadness=10000 }\fi
\makeatother
\begin{document}
\frontmatter
\tableofcontents
\mainmatter
\chapter{Deficiência psicossocial e}
\chapter{Deficiência psicossocial e transtorno}
\chapter{Deficiência psicossocial e transtorno do}
\chapter{Deficiência psicossocial e transtorno do espectro}
\chapter{Deficiência psicossocial e transtorno do espectro autista:}
\chapter{Deficiência psicossocial e transtorno do espectro autista: orientações}
\chapter{Deficiência psicossocial e transtorno do espectro autista: orientações para}
\chapter{Deficiência psicossocial e transtorno do espectro autista: orientações para a}
\chapter{Deficiência psicossocial e transtorno do espectro autista: orientações para a realização}
\chapter{Deficiência psicossocial e transtorno do espectro autista: orientações para a realização de}
\chapter{Deficiência psicossocial e transtorno do espectro autista: orientações para a realização de eventos}
\chapter{Deficiência psicossocial e transtorno do espectro autista: orientações para a realização de eventos e}
\chapter{Deficiência psicossocial e transtorno do espectro autista: orientações para a realização de eventos e apresentações}
\chapter{Deficiência psicossocial e transtorno do espectro autista: orientações para a realização de eventos e apresentações acessíveis}
\chapter{Deficiência psicossocial e transtorno do espectro autista: orientações para a realização de eventos e apresentações acessíveis em}
\chapter{Deficiência psicossocial e transtorno do espectro autista: orientações para a realização de eventos e apresentações acessíveis em todo}
\chapter{Comunicação alternativa e}
\chapter{Comunicação alternativa e aumentativa}
\chapter{Comunicação alternativa e aumentativa na}
\chapter{Comunicação alternativa e aumentativa na participação}
\chapter{Comunicação alternativa e aumentativa na participação social}
\chapter{Comunicação alternativa e aumentativa na participação social de}
\chapter{Comunicação alternativa e aumentativa na participação social de pessoas}
\chapter{Comunicação alternativa e aumentativa na participação social de pessoas com}
\chapter{Comunicação alternativa e aumentativa na participação social de pessoas com deficiência}
\chapter{Comunicação alternativa e aumentativa na participação social de pessoas com deficiência intelectual}
\chapter{Comunicação alternativa e aumentativa na participação social de pessoas com deficiência intelectual em}
\chapter{Comunicação alternativa e aumentativa na participação social de pessoas com deficiência intelectual em eventos}
\chapter{Comunicação alternativa e aumentativa na participação social de pessoas com deficiência intelectual em eventos públicos}
\chapter{Comunicação alternativa e aumentativa na participação social de pessoas com deficiência intelectual em eventos públicos de}
\chapter{Comunicação alternativa e aumentativa na participação social de pessoas com deficiência intelectual em eventos públicos de grande}
\chapter{Comunicação alternativa e aumentativa na participação social de pessoas com deficiência intelectual em eventos públicos de grande porte}
\chapter{Psicopedagogicamente inconstitucionalíssimo otorrinolaringologista}
\chapter{Psicopedagogicamente inconstitucionalíssimo otorrinolaringologista socioeducativo}
\chapter{Psicopedagogicamente inconstitucionalíssimo otorrinolaringologista socioeducativo ético-profissional}
\chapter{Psicopedagogicamente inconstitucionalíssimo otorrinolaringologista socioeducativo ético-profissional sociodemográfico}
\chapter{Psicopedagogicamente inconstitucionalíssimo otorrinolaringologista socioeducativo ético-profissional sociodemográfico interdisciplinaridade}
\chapter{Psicopedagogicamente inconstitucionalíssimo otorrinolaringologista socioeducativo ético-profissional sociodemográfico interdisciplinaridade neuropsicológica}
\chapter{Psicopedagogicamente inconstitucionalíssimo otorrinolaringologista socioeducativo ético-profissional sociodemográfico interdisciplinaridade neuropsicológica avaliação}
\chapter{Psicopedagogicamente inconstitucionalíssimo otorrinolaringologista socioeducativo ético-profissional sociodemográfico interdisciplinaridade neuropsicológica avaliação}
\chapter{Psicopedagogicamente inconstitucionalíssimo otorrinolaringologista socioeducativo ético-profissional sociodemográfico interdisciplinaridade neuropsicológica avaliação}
\chapter{Psicopedagogicamente inconstitucionalíssimo otorrinolaringologista socioeducativo ético-profissional sociodemográfico interdisciplinaridade neuropsicológica avaliação}
\chapter{Psicopedagogicamente inconstitucionalíssimo otorrinolaringologista socioeducativo ético-profissional sociodemográfico interdisciplinaridade neuropsicológica avaliação}
\chapter{Psicopedagogicamente inconstitucionalíssimo otorrinolaringologista socioeducativo ético-profissional sociodemográfico interdisciplinaridade neuropsicológica avaliação}
\chapter{Psicopedagogicamente inconstitucionalíssimo otorrinolaringologista socioeducativo ético-profissional sociodemográfico interdisciplinaridade neuropsicológica avaliação}
\chapter{Psicopedagogicamente inconstitucionalíssimo otorrinolaringologista socioeducativo ético-profissional sociodemográfico interdisciplinaridade neuropsicológica avaliação}
\chapter{Psicopedagogicamente inconstitucionalíssimo otorrinolaringologista socioeducativo ético-profissional sociodemográfico interdisciplinaridade neuropsicológica avaliação}
\chapter{Psicopedagogicamente inconstitucionalíssimo otorrinolaringologista socioeducativo ético-profissional sociodemográfico interdisciplinaridade neuropsicológica avaliação}
\end{document}
'   % C: plus 4em + \tolerance 10000
%   lualatex '\def\var{S}% !TeX program = lualatex
% ===================================================================
% sumario-longo.tex — mede o sumário sem hifenização (issue #35, "Fora do
% escopo"): com a hifenização desligada, a entrada longa de capítulo pode
% bater no número da página. Gera 3 x 16 capítulos, cada um com um prefixo
% (de 3 a 18 palavras) de uma frase longa: frase 1 de palavras médias,
% frase 2 com números, frase 3 só de palavras muito longas. Conta no log
% os "Overfull \hbox" do sumário: o sumário é a única parte do documento
% sem parágrafo de corpo, então todo Overfull "in paragraph" antes do
% primeiro \chapter é entrada que passou da margem.
%
% Uso (de uma pasta com livros_crp_acessivel_book.sty e crpsp-leg.sty no
% caminho; as fontes NEWJUNE por OSFONTDIR):
%   lualatex '\def\var{R}% !TeX program = lualatex
% ===================================================================
% sumario-longo.tex — mede o sumário sem hifenização (issue #35, "Fora do
% escopo"): com a hifenização desligada, a entrada longa de capítulo pode
% bater no número da página. Gera 3 x 16 capítulos, cada um com um prefixo
% (de 3 a 18 palavras) de uma frase longa: frase 1 de palavras médias,
% frase 2 com números, frase 3 só de palavras muito longas. Conta no log
% os "Overfull \hbox" do sumário: o sumário é a única parte do documento
% sem parágrafo de corpo, então todo Overfull "in paragraph" antes do
% primeiro \chapter é entrada que passou da margem.
%
% Uso (de uma pasta com livros_crp_acessivel_book.sty e crpsp-leg.sty no
% caminho; as fontes NEWJUNE por OSFONTDIR):
%   lualatex '\def\var{R}% !TeX program = lualatex
% ===================================================================
% sumario-longo.tex — mede o sumário sem hifenização (issue #35, "Fora do
% escopo"): com a hifenização desligada, a entrada longa de capítulo pode
% bater no número da página. Gera 3 x 16 capítulos, cada um com um prefixo
% (de 3 a 18 palavras) de uma frase longa: frase 1 de palavras médias,
% frase 2 com números, frase 3 só de palavras muito longas. Conta no log
% os "Overfull \hbox" do sumário: o sumário é a única parte do documento
% sem parágrafo de corpo, então todo Overfull "in paragraph" antes do
% primeiro \chapter é entrada que passou da margem.
%
% Uso (de uma pasta com livros_crp_acessivel_book.sty e crpsp-leg.sty no
% caminho; as fontes NEWJUNE por OSFONTDIR):
%   lualatex '\def\var{R}\input{sumario-longo}'   % R: rígido (até 0.5.x)
%   lualatex '\def\var{B}\input{sumario-longo}'   % B: plus 4em (o do guia)
%   lualatex '\def\var{C}\input{sumario-longo}'   % C: plus 4em + \tolerance 10000
%   lualatex '\def\var{S}\input{sumario-longo}'   % S: o que a .sty define
% Cada variante só redefine o \l@chapter depois da carga da .sty.
% ===================================================================
\DocumentMetadata{lang=pt-BR, pdfstandard=ua-2, pdfversion=2.0, tagging=on}
\documentclass[11pt,a5paper,twoside,openright]{book}
\usepackage{livros_crp_acessivel_book}
\hypersetup{pdftitle={Sumario longo}, pdfauthor={CRP SP}}
\makeatletter
\providecommand\var{S}
\def\crp@vR{R}\def\crp@vB{B}\def\crp@vC{C}
\def\crp@lchap#1{%  #1 = acréscimo à margem direita (rightskip)
  \renewcommand*\l@chapter[2]{%
    \ifnum \c@tocdepth >\m@ne
        \addpenalty{-\@highpenalty}\addvspace{0.5em \@plus\p@}%
        \setlength\@tempdima{1.5em}%
        \begingroup
            \parindent \z@ \rightskip \@pnumwidth #1
            \parfillskip -\@pnumwidth
            \leavevmode \bfseries
            \advance\leftskip\@tempdima \hskip -\leftskip
            ##1\nobreak\hfil
            \nobreak\hb@xt@\@pnumwidth{\hss ##2\kern-\p@\kern\p@}\par
            \penalty\@highpenalty
        \endgroup
    \fi}}
\ifx\var\crp@vR \crp@lchap{}\fi
\ifx\var\crp@vB \crp@lchap{\@plus 4em}\fi
\ifx\var\crp@vC \crp@lchap{\@plus 4em \tolerance=10000 \hbadness=10000 }\fi
\makeatother
\begin{document}
\frontmatter
\tableofcontents
\mainmatter
\chapter{Deficiência psicossocial e}
\chapter{Deficiência psicossocial e transtorno}
\chapter{Deficiência psicossocial e transtorno do}
\chapter{Deficiência psicossocial e transtorno do espectro}
\chapter{Deficiência psicossocial e transtorno do espectro autista:}
\chapter{Deficiência psicossocial e transtorno do espectro autista: orientações}
\chapter{Deficiência psicossocial e transtorno do espectro autista: orientações para}
\chapter{Deficiência psicossocial e transtorno do espectro autista: orientações para a}
\chapter{Deficiência psicossocial e transtorno do espectro autista: orientações para a realização}
\chapter{Deficiência psicossocial e transtorno do espectro autista: orientações para a realização de}
\chapter{Deficiência psicossocial e transtorno do espectro autista: orientações para a realização de eventos}
\chapter{Deficiência psicossocial e transtorno do espectro autista: orientações para a realização de eventos e}
\chapter{Deficiência psicossocial e transtorno do espectro autista: orientações para a realização de eventos e apresentações}
\chapter{Deficiência psicossocial e transtorno do espectro autista: orientações para a realização de eventos e apresentações acessíveis}
\chapter{Deficiência psicossocial e transtorno do espectro autista: orientações para a realização de eventos e apresentações acessíveis em}
\chapter{Deficiência psicossocial e transtorno do espectro autista: orientações para a realização de eventos e apresentações acessíveis em todo}
\chapter{Comunicação alternativa e}
\chapter{Comunicação alternativa e aumentativa}
\chapter{Comunicação alternativa e aumentativa na}
\chapter{Comunicação alternativa e aumentativa na participação}
\chapter{Comunicação alternativa e aumentativa na participação social}
\chapter{Comunicação alternativa e aumentativa na participação social de}
\chapter{Comunicação alternativa e aumentativa na participação social de pessoas}
\chapter{Comunicação alternativa e aumentativa na participação social de pessoas com}
\chapter{Comunicação alternativa e aumentativa na participação social de pessoas com deficiência}
\chapter{Comunicação alternativa e aumentativa na participação social de pessoas com deficiência intelectual}
\chapter{Comunicação alternativa e aumentativa na participação social de pessoas com deficiência intelectual em}
\chapter{Comunicação alternativa e aumentativa na participação social de pessoas com deficiência intelectual em eventos}
\chapter{Comunicação alternativa e aumentativa na participação social de pessoas com deficiência intelectual em eventos públicos}
\chapter{Comunicação alternativa e aumentativa na participação social de pessoas com deficiência intelectual em eventos públicos de}
\chapter{Comunicação alternativa e aumentativa na participação social de pessoas com deficiência intelectual em eventos públicos de grande}
\chapter{Comunicação alternativa e aumentativa na participação social de pessoas com deficiência intelectual em eventos públicos de grande porte}
\chapter{Psicopedagogicamente inconstitucionalíssimo otorrinolaringologista}
\chapter{Psicopedagogicamente inconstitucionalíssimo otorrinolaringologista socioeducativo}
\chapter{Psicopedagogicamente inconstitucionalíssimo otorrinolaringologista socioeducativo ético-profissional}
\chapter{Psicopedagogicamente inconstitucionalíssimo otorrinolaringologista socioeducativo ético-profissional sociodemográfico}
\chapter{Psicopedagogicamente inconstitucionalíssimo otorrinolaringologista socioeducativo ético-profissional sociodemográfico interdisciplinaridade}
\chapter{Psicopedagogicamente inconstitucionalíssimo otorrinolaringologista socioeducativo ético-profissional sociodemográfico interdisciplinaridade neuropsicológica}
\chapter{Psicopedagogicamente inconstitucionalíssimo otorrinolaringologista socioeducativo ético-profissional sociodemográfico interdisciplinaridade neuropsicológica avaliação}
\chapter{Psicopedagogicamente inconstitucionalíssimo otorrinolaringologista socioeducativo ético-profissional sociodemográfico interdisciplinaridade neuropsicológica avaliação}
\chapter{Psicopedagogicamente inconstitucionalíssimo otorrinolaringologista socioeducativo ético-profissional sociodemográfico interdisciplinaridade neuropsicológica avaliação}
\chapter{Psicopedagogicamente inconstitucionalíssimo otorrinolaringologista socioeducativo ético-profissional sociodemográfico interdisciplinaridade neuropsicológica avaliação}
\chapter{Psicopedagogicamente inconstitucionalíssimo otorrinolaringologista socioeducativo ético-profissional sociodemográfico interdisciplinaridade neuropsicológica avaliação}
\chapter{Psicopedagogicamente inconstitucionalíssimo otorrinolaringologista socioeducativo ético-profissional sociodemográfico interdisciplinaridade neuropsicológica avaliação}
\chapter{Psicopedagogicamente inconstitucionalíssimo otorrinolaringologista socioeducativo ético-profissional sociodemográfico interdisciplinaridade neuropsicológica avaliação}
\chapter{Psicopedagogicamente inconstitucionalíssimo otorrinolaringologista socioeducativo ético-profissional sociodemográfico interdisciplinaridade neuropsicológica avaliação}
\chapter{Psicopedagogicamente inconstitucionalíssimo otorrinolaringologista socioeducativo ético-profissional sociodemográfico interdisciplinaridade neuropsicológica avaliação}
\chapter{Psicopedagogicamente inconstitucionalíssimo otorrinolaringologista socioeducativo ético-profissional sociodemográfico interdisciplinaridade neuropsicológica avaliação}
\end{document}
'   % R: rígido (até 0.5.x)
%   lualatex '\def\var{B}% !TeX program = lualatex
% ===================================================================
% sumario-longo.tex — mede o sumário sem hifenização (issue #35, "Fora do
% escopo"): com a hifenização desligada, a entrada longa de capítulo pode
% bater no número da página. Gera 3 x 16 capítulos, cada um com um prefixo
% (de 3 a 18 palavras) de uma frase longa: frase 1 de palavras médias,
% frase 2 com números, frase 3 só de palavras muito longas. Conta no log
% os "Overfull \hbox" do sumário: o sumário é a única parte do documento
% sem parágrafo de corpo, então todo Overfull "in paragraph" antes do
% primeiro \chapter é entrada que passou da margem.
%
% Uso (de uma pasta com livros_crp_acessivel_book.sty e crpsp-leg.sty no
% caminho; as fontes NEWJUNE por OSFONTDIR):
%   lualatex '\def\var{R}\input{sumario-longo}'   % R: rígido (até 0.5.x)
%   lualatex '\def\var{B}\input{sumario-longo}'   % B: plus 4em (o do guia)
%   lualatex '\def\var{C}\input{sumario-longo}'   % C: plus 4em + \tolerance 10000
%   lualatex '\def\var{S}\input{sumario-longo}'   % S: o que a .sty define
% Cada variante só redefine o \l@chapter depois da carga da .sty.
% ===================================================================
\DocumentMetadata{lang=pt-BR, pdfstandard=ua-2, pdfversion=2.0, tagging=on}
\documentclass[11pt,a5paper,twoside,openright]{book}
\usepackage{livros_crp_acessivel_book}
\hypersetup{pdftitle={Sumario longo}, pdfauthor={CRP SP}}
\makeatletter
\providecommand\var{S}
\def\crp@vR{R}\def\crp@vB{B}\def\crp@vC{C}
\def\crp@lchap#1{%  #1 = acréscimo à margem direita (rightskip)
  \renewcommand*\l@chapter[2]{%
    \ifnum \c@tocdepth >\m@ne
        \addpenalty{-\@highpenalty}\addvspace{0.5em \@plus\p@}%
        \setlength\@tempdima{1.5em}%
        \begingroup
            \parindent \z@ \rightskip \@pnumwidth #1
            \parfillskip -\@pnumwidth
            \leavevmode \bfseries
            \advance\leftskip\@tempdima \hskip -\leftskip
            ##1\nobreak\hfil
            \nobreak\hb@xt@\@pnumwidth{\hss ##2\kern-\p@\kern\p@}\par
            \penalty\@highpenalty
        \endgroup
    \fi}}
\ifx\var\crp@vR \crp@lchap{}\fi
\ifx\var\crp@vB \crp@lchap{\@plus 4em}\fi
\ifx\var\crp@vC \crp@lchap{\@plus 4em \tolerance=10000 \hbadness=10000 }\fi
\makeatother
\begin{document}
\frontmatter
\tableofcontents
\mainmatter
\chapter{Deficiência psicossocial e}
\chapter{Deficiência psicossocial e transtorno}
\chapter{Deficiência psicossocial e transtorno do}
\chapter{Deficiência psicossocial e transtorno do espectro}
\chapter{Deficiência psicossocial e transtorno do espectro autista:}
\chapter{Deficiência psicossocial e transtorno do espectro autista: orientações}
\chapter{Deficiência psicossocial e transtorno do espectro autista: orientações para}
\chapter{Deficiência psicossocial e transtorno do espectro autista: orientações para a}
\chapter{Deficiência psicossocial e transtorno do espectro autista: orientações para a realização}
\chapter{Deficiência psicossocial e transtorno do espectro autista: orientações para a realização de}
\chapter{Deficiência psicossocial e transtorno do espectro autista: orientações para a realização de eventos}
\chapter{Deficiência psicossocial e transtorno do espectro autista: orientações para a realização de eventos e}
\chapter{Deficiência psicossocial e transtorno do espectro autista: orientações para a realização de eventos e apresentações}
\chapter{Deficiência psicossocial e transtorno do espectro autista: orientações para a realização de eventos e apresentações acessíveis}
\chapter{Deficiência psicossocial e transtorno do espectro autista: orientações para a realização de eventos e apresentações acessíveis em}
\chapter{Deficiência psicossocial e transtorno do espectro autista: orientações para a realização de eventos e apresentações acessíveis em todo}
\chapter{Comunicação alternativa e}
\chapter{Comunicação alternativa e aumentativa}
\chapter{Comunicação alternativa e aumentativa na}
\chapter{Comunicação alternativa e aumentativa na participação}
\chapter{Comunicação alternativa e aumentativa na participação social}
\chapter{Comunicação alternativa e aumentativa na participação social de}
\chapter{Comunicação alternativa e aumentativa na participação social de pessoas}
\chapter{Comunicação alternativa e aumentativa na participação social de pessoas com}
\chapter{Comunicação alternativa e aumentativa na participação social de pessoas com deficiência}
\chapter{Comunicação alternativa e aumentativa na participação social de pessoas com deficiência intelectual}
\chapter{Comunicação alternativa e aumentativa na participação social de pessoas com deficiência intelectual em}
\chapter{Comunicação alternativa e aumentativa na participação social de pessoas com deficiência intelectual em eventos}
\chapter{Comunicação alternativa e aumentativa na participação social de pessoas com deficiência intelectual em eventos públicos}
\chapter{Comunicação alternativa e aumentativa na participação social de pessoas com deficiência intelectual em eventos públicos de}
\chapter{Comunicação alternativa e aumentativa na participação social de pessoas com deficiência intelectual em eventos públicos de grande}
\chapter{Comunicação alternativa e aumentativa na participação social de pessoas com deficiência intelectual em eventos públicos de grande porte}
\chapter{Psicopedagogicamente inconstitucionalíssimo otorrinolaringologista}
\chapter{Psicopedagogicamente inconstitucionalíssimo otorrinolaringologista socioeducativo}
\chapter{Psicopedagogicamente inconstitucionalíssimo otorrinolaringologista socioeducativo ético-profissional}
\chapter{Psicopedagogicamente inconstitucionalíssimo otorrinolaringologista socioeducativo ético-profissional sociodemográfico}
\chapter{Psicopedagogicamente inconstitucionalíssimo otorrinolaringologista socioeducativo ético-profissional sociodemográfico interdisciplinaridade}
\chapter{Psicopedagogicamente inconstitucionalíssimo otorrinolaringologista socioeducativo ético-profissional sociodemográfico interdisciplinaridade neuropsicológica}
\chapter{Psicopedagogicamente inconstitucionalíssimo otorrinolaringologista socioeducativo ético-profissional sociodemográfico interdisciplinaridade neuropsicológica avaliação}
\chapter{Psicopedagogicamente inconstitucionalíssimo otorrinolaringologista socioeducativo ético-profissional sociodemográfico interdisciplinaridade neuropsicológica avaliação}
\chapter{Psicopedagogicamente inconstitucionalíssimo otorrinolaringologista socioeducativo ético-profissional sociodemográfico interdisciplinaridade neuropsicológica avaliação}
\chapter{Psicopedagogicamente inconstitucionalíssimo otorrinolaringologista socioeducativo ético-profissional sociodemográfico interdisciplinaridade neuropsicológica avaliação}
\chapter{Psicopedagogicamente inconstitucionalíssimo otorrinolaringologista socioeducativo ético-profissional sociodemográfico interdisciplinaridade neuropsicológica avaliação}
\chapter{Psicopedagogicamente inconstitucionalíssimo otorrinolaringologista socioeducativo ético-profissional sociodemográfico interdisciplinaridade neuropsicológica avaliação}
\chapter{Psicopedagogicamente inconstitucionalíssimo otorrinolaringologista socioeducativo ético-profissional sociodemográfico interdisciplinaridade neuropsicológica avaliação}
\chapter{Psicopedagogicamente inconstitucionalíssimo otorrinolaringologista socioeducativo ético-profissional sociodemográfico interdisciplinaridade neuropsicológica avaliação}
\chapter{Psicopedagogicamente inconstitucionalíssimo otorrinolaringologista socioeducativo ético-profissional sociodemográfico interdisciplinaridade neuropsicológica avaliação}
\chapter{Psicopedagogicamente inconstitucionalíssimo otorrinolaringologista socioeducativo ético-profissional sociodemográfico interdisciplinaridade neuropsicológica avaliação}
\end{document}
'   % B: plus 4em (o do guia)
%   lualatex '\def\var{C}% !TeX program = lualatex
% ===================================================================
% sumario-longo.tex — mede o sumário sem hifenização (issue #35, "Fora do
% escopo"): com a hifenização desligada, a entrada longa de capítulo pode
% bater no número da página. Gera 3 x 16 capítulos, cada um com um prefixo
% (de 3 a 18 palavras) de uma frase longa: frase 1 de palavras médias,
% frase 2 com números, frase 3 só de palavras muito longas. Conta no log
% os "Overfull \hbox" do sumário: o sumário é a única parte do documento
% sem parágrafo de corpo, então todo Overfull "in paragraph" antes do
% primeiro \chapter é entrada que passou da margem.
%
% Uso (de uma pasta com livros_crp_acessivel_book.sty e crpsp-leg.sty no
% caminho; as fontes NEWJUNE por OSFONTDIR):
%   lualatex '\def\var{R}\input{sumario-longo}'   % R: rígido (até 0.5.x)
%   lualatex '\def\var{B}\input{sumario-longo}'   % B: plus 4em (o do guia)
%   lualatex '\def\var{C}\input{sumario-longo}'   % C: plus 4em + \tolerance 10000
%   lualatex '\def\var{S}\input{sumario-longo}'   % S: o que a .sty define
% Cada variante só redefine o \l@chapter depois da carga da .sty.
% ===================================================================
\DocumentMetadata{lang=pt-BR, pdfstandard=ua-2, pdfversion=2.0, tagging=on}
\documentclass[11pt,a5paper,twoside,openright]{book}
\usepackage{livros_crp_acessivel_book}
\hypersetup{pdftitle={Sumario longo}, pdfauthor={CRP SP}}
\makeatletter
\providecommand\var{S}
\def\crp@vR{R}\def\crp@vB{B}\def\crp@vC{C}
\def\crp@lchap#1{%  #1 = acréscimo à margem direita (rightskip)
  \renewcommand*\l@chapter[2]{%
    \ifnum \c@tocdepth >\m@ne
        \addpenalty{-\@highpenalty}\addvspace{0.5em \@plus\p@}%
        \setlength\@tempdima{1.5em}%
        \begingroup
            \parindent \z@ \rightskip \@pnumwidth #1
            \parfillskip -\@pnumwidth
            \leavevmode \bfseries
            \advance\leftskip\@tempdima \hskip -\leftskip
            ##1\nobreak\hfil
            \nobreak\hb@xt@\@pnumwidth{\hss ##2\kern-\p@\kern\p@}\par
            \penalty\@highpenalty
        \endgroup
    \fi}}
\ifx\var\crp@vR \crp@lchap{}\fi
\ifx\var\crp@vB \crp@lchap{\@plus 4em}\fi
\ifx\var\crp@vC \crp@lchap{\@plus 4em \tolerance=10000 \hbadness=10000 }\fi
\makeatother
\begin{document}
\frontmatter
\tableofcontents
\mainmatter
\chapter{Deficiência psicossocial e}
\chapter{Deficiência psicossocial e transtorno}
\chapter{Deficiência psicossocial e transtorno do}
\chapter{Deficiência psicossocial e transtorno do espectro}
\chapter{Deficiência psicossocial e transtorno do espectro autista:}
\chapter{Deficiência psicossocial e transtorno do espectro autista: orientações}
\chapter{Deficiência psicossocial e transtorno do espectro autista: orientações para}
\chapter{Deficiência psicossocial e transtorno do espectro autista: orientações para a}
\chapter{Deficiência psicossocial e transtorno do espectro autista: orientações para a realização}
\chapter{Deficiência psicossocial e transtorno do espectro autista: orientações para a realização de}
\chapter{Deficiência psicossocial e transtorno do espectro autista: orientações para a realização de eventos}
\chapter{Deficiência psicossocial e transtorno do espectro autista: orientações para a realização de eventos e}
\chapter{Deficiência psicossocial e transtorno do espectro autista: orientações para a realização de eventos e apresentações}
\chapter{Deficiência psicossocial e transtorno do espectro autista: orientações para a realização de eventos e apresentações acessíveis}
\chapter{Deficiência psicossocial e transtorno do espectro autista: orientações para a realização de eventos e apresentações acessíveis em}
\chapter{Deficiência psicossocial e transtorno do espectro autista: orientações para a realização de eventos e apresentações acessíveis em todo}
\chapter{Comunicação alternativa e}
\chapter{Comunicação alternativa e aumentativa}
\chapter{Comunicação alternativa e aumentativa na}
\chapter{Comunicação alternativa e aumentativa na participação}
\chapter{Comunicação alternativa e aumentativa na participação social}
\chapter{Comunicação alternativa e aumentativa na participação social de}
\chapter{Comunicação alternativa e aumentativa na participação social de pessoas}
\chapter{Comunicação alternativa e aumentativa na participação social de pessoas com}
\chapter{Comunicação alternativa e aumentativa na participação social de pessoas com deficiência}
\chapter{Comunicação alternativa e aumentativa na participação social de pessoas com deficiência intelectual}
\chapter{Comunicação alternativa e aumentativa na participação social de pessoas com deficiência intelectual em}
\chapter{Comunicação alternativa e aumentativa na participação social de pessoas com deficiência intelectual em eventos}
\chapter{Comunicação alternativa e aumentativa na participação social de pessoas com deficiência intelectual em eventos públicos}
\chapter{Comunicação alternativa e aumentativa na participação social de pessoas com deficiência intelectual em eventos públicos de}
\chapter{Comunicação alternativa e aumentativa na participação social de pessoas com deficiência intelectual em eventos públicos de grande}
\chapter{Comunicação alternativa e aumentativa na participação social de pessoas com deficiência intelectual em eventos públicos de grande porte}
\chapter{Psicopedagogicamente inconstitucionalíssimo otorrinolaringologista}
\chapter{Psicopedagogicamente inconstitucionalíssimo otorrinolaringologista socioeducativo}
\chapter{Psicopedagogicamente inconstitucionalíssimo otorrinolaringologista socioeducativo ético-profissional}
\chapter{Psicopedagogicamente inconstitucionalíssimo otorrinolaringologista socioeducativo ético-profissional sociodemográfico}
\chapter{Psicopedagogicamente inconstitucionalíssimo otorrinolaringologista socioeducativo ético-profissional sociodemográfico interdisciplinaridade}
\chapter{Psicopedagogicamente inconstitucionalíssimo otorrinolaringologista socioeducativo ético-profissional sociodemográfico interdisciplinaridade neuropsicológica}
\chapter{Psicopedagogicamente inconstitucionalíssimo otorrinolaringologista socioeducativo ético-profissional sociodemográfico interdisciplinaridade neuropsicológica avaliação}
\chapter{Psicopedagogicamente inconstitucionalíssimo otorrinolaringologista socioeducativo ético-profissional sociodemográfico interdisciplinaridade neuropsicológica avaliação}
\chapter{Psicopedagogicamente inconstitucionalíssimo otorrinolaringologista socioeducativo ético-profissional sociodemográfico interdisciplinaridade neuropsicológica avaliação}
\chapter{Psicopedagogicamente inconstitucionalíssimo otorrinolaringologista socioeducativo ético-profissional sociodemográfico interdisciplinaridade neuropsicológica avaliação}
\chapter{Psicopedagogicamente inconstitucionalíssimo otorrinolaringologista socioeducativo ético-profissional sociodemográfico interdisciplinaridade neuropsicológica avaliação}
\chapter{Psicopedagogicamente inconstitucionalíssimo otorrinolaringologista socioeducativo ético-profissional sociodemográfico interdisciplinaridade neuropsicológica avaliação}
\chapter{Psicopedagogicamente inconstitucionalíssimo otorrinolaringologista socioeducativo ético-profissional sociodemográfico interdisciplinaridade neuropsicológica avaliação}
\chapter{Psicopedagogicamente inconstitucionalíssimo otorrinolaringologista socioeducativo ético-profissional sociodemográfico interdisciplinaridade neuropsicológica avaliação}
\chapter{Psicopedagogicamente inconstitucionalíssimo otorrinolaringologista socioeducativo ético-profissional sociodemográfico interdisciplinaridade neuropsicológica avaliação}
\chapter{Psicopedagogicamente inconstitucionalíssimo otorrinolaringologista socioeducativo ético-profissional sociodemográfico interdisciplinaridade neuropsicológica avaliação}
\end{document}
'   % C: plus 4em + \tolerance 10000
%   lualatex '\def\var{S}% !TeX program = lualatex
% ===================================================================
% sumario-longo.tex — mede o sumário sem hifenização (issue #35, "Fora do
% escopo"): com a hifenização desligada, a entrada longa de capítulo pode
% bater no número da página. Gera 3 x 16 capítulos, cada um com um prefixo
% (de 3 a 18 palavras) de uma frase longa: frase 1 de palavras médias,
% frase 2 com números, frase 3 só de palavras muito longas. Conta no log
% os "Overfull \hbox" do sumário: o sumário é a única parte do documento
% sem parágrafo de corpo, então todo Overfull "in paragraph" antes do
% primeiro \chapter é entrada que passou da margem.
%
% Uso (de uma pasta com livros_crp_acessivel_book.sty e crpsp-leg.sty no
% caminho; as fontes NEWJUNE por OSFONTDIR):
%   lualatex '\def\var{R}\input{sumario-longo}'   % R: rígido (até 0.5.x)
%   lualatex '\def\var{B}\input{sumario-longo}'   % B: plus 4em (o do guia)
%   lualatex '\def\var{C}\input{sumario-longo}'   % C: plus 4em + \tolerance 10000
%   lualatex '\def\var{S}\input{sumario-longo}'   % S: o que a .sty define
% Cada variante só redefine o \l@chapter depois da carga da .sty.
% ===================================================================
\DocumentMetadata{lang=pt-BR, pdfstandard=ua-2, pdfversion=2.0, tagging=on}
\documentclass[11pt,a5paper,twoside,openright]{book}
\usepackage{livros_crp_acessivel_book}
\hypersetup{pdftitle={Sumario longo}, pdfauthor={CRP SP}}
\makeatletter
\providecommand\var{S}
\def\crp@vR{R}\def\crp@vB{B}\def\crp@vC{C}
\def\crp@lchap#1{%  #1 = acréscimo à margem direita (rightskip)
  \renewcommand*\l@chapter[2]{%
    \ifnum \c@tocdepth >\m@ne
        \addpenalty{-\@highpenalty}\addvspace{0.5em \@plus\p@}%
        \setlength\@tempdima{1.5em}%
        \begingroup
            \parindent \z@ \rightskip \@pnumwidth #1
            \parfillskip -\@pnumwidth
            \leavevmode \bfseries
            \advance\leftskip\@tempdima \hskip -\leftskip
            ##1\nobreak\hfil
            \nobreak\hb@xt@\@pnumwidth{\hss ##2\kern-\p@\kern\p@}\par
            \penalty\@highpenalty
        \endgroup
    \fi}}
\ifx\var\crp@vR \crp@lchap{}\fi
\ifx\var\crp@vB \crp@lchap{\@plus 4em}\fi
\ifx\var\crp@vC \crp@lchap{\@plus 4em \tolerance=10000 \hbadness=10000 }\fi
\makeatother
\begin{document}
\frontmatter
\tableofcontents
\mainmatter
\chapter{Deficiência psicossocial e}
\chapter{Deficiência psicossocial e transtorno}
\chapter{Deficiência psicossocial e transtorno do}
\chapter{Deficiência psicossocial e transtorno do espectro}
\chapter{Deficiência psicossocial e transtorno do espectro autista:}
\chapter{Deficiência psicossocial e transtorno do espectro autista: orientações}
\chapter{Deficiência psicossocial e transtorno do espectro autista: orientações para}
\chapter{Deficiência psicossocial e transtorno do espectro autista: orientações para a}
\chapter{Deficiência psicossocial e transtorno do espectro autista: orientações para a realização}
\chapter{Deficiência psicossocial e transtorno do espectro autista: orientações para a realização de}
\chapter{Deficiência psicossocial e transtorno do espectro autista: orientações para a realização de eventos}
\chapter{Deficiência psicossocial e transtorno do espectro autista: orientações para a realização de eventos e}
\chapter{Deficiência psicossocial e transtorno do espectro autista: orientações para a realização de eventos e apresentações}
\chapter{Deficiência psicossocial e transtorno do espectro autista: orientações para a realização de eventos e apresentações acessíveis}
\chapter{Deficiência psicossocial e transtorno do espectro autista: orientações para a realização de eventos e apresentações acessíveis em}
\chapter{Deficiência psicossocial e transtorno do espectro autista: orientações para a realização de eventos e apresentações acessíveis em todo}
\chapter{Comunicação alternativa e}
\chapter{Comunicação alternativa e aumentativa}
\chapter{Comunicação alternativa e aumentativa na}
\chapter{Comunicação alternativa e aumentativa na participação}
\chapter{Comunicação alternativa e aumentativa na participação social}
\chapter{Comunicação alternativa e aumentativa na participação social de}
\chapter{Comunicação alternativa e aumentativa na participação social de pessoas}
\chapter{Comunicação alternativa e aumentativa na participação social de pessoas com}
\chapter{Comunicação alternativa e aumentativa na participação social de pessoas com deficiência}
\chapter{Comunicação alternativa e aumentativa na participação social de pessoas com deficiência intelectual}
\chapter{Comunicação alternativa e aumentativa na participação social de pessoas com deficiência intelectual em}
\chapter{Comunicação alternativa e aumentativa na participação social de pessoas com deficiência intelectual em eventos}
\chapter{Comunicação alternativa e aumentativa na participação social de pessoas com deficiência intelectual em eventos públicos}
\chapter{Comunicação alternativa e aumentativa na participação social de pessoas com deficiência intelectual em eventos públicos de}
\chapter{Comunicação alternativa e aumentativa na participação social de pessoas com deficiência intelectual em eventos públicos de grande}
\chapter{Comunicação alternativa e aumentativa na participação social de pessoas com deficiência intelectual em eventos públicos de grande porte}
\chapter{Psicopedagogicamente inconstitucionalíssimo otorrinolaringologista}
\chapter{Psicopedagogicamente inconstitucionalíssimo otorrinolaringologista socioeducativo}
\chapter{Psicopedagogicamente inconstitucionalíssimo otorrinolaringologista socioeducativo ético-profissional}
\chapter{Psicopedagogicamente inconstitucionalíssimo otorrinolaringologista socioeducativo ético-profissional sociodemográfico}
\chapter{Psicopedagogicamente inconstitucionalíssimo otorrinolaringologista socioeducativo ético-profissional sociodemográfico interdisciplinaridade}
\chapter{Psicopedagogicamente inconstitucionalíssimo otorrinolaringologista socioeducativo ético-profissional sociodemográfico interdisciplinaridade neuropsicológica}
\chapter{Psicopedagogicamente inconstitucionalíssimo otorrinolaringologista socioeducativo ético-profissional sociodemográfico interdisciplinaridade neuropsicológica avaliação}
\chapter{Psicopedagogicamente inconstitucionalíssimo otorrinolaringologista socioeducativo ético-profissional sociodemográfico interdisciplinaridade neuropsicológica avaliação}
\chapter{Psicopedagogicamente inconstitucionalíssimo otorrinolaringologista socioeducativo ético-profissional sociodemográfico interdisciplinaridade neuropsicológica avaliação}
\chapter{Psicopedagogicamente inconstitucionalíssimo otorrinolaringologista socioeducativo ético-profissional sociodemográfico interdisciplinaridade neuropsicológica avaliação}
\chapter{Psicopedagogicamente inconstitucionalíssimo otorrinolaringologista socioeducativo ético-profissional sociodemográfico interdisciplinaridade neuropsicológica avaliação}
\chapter{Psicopedagogicamente inconstitucionalíssimo otorrinolaringologista socioeducativo ético-profissional sociodemográfico interdisciplinaridade neuropsicológica avaliação}
\chapter{Psicopedagogicamente inconstitucionalíssimo otorrinolaringologista socioeducativo ético-profissional sociodemográfico interdisciplinaridade neuropsicológica avaliação}
\chapter{Psicopedagogicamente inconstitucionalíssimo otorrinolaringologista socioeducativo ético-profissional sociodemográfico interdisciplinaridade neuropsicológica avaliação}
\chapter{Psicopedagogicamente inconstitucionalíssimo otorrinolaringologista socioeducativo ético-profissional sociodemográfico interdisciplinaridade neuropsicológica avaliação}
\chapter{Psicopedagogicamente inconstitucionalíssimo otorrinolaringologista socioeducativo ético-profissional sociodemográfico interdisciplinaridade neuropsicológica avaliação}
\end{document}
'   % S: o que a .sty define
% Cada variante só redefine o \l@chapter depois da carga da .sty.
% ===================================================================
\DocumentMetadata{lang=pt-BR, pdfstandard=ua-2, pdfversion=2.0, tagging=on}
\documentclass[11pt,a5paper,twoside,openright]{book}
\usepackage{livros_crp_acessivel_book}
\hypersetup{pdftitle={Sumario longo}, pdfauthor={CRP SP}}
\makeatletter
\providecommand\var{S}
\def\crp@vR{R}\def\crp@vB{B}\def\crp@vC{C}
\def\crp@lchap#1{%  #1 = acréscimo à margem direita (rightskip)
  \renewcommand*\l@chapter[2]{%
    \ifnum \c@tocdepth >\m@ne
        \addpenalty{-\@highpenalty}\addvspace{0.5em \@plus\p@}%
        \setlength\@tempdima{1.5em}%
        \begingroup
            \parindent \z@ \rightskip \@pnumwidth #1
            \parfillskip -\@pnumwidth
            \leavevmode \bfseries
            \advance\leftskip\@tempdima \hskip -\leftskip
            ##1\nobreak\hfil
            \nobreak\hb@xt@\@pnumwidth{\hss ##2\kern-\p@\kern\p@}\par
            \penalty\@highpenalty
        \endgroup
    \fi}}
\ifx\var\crp@vR \crp@lchap{}\fi
\ifx\var\crp@vB \crp@lchap{\@plus 4em}\fi
\ifx\var\crp@vC \crp@lchap{\@plus 4em \tolerance=10000 \hbadness=10000 }\fi
\makeatother
\begin{document}
\frontmatter
\tableofcontents
\mainmatter
\chapter{Deficiência psicossocial e}
\chapter{Deficiência psicossocial e transtorno}
\chapter{Deficiência psicossocial e transtorno do}
\chapter{Deficiência psicossocial e transtorno do espectro}
\chapter{Deficiência psicossocial e transtorno do espectro autista:}
\chapter{Deficiência psicossocial e transtorno do espectro autista: orientações}
\chapter{Deficiência psicossocial e transtorno do espectro autista: orientações para}
\chapter{Deficiência psicossocial e transtorno do espectro autista: orientações para a}
\chapter{Deficiência psicossocial e transtorno do espectro autista: orientações para a realização}
\chapter{Deficiência psicossocial e transtorno do espectro autista: orientações para a realização de}
\chapter{Deficiência psicossocial e transtorno do espectro autista: orientações para a realização de eventos}
\chapter{Deficiência psicossocial e transtorno do espectro autista: orientações para a realização de eventos e}
\chapter{Deficiência psicossocial e transtorno do espectro autista: orientações para a realização de eventos e apresentações}
\chapter{Deficiência psicossocial e transtorno do espectro autista: orientações para a realização de eventos e apresentações acessíveis}
\chapter{Deficiência psicossocial e transtorno do espectro autista: orientações para a realização de eventos e apresentações acessíveis em}
\chapter{Deficiência psicossocial e transtorno do espectro autista: orientações para a realização de eventos e apresentações acessíveis em todo}
\chapter{Comunicação alternativa e}
\chapter{Comunicação alternativa e aumentativa}
\chapter{Comunicação alternativa e aumentativa na}
\chapter{Comunicação alternativa e aumentativa na participação}
\chapter{Comunicação alternativa e aumentativa na participação social}
\chapter{Comunicação alternativa e aumentativa na participação social de}
\chapter{Comunicação alternativa e aumentativa na participação social de pessoas}
\chapter{Comunicação alternativa e aumentativa na participação social de pessoas com}
\chapter{Comunicação alternativa e aumentativa na participação social de pessoas com deficiência}
\chapter{Comunicação alternativa e aumentativa na participação social de pessoas com deficiência intelectual}
\chapter{Comunicação alternativa e aumentativa na participação social de pessoas com deficiência intelectual em}
\chapter{Comunicação alternativa e aumentativa na participação social de pessoas com deficiência intelectual em eventos}
\chapter{Comunicação alternativa e aumentativa na participação social de pessoas com deficiência intelectual em eventos públicos}
\chapter{Comunicação alternativa e aumentativa na participação social de pessoas com deficiência intelectual em eventos públicos de}
\chapter{Comunicação alternativa e aumentativa na participação social de pessoas com deficiência intelectual em eventos públicos de grande}
\chapter{Comunicação alternativa e aumentativa na participação social de pessoas com deficiência intelectual em eventos públicos de grande porte}
\chapter{Psicopedagogicamente inconstitucionalíssimo otorrinolaringologista}
\chapter{Psicopedagogicamente inconstitucionalíssimo otorrinolaringologista socioeducativo}
\chapter{Psicopedagogicamente inconstitucionalíssimo otorrinolaringologista socioeducativo ético-profissional}
\chapter{Psicopedagogicamente inconstitucionalíssimo otorrinolaringologista socioeducativo ético-profissional sociodemográfico}
\chapter{Psicopedagogicamente inconstitucionalíssimo otorrinolaringologista socioeducativo ético-profissional sociodemográfico interdisciplinaridade}
\chapter{Psicopedagogicamente inconstitucionalíssimo otorrinolaringologista socioeducativo ético-profissional sociodemográfico interdisciplinaridade neuropsicológica}
\chapter{Psicopedagogicamente inconstitucionalíssimo otorrinolaringologista socioeducativo ético-profissional sociodemográfico interdisciplinaridade neuropsicológica avaliação}
\chapter{Psicopedagogicamente inconstitucionalíssimo otorrinolaringologista socioeducativo ético-profissional sociodemográfico interdisciplinaridade neuropsicológica avaliação}
\chapter{Psicopedagogicamente inconstitucionalíssimo otorrinolaringologista socioeducativo ético-profissional sociodemográfico interdisciplinaridade neuropsicológica avaliação}
\chapter{Psicopedagogicamente inconstitucionalíssimo otorrinolaringologista socioeducativo ético-profissional sociodemográfico interdisciplinaridade neuropsicológica avaliação}
\chapter{Psicopedagogicamente inconstitucionalíssimo otorrinolaringologista socioeducativo ético-profissional sociodemográfico interdisciplinaridade neuropsicológica avaliação}
\chapter{Psicopedagogicamente inconstitucionalíssimo otorrinolaringologista socioeducativo ético-profissional sociodemográfico interdisciplinaridade neuropsicológica avaliação}
\chapter{Psicopedagogicamente inconstitucionalíssimo otorrinolaringologista socioeducativo ético-profissional sociodemográfico interdisciplinaridade neuropsicológica avaliação}
\chapter{Psicopedagogicamente inconstitucionalíssimo otorrinolaringologista socioeducativo ético-profissional sociodemográfico interdisciplinaridade neuropsicológica avaliação}
\chapter{Psicopedagogicamente inconstitucionalíssimo otorrinolaringologista socioeducativo ético-profissional sociodemográfico interdisciplinaridade neuropsicológica avaliação}
\chapter{Psicopedagogicamente inconstitucionalíssimo otorrinolaringologista socioeducativo ético-profissional sociodemográfico interdisciplinaridade neuropsicológica avaliação}
\end{document}
'   % R: rígido (até 0.5.x)
%   lualatex '\def\var{B}% !TeX program = lualatex
% ===================================================================
% sumario-longo.tex — mede o sumário sem hifenização (issue #35, "Fora do
% escopo"): com a hifenização desligada, a entrada longa de capítulo pode
% bater no número da página. Gera 3 x 16 capítulos, cada um com um prefixo
% (de 3 a 18 palavras) de uma frase longa: frase 1 de palavras médias,
% frase 2 com números, frase 3 só de palavras muito longas. Conta no log
% os "Overfull \hbox" do sumário: o sumário é a única parte do documento
% sem parágrafo de corpo, então todo Overfull "in paragraph" antes do
% primeiro \chapter é entrada que passou da margem.
%
% Uso (de uma pasta com livros_crp_acessivel_book.sty e crpsp-leg.sty no
% caminho; as fontes NEWJUNE por OSFONTDIR):
%   lualatex '\def\var{R}% !TeX program = lualatex
% ===================================================================
% sumario-longo.tex — mede o sumário sem hifenização (issue #35, "Fora do
% escopo"): com a hifenização desligada, a entrada longa de capítulo pode
% bater no número da página. Gera 3 x 16 capítulos, cada um com um prefixo
% (de 3 a 18 palavras) de uma frase longa: frase 1 de palavras médias,
% frase 2 com números, frase 3 só de palavras muito longas. Conta no log
% os "Overfull \hbox" do sumário: o sumário é a única parte do documento
% sem parágrafo de corpo, então todo Overfull "in paragraph" antes do
% primeiro \chapter é entrada que passou da margem.
%
% Uso (de uma pasta com livros_crp_acessivel_book.sty e crpsp-leg.sty no
% caminho; as fontes NEWJUNE por OSFONTDIR):
%   lualatex '\def\var{R}\input{sumario-longo}'   % R: rígido (até 0.5.x)
%   lualatex '\def\var{B}\input{sumario-longo}'   % B: plus 4em (o do guia)
%   lualatex '\def\var{C}\input{sumario-longo}'   % C: plus 4em + \tolerance 10000
%   lualatex '\def\var{S}\input{sumario-longo}'   % S: o que a .sty define
% Cada variante só redefine o \l@chapter depois da carga da .sty.
% ===================================================================
\DocumentMetadata{lang=pt-BR, pdfstandard=ua-2, pdfversion=2.0, tagging=on}
\documentclass[11pt,a5paper,twoside,openright]{book}
\usepackage{livros_crp_acessivel_book}
\hypersetup{pdftitle={Sumario longo}, pdfauthor={CRP SP}}
\makeatletter
\providecommand\var{S}
\def\crp@vR{R}\def\crp@vB{B}\def\crp@vC{C}
\def\crp@lchap#1{%  #1 = acréscimo à margem direita (rightskip)
  \renewcommand*\l@chapter[2]{%
    \ifnum \c@tocdepth >\m@ne
        \addpenalty{-\@highpenalty}\addvspace{0.5em \@plus\p@}%
        \setlength\@tempdima{1.5em}%
        \begingroup
            \parindent \z@ \rightskip \@pnumwidth #1
            \parfillskip -\@pnumwidth
            \leavevmode \bfseries
            \advance\leftskip\@tempdima \hskip -\leftskip
            ##1\nobreak\hfil
            \nobreak\hb@xt@\@pnumwidth{\hss ##2\kern-\p@\kern\p@}\par
            \penalty\@highpenalty
        \endgroup
    \fi}}
\ifx\var\crp@vR \crp@lchap{}\fi
\ifx\var\crp@vB \crp@lchap{\@plus 4em}\fi
\ifx\var\crp@vC \crp@lchap{\@plus 4em \tolerance=10000 \hbadness=10000 }\fi
\makeatother
\begin{document}
\frontmatter
\tableofcontents
\mainmatter
\chapter{Deficiência psicossocial e}
\chapter{Deficiência psicossocial e transtorno}
\chapter{Deficiência psicossocial e transtorno do}
\chapter{Deficiência psicossocial e transtorno do espectro}
\chapter{Deficiência psicossocial e transtorno do espectro autista:}
\chapter{Deficiência psicossocial e transtorno do espectro autista: orientações}
\chapter{Deficiência psicossocial e transtorno do espectro autista: orientações para}
\chapter{Deficiência psicossocial e transtorno do espectro autista: orientações para a}
\chapter{Deficiência psicossocial e transtorno do espectro autista: orientações para a realização}
\chapter{Deficiência psicossocial e transtorno do espectro autista: orientações para a realização de}
\chapter{Deficiência psicossocial e transtorno do espectro autista: orientações para a realização de eventos}
\chapter{Deficiência psicossocial e transtorno do espectro autista: orientações para a realização de eventos e}
\chapter{Deficiência psicossocial e transtorno do espectro autista: orientações para a realização de eventos e apresentações}
\chapter{Deficiência psicossocial e transtorno do espectro autista: orientações para a realização de eventos e apresentações acessíveis}
\chapter{Deficiência psicossocial e transtorno do espectro autista: orientações para a realização de eventos e apresentações acessíveis em}
\chapter{Deficiência psicossocial e transtorno do espectro autista: orientações para a realização de eventos e apresentações acessíveis em todo}
\chapter{Comunicação alternativa e}
\chapter{Comunicação alternativa e aumentativa}
\chapter{Comunicação alternativa e aumentativa na}
\chapter{Comunicação alternativa e aumentativa na participação}
\chapter{Comunicação alternativa e aumentativa na participação social}
\chapter{Comunicação alternativa e aumentativa na participação social de}
\chapter{Comunicação alternativa e aumentativa na participação social de pessoas}
\chapter{Comunicação alternativa e aumentativa na participação social de pessoas com}
\chapter{Comunicação alternativa e aumentativa na participação social de pessoas com deficiência}
\chapter{Comunicação alternativa e aumentativa na participação social de pessoas com deficiência intelectual}
\chapter{Comunicação alternativa e aumentativa na participação social de pessoas com deficiência intelectual em}
\chapter{Comunicação alternativa e aumentativa na participação social de pessoas com deficiência intelectual em eventos}
\chapter{Comunicação alternativa e aumentativa na participação social de pessoas com deficiência intelectual em eventos públicos}
\chapter{Comunicação alternativa e aumentativa na participação social de pessoas com deficiência intelectual em eventos públicos de}
\chapter{Comunicação alternativa e aumentativa na participação social de pessoas com deficiência intelectual em eventos públicos de grande}
\chapter{Comunicação alternativa e aumentativa na participação social de pessoas com deficiência intelectual em eventos públicos de grande porte}
\chapter{Psicopedagogicamente inconstitucionalíssimo otorrinolaringologista}
\chapter{Psicopedagogicamente inconstitucionalíssimo otorrinolaringologista socioeducativo}
\chapter{Psicopedagogicamente inconstitucionalíssimo otorrinolaringologista socioeducativo ético-profissional}
\chapter{Psicopedagogicamente inconstitucionalíssimo otorrinolaringologista socioeducativo ético-profissional sociodemográfico}
\chapter{Psicopedagogicamente inconstitucionalíssimo otorrinolaringologista socioeducativo ético-profissional sociodemográfico interdisciplinaridade}
\chapter{Psicopedagogicamente inconstitucionalíssimo otorrinolaringologista socioeducativo ético-profissional sociodemográfico interdisciplinaridade neuropsicológica}
\chapter{Psicopedagogicamente inconstitucionalíssimo otorrinolaringologista socioeducativo ético-profissional sociodemográfico interdisciplinaridade neuropsicológica avaliação}
\chapter{Psicopedagogicamente inconstitucionalíssimo otorrinolaringologista socioeducativo ético-profissional sociodemográfico interdisciplinaridade neuropsicológica avaliação}
\chapter{Psicopedagogicamente inconstitucionalíssimo otorrinolaringologista socioeducativo ético-profissional sociodemográfico interdisciplinaridade neuropsicológica avaliação}
\chapter{Psicopedagogicamente inconstitucionalíssimo otorrinolaringologista socioeducativo ético-profissional sociodemográfico interdisciplinaridade neuropsicológica avaliação}
\chapter{Psicopedagogicamente inconstitucionalíssimo otorrinolaringologista socioeducativo ético-profissional sociodemográfico interdisciplinaridade neuropsicológica avaliação}
\chapter{Psicopedagogicamente inconstitucionalíssimo otorrinolaringologista socioeducativo ético-profissional sociodemográfico interdisciplinaridade neuropsicológica avaliação}
\chapter{Psicopedagogicamente inconstitucionalíssimo otorrinolaringologista socioeducativo ético-profissional sociodemográfico interdisciplinaridade neuropsicológica avaliação}
\chapter{Psicopedagogicamente inconstitucionalíssimo otorrinolaringologista socioeducativo ético-profissional sociodemográfico interdisciplinaridade neuropsicológica avaliação}
\chapter{Psicopedagogicamente inconstitucionalíssimo otorrinolaringologista socioeducativo ético-profissional sociodemográfico interdisciplinaridade neuropsicológica avaliação}
\chapter{Psicopedagogicamente inconstitucionalíssimo otorrinolaringologista socioeducativo ético-profissional sociodemográfico interdisciplinaridade neuropsicológica avaliação}
\end{document}
'   % R: rígido (até 0.5.x)
%   lualatex '\def\var{B}% !TeX program = lualatex
% ===================================================================
% sumario-longo.tex — mede o sumário sem hifenização (issue #35, "Fora do
% escopo"): com a hifenização desligada, a entrada longa de capítulo pode
% bater no número da página. Gera 3 x 16 capítulos, cada um com um prefixo
% (de 3 a 18 palavras) de uma frase longa: frase 1 de palavras médias,
% frase 2 com números, frase 3 só de palavras muito longas. Conta no log
% os "Overfull \hbox" do sumário: o sumário é a única parte do documento
% sem parágrafo de corpo, então todo Overfull "in paragraph" antes do
% primeiro \chapter é entrada que passou da margem.
%
% Uso (de uma pasta com livros_crp_acessivel_book.sty e crpsp-leg.sty no
% caminho; as fontes NEWJUNE por OSFONTDIR):
%   lualatex '\def\var{R}\input{sumario-longo}'   % R: rígido (até 0.5.x)
%   lualatex '\def\var{B}\input{sumario-longo}'   % B: plus 4em (o do guia)
%   lualatex '\def\var{C}\input{sumario-longo}'   % C: plus 4em + \tolerance 10000
%   lualatex '\def\var{S}\input{sumario-longo}'   % S: o que a .sty define
% Cada variante só redefine o \l@chapter depois da carga da .sty.
% ===================================================================
\DocumentMetadata{lang=pt-BR, pdfstandard=ua-2, pdfversion=2.0, tagging=on}
\documentclass[11pt,a5paper,twoside,openright]{book}
\usepackage{livros_crp_acessivel_book}
\hypersetup{pdftitle={Sumario longo}, pdfauthor={CRP SP}}
\makeatletter
\providecommand\var{S}
\def\crp@vR{R}\def\crp@vB{B}\def\crp@vC{C}
\def\crp@lchap#1{%  #1 = acréscimo à margem direita (rightskip)
  \renewcommand*\l@chapter[2]{%
    \ifnum \c@tocdepth >\m@ne
        \addpenalty{-\@highpenalty}\addvspace{0.5em \@plus\p@}%
        \setlength\@tempdima{1.5em}%
        \begingroup
            \parindent \z@ \rightskip \@pnumwidth #1
            \parfillskip -\@pnumwidth
            \leavevmode \bfseries
            \advance\leftskip\@tempdima \hskip -\leftskip
            ##1\nobreak\hfil
            \nobreak\hb@xt@\@pnumwidth{\hss ##2\kern-\p@\kern\p@}\par
            \penalty\@highpenalty
        \endgroup
    \fi}}
\ifx\var\crp@vR \crp@lchap{}\fi
\ifx\var\crp@vB \crp@lchap{\@plus 4em}\fi
\ifx\var\crp@vC \crp@lchap{\@plus 4em \tolerance=10000 \hbadness=10000 }\fi
\makeatother
\begin{document}
\frontmatter
\tableofcontents
\mainmatter
\chapter{Deficiência psicossocial e}
\chapter{Deficiência psicossocial e transtorno}
\chapter{Deficiência psicossocial e transtorno do}
\chapter{Deficiência psicossocial e transtorno do espectro}
\chapter{Deficiência psicossocial e transtorno do espectro autista:}
\chapter{Deficiência psicossocial e transtorno do espectro autista: orientações}
\chapter{Deficiência psicossocial e transtorno do espectro autista: orientações para}
\chapter{Deficiência psicossocial e transtorno do espectro autista: orientações para a}
\chapter{Deficiência psicossocial e transtorno do espectro autista: orientações para a realização}
\chapter{Deficiência psicossocial e transtorno do espectro autista: orientações para a realização de}
\chapter{Deficiência psicossocial e transtorno do espectro autista: orientações para a realização de eventos}
\chapter{Deficiência psicossocial e transtorno do espectro autista: orientações para a realização de eventos e}
\chapter{Deficiência psicossocial e transtorno do espectro autista: orientações para a realização de eventos e apresentações}
\chapter{Deficiência psicossocial e transtorno do espectro autista: orientações para a realização de eventos e apresentações acessíveis}
\chapter{Deficiência psicossocial e transtorno do espectro autista: orientações para a realização de eventos e apresentações acessíveis em}
\chapter{Deficiência psicossocial e transtorno do espectro autista: orientações para a realização de eventos e apresentações acessíveis em todo}
\chapter{Comunicação alternativa e}
\chapter{Comunicação alternativa e aumentativa}
\chapter{Comunicação alternativa e aumentativa na}
\chapter{Comunicação alternativa e aumentativa na participação}
\chapter{Comunicação alternativa e aumentativa na participação social}
\chapter{Comunicação alternativa e aumentativa na participação social de}
\chapter{Comunicação alternativa e aumentativa na participação social de pessoas}
\chapter{Comunicação alternativa e aumentativa na participação social de pessoas com}
\chapter{Comunicação alternativa e aumentativa na participação social de pessoas com deficiência}
\chapter{Comunicação alternativa e aumentativa na participação social de pessoas com deficiência intelectual}
\chapter{Comunicação alternativa e aumentativa na participação social de pessoas com deficiência intelectual em}
\chapter{Comunicação alternativa e aumentativa na participação social de pessoas com deficiência intelectual em eventos}
\chapter{Comunicação alternativa e aumentativa na participação social de pessoas com deficiência intelectual em eventos públicos}
\chapter{Comunicação alternativa e aumentativa na participação social de pessoas com deficiência intelectual em eventos públicos de}
\chapter{Comunicação alternativa e aumentativa na participação social de pessoas com deficiência intelectual em eventos públicos de grande}
\chapter{Comunicação alternativa e aumentativa na participação social de pessoas com deficiência intelectual em eventos públicos de grande porte}
\chapter{Psicopedagogicamente inconstitucionalíssimo otorrinolaringologista}
\chapter{Psicopedagogicamente inconstitucionalíssimo otorrinolaringologista socioeducativo}
\chapter{Psicopedagogicamente inconstitucionalíssimo otorrinolaringologista socioeducativo ético-profissional}
\chapter{Psicopedagogicamente inconstitucionalíssimo otorrinolaringologista socioeducativo ético-profissional sociodemográfico}
\chapter{Psicopedagogicamente inconstitucionalíssimo otorrinolaringologista socioeducativo ético-profissional sociodemográfico interdisciplinaridade}
\chapter{Psicopedagogicamente inconstitucionalíssimo otorrinolaringologista socioeducativo ético-profissional sociodemográfico interdisciplinaridade neuropsicológica}
\chapter{Psicopedagogicamente inconstitucionalíssimo otorrinolaringologista socioeducativo ético-profissional sociodemográfico interdisciplinaridade neuropsicológica avaliação}
\chapter{Psicopedagogicamente inconstitucionalíssimo otorrinolaringologista socioeducativo ético-profissional sociodemográfico interdisciplinaridade neuropsicológica avaliação}
\chapter{Psicopedagogicamente inconstitucionalíssimo otorrinolaringologista socioeducativo ético-profissional sociodemográfico interdisciplinaridade neuropsicológica avaliação}
\chapter{Psicopedagogicamente inconstitucionalíssimo otorrinolaringologista socioeducativo ético-profissional sociodemográfico interdisciplinaridade neuropsicológica avaliação}
\chapter{Psicopedagogicamente inconstitucionalíssimo otorrinolaringologista socioeducativo ético-profissional sociodemográfico interdisciplinaridade neuropsicológica avaliação}
\chapter{Psicopedagogicamente inconstitucionalíssimo otorrinolaringologista socioeducativo ético-profissional sociodemográfico interdisciplinaridade neuropsicológica avaliação}
\chapter{Psicopedagogicamente inconstitucionalíssimo otorrinolaringologista socioeducativo ético-profissional sociodemográfico interdisciplinaridade neuropsicológica avaliação}
\chapter{Psicopedagogicamente inconstitucionalíssimo otorrinolaringologista socioeducativo ético-profissional sociodemográfico interdisciplinaridade neuropsicológica avaliação}
\chapter{Psicopedagogicamente inconstitucionalíssimo otorrinolaringologista socioeducativo ético-profissional sociodemográfico interdisciplinaridade neuropsicológica avaliação}
\chapter{Psicopedagogicamente inconstitucionalíssimo otorrinolaringologista socioeducativo ético-profissional sociodemográfico interdisciplinaridade neuropsicológica avaliação}
\end{document}
'   % B: plus 4em (o do guia)
%   lualatex '\def\var{C}% !TeX program = lualatex
% ===================================================================
% sumario-longo.tex — mede o sumário sem hifenização (issue #35, "Fora do
% escopo"): com a hifenização desligada, a entrada longa de capítulo pode
% bater no número da página. Gera 3 x 16 capítulos, cada um com um prefixo
% (de 3 a 18 palavras) de uma frase longa: frase 1 de palavras médias,
% frase 2 com números, frase 3 só de palavras muito longas. Conta no log
% os "Overfull \hbox" do sumário: o sumário é a única parte do documento
% sem parágrafo de corpo, então todo Overfull "in paragraph" antes do
% primeiro \chapter é entrada que passou da margem.
%
% Uso (de uma pasta com livros_crp_acessivel_book.sty e crpsp-leg.sty no
% caminho; as fontes NEWJUNE por OSFONTDIR):
%   lualatex '\def\var{R}\input{sumario-longo}'   % R: rígido (até 0.5.x)
%   lualatex '\def\var{B}\input{sumario-longo}'   % B: plus 4em (o do guia)
%   lualatex '\def\var{C}\input{sumario-longo}'   % C: plus 4em + \tolerance 10000
%   lualatex '\def\var{S}\input{sumario-longo}'   % S: o que a .sty define
% Cada variante só redefine o \l@chapter depois da carga da .sty.
% ===================================================================
\DocumentMetadata{lang=pt-BR, pdfstandard=ua-2, pdfversion=2.0, tagging=on}
\documentclass[11pt,a5paper,twoside,openright]{book}
\usepackage{livros_crp_acessivel_book}
\hypersetup{pdftitle={Sumario longo}, pdfauthor={CRP SP}}
\makeatletter
\providecommand\var{S}
\def\crp@vR{R}\def\crp@vB{B}\def\crp@vC{C}
\def\crp@lchap#1{%  #1 = acréscimo à margem direita (rightskip)
  \renewcommand*\l@chapter[2]{%
    \ifnum \c@tocdepth >\m@ne
        \addpenalty{-\@highpenalty}\addvspace{0.5em \@plus\p@}%
        \setlength\@tempdima{1.5em}%
        \begingroup
            \parindent \z@ \rightskip \@pnumwidth #1
            \parfillskip -\@pnumwidth
            \leavevmode \bfseries
            \advance\leftskip\@tempdima \hskip -\leftskip
            ##1\nobreak\hfil
            \nobreak\hb@xt@\@pnumwidth{\hss ##2\kern-\p@\kern\p@}\par
            \penalty\@highpenalty
        \endgroup
    \fi}}
\ifx\var\crp@vR \crp@lchap{}\fi
\ifx\var\crp@vB \crp@lchap{\@plus 4em}\fi
\ifx\var\crp@vC \crp@lchap{\@plus 4em \tolerance=10000 \hbadness=10000 }\fi
\makeatother
\begin{document}
\frontmatter
\tableofcontents
\mainmatter
\chapter{Deficiência psicossocial e}
\chapter{Deficiência psicossocial e transtorno}
\chapter{Deficiência psicossocial e transtorno do}
\chapter{Deficiência psicossocial e transtorno do espectro}
\chapter{Deficiência psicossocial e transtorno do espectro autista:}
\chapter{Deficiência psicossocial e transtorno do espectro autista: orientações}
\chapter{Deficiência psicossocial e transtorno do espectro autista: orientações para}
\chapter{Deficiência psicossocial e transtorno do espectro autista: orientações para a}
\chapter{Deficiência psicossocial e transtorno do espectro autista: orientações para a realização}
\chapter{Deficiência psicossocial e transtorno do espectro autista: orientações para a realização de}
\chapter{Deficiência psicossocial e transtorno do espectro autista: orientações para a realização de eventos}
\chapter{Deficiência psicossocial e transtorno do espectro autista: orientações para a realização de eventos e}
\chapter{Deficiência psicossocial e transtorno do espectro autista: orientações para a realização de eventos e apresentações}
\chapter{Deficiência psicossocial e transtorno do espectro autista: orientações para a realização de eventos e apresentações acessíveis}
\chapter{Deficiência psicossocial e transtorno do espectro autista: orientações para a realização de eventos e apresentações acessíveis em}
\chapter{Deficiência psicossocial e transtorno do espectro autista: orientações para a realização de eventos e apresentações acessíveis em todo}
\chapter{Comunicação alternativa e}
\chapter{Comunicação alternativa e aumentativa}
\chapter{Comunicação alternativa e aumentativa na}
\chapter{Comunicação alternativa e aumentativa na participação}
\chapter{Comunicação alternativa e aumentativa na participação social}
\chapter{Comunicação alternativa e aumentativa na participação social de}
\chapter{Comunicação alternativa e aumentativa na participação social de pessoas}
\chapter{Comunicação alternativa e aumentativa na participação social de pessoas com}
\chapter{Comunicação alternativa e aumentativa na participação social de pessoas com deficiência}
\chapter{Comunicação alternativa e aumentativa na participação social de pessoas com deficiência intelectual}
\chapter{Comunicação alternativa e aumentativa na participação social de pessoas com deficiência intelectual em}
\chapter{Comunicação alternativa e aumentativa na participação social de pessoas com deficiência intelectual em eventos}
\chapter{Comunicação alternativa e aumentativa na participação social de pessoas com deficiência intelectual em eventos públicos}
\chapter{Comunicação alternativa e aumentativa na participação social de pessoas com deficiência intelectual em eventos públicos de}
\chapter{Comunicação alternativa e aumentativa na participação social de pessoas com deficiência intelectual em eventos públicos de grande}
\chapter{Comunicação alternativa e aumentativa na participação social de pessoas com deficiência intelectual em eventos públicos de grande porte}
\chapter{Psicopedagogicamente inconstitucionalíssimo otorrinolaringologista}
\chapter{Psicopedagogicamente inconstitucionalíssimo otorrinolaringologista socioeducativo}
\chapter{Psicopedagogicamente inconstitucionalíssimo otorrinolaringologista socioeducativo ético-profissional}
\chapter{Psicopedagogicamente inconstitucionalíssimo otorrinolaringologista socioeducativo ético-profissional sociodemográfico}
\chapter{Psicopedagogicamente inconstitucionalíssimo otorrinolaringologista socioeducativo ético-profissional sociodemográfico interdisciplinaridade}
\chapter{Psicopedagogicamente inconstitucionalíssimo otorrinolaringologista socioeducativo ético-profissional sociodemográfico interdisciplinaridade neuropsicológica}
\chapter{Psicopedagogicamente inconstitucionalíssimo otorrinolaringologista socioeducativo ético-profissional sociodemográfico interdisciplinaridade neuropsicológica avaliação}
\chapter{Psicopedagogicamente inconstitucionalíssimo otorrinolaringologista socioeducativo ético-profissional sociodemográfico interdisciplinaridade neuropsicológica avaliação}
\chapter{Psicopedagogicamente inconstitucionalíssimo otorrinolaringologista socioeducativo ético-profissional sociodemográfico interdisciplinaridade neuropsicológica avaliação}
\chapter{Psicopedagogicamente inconstitucionalíssimo otorrinolaringologista socioeducativo ético-profissional sociodemográfico interdisciplinaridade neuropsicológica avaliação}
\chapter{Psicopedagogicamente inconstitucionalíssimo otorrinolaringologista socioeducativo ético-profissional sociodemográfico interdisciplinaridade neuropsicológica avaliação}
\chapter{Psicopedagogicamente inconstitucionalíssimo otorrinolaringologista socioeducativo ético-profissional sociodemográfico interdisciplinaridade neuropsicológica avaliação}
\chapter{Psicopedagogicamente inconstitucionalíssimo otorrinolaringologista socioeducativo ético-profissional sociodemográfico interdisciplinaridade neuropsicológica avaliação}
\chapter{Psicopedagogicamente inconstitucionalíssimo otorrinolaringologista socioeducativo ético-profissional sociodemográfico interdisciplinaridade neuropsicológica avaliação}
\chapter{Psicopedagogicamente inconstitucionalíssimo otorrinolaringologista socioeducativo ético-profissional sociodemográfico interdisciplinaridade neuropsicológica avaliação}
\chapter{Psicopedagogicamente inconstitucionalíssimo otorrinolaringologista socioeducativo ético-profissional sociodemográfico interdisciplinaridade neuropsicológica avaliação}
\end{document}
'   % C: plus 4em + \tolerance 10000
%   lualatex '\def\var{S}% !TeX program = lualatex
% ===================================================================
% sumario-longo.tex — mede o sumário sem hifenização (issue #35, "Fora do
% escopo"): com a hifenização desligada, a entrada longa de capítulo pode
% bater no número da página. Gera 3 x 16 capítulos, cada um com um prefixo
% (de 3 a 18 palavras) de uma frase longa: frase 1 de palavras médias,
% frase 2 com números, frase 3 só de palavras muito longas. Conta no log
% os "Overfull \hbox" do sumário: o sumário é a única parte do documento
% sem parágrafo de corpo, então todo Overfull "in paragraph" antes do
% primeiro \chapter é entrada que passou da margem.
%
% Uso (de uma pasta com livros_crp_acessivel_book.sty e crpsp-leg.sty no
% caminho; as fontes NEWJUNE por OSFONTDIR):
%   lualatex '\def\var{R}\input{sumario-longo}'   % R: rígido (até 0.5.x)
%   lualatex '\def\var{B}\input{sumario-longo}'   % B: plus 4em (o do guia)
%   lualatex '\def\var{C}\input{sumario-longo}'   % C: plus 4em + \tolerance 10000
%   lualatex '\def\var{S}\input{sumario-longo}'   % S: o que a .sty define
% Cada variante só redefine o \l@chapter depois da carga da .sty.
% ===================================================================
\DocumentMetadata{lang=pt-BR, pdfstandard=ua-2, pdfversion=2.0, tagging=on}
\documentclass[11pt,a5paper,twoside,openright]{book}
\usepackage{livros_crp_acessivel_book}
\hypersetup{pdftitle={Sumario longo}, pdfauthor={CRP SP}}
\makeatletter
\providecommand\var{S}
\def\crp@vR{R}\def\crp@vB{B}\def\crp@vC{C}
\def\crp@lchap#1{%  #1 = acréscimo à margem direita (rightskip)
  \renewcommand*\l@chapter[2]{%
    \ifnum \c@tocdepth >\m@ne
        \addpenalty{-\@highpenalty}\addvspace{0.5em \@plus\p@}%
        \setlength\@tempdima{1.5em}%
        \begingroup
            \parindent \z@ \rightskip \@pnumwidth #1
            \parfillskip -\@pnumwidth
            \leavevmode \bfseries
            \advance\leftskip\@tempdima \hskip -\leftskip
            ##1\nobreak\hfil
            \nobreak\hb@xt@\@pnumwidth{\hss ##2\kern-\p@\kern\p@}\par
            \penalty\@highpenalty
        \endgroup
    \fi}}
\ifx\var\crp@vR \crp@lchap{}\fi
\ifx\var\crp@vB \crp@lchap{\@plus 4em}\fi
\ifx\var\crp@vC \crp@lchap{\@plus 4em \tolerance=10000 \hbadness=10000 }\fi
\makeatother
\begin{document}
\frontmatter
\tableofcontents
\mainmatter
\chapter{Deficiência psicossocial e}
\chapter{Deficiência psicossocial e transtorno}
\chapter{Deficiência psicossocial e transtorno do}
\chapter{Deficiência psicossocial e transtorno do espectro}
\chapter{Deficiência psicossocial e transtorno do espectro autista:}
\chapter{Deficiência psicossocial e transtorno do espectro autista: orientações}
\chapter{Deficiência psicossocial e transtorno do espectro autista: orientações para}
\chapter{Deficiência psicossocial e transtorno do espectro autista: orientações para a}
\chapter{Deficiência psicossocial e transtorno do espectro autista: orientações para a realização}
\chapter{Deficiência psicossocial e transtorno do espectro autista: orientações para a realização de}
\chapter{Deficiência psicossocial e transtorno do espectro autista: orientações para a realização de eventos}
\chapter{Deficiência psicossocial e transtorno do espectro autista: orientações para a realização de eventos e}
\chapter{Deficiência psicossocial e transtorno do espectro autista: orientações para a realização de eventos e apresentações}
\chapter{Deficiência psicossocial e transtorno do espectro autista: orientações para a realização de eventos e apresentações acessíveis}
\chapter{Deficiência psicossocial e transtorno do espectro autista: orientações para a realização de eventos e apresentações acessíveis em}
\chapter{Deficiência psicossocial e transtorno do espectro autista: orientações para a realização de eventos e apresentações acessíveis em todo}
\chapter{Comunicação alternativa e}
\chapter{Comunicação alternativa e aumentativa}
\chapter{Comunicação alternativa e aumentativa na}
\chapter{Comunicação alternativa e aumentativa na participação}
\chapter{Comunicação alternativa e aumentativa na participação social}
\chapter{Comunicação alternativa e aumentativa na participação social de}
\chapter{Comunicação alternativa e aumentativa na participação social de pessoas}
\chapter{Comunicação alternativa e aumentativa na participação social de pessoas com}
\chapter{Comunicação alternativa e aumentativa na participação social de pessoas com deficiência}
\chapter{Comunicação alternativa e aumentativa na participação social de pessoas com deficiência intelectual}
\chapter{Comunicação alternativa e aumentativa na participação social de pessoas com deficiência intelectual em}
\chapter{Comunicação alternativa e aumentativa na participação social de pessoas com deficiência intelectual em eventos}
\chapter{Comunicação alternativa e aumentativa na participação social de pessoas com deficiência intelectual em eventos públicos}
\chapter{Comunicação alternativa e aumentativa na participação social de pessoas com deficiência intelectual em eventos públicos de}
\chapter{Comunicação alternativa e aumentativa na participação social de pessoas com deficiência intelectual em eventos públicos de grande}
\chapter{Comunicação alternativa e aumentativa na participação social de pessoas com deficiência intelectual em eventos públicos de grande porte}
\chapter{Psicopedagogicamente inconstitucionalíssimo otorrinolaringologista}
\chapter{Psicopedagogicamente inconstitucionalíssimo otorrinolaringologista socioeducativo}
\chapter{Psicopedagogicamente inconstitucionalíssimo otorrinolaringologista socioeducativo ético-profissional}
\chapter{Psicopedagogicamente inconstitucionalíssimo otorrinolaringologista socioeducativo ético-profissional sociodemográfico}
\chapter{Psicopedagogicamente inconstitucionalíssimo otorrinolaringologista socioeducativo ético-profissional sociodemográfico interdisciplinaridade}
\chapter{Psicopedagogicamente inconstitucionalíssimo otorrinolaringologista socioeducativo ético-profissional sociodemográfico interdisciplinaridade neuropsicológica}
\chapter{Psicopedagogicamente inconstitucionalíssimo otorrinolaringologista socioeducativo ético-profissional sociodemográfico interdisciplinaridade neuropsicológica avaliação}
\chapter{Psicopedagogicamente inconstitucionalíssimo otorrinolaringologista socioeducativo ético-profissional sociodemográfico interdisciplinaridade neuropsicológica avaliação}
\chapter{Psicopedagogicamente inconstitucionalíssimo otorrinolaringologista socioeducativo ético-profissional sociodemográfico interdisciplinaridade neuropsicológica avaliação}
\chapter{Psicopedagogicamente inconstitucionalíssimo otorrinolaringologista socioeducativo ético-profissional sociodemográfico interdisciplinaridade neuropsicológica avaliação}
\chapter{Psicopedagogicamente inconstitucionalíssimo otorrinolaringologista socioeducativo ético-profissional sociodemográfico interdisciplinaridade neuropsicológica avaliação}
\chapter{Psicopedagogicamente inconstitucionalíssimo otorrinolaringologista socioeducativo ético-profissional sociodemográfico interdisciplinaridade neuropsicológica avaliação}
\chapter{Psicopedagogicamente inconstitucionalíssimo otorrinolaringologista socioeducativo ético-profissional sociodemográfico interdisciplinaridade neuropsicológica avaliação}
\chapter{Psicopedagogicamente inconstitucionalíssimo otorrinolaringologista socioeducativo ético-profissional sociodemográfico interdisciplinaridade neuropsicológica avaliação}
\chapter{Psicopedagogicamente inconstitucionalíssimo otorrinolaringologista socioeducativo ético-profissional sociodemográfico interdisciplinaridade neuropsicológica avaliação}
\chapter{Psicopedagogicamente inconstitucionalíssimo otorrinolaringologista socioeducativo ético-profissional sociodemográfico interdisciplinaridade neuropsicológica avaliação}
\end{document}
'   % S: o que a .sty define
% Cada variante só redefine o \l@chapter depois da carga da .sty.
% ===================================================================
\DocumentMetadata{lang=pt-BR, pdfstandard=ua-2, pdfversion=2.0, tagging=on}
\documentclass[11pt,a5paper,twoside,openright]{book}
\usepackage{livros_crp_acessivel_book}
\hypersetup{pdftitle={Sumario longo}, pdfauthor={CRP SP}}
\makeatletter
\providecommand\var{S}
\def\crp@vR{R}\def\crp@vB{B}\def\crp@vC{C}
\def\crp@lchap#1{%  #1 = acréscimo à margem direita (rightskip)
  \renewcommand*\l@chapter[2]{%
    \ifnum \c@tocdepth >\m@ne
        \addpenalty{-\@highpenalty}\addvspace{0.5em \@plus\p@}%
        \setlength\@tempdima{1.5em}%
        \begingroup
            \parindent \z@ \rightskip \@pnumwidth #1
            \parfillskip -\@pnumwidth
            \leavevmode \bfseries
            \advance\leftskip\@tempdima \hskip -\leftskip
            ##1\nobreak\hfil
            \nobreak\hb@xt@\@pnumwidth{\hss ##2\kern-\p@\kern\p@}\par
            \penalty\@highpenalty
        \endgroup
    \fi}}
\ifx\var\crp@vR \crp@lchap{}\fi
\ifx\var\crp@vB \crp@lchap{\@plus 4em}\fi
\ifx\var\crp@vC \crp@lchap{\@plus 4em \tolerance=10000 \hbadness=10000 }\fi
\makeatother
\begin{document}
\frontmatter
\tableofcontents
\mainmatter
\chapter{Deficiência psicossocial e}
\chapter{Deficiência psicossocial e transtorno}
\chapter{Deficiência psicossocial e transtorno do}
\chapter{Deficiência psicossocial e transtorno do espectro}
\chapter{Deficiência psicossocial e transtorno do espectro autista:}
\chapter{Deficiência psicossocial e transtorno do espectro autista: orientações}
\chapter{Deficiência psicossocial e transtorno do espectro autista: orientações para}
\chapter{Deficiência psicossocial e transtorno do espectro autista: orientações para a}
\chapter{Deficiência psicossocial e transtorno do espectro autista: orientações para a realização}
\chapter{Deficiência psicossocial e transtorno do espectro autista: orientações para a realização de}
\chapter{Deficiência psicossocial e transtorno do espectro autista: orientações para a realização de eventos}
\chapter{Deficiência psicossocial e transtorno do espectro autista: orientações para a realização de eventos e}
\chapter{Deficiência psicossocial e transtorno do espectro autista: orientações para a realização de eventos e apresentações}
\chapter{Deficiência psicossocial e transtorno do espectro autista: orientações para a realização de eventos e apresentações acessíveis}
\chapter{Deficiência psicossocial e transtorno do espectro autista: orientações para a realização de eventos e apresentações acessíveis em}
\chapter{Deficiência psicossocial e transtorno do espectro autista: orientações para a realização de eventos e apresentações acessíveis em todo}
\chapter{Comunicação alternativa e}
\chapter{Comunicação alternativa e aumentativa}
\chapter{Comunicação alternativa e aumentativa na}
\chapter{Comunicação alternativa e aumentativa na participação}
\chapter{Comunicação alternativa e aumentativa na participação social}
\chapter{Comunicação alternativa e aumentativa na participação social de}
\chapter{Comunicação alternativa e aumentativa na participação social de pessoas}
\chapter{Comunicação alternativa e aumentativa na participação social de pessoas com}
\chapter{Comunicação alternativa e aumentativa na participação social de pessoas com deficiência}
\chapter{Comunicação alternativa e aumentativa na participação social de pessoas com deficiência intelectual}
\chapter{Comunicação alternativa e aumentativa na participação social de pessoas com deficiência intelectual em}
\chapter{Comunicação alternativa e aumentativa na participação social de pessoas com deficiência intelectual em eventos}
\chapter{Comunicação alternativa e aumentativa na participação social de pessoas com deficiência intelectual em eventos públicos}
\chapter{Comunicação alternativa e aumentativa na participação social de pessoas com deficiência intelectual em eventos públicos de}
\chapter{Comunicação alternativa e aumentativa na participação social de pessoas com deficiência intelectual em eventos públicos de grande}
\chapter{Comunicação alternativa e aumentativa na participação social de pessoas com deficiência intelectual em eventos públicos de grande porte}
\chapter{Psicopedagogicamente inconstitucionalíssimo otorrinolaringologista}
\chapter{Psicopedagogicamente inconstitucionalíssimo otorrinolaringologista socioeducativo}
\chapter{Psicopedagogicamente inconstitucionalíssimo otorrinolaringologista socioeducativo ético-profissional}
\chapter{Psicopedagogicamente inconstitucionalíssimo otorrinolaringologista socioeducativo ético-profissional sociodemográfico}
\chapter{Psicopedagogicamente inconstitucionalíssimo otorrinolaringologista socioeducativo ético-profissional sociodemográfico interdisciplinaridade}
\chapter{Psicopedagogicamente inconstitucionalíssimo otorrinolaringologista socioeducativo ético-profissional sociodemográfico interdisciplinaridade neuropsicológica}
\chapter{Psicopedagogicamente inconstitucionalíssimo otorrinolaringologista socioeducativo ético-profissional sociodemográfico interdisciplinaridade neuropsicológica avaliação}
\chapter{Psicopedagogicamente inconstitucionalíssimo otorrinolaringologista socioeducativo ético-profissional sociodemográfico interdisciplinaridade neuropsicológica avaliação}
\chapter{Psicopedagogicamente inconstitucionalíssimo otorrinolaringologista socioeducativo ético-profissional sociodemográfico interdisciplinaridade neuropsicológica avaliação}
\chapter{Psicopedagogicamente inconstitucionalíssimo otorrinolaringologista socioeducativo ético-profissional sociodemográfico interdisciplinaridade neuropsicológica avaliação}
\chapter{Psicopedagogicamente inconstitucionalíssimo otorrinolaringologista socioeducativo ético-profissional sociodemográfico interdisciplinaridade neuropsicológica avaliação}
\chapter{Psicopedagogicamente inconstitucionalíssimo otorrinolaringologista socioeducativo ético-profissional sociodemográfico interdisciplinaridade neuropsicológica avaliação}
\chapter{Psicopedagogicamente inconstitucionalíssimo otorrinolaringologista socioeducativo ético-profissional sociodemográfico interdisciplinaridade neuropsicológica avaliação}
\chapter{Psicopedagogicamente inconstitucionalíssimo otorrinolaringologista socioeducativo ético-profissional sociodemográfico interdisciplinaridade neuropsicológica avaliação}
\chapter{Psicopedagogicamente inconstitucionalíssimo otorrinolaringologista socioeducativo ético-profissional sociodemográfico interdisciplinaridade neuropsicológica avaliação}
\chapter{Psicopedagogicamente inconstitucionalíssimo otorrinolaringologista socioeducativo ético-profissional sociodemográfico interdisciplinaridade neuropsicológica avaliação}
\end{document}
'   % B: plus 4em (o do guia)
%   lualatex '\def\var{C}% !TeX program = lualatex
% ===================================================================
% sumario-longo.tex — mede o sumário sem hifenização (issue #35, "Fora do
% escopo"): com a hifenização desligada, a entrada longa de capítulo pode
% bater no número da página. Gera 3 x 16 capítulos, cada um com um prefixo
% (de 3 a 18 palavras) de uma frase longa: frase 1 de palavras médias,
% frase 2 com números, frase 3 só de palavras muito longas. Conta no log
% os "Overfull \hbox" do sumário: o sumário é a única parte do documento
% sem parágrafo de corpo, então todo Overfull "in paragraph" antes do
% primeiro \chapter é entrada que passou da margem.
%
% Uso (de uma pasta com livros_crp_acessivel_book.sty e crpsp-leg.sty no
% caminho; as fontes NEWJUNE por OSFONTDIR):
%   lualatex '\def\var{R}% !TeX program = lualatex
% ===================================================================
% sumario-longo.tex — mede o sumário sem hifenização (issue #35, "Fora do
% escopo"): com a hifenização desligada, a entrada longa de capítulo pode
% bater no número da página. Gera 3 x 16 capítulos, cada um com um prefixo
% (de 3 a 18 palavras) de uma frase longa: frase 1 de palavras médias,
% frase 2 com números, frase 3 só de palavras muito longas. Conta no log
% os "Overfull \hbox" do sumário: o sumário é a única parte do documento
% sem parágrafo de corpo, então todo Overfull "in paragraph" antes do
% primeiro \chapter é entrada que passou da margem.
%
% Uso (de uma pasta com livros_crp_acessivel_book.sty e crpsp-leg.sty no
% caminho; as fontes NEWJUNE por OSFONTDIR):
%   lualatex '\def\var{R}\input{sumario-longo}'   % R: rígido (até 0.5.x)
%   lualatex '\def\var{B}\input{sumario-longo}'   % B: plus 4em (o do guia)
%   lualatex '\def\var{C}\input{sumario-longo}'   % C: plus 4em + \tolerance 10000
%   lualatex '\def\var{S}\input{sumario-longo}'   % S: o que a .sty define
% Cada variante só redefine o \l@chapter depois da carga da .sty.
% ===================================================================
\DocumentMetadata{lang=pt-BR, pdfstandard=ua-2, pdfversion=2.0, tagging=on}
\documentclass[11pt,a5paper,twoside,openright]{book}
\usepackage{livros_crp_acessivel_book}
\hypersetup{pdftitle={Sumario longo}, pdfauthor={CRP SP}}
\makeatletter
\providecommand\var{S}
\def\crp@vR{R}\def\crp@vB{B}\def\crp@vC{C}
\def\crp@lchap#1{%  #1 = acréscimo à margem direita (rightskip)
  \renewcommand*\l@chapter[2]{%
    \ifnum \c@tocdepth >\m@ne
        \addpenalty{-\@highpenalty}\addvspace{0.5em \@plus\p@}%
        \setlength\@tempdima{1.5em}%
        \begingroup
            \parindent \z@ \rightskip \@pnumwidth #1
            \parfillskip -\@pnumwidth
            \leavevmode \bfseries
            \advance\leftskip\@tempdima \hskip -\leftskip
            ##1\nobreak\hfil
            \nobreak\hb@xt@\@pnumwidth{\hss ##2\kern-\p@\kern\p@}\par
            \penalty\@highpenalty
        \endgroup
    \fi}}
\ifx\var\crp@vR \crp@lchap{}\fi
\ifx\var\crp@vB \crp@lchap{\@plus 4em}\fi
\ifx\var\crp@vC \crp@lchap{\@plus 4em \tolerance=10000 \hbadness=10000 }\fi
\makeatother
\begin{document}
\frontmatter
\tableofcontents
\mainmatter
\chapter{Deficiência psicossocial e}
\chapter{Deficiência psicossocial e transtorno}
\chapter{Deficiência psicossocial e transtorno do}
\chapter{Deficiência psicossocial e transtorno do espectro}
\chapter{Deficiência psicossocial e transtorno do espectro autista:}
\chapter{Deficiência psicossocial e transtorno do espectro autista: orientações}
\chapter{Deficiência psicossocial e transtorno do espectro autista: orientações para}
\chapter{Deficiência psicossocial e transtorno do espectro autista: orientações para a}
\chapter{Deficiência psicossocial e transtorno do espectro autista: orientações para a realização}
\chapter{Deficiência psicossocial e transtorno do espectro autista: orientações para a realização de}
\chapter{Deficiência psicossocial e transtorno do espectro autista: orientações para a realização de eventos}
\chapter{Deficiência psicossocial e transtorno do espectro autista: orientações para a realização de eventos e}
\chapter{Deficiência psicossocial e transtorno do espectro autista: orientações para a realização de eventos e apresentações}
\chapter{Deficiência psicossocial e transtorno do espectro autista: orientações para a realização de eventos e apresentações acessíveis}
\chapter{Deficiência psicossocial e transtorno do espectro autista: orientações para a realização de eventos e apresentações acessíveis em}
\chapter{Deficiência psicossocial e transtorno do espectro autista: orientações para a realização de eventos e apresentações acessíveis em todo}
\chapter{Comunicação alternativa e}
\chapter{Comunicação alternativa e aumentativa}
\chapter{Comunicação alternativa e aumentativa na}
\chapter{Comunicação alternativa e aumentativa na participação}
\chapter{Comunicação alternativa e aumentativa na participação social}
\chapter{Comunicação alternativa e aumentativa na participação social de}
\chapter{Comunicação alternativa e aumentativa na participação social de pessoas}
\chapter{Comunicação alternativa e aumentativa na participação social de pessoas com}
\chapter{Comunicação alternativa e aumentativa na participação social de pessoas com deficiência}
\chapter{Comunicação alternativa e aumentativa na participação social de pessoas com deficiência intelectual}
\chapter{Comunicação alternativa e aumentativa na participação social de pessoas com deficiência intelectual em}
\chapter{Comunicação alternativa e aumentativa na participação social de pessoas com deficiência intelectual em eventos}
\chapter{Comunicação alternativa e aumentativa na participação social de pessoas com deficiência intelectual em eventos públicos}
\chapter{Comunicação alternativa e aumentativa na participação social de pessoas com deficiência intelectual em eventos públicos de}
\chapter{Comunicação alternativa e aumentativa na participação social de pessoas com deficiência intelectual em eventos públicos de grande}
\chapter{Comunicação alternativa e aumentativa na participação social de pessoas com deficiência intelectual em eventos públicos de grande porte}
\chapter{Psicopedagogicamente inconstitucionalíssimo otorrinolaringologista}
\chapter{Psicopedagogicamente inconstitucionalíssimo otorrinolaringologista socioeducativo}
\chapter{Psicopedagogicamente inconstitucionalíssimo otorrinolaringologista socioeducativo ético-profissional}
\chapter{Psicopedagogicamente inconstitucionalíssimo otorrinolaringologista socioeducativo ético-profissional sociodemográfico}
\chapter{Psicopedagogicamente inconstitucionalíssimo otorrinolaringologista socioeducativo ético-profissional sociodemográfico interdisciplinaridade}
\chapter{Psicopedagogicamente inconstitucionalíssimo otorrinolaringologista socioeducativo ético-profissional sociodemográfico interdisciplinaridade neuropsicológica}
\chapter{Psicopedagogicamente inconstitucionalíssimo otorrinolaringologista socioeducativo ético-profissional sociodemográfico interdisciplinaridade neuropsicológica avaliação}
\chapter{Psicopedagogicamente inconstitucionalíssimo otorrinolaringologista socioeducativo ético-profissional sociodemográfico interdisciplinaridade neuropsicológica avaliação}
\chapter{Psicopedagogicamente inconstitucionalíssimo otorrinolaringologista socioeducativo ético-profissional sociodemográfico interdisciplinaridade neuropsicológica avaliação}
\chapter{Psicopedagogicamente inconstitucionalíssimo otorrinolaringologista socioeducativo ético-profissional sociodemográfico interdisciplinaridade neuropsicológica avaliação}
\chapter{Psicopedagogicamente inconstitucionalíssimo otorrinolaringologista socioeducativo ético-profissional sociodemográfico interdisciplinaridade neuropsicológica avaliação}
\chapter{Psicopedagogicamente inconstitucionalíssimo otorrinolaringologista socioeducativo ético-profissional sociodemográfico interdisciplinaridade neuropsicológica avaliação}
\chapter{Psicopedagogicamente inconstitucionalíssimo otorrinolaringologista socioeducativo ético-profissional sociodemográfico interdisciplinaridade neuropsicológica avaliação}
\chapter{Psicopedagogicamente inconstitucionalíssimo otorrinolaringologista socioeducativo ético-profissional sociodemográfico interdisciplinaridade neuropsicológica avaliação}
\chapter{Psicopedagogicamente inconstitucionalíssimo otorrinolaringologista socioeducativo ético-profissional sociodemográfico interdisciplinaridade neuropsicológica avaliação}
\chapter{Psicopedagogicamente inconstitucionalíssimo otorrinolaringologista socioeducativo ético-profissional sociodemográfico interdisciplinaridade neuropsicológica avaliação}
\end{document}
'   % R: rígido (até 0.5.x)
%   lualatex '\def\var{B}% !TeX program = lualatex
% ===================================================================
% sumario-longo.tex — mede o sumário sem hifenização (issue #35, "Fora do
% escopo"): com a hifenização desligada, a entrada longa de capítulo pode
% bater no número da página. Gera 3 x 16 capítulos, cada um com um prefixo
% (de 3 a 18 palavras) de uma frase longa: frase 1 de palavras médias,
% frase 2 com números, frase 3 só de palavras muito longas. Conta no log
% os "Overfull \hbox" do sumário: o sumário é a única parte do documento
% sem parágrafo de corpo, então todo Overfull "in paragraph" antes do
% primeiro \chapter é entrada que passou da margem.
%
% Uso (de uma pasta com livros_crp_acessivel_book.sty e crpsp-leg.sty no
% caminho; as fontes NEWJUNE por OSFONTDIR):
%   lualatex '\def\var{R}\input{sumario-longo}'   % R: rígido (até 0.5.x)
%   lualatex '\def\var{B}\input{sumario-longo}'   % B: plus 4em (o do guia)
%   lualatex '\def\var{C}\input{sumario-longo}'   % C: plus 4em + \tolerance 10000
%   lualatex '\def\var{S}\input{sumario-longo}'   % S: o que a .sty define
% Cada variante só redefine o \l@chapter depois da carga da .sty.
% ===================================================================
\DocumentMetadata{lang=pt-BR, pdfstandard=ua-2, pdfversion=2.0, tagging=on}
\documentclass[11pt,a5paper,twoside,openright]{book}
\usepackage{livros_crp_acessivel_book}
\hypersetup{pdftitle={Sumario longo}, pdfauthor={CRP SP}}
\makeatletter
\providecommand\var{S}
\def\crp@vR{R}\def\crp@vB{B}\def\crp@vC{C}
\def\crp@lchap#1{%  #1 = acréscimo à margem direita (rightskip)
  \renewcommand*\l@chapter[2]{%
    \ifnum \c@tocdepth >\m@ne
        \addpenalty{-\@highpenalty}\addvspace{0.5em \@plus\p@}%
        \setlength\@tempdima{1.5em}%
        \begingroup
            \parindent \z@ \rightskip \@pnumwidth #1
            \parfillskip -\@pnumwidth
            \leavevmode \bfseries
            \advance\leftskip\@tempdima \hskip -\leftskip
            ##1\nobreak\hfil
            \nobreak\hb@xt@\@pnumwidth{\hss ##2\kern-\p@\kern\p@}\par
            \penalty\@highpenalty
        \endgroup
    \fi}}
\ifx\var\crp@vR \crp@lchap{}\fi
\ifx\var\crp@vB \crp@lchap{\@plus 4em}\fi
\ifx\var\crp@vC \crp@lchap{\@plus 4em \tolerance=10000 \hbadness=10000 }\fi
\makeatother
\begin{document}
\frontmatter
\tableofcontents
\mainmatter
\chapter{Deficiência psicossocial e}
\chapter{Deficiência psicossocial e transtorno}
\chapter{Deficiência psicossocial e transtorno do}
\chapter{Deficiência psicossocial e transtorno do espectro}
\chapter{Deficiência psicossocial e transtorno do espectro autista:}
\chapter{Deficiência psicossocial e transtorno do espectro autista: orientações}
\chapter{Deficiência psicossocial e transtorno do espectro autista: orientações para}
\chapter{Deficiência psicossocial e transtorno do espectro autista: orientações para a}
\chapter{Deficiência psicossocial e transtorno do espectro autista: orientações para a realização}
\chapter{Deficiência psicossocial e transtorno do espectro autista: orientações para a realização de}
\chapter{Deficiência psicossocial e transtorno do espectro autista: orientações para a realização de eventos}
\chapter{Deficiência psicossocial e transtorno do espectro autista: orientações para a realização de eventos e}
\chapter{Deficiência psicossocial e transtorno do espectro autista: orientações para a realização de eventos e apresentações}
\chapter{Deficiência psicossocial e transtorno do espectro autista: orientações para a realização de eventos e apresentações acessíveis}
\chapter{Deficiência psicossocial e transtorno do espectro autista: orientações para a realização de eventos e apresentações acessíveis em}
\chapter{Deficiência psicossocial e transtorno do espectro autista: orientações para a realização de eventos e apresentações acessíveis em todo}
\chapter{Comunicação alternativa e}
\chapter{Comunicação alternativa e aumentativa}
\chapter{Comunicação alternativa e aumentativa na}
\chapter{Comunicação alternativa e aumentativa na participação}
\chapter{Comunicação alternativa e aumentativa na participação social}
\chapter{Comunicação alternativa e aumentativa na participação social de}
\chapter{Comunicação alternativa e aumentativa na participação social de pessoas}
\chapter{Comunicação alternativa e aumentativa na participação social de pessoas com}
\chapter{Comunicação alternativa e aumentativa na participação social de pessoas com deficiência}
\chapter{Comunicação alternativa e aumentativa na participação social de pessoas com deficiência intelectual}
\chapter{Comunicação alternativa e aumentativa na participação social de pessoas com deficiência intelectual em}
\chapter{Comunicação alternativa e aumentativa na participação social de pessoas com deficiência intelectual em eventos}
\chapter{Comunicação alternativa e aumentativa na participação social de pessoas com deficiência intelectual em eventos públicos}
\chapter{Comunicação alternativa e aumentativa na participação social de pessoas com deficiência intelectual em eventos públicos de}
\chapter{Comunicação alternativa e aumentativa na participação social de pessoas com deficiência intelectual em eventos públicos de grande}
\chapter{Comunicação alternativa e aumentativa na participação social de pessoas com deficiência intelectual em eventos públicos de grande porte}
\chapter{Psicopedagogicamente inconstitucionalíssimo otorrinolaringologista}
\chapter{Psicopedagogicamente inconstitucionalíssimo otorrinolaringologista socioeducativo}
\chapter{Psicopedagogicamente inconstitucionalíssimo otorrinolaringologista socioeducativo ético-profissional}
\chapter{Psicopedagogicamente inconstitucionalíssimo otorrinolaringologista socioeducativo ético-profissional sociodemográfico}
\chapter{Psicopedagogicamente inconstitucionalíssimo otorrinolaringologista socioeducativo ético-profissional sociodemográfico interdisciplinaridade}
\chapter{Psicopedagogicamente inconstitucionalíssimo otorrinolaringologista socioeducativo ético-profissional sociodemográfico interdisciplinaridade neuropsicológica}
\chapter{Psicopedagogicamente inconstitucionalíssimo otorrinolaringologista socioeducativo ético-profissional sociodemográfico interdisciplinaridade neuropsicológica avaliação}
\chapter{Psicopedagogicamente inconstitucionalíssimo otorrinolaringologista socioeducativo ético-profissional sociodemográfico interdisciplinaridade neuropsicológica avaliação}
\chapter{Psicopedagogicamente inconstitucionalíssimo otorrinolaringologista socioeducativo ético-profissional sociodemográfico interdisciplinaridade neuropsicológica avaliação}
\chapter{Psicopedagogicamente inconstitucionalíssimo otorrinolaringologista socioeducativo ético-profissional sociodemográfico interdisciplinaridade neuropsicológica avaliação}
\chapter{Psicopedagogicamente inconstitucionalíssimo otorrinolaringologista socioeducativo ético-profissional sociodemográfico interdisciplinaridade neuropsicológica avaliação}
\chapter{Psicopedagogicamente inconstitucionalíssimo otorrinolaringologista socioeducativo ético-profissional sociodemográfico interdisciplinaridade neuropsicológica avaliação}
\chapter{Psicopedagogicamente inconstitucionalíssimo otorrinolaringologista socioeducativo ético-profissional sociodemográfico interdisciplinaridade neuropsicológica avaliação}
\chapter{Psicopedagogicamente inconstitucionalíssimo otorrinolaringologista socioeducativo ético-profissional sociodemográfico interdisciplinaridade neuropsicológica avaliação}
\chapter{Psicopedagogicamente inconstitucionalíssimo otorrinolaringologista socioeducativo ético-profissional sociodemográfico interdisciplinaridade neuropsicológica avaliação}
\chapter{Psicopedagogicamente inconstitucionalíssimo otorrinolaringologista socioeducativo ético-profissional sociodemográfico interdisciplinaridade neuropsicológica avaliação}
\end{document}
'   % B: plus 4em (o do guia)
%   lualatex '\def\var{C}% !TeX program = lualatex
% ===================================================================
% sumario-longo.tex — mede o sumário sem hifenização (issue #35, "Fora do
% escopo"): com a hifenização desligada, a entrada longa de capítulo pode
% bater no número da página. Gera 3 x 16 capítulos, cada um com um prefixo
% (de 3 a 18 palavras) de uma frase longa: frase 1 de palavras médias,
% frase 2 com números, frase 3 só de palavras muito longas. Conta no log
% os "Overfull \hbox" do sumário: o sumário é a única parte do documento
% sem parágrafo de corpo, então todo Overfull "in paragraph" antes do
% primeiro \chapter é entrada que passou da margem.
%
% Uso (de uma pasta com livros_crp_acessivel_book.sty e crpsp-leg.sty no
% caminho; as fontes NEWJUNE por OSFONTDIR):
%   lualatex '\def\var{R}\input{sumario-longo}'   % R: rígido (até 0.5.x)
%   lualatex '\def\var{B}\input{sumario-longo}'   % B: plus 4em (o do guia)
%   lualatex '\def\var{C}\input{sumario-longo}'   % C: plus 4em + \tolerance 10000
%   lualatex '\def\var{S}\input{sumario-longo}'   % S: o que a .sty define
% Cada variante só redefine o \l@chapter depois da carga da .sty.
% ===================================================================
\DocumentMetadata{lang=pt-BR, pdfstandard=ua-2, pdfversion=2.0, tagging=on}
\documentclass[11pt,a5paper,twoside,openright]{book}
\usepackage{livros_crp_acessivel_book}
\hypersetup{pdftitle={Sumario longo}, pdfauthor={CRP SP}}
\makeatletter
\providecommand\var{S}
\def\crp@vR{R}\def\crp@vB{B}\def\crp@vC{C}
\def\crp@lchap#1{%  #1 = acréscimo à margem direita (rightskip)
  \renewcommand*\l@chapter[2]{%
    \ifnum \c@tocdepth >\m@ne
        \addpenalty{-\@highpenalty}\addvspace{0.5em \@plus\p@}%
        \setlength\@tempdima{1.5em}%
        \begingroup
            \parindent \z@ \rightskip \@pnumwidth #1
            \parfillskip -\@pnumwidth
            \leavevmode \bfseries
            \advance\leftskip\@tempdima \hskip -\leftskip
            ##1\nobreak\hfil
            \nobreak\hb@xt@\@pnumwidth{\hss ##2\kern-\p@\kern\p@}\par
            \penalty\@highpenalty
        \endgroup
    \fi}}
\ifx\var\crp@vR \crp@lchap{}\fi
\ifx\var\crp@vB \crp@lchap{\@plus 4em}\fi
\ifx\var\crp@vC \crp@lchap{\@plus 4em \tolerance=10000 \hbadness=10000 }\fi
\makeatother
\begin{document}
\frontmatter
\tableofcontents
\mainmatter
\chapter{Deficiência psicossocial e}
\chapter{Deficiência psicossocial e transtorno}
\chapter{Deficiência psicossocial e transtorno do}
\chapter{Deficiência psicossocial e transtorno do espectro}
\chapter{Deficiência psicossocial e transtorno do espectro autista:}
\chapter{Deficiência psicossocial e transtorno do espectro autista: orientações}
\chapter{Deficiência psicossocial e transtorno do espectro autista: orientações para}
\chapter{Deficiência psicossocial e transtorno do espectro autista: orientações para a}
\chapter{Deficiência psicossocial e transtorno do espectro autista: orientações para a realização}
\chapter{Deficiência psicossocial e transtorno do espectro autista: orientações para a realização de}
\chapter{Deficiência psicossocial e transtorno do espectro autista: orientações para a realização de eventos}
\chapter{Deficiência psicossocial e transtorno do espectro autista: orientações para a realização de eventos e}
\chapter{Deficiência psicossocial e transtorno do espectro autista: orientações para a realização de eventos e apresentações}
\chapter{Deficiência psicossocial e transtorno do espectro autista: orientações para a realização de eventos e apresentações acessíveis}
\chapter{Deficiência psicossocial e transtorno do espectro autista: orientações para a realização de eventos e apresentações acessíveis em}
\chapter{Deficiência psicossocial e transtorno do espectro autista: orientações para a realização de eventos e apresentações acessíveis em todo}
\chapter{Comunicação alternativa e}
\chapter{Comunicação alternativa e aumentativa}
\chapter{Comunicação alternativa e aumentativa na}
\chapter{Comunicação alternativa e aumentativa na participação}
\chapter{Comunicação alternativa e aumentativa na participação social}
\chapter{Comunicação alternativa e aumentativa na participação social de}
\chapter{Comunicação alternativa e aumentativa na participação social de pessoas}
\chapter{Comunicação alternativa e aumentativa na participação social de pessoas com}
\chapter{Comunicação alternativa e aumentativa na participação social de pessoas com deficiência}
\chapter{Comunicação alternativa e aumentativa na participação social de pessoas com deficiência intelectual}
\chapter{Comunicação alternativa e aumentativa na participação social de pessoas com deficiência intelectual em}
\chapter{Comunicação alternativa e aumentativa na participação social de pessoas com deficiência intelectual em eventos}
\chapter{Comunicação alternativa e aumentativa na participação social de pessoas com deficiência intelectual em eventos públicos}
\chapter{Comunicação alternativa e aumentativa na participação social de pessoas com deficiência intelectual em eventos públicos de}
\chapter{Comunicação alternativa e aumentativa na participação social de pessoas com deficiência intelectual em eventos públicos de grande}
\chapter{Comunicação alternativa e aumentativa na participação social de pessoas com deficiência intelectual em eventos públicos de grande porte}
\chapter{Psicopedagogicamente inconstitucionalíssimo otorrinolaringologista}
\chapter{Psicopedagogicamente inconstitucionalíssimo otorrinolaringologista socioeducativo}
\chapter{Psicopedagogicamente inconstitucionalíssimo otorrinolaringologista socioeducativo ético-profissional}
\chapter{Psicopedagogicamente inconstitucionalíssimo otorrinolaringologista socioeducativo ético-profissional sociodemográfico}
\chapter{Psicopedagogicamente inconstitucionalíssimo otorrinolaringologista socioeducativo ético-profissional sociodemográfico interdisciplinaridade}
\chapter{Psicopedagogicamente inconstitucionalíssimo otorrinolaringologista socioeducativo ético-profissional sociodemográfico interdisciplinaridade neuropsicológica}
\chapter{Psicopedagogicamente inconstitucionalíssimo otorrinolaringologista socioeducativo ético-profissional sociodemográfico interdisciplinaridade neuropsicológica avaliação}
\chapter{Psicopedagogicamente inconstitucionalíssimo otorrinolaringologista socioeducativo ético-profissional sociodemográfico interdisciplinaridade neuropsicológica avaliação}
\chapter{Psicopedagogicamente inconstitucionalíssimo otorrinolaringologista socioeducativo ético-profissional sociodemográfico interdisciplinaridade neuropsicológica avaliação}
\chapter{Psicopedagogicamente inconstitucionalíssimo otorrinolaringologista socioeducativo ético-profissional sociodemográfico interdisciplinaridade neuropsicológica avaliação}
\chapter{Psicopedagogicamente inconstitucionalíssimo otorrinolaringologista socioeducativo ético-profissional sociodemográfico interdisciplinaridade neuropsicológica avaliação}
\chapter{Psicopedagogicamente inconstitucionalíssimo otorrinolaringologista socioeducativo ético-profissional sociodemográfico interdisciplinaridade neuropsicológica avaliação}
\chapter{Psicopedagogicamente inconstitucionalíssimo otorrinolaringologista socioeducativo ético-profissional sociodemográfico interdisciplinaridade neuropsicológica avaliação}
\chapter{Psicopedagogicamente inconstitucionalíssimo otorrinolaringologista socioeducativo ético-profissional sociodemográfico interdisciplinaridade neuropsicológica avaliação}
\chapter{Psicopedagogicamente inconstitucionalíssimo otorrinolaringologista socioeducativo ético-profissional sociodemográfico interdisciplinaridade neuropsicológica avaliação}
\chapter{Psicopedagogicamente inconstitucionalíssimo otorrinolaringologista socioeducativo ético-profissional sociodemográfico interdisciplinaridade neuropsicológica avaliação}
\end{document}
'   % C: plus 4em + \tolerance 10000
%   lualatex '\def\var{S}% !TeX program = lualatex
% ===================================================================
% sumario-longo.tex — mede o sumário sem hifenização (issue #35, "Fora do
% escopo"): com a hifenização desligada, a entrada longa de capítulo pode
% bater no número da página. Gera 3 x 16 capítulos, cada um com um prefixo
% (de 3 a 18 palavras) de uma frase longa: frase 1 de palavras médias,
% frase 2 com números, frase 3 só de palavras muito longas. Conta no log
% os "Overfull \hbox" do sumário: o sumário é a única parte do documento
% sem parágrafo de corpo, então todo Overfull "in paragraph" antes do
% primeiro \chapter é entrada que passou da margem.
%
% Uso (de uma pasta com livros_crp_acessivel_book.sty e crpsp-leg.sty no
% caminho; as fontes NEWJUNE por OSFONTDIR):
%   lualatex '\def\var{R}\input{sumario-longo}'   % R: rígido (até 0.5.x)
%   lualatex '\def\var{B}\input{sumario-longo}'   % B: plus 4em (o do guia)
%   lualatex '\def\var{C}\input{sumario-longo}'   % C: plus 4em + \tolerance 10000
%   lualatex '\def\var{S}\input{sumario-longo}'   % S: o que a .sty define
% Cada variante só redefine o \l@chapter depois da carga da .sty.
% ===================================================================
\DocumentMetadata{lang=pt-BR, pdfstandard=ua-2, pdfversion=2.0, tagging=on}
\documentclass[11pt,a5paper,twoside,openright]{book}
\usepackage{livros_crp_acessivel_book}
\hypersetup{pdftitle={Sumario longo}, pdfauthor={CRP SP}}
\makeatletter
\providecommand\var{S}
\def\crp@vR{R}\def\crp@vB{B}\def\crp@vC{C}
\def\crp@lchap#1{%  #1 = acréscimo à margem direita (rightskip)
  \renewcommand*\l@chapter[2]{%
    \ifnum \c@tocdepth >\m@ne
        \addpenalty{-\@highpenalty}\addvspace{0.5em \@plus\p@}%
        \setlength\@tempdima{1.5em}%
        \begingroup
            \parindent \z@ \rightskip \@pnumwidth #1
            \parfillskip -\@pnumwidth
            \leavevmode \bfseries
            \advance\leftskip\@tempdima \hskip -\leftskip
            ##1\nobreak\hfil
            \nobreak\hb@xt@\@pnumwidth{\hss ##2\kern-\p@\kern\p@}\par
            \penalty\@highpenalty
        \endgroup
    \fi}}
\ifx\var\crp@vR \crp@lchap{}\fi
\ifx\var\crp@vB \crp@lchap{\@plus 4em}\fi
\ifx\var\crp@vC \crp@lchap{\@plus 4em \tolerance=10000 \hbadness=10000 }\fi
\makeatother
\begin{document}
\frontmatter
\tableofcontents
\mainmatter
\chapter{Deficiência psicossocial e}
\chapter{Deficiência psicossocial e transtorno}
\chapter{Deficiência psicossocial e transtorno do}
\chapter{Deficiência psicossocial e transtorno do espectro}
\chapter{Deficiência psicossocial e transtorno do espectro autista:}
\chapter{Deficiência psicossocial e transtorno do espectro autista: orientações}
\chapter{Deficiência psicossocial e transtorno do espectro autista: orientações para}
\chapter{Deficiência psicossocial e transtorno do espectro autista: orientações para a}
\chapter{Deficiência psicossocial e transtorno do espectro autista: orientações para a realização}
\chapter{Deficiência psicossocial e transtorno do espectro autista: orientações para a realização de}
\chapter{Deficiência psicossocial e transtorno do espectro autista: orientações para a realização de eventos}
\chapter{Deficiência psicossocial e transtorno do espectro autista: orientações para a realização de eventos e}
\chapter{Deficiência psicossocial e transtorno do espectro autista: orientações para a realização de eventos e apresentações}
\chapter{Deficiência psicossocial e transtorno do espectro autista: orientações para a realização de eventos e apresentações acessíveis}
\chapter{Deficiência psicossocial e transtorno do espectro autista: orientações para a realização de eventos e apresentações acessíveis em}
\chapter{Deficiência psicossocial e transtorno do espectro autista: orientações para a realização de eventos e apresentações acessíveis em todo}
\chapter{Comunicação alternativa e}
\chapter{Comunicação alternativa e aumentativa}
\chapter{Comunicação alternativa e aumentativa na}
\chapter{Comunicação alternativa e aumentativa na participação}
\chapter{Comunicação alternativa e aumentativa na participação social}
\chapter{Comunicação alternativa e aumentativa na participação social de}
\chapter{Comunicação alternativa e aumentativa na participação social de pessoas}
\chapter{Comunicação alternativa e aumentativa na participação social de pessoas com}
\chapter{Comunicação alternativa e aumentativa na participação social de pessoas com deficiência}
\chapter{Comunicação alternativa e aumentativa na participação social de pessoas com deficiência intelectual}
\chapter{Comunicação alternativa e aumentativa na participação social de pessoas com deficiência intelectual em}
\chapter{Comunicação alternativa e aumentativa na participação social de pessoas com deficiência intelectual em eventos}
\chapter{Comunicação alternativa e aumentativa na participação social de pessoas com deficiência intelectual em eventos públicos}
\chapter{Comunicação alternativa e aumentativa na participação social de pessoas com deficiência intelectual em eventos públicos de}
\chapter{Comunicação alternativa e aumentativa na participação social de pessoas com deficiência intelectual em eventos públicos de grande}
\chapter{Comunicação alternativa e aumentativa na participação social de pessoas com deficiência intelectual em eventos públicos de grande porte}
\chapter{Psicopedagogicamente inconstitucionalíssimo otorrinolaringologista}
\chapter{Psicopedagogicamente inconstitucionalíssimo otorrinolaringologista socioeducativo}
\chapter{Psicopedagogicamente inconstitucionalíssimo otorrinolaringologista socioeducativo ético-profissional}
\chapter{Psicopedagogicamente inconstitucionalíssimo otorrinolaringologista socioeducativo ético-profissional sociodemográfico}
\chapter{Psicopedagogicamente inconstitucionalíssimo otorrinolaringologista socioeducativo ético-profissional sociodemográfico interdisciplinaridade}
\chapter{Psicopedagogicamente inconstitucionalíssimo otorrinolaringologista socioeducativo ético-profissional sociodemográfico interdisciplinaridade neuropsicológica}
\chapter{Psicopedagogicamente inconstitucionalíssimo otorrinolaringologista socioeducativo ético-profissional sociodemográfico interdisciplinaridade neuropsicológica avaliação}
\chapter{Psicopedagogicamente inconstitucionalíssimo otorrinolaringologista socioeducativo ético-profissional sociodemográfico interdisciplinaridade neuropsicológica avaliação}
\chapter{Psicopedagogicamente inconstitucionalíssimo otorrinolaringologista socioeducativo ético-profissional sociodemográfico interdisciplinaridade neuropsicológica avaliação}
\chapter{Psicopedagogicamente inconstitucionalíssimo otorrinolaringologista socioeducativo ético-profissional sociodemográfico interdisciplinaridade neuropsicológica avaliação}
\chapter{Psicopedagogicamente inconstitucionalíssimo otorrinolaringologista socioeducativo ético-profissional sociodemográfico interdisciplinaridade neuropsicológica avaliação}
\chapter{Psicopedagogicamente inconstitucionalíssimo otorrinolaringologista socioeducativo ético-profissional sociodemográfico interdisciplinaridade neuropsicológica avaliação}
\chapter{Psicopedagogicamente inconstitucionalíssimo otorrinolaringologista socioeducativo ético-profissional sociodemográfico interdisciplinaridade neuropsicológica avaliação}
\chapter{Psicopedagogicamente inconstitucionalíssimo otorrinolaringologista socioeducativo ético-profissional sociodemográfico interdisciplinaridade neuropsicológica avaliação}
\chapter{Psicopedagogicamente inconstitucionalíssimo otorrinolaringologista socioeducativo ético-profissional sociodemográfico interdisciplinaridade neuropsicológica avaliação}
\chapter{Psicopedagogicamente inconstitucionalíssimo otorrinolaringologista socioeducativo ético-profissional sociodemográfico interdisciplinaridade neuropsicológica avaliação}
\end{document}
'   % S: o que a .sty define
% Cada variante só redefine o \l@chapter depois da carga da .sty.
% ===================================================================
\DocumentMetadata{lang=pt-BR, pdfstandard=ua-2, pdfversion=2.0, tagging=on}
\documentclass[11pt,a5paper,twoside,openright]{book}
\usepackage{livros_crp_acessivel_book}
\hypersetup{pdftitle={Sumario longo}, pdfauthor={CRP SP}}
\makeatletter
\providecommand\var{S}
\def\crp@vR{R}\def\crp@vB{B}\def\crp@vC{C}
\def\crp@lchap#1{%  #1 = acréscimo à margem direita (rightskip)
  \renewcommand*\l@chapter[2]{%
    \ifnum \c@tocdepth >\m@ne
        \addpenalty{-\@highpenalty}\addvspace{0.5em \@plus\p@}%
        \setlength\@tempdima{1.5em}%
        \begingroup
            \parindent \z@ \rightskip \@pnumwidth #1
            \parfillskip -\@pnumwidth
            \leavevmode \bfseries
            \advance\leftskip\@tempdima \hskip -\leftskip
            ##1\nobreak\hfil
            \nobreak\hb@xt@\@pnumwidth{\hss ##2\kern-\p@\kern\p@}\par
            \penalty\@highpenalty
        \endgroup
    \fi}}
\ifx\var\crp@vR \crp@lchap{}\fi
\ifx\var\crp@vB \crp@lchap{\@plus 4em}\fi
\ifx\var\crp@vC \crp@lchap{\@plus 4em \tolerance=10000 \hbadness=10000 }\fi
\makeatother
\begin{document}
\frontmatter
\tableofcontents
\mainmatter
\chapter{Deficiência psicossocial e}
\chapter{Deficiência psicossocial e transtorno}
\chapter{Deficiência psicossocial e transtorno do}
\chapter{Deficiência psicossocial e transtorno do espectro}
\chapter{Deficiência psicossocial e transtorno do espectro autista:}
\chapter{Deficiência psicossocial e transtorno do espectro autista: orientações}
\chapter{Deficiência psicossocial e transtorno do espectro autista: orientações para}
\chapter{Deficiência psicossocial e transtorno do espectro autista: orientações para a}
\chapter{Deficiência psicossocial e transtorno do espectro autista: orientações para a realização}
\chapter{Deficiência psicossocial e transtorno do espectro autista: orientações para a realização de}
\chapter{Deficiência psicossocial e transtorno do espectro autista: orientações para a realização de eventos}
\chapter{Deficiência psicossocial e transtorno do espectro autista: orientações para a realização de eventos e}
\chapter{Deficiência psicossocial e transtorno do espectro autista: orientações para a realização de eventos e apresentações}
\chapter{Deficiência psicossocial e transtorno do espectro autista: orientações para a realização de eventos e apresentações acessíveis}
\chapter{Deficiência psicossocial e transtorno do espectro autista: orientações para a realização de eventos e apresentações acessíveis em}
\chapter{Deficiência psicossocial e transtorno do espectro autista: orientações para a realização de eventos e apresentações acessíveis em todo}
\chapter{Comunicação alternativa e}
\chapter{Comunicação alternativa e aumentativa}
\chapter{Comunicação alternativa e aumentativa na}
\chapter{Comunicação alternativa e aumentativa na participação}
\chapter{Comunicação alternativa e aumentativa na participação social}
\chapter{Comunicação alternativa e aumentativa na participação social de}
\chapter{Comunicação alternativa e aumentativa na participação social de pessoas}
\chapter{Comunicação alternativa e aumentativa na participação social de pessoas com}
\chapter{Comunicação alternativa e aumentativa na participação social de pessoas com deficiência}
\chapter{Comunicação alternativa e aumentativa na participação social de pessoas com deficiência intelectual}
\chapter{Comunicação alternativa e aumentativa na participação social de pessoas com deficiência intelectual em}
\chapter{Comunicação alternativa e aumentativa na participação social de pessoas com deficiência intelectual em eventos}
\chapter{Comunicação alternativa e aumentativa na participação social de pessoas com deficiência intelectual em eventos públicos}
\chapter{Comunicação alternativa e aumentativa na participação social de pessoas com deficiência intelectual em eventos públicos de}
\chapter{Comunicação alternativa e aumentativa na participação social de pessoas com deficiência intelectual em eventos públicos de grande}
\chapter{Comunicação alternativa e aumentativa na participação social de pessoas com deficiência intelectual em eventos públicos de grande porte}
\chapter{Psicopedagogicamente inconstitucionalíssimo otorrinolaringologista}
\chapter{Psicopedagogicamente inconstitucionalíssimo otorrinolaringologista socioeducativo}
\chapter{Psicopedagogicamente inconstitucionalíssimo otorrinolaringologista socioeducativo ético-profissional}
\chapter{Psicopedagogicamente inconstitucionalíssimo otorrinolaringologista socioeducativo ético-profissional sociodemográfico}
\chapter{Psicopedagogicamente inconstitucionalíssimo otorrinolaringologista socioeducativo ético-profissional sociodemográfico interdisciplinaridade}
\chapter{Psicopedagogicamente inconstitucionalíssimo otorrinolaringologista socioeducativo ético-profissional sociodemográfico interdisciplinaridade neuropsicológica}
\chapter{Psicopedagogicamente inconstitucionalíssimo otorrinolaringologista socioeducativo ético-profissional sociodemográfico interdisciplinaridade neuropsicológica avaliação}
\chapter{Psicopedagogicamente inconstitucionalíssimo otorrinolaringologista socioeducativo ético-profissional sociodemográfico interdisciplinaridade neuropsicológica avaliação}
\chapter{Psicopedagogicamente inconstitucionalíssimo otorrinolaringologista socioeducativo ético-profissional sociodemográfico interdisciplinaridade neuropsicológica avaliação}
\chapter{Psicopedagogicamente inconstitucionalíssimo otorrinolaringologista socioeducativo ético-profissional sociodemográfico interdisciplinaridade neuropsicológica avaliação}
\chapter{Psicopedagogicamente inconstitucionalíssimo otorrinolaringologista socioeducativo ético-profissional sociodemográfico interdisciplinaridade neuropsicológica avaliação}
\chapter{Psicopedagogicamente inconstitucionalíssimo otorrinolaringologista socioeducativo ético-profissional sociodemográfico interdisciplinaridade neuropsicológica avaliação}
\chapter{Psicopedagogicamente inconstitucionalíssimo otorrinolaringologista socioeducativo ético-profissional sociodemográfico interdisciplinaridade neuropsicológica avaliação}
\chapter{Psicopedagogicamente inconstitucionalíssimo otorrinolaringologista socioeducativo ético-profissional sociodemográfico interdisciplinaridade neuropsicológica avaliação}
\chapter{Psicopedagogicamente inconstitucionalíssimo otorrinolaringologista socioeducativo ético-profissional sociodemográfico interdisciplinaridade neuropsicológica avaliação}
\chapter{Psicopedagogicamente inconstitucionalíssimo otorrinolaringologista socioeducativo ético-profissional sociodemográfico interdisciplinaridade neuropsicológica avaliação}
\end{document}
'   % C: plus 4em + \tolerance 10000
%   lualatex '\def\var{S}% !TeX program = lualatex
% ===================================================================
% sumario-longo.tex — mede o sumário sem hifenização (issue #35, "Fora do
% escopo"): com a hifenização desligada, a entrada longa de capítulo pode
% bater no número da página. Gera 3 x 16 capítulos, cada um com um prefixo
% (de 3 a 18 palavras) de uma frase longa: frase 1 de palavras médias,
% frase 2 com números, frase 3 só de palavras muito longas. Conta no log
% os "Overfull \hbox" do sumário: o sumário é a única parte do documento
% sem parágrafo de corpo, então todo Overfull "in paragraph" antes do
% primeiro \chapter é entrada que passou da margem.
%
% Uso (de uma pasta com livros_crp_acessivel_book.sty e crpsp-leg.sty no
% caminho; as fontes NEWJUNE por OSFONTDIR):
%   lualatex '\def\var{R}% !TeX program = lualatex
% ===================================================================
% sumario-longo.tex — mede o sumário sem hifenização (issue #35, "Fora do
% escopo"): com a hifenização desligada, a entrada longa de capítulo pode
% bater no número da página. Gera 3 x 16 capítulos, cada um com um prefixo
% (de 3 a 18 palavras) de uma frase longa: frase 1 de palavras médias,
% frase 2 com números, frase 3 só de palavras muito longas. Conta no log
% os "Overfull \hbox" do sumário: o sumário é a única parte do documento
% sem parágrafo de corpo, então todo Overfull "in paragraph" antes do
% primeiro \chapter é entrada que passou da margem.
%
% Uso (de uma pasta com livros_crp_acessivel_book.sty e crpsp-leg.sty no
% caminho; as fontes NEWJUNE por OSFONTDIR):
%   lualatex '\def\var{R}\input{sumario-longo}'   % R: rígido (até 0.5.x)
%   lualatex '\def\var{B}\input{sumario-longo}'   % B: plus 4em (o do guia)
%   lualatex '\def\var{C}\input{sumario-longo}'   % C: plus 4em + \tolerance 10000
%   lualatex '\def\var{S}\input{sumario-longo}'   % S: o que a .sty define
% Cada variante só redefine o \l@chapter depois da carga da .sty.
% ===================================================================
\DocumentMetadata{lang=pt-BR, pdfstandard=ua-2, pdfversion=2.0, tagging=on}
\documentclass[11pt,a5paper,twoside,openright]{book}
\usepackage{livros_crp_acessivel_book}
\hypersetup{pdftitle={Sumario longo}, pdfauthor={CRP SP}}
\makeatletter
\providecommand\var{S}
\def\crp@vR{R}\def\crp@vB{B}\def\crp@vC{C}
\def\crp@lchap#1{%  #1 = acréscimo à margem direita (rightskip)
  \renewcommand*\l@chapter[2]{%
    \ifnum \c@tocdepth >\m@ne
        \addpenalty{-\@highpenalty}\addvspace{0.5em \@plus\p@}%
        \setlength\@tempdima{1.5em}%
        \begingroup
            \parindent \z@ \rightskip \@pnumwidth #1
            \parfillskip -\@pnumwidth
            \leavevmode \bfseries
            \advance\leftskip\@tempdima \hskip -\leftskip
            ##1\nobreak\hfil
            \nobreak\hb@xt@\@pnumwidth{\hss ##2\kern-\p@\kern\p@}\par
            \penalty\@highpenalty
        \endgroup
    \fi}}
\ifx\var\crp@vR \crp@lchap{}\fi
\ifx\var\crp@vB \crp@lchap{\@plus 4em}\fi
\ifx\var\crp@vC \crp@lchap{\@plus 4em \tolerance=10000 \hbadness=10000 }\fi
\makeatother
\begin{document}
\frontmatter
\tableofcontents
\mainmatter
\chapter{Deficiência psicossocial e}
\chapter{Deficiência psicossocial e transtorno}
\chapter{Deficiência psicossocial e transtorno do}
\chapter{Deficiência psicossocial e transtorno do espectro}
\chapter{Deficiência psicossocial e transtorno do espectro autista:}
\chapter{Deficiência psicossocial e transtorno do espectro autista: orientações}
\chapter{Deficiência psicossocial e transtorno do espectro autista: orientações para}
\chapter{Deficiência psicossocial e transtorno do espectro autista: orientações para a}
\chapter{Deficiência psicossocial e transtorno do espectro autista: orientações para a realização}
\chapter{Deficiência psicossocial e transtorno do espectro autista: orientações para a realização de}
\chapter{Deficiência psicossocial e transtorno do espectro autista: orientações para a realização de eventos}
\chapter{Deficiência psicossocial e transtorno do espectro autista: orientações para a realização de eventos e}
\chapter{Deficiência psicossocial e transtorno do espectro autista: orientações para a realização de eventos e apresentações}
\chapter{Deficiência psicossocial e transtorno do espectro autista: orientações para a realização de eventos e apresentações acessíveis}
\chapter{Deficiência psicossocial e transtorno do espectro autista: orientações para a realização de eventos e apresentações acessíveis em}
\chapter{Deficiência psicossocial e transtorno do espectro autista: orientações para a realização de eventos e apresentações acessíveis em todo}
\chapter{Comunicação alternativa e}
\chapter{Comunicação alternativa e aumentativa}
\chapter{Comunicação alternativa e aumentativa na}
\chapter{Comunicação alternativa e aumentativa na participação}
\chapter{Comunicação alternativa e aumentativa na participação social}
\chapter{Comunicação alternativa e aumentativa na participação social de}
\chapter{Comunicação alternativa e aumentativa na participação social de pessoas}
\chapter{Comunicação alternativa e aumentativa na participação social de pessoas com}
\chapter{Comunicação alternativa e aumentativa na participação social de pessoas com deficiência}
\chapter{Comunicação alternativa e aumentativa na participação social de pessoas com deficiência intelectual}
\chapter{Comunicação alternativa e aumentativa na participação social de pessoas com deficiência intelectual em}
\chapter{Comunicação alternativa e aumentativa na participação social de pessoas com deficiência intelectual em eventos}
\chapter{Comunicação alternativa e aumentativa na participação social de pessoas com deficiência intelectual em eventos públicos}
\chapter{Comunicação alternativa e aumentativa na participação social de pessoas com deficiência intelectual em eventos públicos de}
\chapter{Comunicação alternativa e aumentativa na participação social de pessoas com deficiência intelectual em eventos públicos de grande}
\chapter{Comunicação alternativa e aumentativa na participação social de pessoas com deficiência intelectual em eventos públicos de grande porte}
\chapter{Psicopedagogicamente inconstitucionalíssimo otorrinolaringologista}
\chapter{Psicopedagogicamente inconstitucionalíssimo otorrinolaringologista socioeducativo}
\chapter{Psicopedagogicamente inconstitucionalíssimo otorrinolaringologista socioeducativo ético-profissional}
\chapter{Psicopedagogicamente inconstitucionalíssimo otorrinolaringologista socioeducativo ético-profissional sociodemográfico}
\chapter{Psicopedagogicamente inconstitucionalíssimo otorrinolaringologista socioeducativo ético-profissional sociodemográfico interdisciplinaridade}
\chapter{Psicopedagogicamente inconstitucionalíssimo otorrinolaringologista socioeducativo ético-profissional sociodemográfico interdisciplinaridade neuropsicológica}
\chapter{Psicopedagogicamente inconstitucionalíssimo otorrinolaringologista socioeducativo ético-profissional sociodemográfico interdisciplinaridade neuropsicológica avaliação}
\chapter{Psicopedagogicamente inconstitucionalíssimo otorrinolaringologista socioeducativo ético-profissional sociodemográfico interdisciplinaridade neuropsicológica avaliação}
\chapter{Psicopedagogicamente inconstitucionalíssimo otorrinolaringologista socioeducativo ético-profissional sociodemográfico interdisciplinaridade neuropsicológica avaliação}
\chapter{Psicopedagogicamente inconstitucionalíssimo otorrinolaringologista socioeducativo ético-profissional sociodemográfico interdisciplinaridade neuropsicológica avaliação}
\chapter{Psicopedagogicamente inconstitucionalíssimo otorrinolaringologista socioeducativo ético-profissional sociodemográfico interdisciplinaridade neuropsicológica avaliação}
\chapter{Psicopedagogicamente inconstitucionalíssimo otorrinolaringologista socioeducativo ético-profissional sociodemográfico interdisciplinaridade neuropsicológica avaliação}
\chapter{Psicopedagogicamente inconstitucionalíssimo otorrinolaringologista socioeducativo ético-profissional sociodemográfico interdisciplinaridade neuropsicológica avaliação}
\chapter{Psicopedagogicamente inconstitucionalíssimo otorrinolaringologista socioeducativo ético-profissional sociodemográfico interdisciplinaridade neuropsicológica avaliação}
\chapter{Psicopedagogicamente inconstitucionalíssimo otorrinolaringologista socioeducativo ético-profissional sociodemográfico interdisciplinaridade neuropsicológica avaliação}
\chapter{Psicopedagogicamente inconstitucionalíssimo otorrinolaringologista socioeducativo ético-profissional sociodemográfico interdisciplinaridade neuropsicológica avaliação}
\end{document}
'   % R: rígido (até 0.5.x)
%   lualatex '\def\var{B}% !TeX program = lualatex
% ===================================================================
% sumario-longo.tex — mede o sumário sem hifenização (issue #35, "Fora do
% escopo"): com a hifenização desligada, a entrada longa de capítulo pode
% bater no número da página. Gera 3 x 16 capítulos, cada um com um prefixo
% (de 3 a 18 palavras) de uma frase longa: frase 1 de palavras médias,
% frase 2 com números, frase 3 só de palavras muito longas. Conta no log
% os "Overfull \hbox" do sumário: o sumário é a única parte do documento
% sem parágrafo de corpo, então todo Overfull "in paragraph" antes do
% primeiro \chapter é entrada que passou da margem.
%
% Uso (de uma pasta com livros_crp_acessivel_book.sty e crpsp-leg.sty no
% caminho; as fontes NEWJUNE por OSFONTDIR):
%   lualatex '\def\var{R}\input{sumario-longo}'   % R: rígido (até 0.5.x)
%   lualatex '\def\var{B}\input{sumario-longo}'   % B: plus 4em (o do guia)
%   lualatex '\def\var{C}\input{sumario-longo}'   % C: plus 4em + \tolerance 10000
%   lualatex '\def\var{S}\input{sumario-longo}'   % S: o que a .sty define
% Cada variante só redefine o \l@chapter depois da carga da .sty.
% ===================================================================
\DocumentMetadata{lang=pt-BR, pdfstandard=ua-2, pdfversion=2.0, tagging=on}
\documentclass[11pt,a5paper,twoside,openright]{book}
\usepackage{livros_crp_acessivel_book}
\hypersetup{pdftitle={Sumario longo}, pdfauthor={CRP SP}}
\makeatletter
\providecommand\var{S}
\def\crp@vR{R}\def\crp@vB{B}\def\crp@vC{C}
\def\crp@lchap#1{%  #1 = acréscimo à margem direita (rightskip)
  \renewcommand*\l@chapter[2]{%
    \ifnum \c@tocdepth >\m@ne
        \addpenalty{-\@highpenalty}\addvspace{0.5em \@plus\p@}%
        \setlength\@tempdima{1.5em}%
        \begingroup
            \parindent \z@ \rightskip \@pnumwidth #1
            \parfillskip -\@pnumwidth
            \leavevmode \bfseries
            \advance\leftskip\@tempdima \hskip -\leftskip
            ##1\nobreak\hfil
            \nobreak\hb@xt@\@pnumwidth{\hss ##2\kern-\p@\kern\p@}\par
            \penalty\@highpenalty
        \endgroup
    \fi}}
\ifx\var\crp@vR \crp@lchap{}\fi
\ifx\var\crp@vB \crp@lchap{\@plus 4em}\fi
\ifx\var\crp@vC \crp@lchap{\@plus 4em \tolerance=10000 \hbadness=10000 }\fi
\makeatother
\begin{document}
\frontmatter
\tableofcontents
\mainmatter
\chapter{Deficiência psicossocial e}
\chapter{Deficiência psicossocial e transtorno}
\chapter{Deficiência psicossocial e transtorno do}
\chapter{Deficiência psicossocial e transtorno do espectro}
\chapter{Deficiência psicossocial e transtorno do espectro autista:}
\chapter{Deficiência psicossocial e transtorno do espectro autista: orientações}
\chapter{Deficiência psicossocial e transtorno do espectro autista: orientações para}
\chapter{Deficiência psicossocial e transtorno do espectro autista: orientações para a}
\chapter{Deficiência psicossocial e transtorno do espectro autista: orientações para a realização}
\chapter{Deficiência psicossocial e transtorno do espectro autista: orientações para a realização de}
\chapter{Deficiência psicossocial e transtorno do espectro autista: orientações para a realização de eventos}
\chapter{Deficiência psicossocial e transtorno do espectro autista: orientações para a realização de eventos e}
\chapter{Deficiência psicossocial e transtorno do espectro autista: orientações para a realização de eventos e apresentações}
\chapter{Deficiência psicossocial e transtorno do espectro autista: orientações para a realização de eventos e apresentações acessíveis}
\chapter{Deficiência psicossocial e transtorno do espectro autista: orientações para a realização de eventos e apresentações acessíveis em}
\chapter{Deficiência psicossocial e transtorno do espectro autista: orientações para a realização de eventos e apresentações acessíveis em todo}
\chapter{Comunicação alternativa e}
\chapter{Comunicação alternativa e aumentativa}
\chapter{Comunicação alternativa e aumentativa na}
\chapter{Comunicação alternativa e aumentativa na participação}
\chapter{Comunicação alternativa e aumentativa na participação social}
\chapter{Comunicação alternativa e aumentativa na participação social de}
\chapter{Comunicação alternativa e aumentativa na participação social de pessoas}
\chapter{Comunicação alternativa e aumentativa na participação social de pessoas com}
\chapter{Comunicação alternativa e aumentativa na participação social de pessoas com deficiência}
\chapter{Comunicação alternativa e aumentativa na participação social de pessoas com deficiência intelectual}
\chapter{Comunicação alternativa e aumentativa na participação social de pessoas com deficiência intelectual em}
\chapter{Comunicação alternativa e aumentativa na participação social de pessoas com deficiência intelectual em eventos}
\chapter{Comunicação alternativa e aumentativa na participação social de pessoas com deficiência intelectual em eventos públicos}
\chapter{Comunicação alternativa e aumentativa na participação social de pessoas com deficiência intelectual em eventos públicos de}
\chapter{Comunicação alternativa e aumentativa na participação social de pessoas com deficiência intelectual em eventos públicos de grande}
\chapter{Comunicação alternativa e aumentativa na participação social de pessoas com deficiência intelectual em eventos públicos de grande porte}
\chapter{Psicopedagogicamente inconstitucionalíssimo otorrinolaringologista}
\chapter{Psicopedagogicamente inconstitucionalíssimo otorrinolaringologista socioeducativo}
\chapter{Psicopedagogicamente inconstitucionalíssimo otorrinolaringologista socioeducativo ético-profissional}
\chapter{Psicopedagogicamente inconstitucionalíssimo otorrinolaringologista socioeducativo ético-profissional sociodemográfico}
\chapter{Psicopedagogicamente inconstitucionalíssimo otorrinolaringologista socioeducativo ético-profissional sociodemográfico interdisciplinaridade}
\chapter{Psicopedagogicamente inconstitucionalíssimo otorrinolaringologista socioeducativo ético-profissional sociodemográfico interdisciplinaridade neuropsicológica}
\chapter{Psicopedagogicamente inconstitucionalíssimo otorrinolaringologista socioeducativo ético-profissional sociodemográfico interdisciplinaridade neuropsicológica avaliação}
\chapter{Psicopedagogicamente inconstitucionalíssimo otorrinolaringologista socioeducativo ético-profissional sociodemográfico interdisciplinaridade neuropsicológica avaliação}
\chapter{Psicopedagogicamente inconstitucionalíssimo otorrinolaringologista socioeducativo ético-profissional sociodemográfico interdisciplinaridade neuropsicológica avaliação}
\chapter{Psicopedagogicamente inconstitucionalíssimo otorrinolaringologista socioeducativo ético-profissional sociodemográfico interdisciplinaridade neuropsicológica avaliação}
\chapter{Psicopedagogicamente inconstitucionalíssimo otorrinolaringologista socioeducativo ético-profissional sociodemográfico interdisciplinaridade neuropsicológica avaliação}
\chapter{Psicopedagogicamente inconstitucionalíssimo otorrinolaringologista socioeducativo ético-profissional sociodemográfico interdisciplinaridade neuropsicológica avaliação}
\chapter{Psicopedagogicamente inconstitucionalíssimo otorrinolaringologista socioeducativo ético-profissional sociodemográfico interdisciplinaridade neuropsicológica avaliação}
\chapter{Psicopedagogicamente inconstitucionalíssimo otorrinolaringologista socioeducativo ético-profissional sociodemográfico interdisciplinaridade neuropsicológica avaliação}
\chapter{Psicopedagogicamente inconstitucionalíssimo otorrinolaringologista socioeducativo ético-profissional sociodemográfico interdisciplinaridade neuropsicológica avaliação}
\chapter{Psicopedagogicamente inconstitucionalíssimo otorrinolaringologista socioeducativo ético-profissional sociodemográfico interdisciplinaridade neuropsicológica avaliação}
\end{document}
'   % B: plus 4em (o do guia)
%   lualatex '\def\var{C}% !TeX program = lualatex
% ===================================================================
% sumario-longo.tex — mede o sumário sem hifenização (issue #35, "Fora do
% escopo"): com a hifenização desligada, a entrada longa de capítulo pode
% bater no número da página. Gera 3 x 16 capítulos, cada um com um prefixo
% (de 3 a 18 palavras) de uma frase longa: frase 1 de palavras médias,
% frase 2 com números, frase 3 só de palavras muito longas. Conta no log
% os "Overfull \hbox" do sumário: o sumário é a única parte do documento
% sem parágrafo de corpo, então todo Overfull "in paragraph" antes do
% primeiro \chapter é entrada que passou da margem.
%
% Uso (de uma pasta com livros_crp_acessivel_book.sty e crpsp-leg.sty no
% caminho; as fontes NEWJUNE por OSFONTDIR):
%   lualatex '\def\var{R}\input{sumario-longo}'   % R: rígido (até 0.5.x)
%   lualatex '\def\var{B}\input{sumario-longo}'   % B: plus 4em (o do guia)
%   lualatex '\def\var{C}\input{sumario-longo}'   % C: plus 4em + \tolerance 10000
%   lualatex '\def\var{S}\input{sumario-longo}'   % S: o que a .sty define
% Cada variante só redefine o \l@chapter depois da carga da .sty.
% ===================================================================
\DocumentMetadata{lang=pt-BR, pdfstandard=ua-2, pdfversion=2.0, tagging=on}
\documentclass[11pt,a5paper,twoside,openright]{book}
\usepackage{livros_crp_acessivel_book}
\hypersetup{pdftitle={Sumario longo}, pdfauthor={CRP SP}}
\makeatletter
\providecommand\var{S}
\def\crp@vR{R}\def\crp@vB{B}\def\crp@vC{C}
\def\crp@lchap#1{%  #1 = acréscimo à margem direita (rightskip)
  \renewcommand*\l@chapter[2]{%
    \ifnum \c@tocdepth >\m@ne
        \addpenalty{-\@highpenalty}\addvspace{0.5em \@plus\p@}%
        \setlength\@tempdima{1.5em}%
        \begingroup
            \parindent \z@ \rightskip \@pnumwidth #1
            \parfillskip -\@pnumwidth
            \leavevmode \bfseries
            \advance\leftskip\@tempdima \hskip -\leftskip
            ##1\nobreak\hfil
            \nobreak\hb@xt@\@pnumwidth{\hss ##2\kern-\p@\kern\p@}\par
            \penalty\@highpenalty
        \endgroup
    \fi}}
\ifx\var\crp@vR \crp@lchap{}\fi
\ifx\var\crp@vB \crp@lchap{\@plus 4em}\fi
\ifx\var\crp@vC \crp@lchap{\@plus 4em \tolerance=10000 \hbadness=10000 }\fi
\makeatother
\begin{document}
\frontmatter
\tableofcontents
\mainmatter
\chapter{Deficiência psicossocial e}
\chapter{Deficiência psicossocial e transtorno}
\chapter{Deficiência psicossocial e transtorno do}
\chapter{Deficiência psicossocial e transtorno do espectro}
\chapter{Deficiência psicossocial e transtorno do espectro autista:}
\chapter{Deficiência psicossocial e transtorno do espectro autista: orientações}
\chapter{Deficiência psicossocial e transtorno do espectro autista: orientações para}
\chapter{Deficiência psicossocial e transtorno do espectro autista: orientações para a}
\chapter{Deficiência psicossocial e transtorno do espectro autista: orientações para a realização}
\chapter{Deficiência psicossocial e transtorno do espectro autista: orientações para a realização de}
\chapter{Deficiência psicossocial e transtorno do espectro autista: orientações para a realização de eventos}
\chapter{Deficiência psicossocial e transtorno do espectro autista: orientações para a realização de eventos e}
\chapter{Deficiência psicossocial e transtorno do espectro autista: orientações para a realização de eventos e apresentações}
\chapter{Deficiência psicossocial e transtorno do espectro autista: orientações para a realização de eventos e apresentações acessíveis}
\chapter{Deficiência psicossocial e transtorno do espectro autista: orientações para a realização de eventos e apresentações acessíveis em}
\chapter{Deficiência psicossocial e transtorno do espectro autista: orientações para a realização de eventos e apresentações acessíveis em todo}
\chapter{Comunicação alternativa e}
\chapter{Comunicação alternativa e aumentativa}
\chapter{Comunicação alternativa e aumentativa na}
\chapter{Comunicação alternativa e aumentativa na participação}
\chapter{Comunicação alternativa e aumentativa na participação social}
\chapter{Comunicação alternativa e aumentativa na participação social de}
\chapter{Comunicação alternativa e aumentativa na participação social de pessoas}
\chapter{Comunicação alternativa e aumentativa na participação social de pessoas com}
\chapter{Comunicação alternativa e aumentativa na participação social de pessoas com deficiência}
\chapter{Comunicação alternativa e aumentativa na participação social de pessoas com deficiência intelectual}
\chapter{Comunicação alternativa e aumentativa na participação social de pessoas com deficiência intelectual em}
\chapter{Comunicação alternativa e aumentativa na participação social de pessoas com deficiência intelectual em eventos}
\chapter{Comunicação alternativa e aumentativa na participação social de pessoas com deficiência intelectual em eventos públicos}
\chapter{Comunicação alternativa e aumentativa na participação social de pessoas com deficiência intelectual em eventos públicos de}
\chapter{Comunicação alternativa e aumentativa na participação social de pessoas com deficiência intelectual em eventos públicos de grande}
\chapter{Comunicação alternativa e aumentativa na participação social de pessoas com deficiência intelectual em eventos públicos de grande porte}
\chapter{Psicopedagogicamente inconstitucionalíssimo otorrinolaringologista}
\chapter{Psicopedagogicamente inconstitucionalíssimo otorrinolaringologista socioeducativo}
\chapter{Psicopedagogicamente inconstitucionalíssimo otorrinolaringologista socioeducativo ético-profissional}
\chapter{Psicopedagogicamente inconstitucionalíssimo otorrinolaringologista socioeducativo ético-profissional sociodemográfico}
\chapter{Psicopedagogicamente inconstitucionalíssimo otorrinolaringologista socioeducativo ético-profissional sociodemográfico interdisciplinaridade}
\chapter{Psicopedagogicamente inconstitucionalíssimo otorrinolaringologista socioeducativo ético-profissional sociodemográfico interdisciplinaridade neuropsicológica}
\chapter{Psicopedagogicamente inconstitucionalíssimo otorrinolaringologista socioeducativo ético-profissional sociodemográfico interdisciplinaridade neuropsicológica avaliação}
\chapter{Psicopedagogicamente inconstitucionalíssimo otorrinolaringologista socioeducativo ético-profissional sociodemográfico interdisciplinaridade neuropsicológica avaliação}
\chapter{Psicopedagogicamente inconstitucionalíssimo otorrinolaringologista socioeducativo ético-profissional sociodemográfico interdisciplinaridade neuropsicológica avaliação}
\chapter{Psicopedagogicamente inconstitucionalíssimo otorrinolaringologista socioeducativo ético-profissional sociodemográfico interdisciplinaridade neuropsicológica avaliação}
\chapter{Psicopedagogicamente inconstitucionalíssimo otorrinolaringologista socioeducativo ético-profissional sociodemográfico interdisciplinaridade neuropsicológica avaliação}
\chapter{Psicopedagogicamente inconstitucionalíssimo otorrinolaringologista socioeducativo ético-profissional sociodemográfico interdisciplinaridade neuropsicológica avaliação}
\chapter{Psicopedagogicamente inconstitucionalíssimo otorrinolaringologista socioeducativo ético-profissional sociodemográfico interdisciplinaridade neuropsicológica avaliação}
\chapter{Psicopedagogicamente inconstitucionalíssimo otorrinolaringologista socioeducativo ético-profissional sociodemográfico interdisciplinaridade neuropsicológica avaliação}
\chapter{Psicopedagogicamente inconstitucionalíssimo otorrinolaringologista socioeducativo ético-profissional sociodemográfico interdisciplinaridade neuropsicológica avaliação}
\chapter{Psicopedagogicamente inconstitucionalíssimo otorrinolaringologista socioeducativo ético-profissional sociodemográfico interdisciplinaridade neuropsicológica avaliação}
\end{document}
'   % C: plus 4em + \tolerance 10000
%   lualatex '\def\var{S}% !TeX program = lualatex
% ===================================================================
% sumario-longo.tex — mede o sumário sem hifenização (issue #35, "Fora do
% escopo"): com a hifenização desligada, a entrada longa de capítulo pode
% bater no número da página. Gera 3 x 16 capítulos, cada um com um prefixo
% (de 3 a 18 palavras) de uma frase longa: frase 1 de palavras médias,
% frase 2 com números, frase 3 só de palavras muito longas. Conta no log
% os "Overfull \hbox" do sumário: o sumário é a única parte do documento
% sem parágrafo de corpo, então todo Overfull "in paragraph" antes do
% primeiro \chapter é entrada que passou da margem.
%
% Uso (de uma pasta com livros_crp_acessivel_book.sty e crpsp-leg.sty no
% caminho; as fontes NEWJUNE por OSFONTDIR):
%   lualatex '\def\var{R}\input{sumario-longo}'   % R: rígido (até 0.5.x)
%   lualatex '\def\var{B}\input{sumario-longo}'   % B: plus 4em (o do guia)
%   lualatex '\def\var{C}\input{sumario-longo}'   % C: plus 4em + \tolerance 10000
%   lualatex '\def\var{S}\input{sumario-longo}'   % S: o que a .sty define
% Cada variante só redefine o \l@chapter depois da carga da .sty.
% ===================================================================
\DocumentMetadata{lang=pt-BR, pdfstandard=ua-2, pdfversion=2.0, tagging=on}
\documentclass[11pt,a5paper,twoside,openright]{book}
\usepackage{livros_crp_acessivel_book}
\hypersetup{pdftitle={Sumario longo}, pdfauthor={CRP SP}}
\makeatletter
\providecommand\var{S}
\def\crp@vR{R}\def\crp@vB{B}\def\crp@vC{C}
\def\crp@lchap#1{%  #1 = acréscimo à margem direita (rightskip)
  \renewcommand*\l@chapter[2]{%
    \ifnum \c@tocdepth >\m@ne
        \addpenalty{-\@highpenalty}\addvspace{0.5em \@plus\p@}%
        \setlength\@tempdima{1.5em}%
        \begingroup
            \parindent \z@ \rightskip \@pnumwidth #1
            \parfillskip -\@pnumwidth
            \leavevmode \bfseries
            \advance\leftskip\@tempdima \hskip -\leftskip
            ##1\nobreak\hfil
            \nobreak\hb@xt@\@pnumwidth{\hss ##2\kern-\p@\kern\p@}\par
            \penalty\@highpenalty
        \endgroup
    \fi}}
\ifx\var\crp@vR \crp@lchap{}\fi
\ifx\var\crp@vB \crp@lchap{\@plus 4em}\fi
\ifx\var\crp@vC \crp@lchap{\@plus 4em \tolerance=10000 \hbadness=10000 }\fi
\makeatother
\begin{document}
\frontmatter
\tableofcontents
\mainmatter
\chapter{Deficiência psicossocial e}
\chapter{Deficiência psicossocial e transtorno}
\chapter{Deficiência psicossocial e transtorno do}
\chapter{Deficiência psicossocial e transtorno do espectro}
\chapter{Deficiência psicossocial e transtorno do espectro autista:}
\chapter{Deficiência psicossocial e transtorno do espectro autista: orientações}
\chapter{Deficiência psicossocial e transtorno do espectro autista: orientações para}
\chapter{Deficiência psicossocial e transtorno do espectro autista: orientações para a}
\chapter{Deficiência psicossocial e transtorno do espectro autista: orientações para a realização}
\chapter{Deficiência psicossocial e transtorno do espectro autista: orientações para a realização de}
\chapter{Deficiência psicossocial e transtorno do espectro autista: orientações para a realização de eventos}
\chapter{Deficiência psicossocial e transtorno do espectro autista: orientações para a realização de eventos e}
\chapter{Deficiência psicossocial e transtorno do espectro autista: orientações para a realização de eventos e apresentações}
\chapter{Deficiência psicossocial e transtorno do espectro autista: orientações para a realização de eventos e apresentações acessíveis}
\chapter{Deficiência psicossocial e transtorno do espectro autista: orientações para a realização de eventos e apresentações acessíveis em}
\chapter{Deficiência psicossocial e transtorno do espectro autista: orientações para a realização de eventos e apresentações acessíveis em todo}
\chapter{Comunicação alternativa e}
\chapter{Comunicação alternativa e aumentativa}
\chapter{Comunicação alternativa e aumentativa na}
\chapter{Comunicação alternativa e aumentativa na participação}
\chapter{Comunicação alternativa e aumentativa na participação social}
\chapter{Comunicação alternativa e aumentativa na participação social de}
\chapter{Comunicação alternativa e aumentativa na participação social de pessoas}
\chapter{Comunicação alternativa e aumentativa na participação social de pessoas com}
\chapter{Comunicação alternativa e aumentativa na participação social de pessoas com deficiência}
\chapter{Comunicação alternativa e aumentativa na participação social de pessoas com deficiência intelectual}
\chapter{Comunicação alternativa e aumentativa na participação social de pessoas com deficiência intelectual em}
\chapter{Comunicação alternativa e aumentativa na participação social de pessoas com deficiência intelectual em eventos}
\chapter{Comunicação alternativa e aumentativa na participação social de pessoas com deficiência intelectual em eventos públicos}
\chapter{Comunicação alternativa e aumentativa na participação social de pessoas com deficiência intelectual em eventos públicos de}
\chapter{Comunicação alternativa e aumentativa na participação social de pessoas com deficiência intelectual em eventos públicos de grande}
\chapter{Comunicação alternativa e aumentativa na participação social de pessoas com deficiência intelectual em eventos públicos de grande porte}
\chapter{Psicopedagogicamente inconstitucionalíssimo otorrinolaringologista}
\chapter{Psicopedagogicamente inconstitucionalíssimo otorrinolaringologista socioeducativo}
\chapter{Psicopedagogicamente inconstitucionalíssimo otorrinolaringologista socioeducativo ético-profissional}
\chapter{Psicopedagogicamente inconstitucionalíssimo otorrinolaringologista socioeducativo ético-profissional sociodemográfico}
\chapter{Psicopedagogicamente inconstitucionalíssimo otorrinolaringologista socioeducativo ético-profissional sociodemográfico interdisciplinaridade}
\chapter{Psicopedagogicamente inconstitucionalíssimo otorrinolaringologista socioeducativo ético-profissional sociodemográfico interdisciplinaridade neuropsicológica}
\chapter{Psicopedagogicamente inconstitucionalíssimo otorrinolaringologista socioeducativo ético-profissional sociodemográfico interdisciplinaridade neuropsicológica avaliação}
\chapter{Psicopedagogicamente inconstitucionalíssimo otorrinolaringologista socioeducativo ético-profissional sociodemográfico interdisciplinaridade neuropsicológica avaliação}
\chapter{Psicopedagogicamente inconstitucionalíssimo otorrinolaringologista socioeducativo ético-profissional sociodemográfico interdisciplinaridade neuropsicológica avaliação}
\chapter{Psicopedagogicamente inconstitucionalíssimo otorrinolaringologista socioeducativo ético-profissional sociodemográfico interdisciplinaridade neuropsicológica avaliação}
\chapter{Psicopedagogicamente inconstitucionalíssimo otorrinolaringologista socioeducativo ético-profissional sociodemográfico interdisciplinaridade neuropsicológica avaliação}
\chapter{Psicopedagogicamente inconstitucionalíssimo otorrinolaringologista socioeducativo ético-profissional sociodemográfico interdisciplinaridade neuropsicológica avaliação}
\chapter{Psicopedagogicamente inconstitucionalíssimo otorrinolaringologista socioeducativo ético-profissional sociodemográfico interdisciplinaridade neuropsicológica avaliação}
\chapter{Psicopedagogicamente inconstitucionalíssimo otorrinolaringologista socioeducativo ético-profissional sociodemográfico interdisciplinaridade neuropsicológica avaliação}
\chapter{Psicopedagogicamente inconstitucionalíssimo otorrinolaringologista socioeducativo ético-profissional sociodemográfico interdisciplinaridade neuropsicológica avaliação}
\chapter{Psicopedagogicamente inconstitucionalíssimo otorrinolaringologista socioeducativo ético-profissional sociodemográfico interdisciplinaridade neuropsicológica avaliação}
\end{document}
'   % S: o que a .sty define
% Cada variante só redefine o \l@chapter depois da carga da .sty.
% ===================================================================
\DocumentMetadata{lang=pt-BR, pdfstandard=ua-2, pdfversion=2.0, tagging=on}
\documentclass[11pt,a5paper,twoside,openright]{book}
\usepackage{livros_crp_acessivel_book}
\hypersetup{pdftitle={Sumario longo}, pdfauthor={CRP SP}}
\makeatletter
\providecommand\var{S}
\def\crp@vR{R}\def\crp@vB{B}\def\crp@vC{C}
\def\crp@lchap#1{%  #1 = acréscimo à margem direita (rightskip)
  \renewcommand*\l@chapter[2]{%
    \ifnum \c@tocdepth >\m@ne
        \addpenalty{-\@highpenalty}\addvspace{0.5em \@plus\p@}%
        \setlength\@tempdima{1.5em}%
        \begingroup
            \parindent \z@ \rightskip \@pnumwidth #1
            \parfillskip -\@pnumwidth
            \leavevmode \bfseries
            \advance\leftskip\@tempdima \hskip -\leftskip
            ##1\nobreak\hfil
            \nobreak\hb@xt@\@pnumwidth{\hss ##2\kern-\p@\kern\p@}\par
            \penalty\@highpenalty
        \endgroup
    \fi}}
\ifx\var\crp@vR \crp@lchap{}\fi
\ifx\var\crp@vB \crp@lchap{\@plus 4em}\fi
\ifx\var\crp@vC \crp@lchap{\@plus 4em \tolerance=10000 \hbadness=10000 }\fi
\makeatother
\begin{document}
\frontmatter
\tableofcontents
\mainmatter
\chapter{Deficiência psicossocial e}
\chapter{Deficiência psicossocial e transtorno}
\chapter{Deficiência psicossocial e transtorno do}
\chapter{Deficiência psicossocial e transtorno do espectro}
\chapter{Deficiência psicossocial e transtorno do espectro autista:}
\chapter{Deficiência psicossocial e transtorno do espectro autista: orientações}
\chapter{Deficiência psicossocial e transtorno do espectro autista: orientações para}
\chapter{Deficiência psicossocial e transtorno do espectro autista: orientações para a}
\chapter{Deficiência psicossocial e transtorno do espectro autista: orientações para a realização}
\chapter{Deficiência psicossocial e transtorno do espectro autista: orientações para a realização de}
\chapter{Deficiência psicossocial e transtorno do espectro autista: orientações para a realização de eventos}
\chapter{Deficiência psicossocial e transtorno do espectro autista: orientações para a realização de eventos e}
\chapter{Deficiência psicossocial e transtorno do espectro autista: orientações para a realização de eventos e apresentações}
\chapter{Deficiência psicossocial e transtorno do espectro autista: orientações para a realização de eventos e apresentações acessíveis}
\chapter{Deficiência psicossocial e transtorno do espectro autista: orientações para a realização de eventos e apresentações acessíveis em}
\chapter{Deficiência psicossocial e transtorno do espectro autista: orientações para a realização de eventos e apresentações acessíveis em todo}
\chapter{Comunicação alternativa e}
\chapter{Comunicação alternativa e aumentativa}
\chapter{Comunicação alternativa e aumentativa na}
\chapter{Comunicação alternativa e aumentativa na participação}
\chapter{Comunicação alternativa e aumentativa na participação social}
\chapter{Comunicação alternativa e aumentativa na participação social de}
\chapter{Comunicação alternativa e aumentativa na participação social de pessoas}
\chapter{Comunicação alternativa e aumentativa na participação social de pessoas com}
\chapter{Comunicação alternativa e aumentativa na participação social de pessoas com deficiência}
\chapter{Comunicação alternativa e aumentativa na participação social de pessoas com deficiência intelectual}
\chapter{Comunicação alternativa e aumentativa na participação social de pessoas com deficiência intelectual em}
\chapter{Comunicação alternativa e aumentativa na participação social de pessoas com deficiência intelectual em eventos}
\chapter{Comunicação alternativa e aumentativa na participação social de pessoas com deficiência intelectual em eventos públicos}
\chapter{Comunicação alternativa e aumentativa na participação social de pessoas com deficiência intelectual em eventos públicos de}
\chapter{Comunicação alternativa e aumentativa na participação social de pessoas com deficiência intelectual em eventos públicos de grande}
\chapter{Comunicação alternativa e aumentativa na participação social de pessoas com deficiência intelectual em eventos públicos de grande porte}
\chapter{Psicopedagogicamente inconstitucionalíssimo otorrinolaringologista}
\chapter{Psicopedagogicamente inconstitucionalíssimo otorrinolaringologista socioeducativo}
\chapter{Psicopedagogicamente inconstitucionalíssimo otorrinolaringologista socioeducativo ético-profissional}
\chapter{Psicopedagogicamente inconstitucionalíssimo otorrinolaringologista socioeducativo ético-profissional sociodemográfico}
\chapter{Psicopedagogicamente inconstitucionalíssimo otorrinolaringologista socioeducativo ético-profissional sociodemográfico interdisciplinaridade}
\chapter{Psicopedagogicamente inconstitucionalíssimo otorrinolaringologista socioeducativo ético-profissional sociodemográfico interdisciplinaridade neuropsicológica}
\chapter{Psicopedagogicamente inconstitucionalíssimo otorrinolaringologista socioeducativo ético-profissional sociodemográfico interdisciplinaridade neuropsicológica avaliação}
\chapter{Psicopedagogicamente inconstitucionalíssimo otorrinolaringologista socioeducativo ético-profissional sociodemográfico interdisciplinaridade neuropsicológica avaliação}
\chapter{Psicopedagogicamente inconstitucionalíssimo otorrinolaringologista socioeducativo ético-profissional sociodemográfico interdisciplinaridade neuropsicológica avaliação}
\chapter{Psicopedagogicamente inconstitucionalíssimo otorrinolaringologista socioeducativo ético-profissional sociodemográfico interdisciplinaridade neuropsicológica avaliação}
\chapter{Psicopedagogicamente inconstitucionalíssimo otorrinolaringologista socioeducativo ético-profissional sociodemográfico interdisciplinaridade neuropsicológica avaliação}
\chapter{Psicopedagogicamente inconstitucionalíssimo otorrinolaringologista socioeducativo ético-profissional sociodemográfico interdisciplinaridade neuropsicológica avaliação}
\chapter{Psicopedagogicamente inconstitucionalíssimo otorrinolaringologista socioeducativo ético-profissional sociodemográfico interdisciplinaridade neuropsicológica avaliação}
\chapter{Psicopedagogicamente inconstitucionalíssimo otorrinolaringologista socioeducativo ético-profissional sociodemográfico interdisciplinaridade neuropsicológica avaliação}
\chapter{Psicopedagogicamente inconstitucionalíssimo otorrinolaringologista socioeducativo ético-profissional sociodemográfico interdisciplinaridade neuropsicológica avaliação}
\chapter{Psicopedagogicamente inconstitucionalíssimo otorrinolaringologista socioeducativo ético-profissional sociodemográfico interdisciplinaridade neuropsicológica avaliação}
\end{document}
'   % S: o que a .sty define
% Cada variante só redefine o \l@chapter depois da carga da .sty.
% ===================================================================
\DocumentMetadata{lang=pt-BR, pdfstandard=ua-2, pdfversion=2.0, tagging=on}
\documentclass[11pt,a5paper,twoside,openright]{book}
\usepackage{livros_crp_acessivel_book}
\hypersetup{pdftitle={Sumario longo}, pdfauthor={CRP SP}}
\makeatletter
\providecommand\var{S}
\def\crp@vR{R}\def\crp@vB{B}\def\crp@vC{C}
\def\crp@lchap#1{%  #1 = acréscimo à margem direita (rightskip)
  \renewcommand*\l@chapter[2]{%
    \ifnum \c@tocdepth >\m@ne
        \addpenalty{-\@highpenalty}\addvspace{0.5em \@plus\p@}%
        \setlength\@tempdima{1.5em}%
        \begingroup
            \parindent \z@ \rightskip \@pnumwidth #1
            \parfillskip -\@pnumwidth
            \leavevmode \bfseries
            \advance\leftskip\@tempdima \hskip -\leftskip
            ##1\nobreak\hfil
            \nobreak\hb@xt@\@pnumwidth{\hss ##2\kern-\p@\kern\p@}\par
            \penalty\@highpenalty
        \endgroup
    \fi}}
\ifx\var\crp@vR \crp@lchap{}\fi
\ifx\var\crp@vB \crp@lchap{\@plus 4em}\fi
\ifx\var\crp@vC \crp@lchap{\@plus 4em \tolerance=10000 \hbadness=10000 }\fi
\makeatother
\begin{document}
\frontmatter
\tableofcontents
\mainmatter
\chapter{Deficiência psicossocial e}
\chapter{Deficiência psicossocial e transtorno}
\chapter{Deficiência psicossocial e transtorno do}
\chapter{Deficiência psicossocial e transtorno do espectro}
\chapter{Deficiência psicossocial e transtorno do espectro autista:}
\chapter{Deficiência psicossocial e transtorno do espectro autista: orientações}
\chapter{Deficiência psicossocial e transtorno do espectro autista: orientações para}
\chapter{Deficiência psicossocial e transtorno do espectro autista: orientações para a}
\chapter{Deficiência psicossocial e transtorno do espectro autista: orientações para a realização}
\chapter{Deficiência psicossocial e transtorno do espectro autista: orientações para a realização de}
\chapter{Deficiência psicossocial e transtorno do espectro autista: orientações para a realização de eventos}
\chapter{Deficiência psicossocial e transtorno do espectro autista: orientações para a realização de eventos e}
\chapter{Deficiência psicossocial e transtorno do espectro autista: orientações para a realização de eventos e apresentações}
\chapter{Deficiência psicossocial e transtorno do espectro autista: orientações para a realização de eventos e apresentações acessíveis}
\chapter{Deficiência psicossocial e transtorno do espectro autista: orientações para a realização de eventos e apresentações acessíveis em}
\chapter{Deficiência psicossocial e transtorno do espectro autista: orientações para a realização de eventos e apresentações acessíveis em todo}
\chapter{Comunicação alternativa e}
\chapter{Comunicação alternativa e aumentativa}
\chapter{Comunicação alternativa e aumentativa na}
\chapter{Comunicação alternativa e aumentativa na participação}
\chapter{Comunicação alternativa e aumentativa na participação social}
\chapter{Comunicação alternativa e aumentativa na participação social de}
\chapter{Comunicação alternativa e aumentativa na participação social de pessoas}
\chapter{Comunicação alternativa e aumentativa na participação social de pessoas com}
\chapter{Comunicação alternativa e aumentativa na participação social de pessoas com deficiência}
\chapter{Comunicação alternativa e aumentativa na participação social de pessoas com deficiência intelectual}
\chapter{Comunicação alternativa e aumentativa na participação social de pessoas com deficiência intelectual em}
\chapter{Comunicação alternativa e aumentativa na participação social de pessoas com deficiência intelectual em eventos}
\chapter{Comunicação alternativa e aumentativa na participação social de pessoas com deficiência intelectual em eventos públicos}
\chapter{Comunicação alternativa e aumentativa na participação social de pessoas com deficiência intelectual em eventos públicos de}
\chapter{Comunicação alternativa e aumentativa na participação social de pessoas com deficiência intelectual em eventos públicos de grande}
\chapter{Comunicação alternativa e aumentativa na participação social de pessoas com deficiência intelectual em eventos públicos de grande porte}
\chapter{Psicopedagogicamente inconstitucionalíssimo otorrinolaringologista}
\chapter{Psicopedagogicamente inconstitucionalíssimo otorrinolaringologista socioeducativo}
\chapter{Psicopedagogicamente inconstitucionalíssimo otorrinolaringologista socioeducativo ético-profissional}
\chapter{Psicopedagogicamente inconstitucionalíssimo otorrinolaringologista socioeducativo ético-profissional sociodemográfico}
\chapter{Psicopedagogicamente inconstitucionalíssimo otorrinolaringologista socioeducativo ético-profissional sociodemográfico interdisciplinaridade}
\chapter{Psicopedagogicamente inconstitucionalíssimo otorrinolaringologista socioeducativo ético-profissional sociodemográfico interdisciplinaridade neuropsicológica}
\chapter{Psicopedagogicamente inconstitucionalíssimo otorrinolaringologista socioeducativo ético-profissional sociodemográfico interdisciplinaridade neuropsicológica avaliação}
\chapter{Psicopedagogicamente inconstitucionalíssimo otorrinolaringologista socioeducativo ético-profissional sociodemográfico interdisciplinaridade neuropsicológica avaliação}
\chapter{Psicopedagogicamente inconstitucionalíssimo otorrinolaringologista socioeducativo ético-profissional sociodemográfico interdisciplinaridade neuropsicológica avaliação}
\chapter{Psicopedagogicamente inconstitucionalíssimo otorrinolaringologista socioeducativo ético-profissional sociodemográfico interdisciplinaridade neuropsicológica avaliação}
\chapter{Psicopedagogicamente inconstitucionalíssimo otorrinolaringologista socioeducativo ético-profissional sociodemográfico interdisciplinaridade neuropsicológica avaliação}
\chapter{Psicopedagogicamente inconstitucionalíssimo otorrinolaringologista socioeducativo ético-profissional sociodemográfico interdisciplinaridade neuropsicológica avaliação}
\chapter{Psicopedagogicamente inconstitucionalíssimo otorrinolaringologista socioeducativo ético-profissional sociodemográfico interdisciplinaridade neuropsicológica avaliação}
\chapter{Psicopedagogicamente inconstitucionalíssimo otorrinolaringologista socioeducativo ético-profissional sociodemográfico interdisciplinaridade neuropsicológica avaliação}
\chapter{Psicopedagogicamente inconstitucionalíssimo otorrinolaringologista socioeducativo ético-profissional sociodemográfico interdisciplinaridade neuropsicológica avaliação}
\chapter{Psicopedagogicamente inconstitucionalíssimo otorrinolaringologista socioeducativo ético-profissional sociodemográfico interdisciplinaridade neuropsicológica avaliação}
\end{document}
'   % S: o que a .sty define
% Cada variante só redefine o \l@chapter depois da carga da .sty.
% ===================================================================
\DocumentMetadata{lang=pt-BR, pdfstandard=ua-2, pdfversion=2.0, tagging=on}
\documentclass[11pt,a5paper,twoside,openright]{book}
\usepackage{livros_crp_acessivel_book}
\hypersetup{pdftitle={Sumario longo}, pdfauthor={CRP SP}}
\makeatletter
\providecommand\var{S}
\def\crp@vR{R}\def\crp@vB{B}\def\crp@vC{C}
\def\crp@lchap#1{%  #1 = acréscimo à margem direita (rightskip)
  \renewcommand*\l@chapter[2]{%
    \ifnum \c@tocdepth >\m@ne
        \addpenalty{-\@highpenalty}\addvspace{0.5em \@plus\p@}%
        \setlength\@tempdima{1.5em}%
        \begingroup
            \parindent \z@ \rightskip \@pnumwidth #1
            \parfillskip -\@pnumwidth
            \leavevmode \bfseries
            \advance\leftskip\@tempdima \hskip -\leftskip
            ##1\nobreak\hfil
            \nobreak\hb@xt@\@pnumwidth{\hss ##2\kern-\p@\kern\p@}\par
            \penalty\@highpenalty
        \endgroup
    \fi}}
\ifx\var\crp@vR \crp@lchap{}\fi
\ifx\var\crp@vB \crp@lchap{\@plus 4em}\fi
\ifx\var\crp@vC \crp@lchap{\@plus 4em \tolerance=10000 \hbadness=10000 }\fi
\makeatother
\begin{document}
\frontmatter
\tableofcontents
\mainmatter
\chapter{Deficiência psicossocial e}
\chapter{Deficiência psicossocial e transtorno}
\chapter{Deficiência psicossocial e transtorno do}
\chapter{Deficiência psicossocial e transtorno do espectro}
\chapter{Deficiência psicossocial e transtorno do espectro autista:}
\chapter{Deficiência psicossocial e transtorno do espectro autista: orientações}
\chapter{Deficiência psicossocial e transtorno do espectro autista: orientações para}
\chapter{Deficiência psicossocial e transtorno do espectro autista: orientações para a}
\chapter{Deficiência psicossocial e transtorno do espectro autista: orientações para a realização}
\chapter{Deficiência psicossocial e transtorno do espectro autista: orientações para a realização de}
\chapter{Deficiência psicossocial e transtorno do espectro autista: orientações para a realização de eventos}
\chapter{Deficiência psicossocial e transtorno do espectro autista: orientações para a realização de eventos e}
\chapter{Deficiência psicossocial e transtorno do espectro autista: orientações para a realização de eventos e apresentações}
\chapter{Deficiência psicossocial e transtorno do espectro autista: orientações para a realização de eventos e apresentações acessíveis}
\chapter{Deficiência psicossocial e transtorno do espectro autista: orientações para a realização de eventos e apresentações acessíveis em}
\chapter{Deficiência psicossocial e transtorno do espectro autista: orientações para a realização de eventos e apresentações acessíveis em todo}
\chapter{Comunicação alternativa e}
\chapter{Comunicação alternativa e aumentativa}
\chapter{Comunicação alternativa e aumentativa na}
\chapter{Comunicação alternativa e aumentativa na participação}
\chapter{Comunicação alternativa e aumentativa na participação social}
\chapter{Comunicação alternativa e aumentativa na participação social de}
\chapter{Comunicação alternativa e aumentativa na participação social de pessoas}
\chapter{Comunicação alternativa e aumentativa na participação social de pessoas com}
\chapter{Comunicação alternativa e aumentativa na participação social de pessoas com deficiência}
\chapter{Comunicação alternativa e aumentativa na participação social de pessoas com deficiência intelectual}
\chapter{Comunicação alternativa e aumentativa na participação social de pessoas com deficiência intelectual em}
\chapter{Comunicação alternativa e aumentativa na participação social de pessoas com deficiência intelectual em eventos}
\chapter{Comunicação alternativa e aumentativa na participação social de pessoas com deficiência intelectual em eventos públicos}
\chapter{Comunicação alternativa e aumentativa na participação social de pessoas com deficiência intelectual em eventos públicos de}
\chapter{Comunicação alternativa e aumentativa na participação social de pessoas com deficiência intelectual em eventos públicos de grande}
\chapter{Comunicação alternativa e aumentativa na participação social de pessoas com deficiência intelectual em eventos públicos de grande porte}
\chapter{Psicopedagogicamente inconstitucionalíssimo otorrinolaringologista}
\chapter{Psicopedagogicamente inconstitucionalíssimo otorrinolaringologista socioeducativo}
\chapter{Psicopedagogicamente inconstitucionalíssimo otorrinolaringologista socioeducativo ético-profissional}
\chapter{Psicopedagogicamente inconstitucionalíssimo otorrinolaringologista socioeducativo ético-profissional sociodemográfico}
\chapter{Psicopedagogicamente inconstitucionalíssimo otorrinolaringologista socioeducativo ético-profissional sociodemográfico interdisciplinaridade}
\chapter{Psicopedagogicamente inconstitucionalíssimo otorrinolaringologista socioeducativo ético-profissional sociodemográfico interdisciplinaridade neuropsicológica}
\chapter{Psicopedagogicamente inconstitucionalíssimo otorrinolaringologista socioeducativo ético-profissional sociodemográfico interdisciplinaridade neuropsicológica avaliação}
\chapter{Psicopedagogicamente inconstitucionalíssimo otorrinolaringologista socioeducativo ético-profissional sociodemográfico interdisciplinaridade neuropsicológica avaliação}
\chapter{Psicopedagogicamente inconstitucionalíssimo otorrinolaringologista socioeducativo ético-profissional sociodemográfico interdisciplinaridade neuropsicológica avaliação}
\chapter{Psicopedagogicamente inconstitucionalíssimo otorrinolaringologista socioeducativo ético-profissional sociodemográfico interdisciplinaridade neuropsicológica avaliação}
\chapter{Psicopedagogicamente inconstitucionalíssimo otorrinolaringologista socioeducativo ético-profissional sociodemográfico interdisciplinaridade neuropsicológica avaliação}
\chapter{Psicopedagogicamente inconstitucionalíssimo otorrinolaringologista socioeducativo ético-profissional sociodemográfico interdisciplinaridade neuropsicológica avaliação}
\chapter{Psicopedagogicamente inconstitucionalíssimo otorrinolaringologista socioeducativo ético-profissional sociodemográfico interdisciplinaridade neuropsicológica avaliação}
\chapter{Psicopedagogicamente inconstitucionalíssimo otorrinolaringologista socioeducativo ético-profissional sociodemográfico interdisciplinaridade neuropsicológica avaliação}
\chapter{Psicopedagogicamente inconstitucionalíssimo otorrinolaringologista socioeducativo ético-profissional sociodemográfico interdisciplinaridade neuropsicológica avaliação}
\chapter{Psicopedagogicamente inconstitucionalíssimo otorrinolaringologista socioeducativo ético-profissional sociodemográfico interdisciplinaridade neuropsicológica avaliação}
\end{document}
